\documentclass[11pt,oneside,leqno]{article}
\usepackage[paperwidth=6.5in,paperheight=9.5in,margin=.72in]{geometry}
\usepackage{fontspec}
\setmainfont{lmroman10-regular.otf}[BoldFont=lmroman10-bold.otf,ItalicFont=lmroman10-italic.otf,BoldItalicFont=lmroman10-bolditalic.otf,SmallCapsFont=lmromancaps10-regular.otf]
\usepackage[german]{babel}
\germansetup{glottonyms=legacy,hyphenrules=legacy}
\usepackage{amsmath,amssymb,microtype,needspace,multirow,booktabs,tikz}
\usepackage[unicode,hidelinks]{hyperref}
\hypersetup{pdftitle={Verwendung complexer Grössen für die Geometrie — Das Dreieck; Pothenots Aufgabe und das Viereck; Der Kreis; Die Kegelschnitte; Projection des Würfels; Geometrische Seite der ternären Formen; Die Kugel; Mutationen des Raumes},pdfauthor={Carl Friedrich Gauss}}
\newcommand{\SourceLeaf}[1]{\hypertarget{native-#1}{}}
\newcommand{\BareRoot}[1]{\mathord{\surd}#1}
\setlength{\parindent}{1.2em}
\setlength{\emergencystretch}{1.4em}
\widowpenalty=10000\clubpenalty=10000\raggedbottom
\newcommand{\HistoricalNote}[3]{\begingroup\renewcommand{\thefootnote}{#1}\footnotemark\renewcommand{\thefootnote}{#2}\footnotetext{#3}\endgroup}
\makeatletter
\long\def\@makefntext#1{\parindent1em\noindent\hb@xt@2.4em{\hss\@thefnmark\ }#1}
\makeatother
\newcommand{\EqualsOverGreater}{\mathrel{\vcenter{\offinterlineskip\hbox{$=$}\kern1pt\hbox{$>$}}}}
\newfontfamily\CyrillicFont{Times New Roman}
\newcommand{\gausscyr}[1]{{\CyrillicFont #1}}
\newcommand{\TripleParallel}{\mathrel{\rotatebox{-12}{$\vert$}\mkern-2mu\rotatebox{-12}{$\vert$}\mkern-2mu\rotatebox{-12}{$\vert$}}}
\begin{document}
\clearpage\SourceLeaf{308}
\thispagestyle{empty}
\vspace*{4cm}
\begin{center}
\textbf{\Large VERWENDUNG COMPLEXER GRÖSSEN}\par\bigskip
\textbf{\Large FÜR DIE GEOMETRIE.}\par\bigskip
NACHTRÄGE ZU BAND IV.
\end{center}
\clearpage
%
\SourceLeaf{310}
\begin{center}
\textbf{\Large NACHLASS UND BRIEFWECHSEL.}\par\bigskip
\textbf{[DAS DREIECK.]}
\end{center}
\bigskip

Die Aufgabe, die Lage eines Punktes aus dem bekannten Verhältniss seiner Abstände von dreien der Lage nach bekannten Punkten zu finden.

Coordinaten der gegebenen Punkte $\left\{\begin{array}{c}a,a',a''\\b,b',b''\end{array}\right\}$; des unbekannten $\left\{\begin{array}{c}x\\y\end{array}\right\}$; Verhältniss der drei Abstände wie $m$, $m'$, $m''$.

\textls{Erste Auflösungsart.}

Man mache
\[
\begin{array}{lll}
a''-a'=h\cos A,&a-a''=h'\cos A',&a'-a=h''\cos A''\\
b''-b'=h\sin A,&b-b''=h'\sin A',&b'-b=h''\sin A''
\end{array}
\]
\noindent und nehme an
\[
\begin{aligned}
x&=a+r\cos\varphi\\
y&=b+r\sin\varphi.
\end{aligned}
\]

So hat man die Gleichungen:
\[
\begin{aligned}
rr-2rh''\cos(\varphi-A'')+h''h''&=\frac{rrm'm'}{mm}\\
rr+2rh'\cos(\varphi-A')+h'h'&=\frac{rrm''m''}{mm}.
\end{aligned}
\]
\begin{center}\resizebox{4.4cm}{!}{%
\begin{tikzpicture}[x=1pt,y=-1pt]
\path[use as bounding box] (0,0) rectangle (850,740);
\path[fill=black,draw=none,even odd rule]
(542,110) -- (535,116) -- (534,119) -- (532,121) -- (530,125) -- (524,131) -- (523,134) -- (516,142) -- (513,148) -- (507,154) -- (507,155) -- (504,158) -- (503,161) -- (495,170) -- (492,176) -- (486,182) -- (485,185) -- (482,188) -- (480,192) -- (466,208) -- (465,211) -- (460,215) -- (460,217) -- (453,227) -- (453,229) -- (451,231) -- (450,234) -- (446,237) -- (446,239) -- (441,244) -- (440,246) -- (440,249) -- (433,254) -- (434,256) -- (432,258) -- (431,258) -- (423,267) -- (423,270) -- (421,272) -- (419,272) -- (418,273) -- (418,275) -- (410,284) -- (409,287) -- (405,292) -- (405,294) -- (401,297) -- (399,301) -- (394,306) -- (393,310) -- (384,318) -- (384,321) -- (381,324) -- (381,326) -- (375,333) -- (373,337) -- (365,345) -- (365,348) -- (364,350) -- (357,356) -- (353,364) -- (348,369) -- (339,383) -- (333,389) -- (333,391) -- (329,395) -- (328,398) -- (321,406) -- (320,409) -- (309,422) -- (309,423) -- (307,425) -- (307,427) -- (303,431) -- (303,433) -- (298,437) -- (297,440) -- (294,443) -- (294,445) -- (282,459) -- (281,462) -- (277,467) -- (276,470) -- (269,478) -- (267,482) -- (264,485) -- (261,491) -- (258,494) -- (256,494) -- (256,498) -- (247,507) -- (247,509) -- (243,513) -- (243,514) -- (241,516) -- (241,518) -- (239,520) -- (238,520) -- (235,523) -- (234,526) -- (229,530) -- (226,537) -- (219,544) -- (219,546) -- (217,548) -- (217,550) -- (213,554) -- (212,557) -- (206,563) -- (204,567) -- (200,571) -- (200,573) -- (192,582) -- (192,583) -- (190,585) -- (190,587) -- (186,591) -- (184,595) -- (180,598) -- (180,600) -- (176,604) -- (175,607) -- (168,615) -- (167,618) -- (164,621) -- (164,623) -- (167,623) -- (180,609) -- (182,605) -- (186,601) -- (188,596) -- (193,591) -- (195,586) -- (201,579) -- (203,575) -- (218,557) -- (219,554) -- (221,552) -- (221,550) -- (231,539) -- (231,538) -- (234,535) -- (236,530) -- (245,520) -- (247,516) -- (254,508) -- (255,505) -- (260,501) -- (260,498) -- (264,495) -- (264,493) -- (270,487) -- (270,485) -- (279,475) -- (279,474) -- (281,472) -- (281,470) -- (285,465) -- (285,463) -- (288,458) -- (295,451) -- (295,449) -- (301,443) -- (302,440) -- (310,431) -- (312,427) -- (319,420) -- (319,418) -- (322,415) -- (324,411) -- (330,405) -- (331,402) -- (335,398) -- (336,394) -- (342,387) -- (344,383) -- (349,378) -- (349,377) -- (351,375) -- (351,373) -- (358,366) -- (359,363) -- (363,359) -- (367,352) -- (370,349) -- (370,347) -- (381,335) -- (381,333) -- (383,331) -- (385,331) -- (388,328) -- (388,325) -- (392,318) -- (397,313) -- (398,310) -- (403,305) -- (408,297) -- (415,289) -- (415,287) -- (418,284) -- (418,283) -- (426,275) -- (426,273) -- (430,269) -- (434,262) -- (447,247) -- (447,245) -- (453,239) -- (453,237) -- (455,235) -- (455,234) -- (459,231) -- (460,227) -- (463,224) -- (465,224) -- (466,223) -- (465,222) -- (465,220) -- (467,218) -- (468,215) -- (472,211) -- (473,208) -- (477,204) -- (477,203) -- (483,197) -- (486,191) -- (491,186) -- (492,183) -- (495,180) -- (495,179) -- (500,174) -- (501,170) -- (511,159) -- (514,153) -- (518,148) -- (519,145) -- (523,142) -- (524,139) -- (536,125) -- (537,125) -- (539,127) -- (539,130) -- (538,131) -- (538,135) -- (537,136) -- (537,139) -- (536,140) -- (537,145) -- (536,146) -- (536,149) -- (535,150) -- (535,157) -- (534,158) -- (534,161) -- (533,162) -- (533,165) -- (534,167) -- (533,168) -- (533,174) -- (532,175) -- (532,187) -- (530,190) -- (530,200) -- (528,203) -- (528,210) -- (527,211) -- (527,224) -- (526,225) -- (526,230) -- (527,231) -- (527,233) -- (526,234) -- (526,242) -- (525,243) -- (525,246) -- (524,247) -- (525,253) -- (523,257) -- (523,265) -- (522,266) -- (522,273) -- (521,274) -- (522,278) -- (521,279) -- (521,285) -- (520,286) -- (520,293) -- (518,295) -- (518,297) -- (519,298) -- (519,303) -- (518,304) -- (518,308) -- (517,309) -- (517,319) -- (516,320) -- (516,326) -- (515,327) -- (515,342) -- (514,343) -- (514,347) -- (515,348) -- (515,350) -- (514,351) -- (514,360) -- (513,361) -- (513,365) -- (511,368) -- (511,383) -- (510,384) -- (510,388) -- (509,389) -- (510,392) -- (509,393) -- (509,397) -- (508,398) -- (509,399) -- (508,400) -- (508,403) -- (509,405) -- (508,406) -- (508,410) -- (507,411) -- (507,414) -- (506,415) -- (506,424) -- (507,426) -- (506,427) -- (506,430) -- (505,431) -- (506,434) -- (504,436) -- (504,442) -- (505,444) -- (503,446) -- (503,452) -- (502,453) -- (503,455) -- (503,458) -- (502,459) -- (502,463) -- (500,465) -- (502,468) -- (502,475) -- (501,476) -- (501,480) -- (499,483) -- (500,487) -- (498,491) -- (499,493) -- (499,496) -- (498,497) -- (498,505) -- (499,506) -- (497,509) -- (497,515) -- (495,518) -- (496,520) -- (496,523) -- (494,526) -- (494,536) -- (492,539) -- (494,547) -- (492,549) -- (492,551) -- (490,553) -- (486,555) -- (476,556) -- (470,559) -- (463,559) -- (460,561) -- (458,561) -- (451,564) -- (446,564) -- (445,565) -- (441,565) -- (440,566) -- (432,566) -- (428,569) -- (416,570) -- (415,571) -- (412,571) -- (411,572) -- (403,572) -- (399,574) -- (394,574) -- (390,576) -- (386,576) -- (385,577) -- (378,577) -- (376,579) -- (371,579) -- (367,581) -- (364,581) -- (361,583) -- (354,583) -- (353,584) -- (350,584) -- (349,585) -- (344,585) -- (343,586) -- (340,586) -- (337,590) -- (327,590) -- (326,591) -- (323,591) -- (319,593) -- (312,593) -- (306,596) -- (302,596) -- (299,598) -- (291,598) -- (290,599) -- (287,599) -- (286,600) -- (281,600) -- (277,602) -- (272,602) -- (271,603) -- (268,603) -- (267,604) -- (263,604) -- (260,606) -- (256,606) -- (255,607) -- (251,607) -- (250,608) -- (247,608) -- (245,607) -- (243,609) -- (233,610) -- (232,611) -- (229,611) -- (225,613) -- (219,613) -- (217,615) -- (214,615) -- (211,617) -- (204,617) -- (203,618) -- (199,618) -- (196,620) -- (191,620) -- (190,621) -- (185,621) -- (184,622) -- (181,622) -- (177,624) -- (170,624) -- (168,626) -- (159,626) -- (158,627) -- (158,631) -- (160,633) -- (162,634) -- (185,634) -- (186,635) -- (188,635) -- (189,634) -- (213,634) -- (214,635) -- (218,635) -- (221,633) -- (226,634) -- (227,633) -- (249,633) -- (250,632) -- (254,634) -- (256,633) -- (293,633) -- (294,634) -- (296,634) -- (297,633) -- (300,633) -- (301,634) -- (305,634) -- (306,633) -- (310,633) -- (311,634) -- (313,633) -- (327,633) -- (328,634) -- (332,634) -- (333,633) -- (338,633) -- (339,632) -- (354,632) -- (355,631) -- (374,631) -- (375,632) -- (384,632) -- (385,631) -- (387,632) -- (391,632) -- (392,631) -- (405,631) -- (407,632) -- (408,631) -- (413,631) -- (414,630) -- (416,630) -- (417,631) -- (439,631) -- (440,632) -- (442,632) -- (443,631) -- (457,631) -- (458,630) -- (463,630) -- (464,631) -- (469,631) -- (471,630) -- (472,631) -- (479,631) -- (480,632) -- (485,632) -- (486,631) -- (500,631) -- (502,632) -- (503,631) -- (508,631) -- (509,630) -- (511,631) -- (527,631) -- (528,630) -- (531,630) -- (532,631) -- (537,631) -- (538,632) -- (540,631) -- (545,631) -- (546,632) -- (548,632) -- (549,631) -- (551,631) -- (552,632) -- (554,631) -- (567,631) -- (569,632) -- (570,631) -- (576,631) -- (577,630) -- (601,630) -- (602,629) -- (604,629) -- (605,630) -- (611,630) -- (612,629) -- (634,629) -- (635,630) -- (637,629) -- (649,629) -- (651,631) -- (654,632) -- (654,622) -- (653,621) -- (653,618) -- (651,615) -- (652,613) -- (652,610) -- (649,604) -- (649,601) -- (651,599) -- (651,598) -- (647,593) -- (647,590) -- (646,589) -- (646,586) -- (645,585) -- (645,579) -- (642,575) -- (641,564) -- (639,562) -- (639,558) -- (641,556) -- (641,552) -- (636,546) -- (637,544) -- (637,541) -- (634,535) -- (634,532) -- (636,530) -- (636,529) -- (632,524) -- (632,519) -- (631,518) -- (630,512) -- (629,511) -- (628,505) -- (627,504) -- (627,498) -- (624,495) -- (625,493) -- (625,488) -- (623,485) -- (623,481) -- (622,480) -- (622,477) -- (621,476) -- (621,473) -- (620,472) -- (620,468) -- (617,465) -- (617,462) -- (616,461) -- (616,453) -- (615,452) -- (615,448) -- (613,446) -- (614,441) -- (612,439) -- (612,431) -- (609,427) -- (609,420) -- (607,416) -- (608,414) -- (608,411) -- (606,408) -- (606,404) -- (604,401) -- (603,390) -- (601,387) -- (601,382) -- (600,381) -- (600,378) -- (602,377) -- (602,376) -- (600,375) -- (598,372) -- (597,356) -- (595,352) -- (595,349) -- (592,345) -- (592,341) -- (590,338) -- (590,331) -- (589,330) -- (589,325) -- (588,324) -- (587,312) -- (585,310) -- (585,307) -- (584,306) -- (584,301) -- (582,298) -- (583,294) -- (581,291) -- (581,285) -- (580,284) -- (580,282) -- (579,281) -- (580,280) -- (578,276) -- (579,273) -- (576,267) -- (577,263) -- (575,259) -- (575,256) -- (574,255) -- (574,252) -- (573,251) -- (573,248) -- (572,247) -- (572,242) -- (571,241) -- (571,238) -- (569,236) -- (569,228) -- (566,225) -- (566,217) -- (564,215) -- (563,204) -- (562,203) -- (563,199) -- (560,196) -- (560,192) -- (559,191) -- (559,183) -- (557,180) -- (557,177) -- (556,176) -- (556,173) -- (555,172) -- (555,169) -- (554,168) -- (554,161) -- (552,158) -- (552,155) -- (551,154) -- (552,151) -- (550,149) -- (550,143) -- (549,142) -- (549,139) -- (547,136) -- (547,133) -- (544,131) -- (544,129) -- (546,127) -- (546,125) -- (545,124) -- (545,120) -- (544,119) -- (545,114) -- cycle
(641,623) -- (641,625) -- (640,626) -- (626,626) -- (625,625) -- (623,626) -- (621,625) -- (619,625) -- (618,626) -- (616,625) -- (595,625) -- (594,626) -- (588,626) -- (587,625) -- (584,625) -- (583,626) -- (577,626) -- (576,627) -- (569,627) -- (568,626) -- (560,626) -- (559,627) -- (534,627) -- (533,626) -- (531,626) -- (530,627) -- (527,627) -- (526,628) -- (518,627) -- (517,626) -- (515,626) -- (514,627) -- (506,627) -- (505,626) -- (503,626) -- (502,627) -- (474,627) -- (473,626) -- (468,627) -- (467,626) -- (461,626) -- (460,627) -- (449,627) -- (448,628) -- (447,627) -- (437,627) -- (436,628) -- (434,628) -- (433,627) -- (423,627) -- (422,628) -- (420,627) -- (418,628) -- (416,627) -- (414,627) -- (413,628) -- (406,628) -- (405,627) -- (395,627) -- (394,628) -- (392,627) -- (381,627) -- (380,628) -- (378,628) -- (377,627) -- (358,627) -- (357,628) -- (353,628) -- (352,627) -- (350,627) -- (349,628) -- (339,628) -- (338,629) -- (336,628) -- (334,628) -- (333,629) -- (330,629) -- (329,628) -- (324,628) -- (323,629) -- (321,629) -- (320,628) -- (311,628) -- (310,629) -- (308,628) -- (303,628) -- (302,629) -- (270,629) -- (269,628) -- (267,629) -- (230,629) -- (229,628) -- (227,628) -- (226,629) -- (212,629) -- (211,630) -- (201,630) -- (200,629) -- (197,629) -- (196,630) -- (184,630) -- (183,629) -- (179,630) -- (177,629) -- (176,630) -- (174,630) -- (173,629) -- (174,627) -- (177,627) -- (178,626) -- (181,626) -- (182,625) -- (193,624) -- (194,623) -- (198,623) -- (199,622) -- (203,622) -- (204,621) -- (210,621) -- (211,620) -- (214,620) -- (219,617) -- (232,616) -- (238,613) -- (245,613) -- (246,612) -- (249,612) -- (250,611) -- (254,611) -- (257,609) -- (265,609) -- (269,607) -- (280,606) -- (281,605) -- (285,605) -- (286,604) -- (295,603) -- (299,601) -- (304,601) -- (311,598) -- (315,598) -- (317,600) -- (318,600) -- (322,596) -- (326,596) -- (327,595) -- (330,595) -- (331,594) -- (336,594) -- (339,592) -- (341,592) -- (343,590) -- (346,590) -- (347,589) -- (350,590) -- (351,589) -- (354,589) -- (355,588) -- (359,588) -- (362,586) -- (368,586) -- (369,585) -- (372,585) -- (375,583) -- (382,583) -- (383,582) -- (392,581) -- (398,578) -- (407,578) -- (410,576) -- (424,575) -- (426,573) -- (431,573) -- (432,572) -- (444,571) -- (447,569) -- (452,569) -- (454,567) -- (457,567) -- (460,565) -- (466,565) -- (470,563) -- (472,563) -- (475,565) -- (478,562) -- (481,562) -- (484,560) -- (487,560) -- (488,559) -- (496,559) -- (497,558) -- (510,565) -- (517,566) -- (522,570) -- (526,571) -- (530,574) -- (535,575) -- (542,580) -- (545,580) -- (548,583) -- (552,583) -- (554,585) -- (556,585) -- (560,589) -- (564,589) -- (567,592) -- (572,593) -- (578,596) -- (581,596) -- (585,599) -- (589,600) -- (596,604) -- (598,604) -- (600,606) -- (602,606) -- (606,609) -- (611,610) -- (615,612) -- (617,614) -- (621,614) -- (623,616) -- (628,617) -- (635,621) -- (639,621) -- cycle
(542,138) -- (544,139) -- (544,143) -- (546,147) -- (545,150) -- (548,156) -- (547,158) -- (547,161) -- (549,165) -- (550,174) -- (552,177) -- (553,189) -- (555,191) -- (554,195) -- (555,196) -- (556,205) -- (558,209) -- (557,213) -- (559,216) -- (558,217) -- (559,218) -- (559,220) -- (561,222) -- (561,227) -- (562,228) -- (563,235) -- (565,238) -- (565,242) -- (566,243) -- (566,247) -- (567,248) -- (567,252) -- (568,253) -- (568,258) -- (569,259) -- (570,266) -- (572,269) -- (572,274) -- (574,278) -- (575,288) -- (577,291) -- (577,299) -- (580,306) -- (580,315) -- (582,318) -- (582,325) -- (584,328) -- (585,337) -- (587,340) -- (587,343) -- (589,347) -- (589,352) -- (591,355) -- (591,361) -- (592,362) -- (592,366) -- (593,367) -- (594,375) -- (595,376) -- (595,379) -- (596,380) -- (596,384) -- (597,385) -- (597,389) -- (598,390) -- (599,399) -- (602,404) -- (601,410) -- (602,411) -- (603,417) -- (604,418) -- (604,423) -- (605,424) -- (605,427) -- (607,430) -- (607,437) -- (609,441) -- (608,442) -- (609,443) -- (609,447) -- (610,448) -- (611,458) -- (613,462) -- (613,465) -- (615,467) -- (616,479) -- (618,483) -- (619,491) -- (621,493) -- (621,498) -- (622,499) -- (622,507) -- (623,508) -- (623,511) -- (625,513) -- (625,519) -- (624,521) -- (626,524) -- (627,531) -- (629,533) -- (629,539) -- (631,543) -- (630,553) -- (633,557) -- (634,567) -- (636,570) -- (636,574) -- (637,575) -- (638,583) -- (640,586) -- (640,588) -- (642,590) -- (642,596) -- (644,598) -- (645,609) -- (647,611) -- (646,614) -- (648,616) -- (648,619) -- (649,620) -- (649,622) -- (647,624) -- (645,624) -- (639,619) -- (635,618) -- (632,615) -- (630,615) -- (626,612) -- (620,609) -- (618,609) -- (614,606) -- (612,606) -- (601,600) -- (596,599) -- (594,597) -- (589,595) -- (587,593) -- (584,593) -- (581,591) -- (579,591) -- (577,589) -- (573,589) -- (569,586) -- (563,585) -- (561,583) -- (558,583) -- (556,581) -- (552,579) -- (545,577) -- (543,575) -- (536,572) -- (532,569) -- (525,567) -- (521,564) -- (515,563) -- (512,561) -- (510,561) -- (504,558) -- (501,555) -- (499,555) -- (497,553) -- (497,546) -- (498,545) -- (498,542) -- (499,541) -- (500,528) -- (501,527) -- (501,523) -- (500,522) -- (500,520) -- (502,516) -- (501,514) -- (502,512) -- (501,507) -- (503,505) -- (504,494) -- (505,493) -- (504,491) -- (504,488) -- (505,487) -- (505,479) -- (509,474) -- (508,472) -- (508,468) -- (506,466) -- (506,462) -- (507,461) -- (507,452) -- (509,448) -- (508,441) -- (510,438) -- (511,432) -- (514,427) -- (514,421) -- (511,418) -- (511,417) -- (514,414) -- (513,412) -- (513,405) -- (514,404) -- (514,399) -- (516,396) -- (516,393) -- (515,392) -- (516,390) -- (516,388) -- (515,387) -- (515,383) -- (518,378) -- (517,376) -- (518,374) -- (517,372) -- (517,369) -- (518,368) -- (519,356) -- (520,355) -- (520,353) -- (519,352) -- (519,350) -- (520,349) -- (520,347) -- (519,346) -- (519,341) -- (520,340) -- (520,336) -- (521,335) -- (521,326) -- (522,325) -- (523,316) -- (527,311) -- (527,309) -- (524,305) -- (524,301) -- (525,300) -- (525,293) -- (528,289) -- (528,284) -- (529,283) -- (529,281) -- (527,279) -- (527,278) -- (530,275) -- (530,273) -- (528,271) -- (528,268) -- (529,267) -- (529,260) -- (530,259) -- (529,253) -- (530,252) -- (530,248) -- (532,246) -- (531,241) -- (535,239) -- (532,233) -- (533,231) -- (533,227) -- (531,224) -- (531,220) -- (533,216) -- (532,210) -- (533,209) -- (533,205) -- (534,204) -- (534,201) -- (535,200) -- (534,198) -- (535,194) -- (533,192) -- (533,190) -- (536,186) -- (536,183) -- (535,182) -- (536,180) -- (536,174) -- (537,173) -- (537,170) -- (536,168) -- (539,164) -- (539,162) -- (538,161) -- (538,159) -- (539,158) -- (539,145) -- (540,144) -- (540,140) -- cycle;
\node[inner sep=0pt,font=\fontsize{10}{10}\selectfont] at (567.5,71.5) {\scalebox{5.6}{$a',b'$}};
\node[inner sep=0pt,font=\fontsize{10}{10}\selectfont] at (329.0,341.0) {\scalebox{5.4}{$h''$}};
\node[inner sep=0pt,font=\fontsize{10}{10}\selectfont] at (621.0,355.5) {\scalebox{6.4}{$h$}};
\node[inner sep=0pt,font=\fontsize{10}{10}\selectfont] at (546.5,537.5) {\scalebox{5.2}{$x,y$}};
\node[inner sep=0pt,font=\fontsize{10}{10}\selectfont] at (315.5,571.5) {\scalebox{4.3}{$r$}};
\node[inner sep=0pt,font=\fontsize{10}{10}\selectfont] at (83.5,640.0) {\scalebox{6.3}{$a,b$}};
\node[inner sep=0pt,font=\fontsize{10}{10}\selectfont] at (720.0,642.0) {\scalebox{5.8}{$a'',b''$}};
\node[inner sep=0pt,font=\fontsize{10}{10}\selectfont] at (437.0,665.0) {\scalebox{5.4}{$h'$}};
\end{tikzpicture}%
}
\end{center}
%
\SourceLeaf{311}

[Multiplicirt man die erste dieser Gleichungen mit
\[
\dfrac{h'h'}{rr},
\]
die zweite mit
\[
\dfrac{h''h''}{rr},
\]
so ergibt sich durch Subtraction, indem man die Identitäten

\[
\dfrac{h}{\sin(A''-A')} = \dfrac{h'}{\sin(A-A'')} = \dfrac{h''}{\sin(A'-A)}
\]

und

\[
\begin{aligned}\sin(A-A')\cos(\varphi-A') &+ \sin(A''-A)\cos(\varphi-A'')\\ &= \sin(A''-A')\cos(\varphi+A-A'-A'')\end{aligned}
\]

zu Hülfe nimmt, die Gleichung:]

\[
\mathrm{I.}\quad \begin{aligned}&\left\{h''h''\left(1-\dfrac{m''m''}{mm}\right) - h'h'\left(1-\dfrac{m'm'}{mm}\right)\right\}r \\
&= 2hh'h''\cos(\varphi+A-A'-A'').\end{aligned}
\]

[Multiplicirt man ferner die erste jener Gleichungen mit
\[
1 - \dfrac{m''m''}{mm},
\]
die zweite mit
\[
1 - \dfrac{m'm'}{mm},
\]
so ergibt sich durch Subtraction die Gleichung:]

\[
\mathrm{II.}\quad 
\begin{aligned}
&\left\{h''h''\left(1-\dfrac{m''m''}{mm}\right) - h'h'\left(1-\dfrac{m'm'}{mm}\right)\right\} \\
&= 2r\left\{h''\left(1-\dfrac{m''m''}{mm}\right)\cos(\varphi-A'')\right.\\ &\qquad\left. + h'\left(1-\dfrac{m'm'}{mm}\right)\cos(\varphi-A')\right\}\\
&= -\dfrac{2r}{mm}\left\{mmh\cos(\varphi-A) + m'm'h'\cos(\varphi-A')\right.\\ &\qquad\left. + m''m''h''\cos(\varphi-A'')\right\}.
\end{aligned}
\]

Durch Multiplication entsteht also:

\[
\mathrm{III.}\quad 
\begin{aligned}
&\left\{h''h''\left(1-\dfrac{m''m''}{mm}\right) - h'h'\left(1-\dfrac{m'm'}{mm}\right)\right\}^2 \\
&- 2hh'h''\left\{h''\left(1-\dfrac{m''m''}{mm}\right)\cos(A-A')\right.\\ &\qquad\left. + h'\left(1-\dfrac{m'm'}{mm}\right)\cos(A-A'')\right\}\\
&= 2hh'h''\left\{h''\left(1-\dfrac{m''m''}{mm}\right)\cos(2\varphi+A-A'-2A'')\right.\\ &\qquad\left. + h'\left(1-\dfrac{m'm'}{mm}\right)\cos(2\varphi+A-2A'-A'')\right\}.
\end{aligned}
\]

Hieraus folgen zwei Werthe für $2\varphi$ und eben so viele für $\varphi$, die in I. ein positives $r$ geben, mithin zwei Auflösungen unsrer Gleichung.
%
\SourceLeaf{312}

\Needspace{8\baselineskip}\textls{Die zweite Auflösungsart.}

Es seien $\xi$, $\eta$ [die] Coordinaten des Mittelpunkts des durch die drei Punkte gelegten Kreises, $\rho$ dessen Halbmesser. Man mache, wie vorhin

\[
a'' - a' = h\cos A, \qquad a - a'' = h'\cos A', \qquad a' - a = h''\cos A'',
\]
\[
b'' - b' = h\sin A, \qquad b - b'' = h'\sin A', \qquad b' - b = h''\sin A''.
\]

Man hat dann

\[
\begin{aligned}
\xi &= a + \rho\sin(A'+A''-A)\\ &= a' + \rho\sin(A''+A-A')\\ &= a'' + \rho\sin(A+A'-A''),
\end{aligned}
\]
\[
\begin{aligned}
\eta &= b - \rho\cos(A'+A''-A)\\ &= b' - \rho\cos(A''+A-A')\\ &= b'' - \rho\cos(A+A'-A''),
\end{aligned}
\]
\[
\rho = \dfrac{h}{\sin(A''-A')} = \dfrac{h'}{\sin(A-A'')} = \dfrac{h''}{\sin(A'-A)}.
\]

\begin{center}\resizebox{5.3cm}{!}{%
\begin{tikzpicture}[x=1pt,y=-1pt]
\path[use as bounding box] (0,0) rectangle (825,765);
\path[fill=black,draw=none,even odd rule]
(303,57) -- (299,59) -- (295,59) -- (294,60) -- (286,61) -- (284,63) -- (278,63) -- (277,64) -- (274,64) -- (273,65) -- (268,64) -- (266,66) -- (262,68) -- (258,68) -- (253,71) -- (249,71) -- (247,73) -- (245,74) -- (242,74) -- (237,77) -- (231,78) -- (227,81) -- (223,82) -- (211,89) -- (203,90) -- (198,95) -- (195,96) -- (191,100) -- (189,100) -- (182,105) -- (179,105) -- (174,111) -- (166,115) -- (159,121) -- (156,122) -- (140,137) -- (138,137) -- (137,139) -- (117,158) -- (117,159) -- (96,181) -- (95,185) -- (91,189) -- (90,192) -- (79,207) -- (79,209) -- (76,216) -- (73,219) -- (71,224) -- (68,227) -- (68,231) -- (64,235) -- (64,239) -- (60,243) -- (60,245) -- (58,247) -- (58,249) -- (56,251) -- (56,254) -- (54,256) -- (54,259) -- (51,263) -- (52,267) -- (49,271) -- (49,274) -- (46,278) -- (46,282) -- (45,283) -- (43,291) -- (41,294) -- (41,298) -- (39,301) -- (39,305) -- (36,308) -- (36,312) -- (35,313) -- (35,316) -- (34,317) -- (35,319) -- (35,321) -- (34,322) -- (34,329) -- (33,330) -- (33,332) -- (30,335) -- (33,342) -- (31,346) -- (31,360) -- (30,361) -- (30,364) -- (31,366) -- (29,368) -- (29,377) -- (30,378) -- (30,383) -- (29,384) -- (29,387) -- (30,388) -- (30,393) -- (31,394) -- (30,403) -- (31,404) -- (31,411) -- (32,412) -- (32,419) -- (33,420) -- (32,428) -- (34,430) -- (35,440) -- (36,441) -- (36,446) -- (38,449) -- (38,454) -- (39,455) -- (40,461) -- (42,465) -- (42,470) -- (45,475) -- (45,478) -- (46,479) -- (46,484) -- (48,487) -- (49,492) -- (53,500) -- (55,508) -- (60,515) -- (60,517) -- (62,520) -- (63,525) -- (67,530) -- (67,533) -- (71,538) -- (75,547) -- (77,549) -- (77,551) -- (84,559) -- (88,567) -- (93,572) -- (93,574) -- (99,580) -- (99,583) -- (102,586) -- (102,587) -- (117,602) -- (120,607) -- (129,616) -- (129,617) -- (151,639) -- (156,641) -- (161,646) -- (164,647) -- (171,653) -- (175,654) -- (182,660) -- (185,660) -- (187,662) -- (190,663) -- (197,669) -- (202,671) -- (206,674) -- (210,675) -- (213,678) -- (216,678) -- (218,680) -- (221,680) -- (223,682) -- (227,683) -- (231,686) -- (234,686) -- (241,691) -- (244,691) -- (245,690) -- (249,691) -- (261,697) -- (263,697) -- (270,700) -- (273,700) -- (277,703) -- (280,703) -- (284,706) -- (291,706) -- (294,709) -- (297,709) -- (298,710) -- (301,709) -- (305,711) -- (307,711) -- (308,710) -- (310,711) -- (316,710) -- (319,712) -- (321,711) -- (323,712) -- (326,712) -- (327,711) -- (334,711) -- (335,712) -- (343,711) -- (345,713) -- (352,713) -- (353,714) -- (355,714) -- (358,712) -- (361,712) -- (364,714) -- (375,714) -- (376,713) -- (406,713) -- (409,711) -- (428,709) -- (429,708) -- (435,707) -- (436,706) -- (442,706) -- (445,704) -- (451,704) -- (454,702) -- (460,702) -- (464,699) -- (467,699) -- (474,696) -- (478,696) -- (483,692) -- (489,691) -- (492,688) -- (499,686) -- (503,683) -- (507,682) -- (509,680) -- (511,680) -- (514,678) -- (516,678) -- (520,676) -- (522,674) -- (525,674) -- (528,670) -- (532,669) -- (537,665) -- (545,661) -- (548,658) -- (550,657) -- (554,657) -- (560,651) -- (564,651) -- (565,650) -- (567,650) -- (571,646) -- (572,643) -- (574,641) -- (574,639) -- (577,636) -- (579,636) -- (591,624) -- (594,624) -- (597,621) -- (597,619) -- (599,617) -- (602,617) -- (607,612) -- (606,608) -- (618,596) -- (619,593) -- (622,590) -- (622,588) -- (630,580) -- (630,579) -- (633,576) -- (634,573) -- (641,565) -- (643,561) -- (650,552) -- (653,546) -- (653,543) -- (656,540) -- (656,538) -- (660,533) -- (661,528) -- (664,525) -- (666,525) -- (667,524) -- (667,520) -- (671,512) -- (671,506) -- (674,501) -- (674,498) -- (677,495) -- (677,491) -- (678,490) -- (678,487) -- (679,486) -- (679,481) -- (682,477) -- (682,471) -- (684,468) -- (684,465) -- (685,464) -- (685,461) -- (686,460) -- (686,455) -- (687,454) -- (688,450) -- (692,447) -- (692,446) -- (689,443) -- (689,441) -- (691,438) -- (691,434) -- (693,432) -- (692,430) -- (692,423) -- (693,422) -- (693,415) -- (694,414) -- (694,410) -- (693,409) -- (693,392) -- (694,391) -- (694,386) -- (695,385) -- (695,378) -- (694,377) -- (694,374) -- (693,373) -- (695,369) -- (693,367) -- (693,365) -- (694,364) -- (694,359) -- (693,358) -- (693,351) -- (692,350) -- (692,342) -- (690,340) -- (690,336) -- (689,335) -- (690,332) -- (687,329) -- (687,321) -- (686,320) -- (686,315) -- (685,314) -- (685,311) -- (683,307) -- (682,299) -- (680,296) -- (680,291) -- (676,283) -- (675,275) -- (672,269) -- (672,265) -- (667,260) -- (667,256) -- (665,254) -- (665,252) -- (664,251) -- (664,247) -- (661,243) -- (660,240) -- (657,237) -- (657,234) -- (655,230) -- (653,228) -- (653,226) -- (649,222) -- (649,220) -- (647,218) -- (642,207) -- (636,201) -- (636,199) -- (625,187) -- (625,184) -- (624,182) -- (620,179) -- (619,176) -- (607,164) -- (607,163) -- (603,159) -- (603,158) -- (583,138) -- (582,138) -- (577,132) -- (573,130) -- (569,126) -- (567,126) -- (555,116) -- (548,112) -- (543,107) -- (537,104) -- (535,104) -- (530,100) -- (522,97) -- (519,94) -- (513,91) -- (506,89) -- (504,87) -- (504,85) -- (503,84) -- (498,85) -- (488,79) -- (485,80) -- (482,77) -- (478,75) -- (472,74) -- (468,71) -- (464,71) -- (462,69) -- (460,69) -- (459,68) -- (451,67) -- (450,66) -- (444,65) -- (443,64) -- (440,64) -- (436,62) -- (431,62) -- (430,61) -- (427,61) -- (423,59) -- (416,59) -- (412,57) -- (409,58) -- (406,56) -- (403,57) -- (402,56) -- (395,56) -- (394,55) -- (380,55) -- (379,54) -- (339,54) -- (337,53) -- (336,54) -- (327,54) -- (326,55) -- (323,55) -- (322,56) -- (311,56) -- (310,57) -- cycle
(607,600) -- (607,601) -- (605,603) -- (604,603) -- (577,629) -- (573,631) -- (567,637) -- (564,638) -- (555,647) -- (547,651) -- (542,656) -- (540,656) -- (535,660) -- (531,660) -- (527,665) -- (518,670) -- (514,671) -- (510,675) -- (502,677) -- (495,681) -- (492,681) -- (489,684) -- (482,686) -- (480,688) -- (476,688) -- (473,690) -- (469,690) -- (468,691) -- (463,692) -- (460,694) -- (453,694) -- (451,693) -- (448,697) -- (446,697) -- (443,699) -- (438,699) -- (437,700) -- (434,700) -- (430,702) -- (417,702) -- (414,704) -- (410,704) -- (407,706) -- (397,706) -- (396,707) -- (387,707) -- (386,708) -- (371,708) -- (370,707) -- (362,707) -- (361,708) -- (359,708) -- (358,707) -- (343,707) -- (342,706) -- (336,706) -- (335,705) -- (312,704) -- (311,703) -- (307,703) -- (304,701) -- (299,701) -- (298,700) -- (294,700) -- (293,699) -- (290,699) -- (289,698) -- (285,698) -- (282,696) -- (273,695) -- (271,693) -- (267,693) -- (264,691) -- (259,690) -- (256,688) -- (253,688) -- (248,685) -- (245,685) -- (244,684) -- (239,683) -- (236,681) -- (231,680) -- (227,677) -- (225,677) -- (223,675) -- (221,675) -- (219,673) -- (217,673) -- (214,671) -- (211,671) -- (202,664) -- (200,664) -- (198,662) -- (195,661) -- (190,656) -- (187,656) -- (184,653) -- (178,650) -- (169,642) -- (161,638) -- (147,625) -- (144,624) -- (141,621) -- (141,620) -- (139,619) -- (136,616) -- (136,615) -- (123,603) -- (123,601) -- (124,600) -- (128,600) -- (129,601) -- (132,601) -- (133,600) -- (135,601) -- (138,601) -- (139,600) -- (141,601) -- (167,601) -- (168,600) -- (170,600) -- (174,602) -- (176,602) -- (177,601) -- (204,601) -- (205,602) -- (223,602) -- (224,601) -- (226,602) -- (243,602) -- (244,601) -- (272,601) -- (273,602) -- (310,602) -- (311,601) -- (313,601) -- (314,602) -- (322,602) -- (323,603) -- (324,603) -- (326,601) -- (335,601) -- (336,600) -- (338,601) -- (340,600) -- (342,601) -- (346,601) -- (347,600) -- (349,600) -- (350,601) -- (356,601) -- (357,600) -- (359,600) -- (360,601) -- (363,601) -- (364,602) -- (370,601) -- (374,603) -- (376,603) -- (380,601) -- (383,601) -- (384,602) -- (392,601) -- (393,602) -- (401,602) -- (403,603) -- (404,602) -- (407,602) -- (408,601) -- (434,601) -- (435,600) -- (437,601) -- (455,601) -- (456,602) -- (459,602) -- (460,601) -- (466,601) -- (467,602) -- (469,601) -- (475,601) -- (476,602) -- (484,601) -- (485,602) -- (489,602) -- (493,600) -- (496,600) -- (497,601) -- (508,601) -- (509,600) -- (512,600) -- (513,601) -- (516,601) -- (517,600) -- (535,600) -- (536,601) -- (539,601) -- (540,600) -- (542,601) -- (545,601) -- (546,600) -- (554,600) -- (555,601) -- (562,601) -- (563,600) -- (566,600) -- (567,601) -- (571,601) -- (572,600) -- (575,600) -- (576,599) -- (587,599) -- (588,598) -- (590,598) -- (591,599) -- (601,599) -- (602,598) -- (605,598) -- cycle
(601,592) -- (601,594) -- (600,595) -- (590,595) -- (589,594) -- (587,594) -- (586,595) -- (577,595) -- (576,594) -- (572,594) -- (571,595) -- (550,595) -- (549,596) -- (547,595) -- (542,595) -- (541,596) -- (539,595) -- (524,595) -- (523,594) -- (521,595) -- (519,594) -- (517,595) -- (501,595) -- (500,594) -- (495,594) -- (494,595) -- (490,595) -- (489,596) -- (485,596) -- (484,595) -- (478,595) -- (477,596) -- (473,596) -- (472,595) -- (467,595) -- (466,596) -- (463,596) -- (462,595) -- (455,595) -- (454,596) -- (447,596) -- (446,595) -- (423,595) -- (422,596) -- (420,595) -- (408,595) -- (407,596) -- (381,596) -- (380,595) -- (377,595) -- (376,596) -- (365,596) -- (364,595) -- (362,595) -- (361,596) -- (356,596) -- (355,595) -- (350,595) -- (349,596) -- (347,596) -- (346,595) -- (326,595) -- (325,596) -- (272,596) -- (270,595) -- (269,596) -- (244,596) -- (243,597) -- (165,597) -- (163,596) -- (160,598) -- (152,598) -- (151,597) -- (144,597) -- (143,598) -- (133,598) -- (131,596) -- (133,594) -- (137,594) -- (138,593) -- (142,593) -- (143,592) -- (149,592) -- (150,591) -- (161,590) -- (163,588) -- (169,587) -- (170,586) -- (176,586) -- (177,585) -- (180,585) -- (181,584) -- (190,584) -- (194,582) -- (200,582) -- (201,581) -- (206,581) -- (210,579) -- (217,579) -- (218,578) -- (222,578) -- (223,577) -- (226,577) -- (227,576) -- (238,575) -- (241,573) -- (243,573) -- (245,571) -- (251,571) -- (252,570) -- (256,570) -- (257,569) -- (261,569) -- (262,568) -- (269,568) -- (276,565) -- (282,565) -- (283,564) -- (294,564) -- (296,562) -- (298,562) -- (299,561) -- (307,561) -- (313,558) -- (324,557) -- (329,554) -- (335,554) -- (336,553) -- (340,553) -- (341,552) -- (347,552) -- (351,550) -- (362,549) -- (365,547) -- (371,547) -- (375,545) -- (382,545) -- (383,544) -- (387,544) -- (388,543) -- (391,543) -- (392,542) -- (401,542) -- (404,540) -- (407,540) -- (413,537) -- (417,537) -- (418,536) -- (427,535) -- (428,534) -- (431,534) -- (435,532) -- (442,531) -- (446,529) -- (450,529) -- (451,528) -- (458,528) -- (459,527) -- (461,527) -- (465,530) -- (469,530) -- (471,532) -- (479,536) -- (481,536) -- (483,538) -- (482,540) -- (485,542) -- (491,541) -- (493,543) -- (497,544) -- (501,546) -- (503,548) -- (511,552) -- (513,552) -- (515,554) -- (518,554) -- (523,558) -- (526,558) -- (528,560) -- (531,560) -- (536,564) -- (538,564) -- (545,568) -- (547,568) -- (552,571) -- (554,571) -- (556,573) -- (563,577) -- (566,577) -- (570,580) -- (576,580) -- (582,585) -- (585,585) -- (592,589) -- (595,589) -- cycle
(502,112) -- (503,112) -- (506,115) -- (506,125) -- (507,126) -- (507,129) -- (508,130) -- (508,134) -- (509,135) -- (510,144) -- (512,147) -- (512,153) -- (513,154) -- (513,158) -- (514,159) -- (514,162) -- (515,163) -- (515,168) -- (517,171) -- (517,177) -- (518,178) -- (518,181) -- (520,185) -- (520,188) -- (522,191) -- (522,201) -- (523,202) -- (523,204) -- (526,207) -- (525,212) -- (527,214) -- (527,216) -- (526,217) -- (526,219) -- (527,220) -- (527,225) -- (529,227) -- (529,234) -- (531,238) -- (531,244) -- (532,245) -- (532,250) -- (533,251) -- (533,254) -- (535,258) -- (535,261) -- (538,267) -- (538,274) -- (540,278) -- (540,286) -- (541,287) -- (541,290) -- (542,291) -- (543,301) -- (545,303) -- (546,313) -- (548,315) -- (549,328) -- (552,335) -- (552,343) -- (555,346) -- (555,350) -- (556,351) -- (557,359) -- (558,360) -- (558,366) -- (561,373) -- (561,379) -- (563,381) -- (562,384) -- (563,385) -- (563,388) -- (564,389) -- (565,400) -- (568,405) -- (569,418) -- (572,424) -- (572,426) -- (574,429) -- (574,437) -- (576,439) -- (576,445) -- (578,448) -- (578,455) -- (580,458) -- (580,465) -- (582,469) -- (582,473) -- (583,474) -- (584,480) -- (585,481) -- (585,486) -- (586,487) -- (586,492) -- (587,493) -- (588,501) -- (589,502) -- (589,507) -- (591,509) -- (591,518) -- (593,521) -- (593,523) -- (595,527) -- (595,530) -- (597,533) -- (598,546) -- (600,549) -- (600,557) -- (602,560) -- (602,562) -- (603,563) -- (603,567) -- (605,570) -- (605,572) -- (606,573) -- (606,577) -- (608,581) -- (608,584) -- (607,585) -- (605,585) -- (601,581) -- (596,579) -- (595,581) -- (601,586) -- (600,588) -- (593,585) -- (591,583) -- (587,582) -- (585,580) -- (580,579) -- (578,577) -- (574,575) -- (569,574) -- (565,571) -- (562,571) -- (560,569) -- (558,569) -- (556,567) -- (546,562) -- (542,561) -- (537,557) -- (534,557) -- (532,555) -- (530,555) -- (523,551) -- (521,551) -- (517,548) -- (512,547) -- (508,543) -- (504,541) -- (501,541) -- (496,537) -- (494,537) -- (491,535) -- (487,535) -- (474,528) -- (472,528) -- (469,525) -- (464,524) -- (460,521) -- (461,519) -- (461,516) -- (462,515) -- (462,502) -- (463,501) -- (463,494) -- (462,492) -- (464,489) -- (465,471) -- (466,470) -- (466,460) -- (468,457) -- (468,453) -- (467,452) -- (468,450) -- (468,444) -- (469,443) -- (468,440) -- (472,435) -- (473,436) -- (475,436) -- (474,433) -- (472,431) -- (470,427) -- (470,423) -- (471,422) -- (471,412) -- (472,411) -- (472,408) -- (471,406) -- (472,405) -- (472,397) -- (474,394) -- (473,386) -- (475,382) -- (475,373) -- (476,372) -- (476,365) -- (477,364) -- (477,355) -- (478,354) -- (478,348) -- (479,347) -- (479,346) -- (477,344) -- (477,342) -- (478,341) -- (478,333) -- (480,329) -- (480,316) -- (481,315) -- (481,311) -- (482,310) -- (483,296) -- (484,295) -- (484,285) -- (485,284) -- (485,276) -- (486,275) -- (486,267) -- (487,266) -- (487,264) -- (486,263) -- (486,259) -- (487,258) -- (487,254) -- (488,253) -- (487,247) -- (491,240) -- (491,238) -- (488,235) -- (488,233) -- (490,231) -- (491,217) -- (492,216) -- (491,211) -- (493,209) -- (492,207) -- (492,205) -- (493,204) -- (493,202) -- (492,201) -- (492,193) -- (494,190) -- (494,187) -- (495,186) -- (494,184) -- (495,165) -- (496,164) -- (496,159) -- (497,158) -- (497,150) -- (499,146) -- (498,140) -- (499,139) -- (499,131) -- (500,130) -- (500,123) -- (501,122) -- (501,118) -- (500,117) -- (500,115) -- cycle
(574,563) -- (574,564) -- (578,568) -- (579,568) -- (586,574) -- (588,575) -- (590,575) -- (591,574) -- (591,573) -- (580,563) -- (578,562) -- (575,562) -- cycle
(555,548) -- (555,550) -- (565,559) -- (570,559) -- (571,558) -- (571,557) -- (563,549) -- (559,547) -- (556,547) -- cycle
(535,530) -- (535,533) -- (546,543) -- (550,543) -- (551,542) -- (551,539) -- (546,536) -- (538,529) -- (536,529) -- cycle
(514,514) -- (515,517) -- (517,518) -- (524,525) -- (528,527) -- (531,527) -- (532,526) -- (532,524) -- (526,518) -- (524,514) -- (520,512) -- (516,512) -- cycle
(493,497) -- (494,501) -- (499,505) -- (504,508) -- (507,508) -- (508,507) -- (508,505) -- (499,497) -- (497,496) -- (494,496) -- cycle
(469,477) -- (469,479) -- (471,481) -- (471,482) -- (482,491) -- (486,492) -- (488,490) -- (488,489) -- (477,479) -- (471,476) -- cycle
(500,101) -- (503,104) -- (503,105) -- (501,107) -- (500,107) -- (498,110) -- (496,129) -- (494,131) -- (494,135) -- (495,137) -- (494,138) -- (494,143) -- (493,144) -- (493,159) -- (491,161) -- (490,185) -- (489,186) -- (490,194) -- (488,196) -- (489,205) -- (486,211) -- (485,229) -- (482,236) -- (483,238) -- (483,244) -- (482,245) -- (482,255) -- (481,256) -- (481,277) -- (480,278) -- (480,281) -- (479,282) -- (478,303) -- (475,307) -- (476,309) -- (476,312) -- (475,313) -- (476,318) -- (475,319) -- (475,324) -- (473,328) -- (473,339) -- (472,340) -- (472,345) -- (473,346) -- (472,348) -- (472,366) -- (471,367) -- (471,371) -- (470,372) -- (470,375) -- (469,376) -- (469,378) -- (470,379) -- (470,383) -- (468,386) -- (468,390) -- (467,391) -- (467,394) -- (468,395) -- (468,397) -- (466,400) -- (466,405) -- (467,407) -- (465,410) -- (465,415) -- (464,416) -- (464,438) -- (463,439) -- (463,457) -- (462,458) -- (463,462) -- (461,465) -- (459,465) -- (458,464) -- (456,464) -- (451,458) -- (448,458) -- (447,459) -- (447,461) -- (448,462) -- (450,462) -- (454,467) -- (455,467) -- (461,472) -- (461,476) -- (460,477) -- (459,485) -- (458,486) -- (459,491) -- (458,492) -- (458,501) -- (457,502) -- (457,515) -- (456,516) -- (454,516) -- (451,513) -- (451,512) -- (447,508) -- (446,505) -- (442,501) -- (435,490) -- (432,487) -- (431,484) -- (428,481) -- (426,477) -- (417,467) -- (417,466) -- (413,461) -- (412,458) -- (408,454) -- (408,453) -- (404,450) -- (403,447) -- (399,443) -- (396,436) -- (397,435) -- (400,435) -- (401,434) -- (401,432) -- (399,430) -- (396,430) -- (396,431) -- (394,433) -- (385,424) -- (385,423) -- (380,417) -- (375,409) -- (369,409) -- (369,410) -- (378,420) -- (378,422) -- (380,424) -- (380,425) -- (387,432) -- (388,435) -- (393,440) -- (397,447) -- (402,453) -- (403,456) -- (412,466) -- (412,468) -- (414,472) -- (418,476) -- (420,481) -- (428,489) -- (430,495) -- (435,500) -- (436,502) -- (436,505) -- (445,512) -- (445,513) -- (450,518) -- (453,523) -- (451,525) -- (441,526) -- (437,528) -- (433,528) -- (432,527) -- (430,527) -- (425,530) -- (418,530) -- (417,529) -- (415,529) -- (412,532) -- (409,532) -- (404,535) -- (397,535) -- (393,537) -- (386,537) -- (383,539) -- (375,539) -- (369,542) -- (358,543) -- (357,544) -- (353,544) -- (352,545) -- (346,545) -- (345,546) -- (342,546) -- (341,547) -- (338,547) -- (337,548) -- (331,548) -- (330,549) -- (327,549) -- (325,551) -- (319,551) -- (318,552) -- (315,552) -- (312,554) -- (303,554) -- (300,556) -- (296,556) -- (292,558) -- (287,558) -- (283,560) -- (278,560) -- (275,562) -- (268,562) -- (264,564) -- (258,564) -- (255,566) -- (248,566) -- (247,567) -- (245,567) -- (243,569) -- (233,570) -- (232,571) -- (228,571) -- (227,572) -- (218,573) -- (214,575) -- (208,575) -- (207,576) -- (203,576) -- (202,577) -- (198,577) -- (197,578) -- (193,578) -- (192,579) -- (187,579) -- (184,581) -- (177,581) -- (176,582) -- (173,582) -- (172,583) -- (169,583) -- (168,584) -- (163,584) -- (160,586) -- (157,586) -- (156,587) -- (151,587) -- (145,590) -- (139,590) -- (136,592) -- (133,592) -- (132,591) -- (132,589) -- (137,585) -- (137,582) -- (136,582) -- (134,580) -- (140,574) -- (142,573) -- (142,571) -- (145,568) -- (148,563) -- (159,551) -- (159,550) -- (161,548) -- (161,546) -- (166,541) -- (168,536) -- (175,530) -- (175,529) -- (179,525) -- (179,524) -- (184,518) -- (184,516) -- (189,512) -- (189,510) -- (191,508) -- (195,501) -- (200,496) -- (201,493) -- (210,484) -- (216,474) -- (222,468) -- (222,466) -- (226,462) -- (229,456) -- (235,450) -- (235,449) -- (237,447) -- (237,445) -- (241,441) -- (242,438) -- (253,428) -- (251,426) -- (251,425) -- (257,418) -- (257,417) -- (262,412) -- (264,407) -- (268,403) -- (269,400) -- (275,394) -- (277,390) -- (280,387) -- (280,386) -- (283,383) -- (286,378) -- (299,363) -- (305,353) -- (310,347) -- (311,344) -- (321,333) -- (321,332) -- (323,330) -- (323,328) -- (328,323) -- (330,318) -- (338,310) -- (341,305) -- (347,299) -- (347,298) -- (349,296) -- (349,294) -- (352,291) -- (352,289) -- (360,281) -- (363,276) -- (367,272) -- (369,268) -- (372,265) -- (373,265) -- (375,259) -- (379,255) -- (381,255) -- (382,254) -- (382,250) -- (384,248) -- (384,247) -- (389,242) -- (390,239) -- (396,233) -- (396,231) -- (399,228) -- (399,227) -- (405,221) -- (406,217) -- (411,212) -- (411,210) -- (413,208) -- (414,208) -- (419,202) -- (420,198) -- (423,195) -- (424,195) -- (426,191) -- (432,186) -- (432,184) -- (434,182) -- (434,180) -- (438,176) -- (440,171) -- (443,168) -- (443,167) -- (449,162) -- (450,159) -- (453,156) -- (454,153) -- (469,135) -- (472,129) -- (475,126) -- (477,122) -- (481,118) -- (486,110) -- (490,106) -- (494,107) -- (495,106) -- (495,103) -- (497,101) -- cycle
(163,560) -- (161,560) -- (155,566) -- (154,566) -- (149,571) -- (148,571) -- (145,574) -- (145,575) -- (146,576) -- (149,576) -- (161,565) -- (161,564) -- (163,562) -- cycle
(186,540) -- (184,539) -- (182,541) -- (181,541) -- (174,548) -- (171,549) -- (170,550) -- (170,552) -- (171,553) -- (174,553) -- (176,552) -- (179,548) -- (183,547) -- (186,544) -- cycle
(210,520) -- (209,519) -- (206,519) -- (202,521) -- (200,525) -- (194,530) -- (193,533) -- (194,534) -- (198,534) -- (209,524) -- cycle
(234,495) -- (231,496) -- (223,504) -- (218,507) -- (217,509) -- (217,512) -- (218,513) -- (221,513) -- (223,512) -- (226,508) -- (227,508) -- (234,502) -- (234,500) -- (235,499) -- (235,496) -- cycle
(261,475) -- (259,473) -- (256,473) -- (248,480) -- (247,480) -- (244,483) -- (243,486) -- (244,487) -- (248,487) -- (252,482) -- (258,479) -- (261,476) -- cycle
(286,450) -- (285,449) -- (281,450) -- (271,460) -- (270,460) -- (266,464) -- (266,465) -- (268,467) -- (271,467) -- (273,466) -- (286,454) -- cycle
(428,443) -- (428,446) -- (437,454) -- (441,455) -- (443,453) -- (443,452) -- (434,444) -- (433,444) -- (431,442) -- (429,442) -- cycle
(312,430) -- (309,427) -- (307,427) -- (300,434) -- (299,434) -- (292,440) -- (291,440) -- (289,438) -- (289,437) -- (287,437) -- (288,441) -- (293,446) -- (295,446) -- (308,433) -- (309,433) -- cycle
(402,424) -- (402,426) -- (403,427) -- (405,427) -- (418,438) -- (420,438) -- (423,436) -- (423,435) -- (419,431) -- (413,428) -- (409,424) -- (407,423) -- (403,423) -- cycle
(332,405) -- (329,405) -- (321,414) -- (319,414) -- (316,417) -- (315,421) -- (317,423) -- (322,423) -- (324,421) -- (324,419) -- (326,417) -- (326,416) -- (334,410) -- (334,407) -- cycle
(380,405) -- (380,406) -- (394,419) -- (396,420) -- (398,420) -- (399,419) -- (399,417) -- (396,414) -- (393,413) -- (387,407) -- (382,404) -- cycle
(338,398) -- (338,399) -- (340,401) -- (344,401) -- (351,394) -- (356,391) -- (363,399) -- (363,400) -- (365,402) -- (365,404) -- (365,399) -- (368,397) -- (369,397) -- (375,402) -- (377,401) -- (377,399) -- (376,399) -- (372,395) -- (371,395) -- (367,391) -- (366,391) -- (365,389) -- (364,389) -- (357,383) -- (349,387) -- cycle
(365,367) -- (363,369) -- (362,373) -- (358,377) -- (358,378) -- (360,380) -- (362,380) -- (364,378) -- (364,376) -- (365,375) -- cycle
(377,338) -- (376,339) -- (375,343) -- (380,343) -- (381,342) -- (381,339) -- (380,338) -- cycle
(395,310) -- (392,310) -- (383,324) -- (383,329) -- (385,331) -- (387,331) -- (389,329) -- (390,323) -- (392,321) -- (396,319) -- (395,317) -- (396,311) -- cycle
(404,286) -- (402,288) -- (402,290) -- (399,294) -- (399,296) -- (397,298) -- (397,302) -- (398,303) -- (399,303) -- (403,298) -- (403,296) -- (404,295) -- cycle
(423,253) -- (420,252) -- (417,255) -- (417,257) -- (415,259) -- (415,261) -- (414,262) -- (414,265) -- (411,268) -- (411,270) -- (417,270) -- (419,267) -- (419,263) -- (422,260) -- (424,256) -- (424,254) -- cycle
(437,222) -- (435,222) -- (431,225) -- (431,228) -- (428,233) -- (428,235) -- (426,239) -- (423,242) -- (423,244) -- (425,246) -- (428,246) -- (430,242) -- (434,238) -- (434,234) -- (437,230) -- (437,228) -- (439,225) -- (439,224) -- cycle
(451,194) -- (449,194) -- (447,196) -- (446,200) -- (444,202) -- (440,210) -- (440,218) -- (442,220) -- (443,220) -- (450,214) -- (449,213) -- (447,213) -- (446,212) -- (446,210) -- (449,204) -- (453,199) -- (453,196) -- cycle
(462,169) -- (461,169) -- (459,171) -- (459,174) -- (455,179) -- (454,183) -- (452,185) -- (452,187) -- (454,189) -- (456,189) -- (459,184) -- (460,179) -- (464,173) -- (464,171) -- cycle
(476,142) -- (475,142) -- (473,144) -- (473,145) -- (471,147) -- (471,149) -- (469,153) -- (466,156) -- (466,158) -- (464,161) -- (465,163) -- (468,163) -- (470,161) -- (471,159) -- (471,156) -- (473,154) -- (474,151) -- (478,147) -- (477,143) -- cycle
(490,115) -- (488,115) -- (487,116) -- (486,120) -- (482,126) -- (482,128) -- (479,131) -- (478,133) -- (478,135) -- (479,136) -- (481,136) -- (482,135) -- (484,130) -- (486,128) -- (486,126) -- (489,122) -- (489,120) -- (490,119) -- cycle
(505,93) -- (515,97) -- (517,99) -- (523,102) -- (527,106) -- (531,108) -- (533,108) -- (535,110) -- (538,111) -- (541,114) -- (545,116) -- (549,120) -- (551,120) -- (562,130) -- (568,133) -- (572,137) -- (573,137) -- (575,139) -- (575,140) -- (577,141) -- (583,147) -- (583,148) -- (585,149) -- (594,158) -- (594,159) -- (601,166) -- (601,167) -- (607,172) -- (607,173) -- (623,190) -- (623,192) -- (626,195) -- (628,199) -- (630,200) -- (635,206) -- (635,208) -- (637,211) -- (637,214) -- (642,220) -- (642,223) -- (650,235) -- (650,237) -- (653,243) -- (655,245) -- (655,247) -- (657,251) -- (659,253) -- (659,257) -- (661,261) -- (665,266) -- (665,268) -- (668,275) -- (668,278) -- (669,279) -- (670,284) -- (672,287) -- (672,290) -- (673,291) -- (673,294) -- (674,295) -- (674,299) -- (676,301) -- (676,306) -- (678,309) -- (679,320) -- (681,323) -- (680,325) -- (681,327) -- (681,333) -- (682,334) -- (681,337) -- (683,340) -- (683,343) -- (684,344) -- (684,352) -- (685,353) -- (685,356) -- (686,357) -- (686,375) -- (687,376) -- (686,378) -- (687,380) -- (686,382) -- (687,384) -- (687,390) -- (686,391) -- (686,393) -- (687,394) -- (687,396) -- (686,397) -- (686,404) -- (687,405) -- (687,408) -- (684,412) -- (684,417) -- (685,419) -- (683,423) -- (683,425) -- (684,426) -- (684,433) -- (683,434) -- (683,438) -- (682,439) -- (682,448) -- (681,449) -- (681,452) -- (680,453) -- (680,456) -- (679,457) -- (679,460) -- (678,461) -- (678,465) -- (676,467) -- (676,472) -- (674,474) -- (674,479) -- (671,486) -- (671,489) -- (668,493) -- (668,497) -- (664,502) -- (664,505) -- (663,506) -- (663,508) -- (661,510) -- (661,513) -- (660,514) -- (660,519) -- (657,526) -- (655,528) -- (655,531) -- (653,535) -- (651,537) -- (650,541) -- (646,546) -- (644,550) -- (644,552) -- (641,555) -- (640,558) -- (635,564) -- (635,565) -- (631,569) -- (635,569) -- (636,570) -- (634,572) -- (634,573) -- (620,587) -- (619,585) -- (616,589) -- (614,589) -- (612,587) -- (612,582) -- (611,581) -- (611,577) -- (610,576) -- (610,569) -- (609,568) -- (608,560) -- (606,556) -- (606,549) -- (607,548) -- (607,546) -- (604,543) -- (604,536) -- (603,535) -- (603,532) -- (605,530) -- (605,529) -- (601,525) -- (601,524) -- (603,522) -- (601,518) -- (601,512) -- (598,508) -- (596,497) -- (595,496) -- (595,494) -- (592,490) -- (592,483) -- (593,481) -- (590,478) -- (590,473) -- (587,467) -- (587,463) -- (588,461) -- (585,458) -- (585,452) -- (584,451) -- (584,442) -- (582,440) -- (582,437) -- (579,433) -- (580,431) -- (580,428) -- (577,423) -- (577,414) -- (578,413) -- (578,411) -- (572,404) -- (572,394) -- (570,391) -- (570,387) -- (571,386) -- (568,383) -- (568,375) -- (566,372) -- (566,370) -- (564,366) -- (564,362) -- (563,361) -- (563,358) -- (562,357) -- (562,354) -- (561,353) -- (561,350) -- (560,349) -- (560,345) -- (558,343) -- (558,340) -- (557,339) -- (557,334) -- (555,332) -- (555,328) -- (556,327) -- (556,325) -- (555,324) -- (555,321) -- (553,317) -- (552,308) -- (550,305) -- (550,302) -- (549,301) -- (550,295) -- (548,291) -- (548,285) -- (545,278) -- (544,268) -- (541,262) -- (540,256) -- (538,254) -- (538,253) -- (541,250) -- (538,247) -- (538,245) -- (537,244) -- (537,239) -- (535,236) -- (535,232) -- (533,229) -- (533,220) -- (531,217) -- (531,214) -- (530,213) -- (530,208) -- (529,207) -- (529,203) -- (528,202) -- (528,198) -- (527,197) -- (527,194) -- (526,193) -- (525,185) -- (523,183) -- (523,179) -- (521,175) -- (522,171) -- (520,169) -- (520,167) -- (519,166) -- (519,160) -- (518,159) -- (518,155) -- (516,151) -- (516,144) -- (515,143) -- (515,139) -- (514,138) -- (514,135) -- (512,132) -- (512,126) -- (510,123) -- (510,119) -- (509,118) -- (508,109) -- (506,106) -- (506,102) -- (503,99) -- (503,95) -- cycle
(493,87) -- (495,89) -- (493,94) -- (489,98) -- (484,108) -- (472,121) -- (472,122) -- (469,125) -- (468,128) -- (459,138) -- (459,139) -- (455,143) -- (454,146) -- (442,160) -- (440,166) -- (426,181) -- (426,183) -- (424,185) -- (423,188) -- (409,204) -- (409,206) -- (403,212) -- (403,214) -- (395,222) -- (392,228) -- (386,233) -- (385,236) -- (382,239) -- (381,242) -- (369,256) -- (369,258) -- (365,262) -- (364,265) -- (361,268) -- (361,269) -- (357,272) -- (357,275) -- (348,284) -- (347,287) -- (339,296) -- (339,298) -- (330,308) -- (330,309) -- (328,311) -- (328,313) -- (324,317) -- (323,320) -- (319,324) -- (318,327) -- (313,332) -- (310,337) -- (303,345) -- (300,352) -- (291,362) -- (288,367) -- (282,373) -- (282,375) -- (279,378) -- (279,379) -- (272,386) -- (270,386) -- (269,387) -- (269,389) -- (266,393) -- (266,396) -- (259,404) -- (255,411) -- (247,420) -- (247,421) -- (243,426) -- (242,429) -- (236,437) -- (235,440) -- (232,443) -- (232,445) -- (225,452) -- (224,455) -- (217,462) -- (216,465) -- (209,472) -- (209,474) -- (201,483) -- (194,494) -- (181,509) -- (181,513) -- (170,525) -- (169,528) -- (163,534) -- (162,538) -- (158,542) -- (156,546) -- (145,559) -- (143,563) -- (139,567) -- (136,572) -- (131,577) -- (131,580) -- (126,585) -- (126,587) -- (118,596) -- (114,594) -- (111,591) -- (110,588) -- (106,584) -- (105,580) -- (98,573) -- (97,570) -- (90,562) -- (90,559) -- (86,555) -- (83,549) -- (80,546) -- (80,543) -- (76,538) -- (75,532) -- (73,530) -- (73,528) -- (71,526) -- (69,520) -- (65,516) -- (65,513) -- (62,510) -- (62,507) -- (60,505) -- (59,500) -- (56,495) -- (56,491) -- (52,484) -- (52,479) -- (50,475) -- (51,471) -- (47,466) -- (47,455) -- (46,454) -- (45,450) -- (43,448) -- (43,444) -- (41,442) -- (41,434) -- (38,429) -- (38,424) -- (37,423) -- (37,421) -- (40,418) -- (40,417) -- (39,416) -- (39,413) -- (37,411) -- (37,406) -- (38,405) -- (38,399) -- (37,398) -- (37,395) -- (38,393) -- (36,391) -- (36,388) -- (35,387) -- (35,385) -- (37,383) -- (35,376) -- (38,373) -- (38,370) -- (36,368) -- (36,365) -- (37,364) -- (36,361) -- (37,360) -- (37,357) -- (38,356) -- (38,353) -- (37,352) -- (38,350) -- (38,348) -- (37,347) -- (37,340) -- (38,339) -- (38,336) -- (39,335) -- (39,331) -- (41,329) -- (40,327) -- (40,324) -- (41,323) -- (42,316) -- (44,312) -- (44,307) -- (45,306) -- (45,303) -- (47,300) -- (48,293) -- (50,291) -- (50,285) -- (52,282) -- (52,279) -- (55,276) -- (54,273) -- (57,267) -- (57,262) -- (60,259) -- (60,257) -- (62,255) -- (62,251) -- (64,249) -- (64,245) -- (66,243) -- (66,239) -- (70,236) -- (74,226) -- (76,224) -- (83,211) -- (87,207) -- (87,206) -- (90,203) -- (90,201) -- (94,197) -- (95,193) -- (103,184) -- (103,183) -- (106,180) -- (109,175) -- (117,167) -- (119,166) -- (121,162) -- (124,159) -- (125,159) -- (129,154) -- (130,154) -- (133,151) -- (135,151) -- (136,150) -- (136,148) -- (139,145) -- (147,144) -- (149,142) -- (149,138) -- (152,135) -- (155,135) -- (159,131) -- (160,131) -- (165,125) -- (165,124) -- (173,117) -- (179,115) -- (183,110) -- (185,110) -- (191,107) -- (193,105) -- (195,105) -- (198,102) -- (202,101) -- (207,97) -- (211,97) -- (213,95) -- (216,94) -- (220,91) -- (224,91) -- (226,90) -- (229,87) -- (233,87) -- (241,83) -- (244,84) -- (251,78) -- (253,78) -- (256,76) -- (261,76) -- (264,73) -- (268,71) -- (271,71) -- (272,70) -- (278,70) -- (279,69) -- (285,68) -- (286,67) -- (292,67) -- (295,65) -- (298,66) -- (302,63) -- (305,63) -- (306,64) -- (309,64) -- (313,62) -- (321,63) -- (323,61) -- (324,61) -- (326,63) -- (331,63) -- (335,60) -- (339,60) -- (340,59) -- (342,59) -- (343,60) -- (361,60) -- (362,59) -- (371,59) -- (372,60) -- (377,60) -- (379,62) -- (382,62) -- (384,61) -- (385,62) -- (393,62) -- (394,63) -- (410,63) -- (411,64) -- (413,64) -- (417,67) -- (422,67) -- (424,66) -- (425,67) -- (428,67) -- (429,68) -- (432,68) -- (433,69) -- (440,69) -- (442,71) -- (445,71) -- (450,74) -- (461,76) -- (464,78) -- (467,78) -- (473,81) -- (479,82) -- (484,85) -- (490,85) -- (492,87) -- cycle;
\node[inner sep=0pt,font=\fontsize{10}{10}\selectfont] at (547.5,51.5) {\scalebox{6.0}{$a',b'$}};
\node[inner sep=0pt,font=\fontsize{10}{10}\selectfont] at (264.5,319.5) {\scalebox{6.0}{$h''$}};
\node[inner sep=0pt,font=\fontsize{10}{10}\selectfont] at (425.0,297.5) {\scalebox{5.9}{$\varrho$}};
\node[inner sep=0pt,font=\fontsize{10}{10}\selectfont] at (415,367) {\scalebox{4.6}{$\xi,\eta$}};
\node[inner sep=0pt,font=\fontsize{10}{10}\selectfont] at (597.0,327.0) {\scalebox{6.2}{$h$}};
\node[inner sep=0pt,font=\fontsize{10}{10}\selectfont] at (373.0,462.0) {\scalebox{5.6}{$R$}};
\node[inner sep=0pt,font=\fontsize{10}{10}\selectfont] at (262.0,503.0) {\scalebox{5.8}{$\varrho$}};
\node[inner sep=0pt,font=\fontsize{10}{10}\selectfont] at (502.0,469.0) {\scalebox{5.8}{$\varrho$}};
\node[inner sep=0pt,font=\fontsize{10}{10}\selectfont] at (458,566) {\scalebox{4.6}{$x,y$}};
\node[inner sep=0pt,font=\fontsize{10}{10}\selectfont] at (67,612) {\scalebox{5.0}{$a,b$}};
\node[inner sep=0pt,font=\fontsize{10}{10}\selectfont] at (701,606) {\scalebox{5.0}{$a'',b''$}};
\node[inner sep=0pt,font=\fontsize{10}{10}\selectfont] at (367.0,643.5) {\scalebox{5.8}{$h'$}};
\end{tikzpicture}%
}
\end{center}

Man nehme ferner an

\[
x = \xi + R\cos\psi, \qquad y = \eta + R\sin\psi;
\]

so findet man:

\[
\begin{aligned}
&\dfrac{RR+\rho\rho-2R\rho\sin(\psi+A-A'-A'')}{mm}\\ &= \dfrac{RR+\rho\rho-2R\rho\sin(\psi+A'-A''-A)}{m'm'}\\ &= \dfrac{RR+\rho\rho-2R\rho\sin(\psi+A''-A-A')}{m''m''} = \Delta.
\end{aligned}
\]

Hieraus wird

\[
\begin{aligned}&\mathrm{I.}\quad RR+\rho\rho\\ &=\dfrac{mm\sin(2A''-2A')+m'm'\sin(2A-2A'')+m''m''\sin(2A'-2A)}{4\sin(A''-A')\sin(A''-A)\sin(A'-A)}\Delta,\end{aligned}
\]

\[
\begin{aligned}\mathrm{II.}\quad 2R\rho\cos\psi&=\\ &\dfrac{mmh\sin A+m'm'h'\sin A'+m''m''h''\sin A''}{\substack{\sin(A'+A''-A)h\sin A+\sin(A''+A-A')h'\sin A'\\+\sin(A+A'-A'')h''\sin A''}}\Delta\end{aligned}
\]

[und

\[
\begin{aligned}\mathrm{III.}\quad -2R\rho\sin\psi&=\\ &\dfrac{mmh\cos A+m'm'h'\cos A'+m''m''h''\cos A''}{\substack{\cos(A'+A''-A)h\cos A+\cos(A''+A-A')h'\cos A'\\+\cos(A+A'-A'')h''\cos A''}}\Delta.\end{aligned}
\]
%
\SourceLeaf{313}
Da nun die Nenner in II. und III. beide denselben Werth, nemlich

\[
\tfrac{1}{2}\rho\left\{\sin(2A-2A')+\sin(2A'-2A'')+\sin(2A''-2A)\right\},
\]

haben, so ergibt sich aus diesen Gleichungen durch Division:

\[
\operatorname{tg}\psi = -\dfrac{mmh\cos A+m'm'h'\cos A'+m''m''h''\cos A''}{mmh\sin A+m'm'h'\sin A'+m''m''h''\sin A''},
\]

wodurch für $\psi$ und damit auch für $R$ zwei Werthe folgen.]

\bigskip
\begin{center}
\hrulefill
\end{center}
\bigskip

\Needspace{9\baselineskip}
\begin{center}
BEMERKUNG.
\end{center}

Die vorstehende Notiz, die sich in einem Handbuche findet und vermuthlich aus dem Jahre 1816 stammt, ist sehr flüchtig geschrieben, sodass zahlreiche Schreibfehler in den Formeln verbessert werden mussten. Sie bricht bei der Formel II der zweiten Auflösungsart ab, bei der, was im Druck nicht ersichtlich gemacht werden konnte, der Nenner und der Factor $\Delta$ ergänzt werden mussten. Der Deutlichkeit wegen sind dem Texte zwei Figuren hinzugefügt worden.

Es ist anzunehmen, dass \textsc{Gauss} seine Auflösungen gefunden hat, indem er von der Darstellung der Ecken des Dreiecks durch complexe Grössen ausging, denn hierbei ergeben sich genau die Formeln des Textes, wenn man darin die Paare von Gleichungen, die den Cosinus und den Sinus desselben Winkels enthalten, mittelst des Symboles $i$ immer in eine einzige Gleichung zusammenfasst. Dass \textsc{Gauss} die Darstellung der Punkte einer Ebene durch complexe Grössen schon früh bei geometrischen Aufgaben benutzt hat, zeigt zum Beispiel die auf der folgenden Seite abgedruckte Notiz über die Construction der \textsc{Pothenot}schen Aufgabe, sowie eine handschriftliche Bemerkung zu \textsc{Schumachers} Übersetzung der \emph{Geometrie der Stellung} von \textsc{Carnot} aus dem Jahre 1810, die sich ebenfalls auf die Geometrie des Dreiecks bezieht und bereits Bd.\ IV.\ S.\ 396 bis 397 Platz gefunden hat.

\hfill \textsc{Stäckel}.

\bigskip
\begin{center}
\hrulefill
\end{center}
%
\clearpage
\clearpage\SourceLeaf{314}

\begin{center}
[POTHENOTS AUFGABE UND DAS VIERECK.]
\end{center}

\begin{center}
[I.]

[CONSTRUCTION DER POTHENOTSCHEN AUFGABE.]

[1.]
\end{center}

Es werden drei Punkte $a$, $b$, $c$ durch die eben so bezeichneten complexen
Zahlen dargestellt, es seien ferner $\alpha$, $\beta$, $\gamma$ drei complexe
Zahlen, so dass
\[
\alpha+\beta+\gamma=0.
\]
Man mache
\[
\alpha A+\beta b+\gamma c=0
\]
\[
\alpha a+\beta B+\gamma c=0
\]
\[
\alpha a+\beta b+\gamma C=0,
\]
und $A$, $B$, $C$ seien die Punkte, denen die complexen Zahlen
$A$, $B$, $C$ entsprechen. Dann schneiden sich $aA$,
$bB$, $cC$ in Einem Punkte $P$ und dieser ist die Auflösung
von \textsc{Pothenots} Aufgabe.

Sind $\mathfrak{A}$, $\mathfrak{B}$, $\mathfrak{C}$ die in $P$ beobachteten
(nicht orientirten) Azimuthe, so verhalten sich
\[
\alpha e^{i\mathfrak{A}}, \qquad \beta e^{i\mathfrak{B}}, \qquad \gamma e^{i\mathfrak{C}}.
\]
wie reelle Zahlen.

\Needspace{8\baselineskip}
\begin{center}
BEMERKUNG.
\end{center}

Diese Notiz hat \textsc{Gauss} auf der innern Seite des hintern Deckels
seines Exemplars des Werkes: \textsc{Carnot}, \emph{Geometrie der
Stellung}. Übersetzt von \textsc{Schumacher}. Bd.\ II.\ Altona 1810
eingetragen.

\par\hfill\textsc{Stäckel}.\par
%
\Needspace{12\baselineskip}\SourceLeaf{315}

\begin{center}
[2.]
\end{center}

\begin{center}
\textsc{Gauss} an \textsc{Gerling}. Göttingen, 7. November 1830.
\end{center}

\indent\mbox{. . . . . . . . . .} Über die vier merkwürdigen Punkte im Dreieck
habe ich selbst nie etwas bekannt gemacht. Was ich jedoch 1808 darüber
\textsc{Schumacher} mitgetheilt habe, hat dieser im Anhange zu seiner
Übersetzung von \textsc{Carnot}, Géométrie de position, zweiter
Bd., welcher 1810 erschienen ist, drucken lassen. Ich zweifle fast, dass
der Punkt, welcher sich nach dem von Ihnen angeführten Satze ergibt,
besonders elegante Relationen zu den übrigen darbieten wird, und zwar
deshalb, weil jenes Theorem nur ein ganz specieller Fall eines viel
allgemeinern ist.

Wenn man nemlich an die Seiten eines gegebenen Dreiecks $ABC$
drei andere $AB\gamma$, $\alpha BC$, $AC\beta$ legt (entweder alle drei
an die äussere oder alle drei an die innere Seite des Dreiecks
$ABC$ gelegt), die jedes einem zweiten gegebenen Dreieck $abc$ ähnlich
sind, und zwar dergestalt, dass sie auch ähnlich liegend und die Winkel an
$\alpha$, $\beta$, $\gamma$ denen an $a$, $b$, $c$ gleich werden, so
schneiden sich die Linien $A\alpha$, $B\beta$, $C\gamma$
(nöthigenfalls vorwärts oder rückwärts fortgesetzt) in Einem Punkte
$D$ und die Winkel an $D$ ($ADB$, $BDC$, $CDA$) sind
denen an $\gamma$, $\alpha$, $\beta$ entweder selbst oder ihren
Nebenwinkeln gleich und ähnlich liegend.

Offenbar ist diess die zierlichste Construction der
\textsc{Pothenots}chen Aufgabe; wo man aber in der Praxis sich begnügen
kann, nur Einen Punkt z.\,B. $\alpha$ zu construiren und den Messtisch nach
der Linie $A\alpha$ zu orientiren. Wie ich von meinem Sohne höre, ist
letztere Manier unter dem Namen der \textsc{Schulz-Montanus}'schen
bekannt. \mbox{. . . . . . .}
%
\clearpage\SourceLeaf{316}

\begin{center}
[II.]

[DIE PHYSISCHE MÖGLICHKEIT DER DATEN IN POTHENOTS AUFGABE.]

[1.]
\end{center}

\begin{center}
\textsc{Gauss} an \textsc{Schumacher}. Göttingen, 13. April 1836.
\end{center}

\indent \mbox{. . . . . . .} Bei dieser Gelegenheit will ich doch auch einer
andern Aufgabe erwähnen, nemlich der sogenannten \textsc{Pothenots}chen.
Wenn drei Punkte $A$, $B$, $C$ gegeben sind, und zugleich die
Winkel, die $XA$, $XB$, $XC$ unter einander machen
(natürlich inclusive des \emph{sens}, d.\,i. ob die Winkel von der linken
nach der rechten [Seite] oder umgekehrt gezählt sind), so ist
$X$ (allgemein zu reden, d.\,i. mit Ausnahme des Falls, wo $A$,
$B$, $C$, $X$ auf Einem Kreise liegen) \textls[120]{bestimmt}. Gründen sich die
Data auf einen wirklichen concreten Fall, so versteht sich von selbst,
dass die Aufgabe, weil wirklich, auch möglich ist. Aber wenn jemand (um
Sie aufs Eis zu führen) Ihnen bloss die Zahlen aufgibt, so kann sich
hinterher ergeben, dass die Aufgabe physisch unmöglich ist, und dass Sie
statt einer der drei Richtungen die entgegengesetzte erhalten
($180^0$ verschieden von der aufgegebenen).

Die Frage ist nun, wie man \textls[120]{am einfachsten} die Bedingung der physischen
Möglichkeit durch eine Gleichung \textls[120]{zwischen den Datis} ausdrücken kann? Ich
habe, mich verwundernd, dass bei einer Aufgabe, worüber so viel
geschrieben ist, doch von allen, so viel ich weiss, ein so wesentlicher
Umstand ganz übersehen ist, — die Auflösung oft in meinen Vorlesungen
vorgetragen. Aber so ganz leicht muss sie doch wohl nicht zu finden sein;
wenigstens habe ich einmal einem sonst sehr geschickten jungen Mann (dem
sie von mir \textls[120]{nicht} vorgetragen war) sie als Thema einer Doctordissertation
vorgeschlagen, und der hatte, obgleich er sich lange damit gequält, die
Auflösung nicht finden können. \mbox{. . . . . . .}
%
\Needspace{12\baselineskip}\SourceLeaf{317}

\begin{center}
\textsc{Gauss} an \textsc{Schumacher}. Göttingen, 21. April 1836.
\end{center}

\indent \mbox{. . . . . .} Um ein Beispiel zu haben, wo \textsc{Pothenots}
Aufgabe physisch (i.\,e. optisch) unmöglich ist, nehmen Sie an, die drei
Richtungen sollen, dem Verlangen nach, Winkel von $120^0$ unter einander
machen, was allemal unmöglich wird, sobald die drei Punkte ein Dreieck
bilden, worin ein stumpfer Winkel grösser als $120^0$ ist.

\begin{center}\resizebox{7.0cm}{!}{%
\begin{tikzpicture}[x=1pt,y=-1pt]
\path[use as bounding box] (0,0) rectangle (1615,720);
\path[fill=black,draw=none,even odd rule]
(1545,658) -- (1545,657) -- (1543,655) -- (1538,653) -- (1532,648) -- (1530,648) -- (1524,642) -- (1520,642) -- (1512,635) -- (1509,635) -- (1507,633) -- (1500,630) -- (1495,624) -- (1490,623) -- (1482,619) -- (1476,611) -- (1474,612) -- (1471,612) -- (1467,610) -- (1464,607) -- (1460,606) -- (1458,604) -- (1448,600) -- (1445,597) -- (1440,595) -- (1438,593) -- (1436,593) -- (1434,591) -- (1432,591) -- (1428,589) -- (1422,584) -- (1420,584) -- (1413,579) -- (1410,578) -- (1407,575) -- (1406,572) -- (1401,572) -- (1400,571) -- (1398,571) -- (1395,568) -- (1393,568) -- (1389,565) -- (1385,564) -- (1382,561) -- (1379,560) -- (1374,554) -- (1372,555) -- (1367,552) -- (1365,552) -- (1360,548) -- (1355,546) -- (1350,541) -- (1346,540) -- (1343,537) -- (1341,537) -- (1336,533) -- (1329,530) -- (1326,526) -- (1322,526) -- (1320,525) -- (1314,520) -- (1312,520) -- (1310,518) -- (1308,518) -- (1305,514) -- (1303,514) -- (1293,509) -- (1290,506) -- (1287,505) -- (1284,502) -- (1282,502) -- (1278,500) -- (1275,497) -- (1272,497) -- (1268,493) -- (1254,486) -- (1247,480) -- (1245,480) -- (1243,478) -- (1239,478) -- (1230,471) -- (1226,470) -- (1221,466) -- (1219,466) -- (1217,464) -- (1214,463) -- (1206,457) -- (1202,456) -- (1198,452) -- (1196,452) -- (1191,448) -- (1189,448) -- (1185,446) -- (1180,442) -- (1173,439) -- (1169,435) -- (1167,435) -- (1163,433) -- (1160,430) -- (1155,428) -- (1152,425) -- (1150,425) -- (1148,423) -- (1144,421) -- (1141,421) -- (1138,419) -- (1138,416) -- (1134,416) -- (1132,415) -- (1130,413) -- (1127,412) -- (1125,410) -- (1119,407) -- (1116,404) -- (1114,404) -- (1112,402) -- (1110,402) -- (1107,399) -- (1105,399) -- (1103,397) -- (1094,393) -- (1089,389) -- (1075,382) -- (1067,376) -- (1065,376) -- (1061,373) -- (1058,372) -- (1055,369) -- (1053,369) -- (1050,366) -- (1049,366) -- (1046,363) -- (1044,363) -- (1041,360) -- (1039,359) -- (1036,359) -- (1034,357) -- (1029,355) -- (1025,351) -- (1018,349) -- (1013,345) -- (1009,343) -- (1007,343) -- (1004,340) -- (995,336) -- (993,334) -- (991,334) -- (987,330) -- (982,328) -- (980,326) -- (976,325) -- (966,317) -- (963,317) -- (960,314) -- (957,314) -- (955,313) -- (952,310) -- (946,308) -- (943,305) -- (939,304) -- (937,302) -- (935,302) -- (932,299) -- (922,294) -- (917,290) -- (915,290) -- (913,288) -- (910,287) -- (903,282) -- (896,279) -- (891,275) -- (889,275) -- (887,273) -- (885,273) -- (883,271) -- (878,269) -- (875,266) -- (871,265) -- (867,262) -- (864,262) -- (860,258) -- (856,258) -- (854,256) -- (851,255) -- (849,253) -- (845,251) -- (837,255) -- (830,257) -- (827,260) -- (825,260) -- (821,263) -- (817,263) -- (815,264) -- (812,267) -- (811,267) -- (803,273) -- (800,273) -- (797,275) -- (795,275) -- (792,278) -- (783,282) -- (781,284) -- (778,285) -- (775,288) -- (771,288) -- (766,291) -- (763,291) -- (761,293) -- (759,293) -- (754,298) -- (751,298) -- (747,302) -- (740,304) -- (737,307) -- (731,310) -- (729,310) -- (726,313) -- (717,317) -- (710,322) -- (707,322) -- (701,325) -- (698,328) -- (695,328) -- (688,333) -- (682,334) -- (680,336) -- (676,337) -- (674,339) -- (669,341) -- (664,346) -- (662,346) -- (656,350) -- (653,350) -- (650,354) -- (645,354) -- (640,359) -- (636,360) -- (634,362) -- (632,362) -- (630,364) -- (626,366) -- (623,366) -- (621,368) -- (619,368) -- (616,371) -- (610,374) -- (608,374) -- (605,377) -- (603,377) -- (599,379) -- (595,383) -- (589,384) -- (585,387) -- (583,387) -- (581,389) -- (578,390) -- (571,395) -- (562,399) -- (560,401) -- (556,403) -- (554,403) -- (546,408) -- (544,408) -- (539,412) -- (535,413) -- (533,415) -- (531,415) -- (529,417) -- (523,420) -- (520,423) -- (513,425) -- (509,428) -- (507,428) -- (505,430) -- (502,430) -- (497,434) -- (494,434) -- (492,436) -- (491,436) -- (480,444) -- (478,444) -- (475,447) -- (472,448) -- (470,450) -- (461,454) -- (459,456) -- (455,457) -- (445,464) -- (443,464) -- (437,468) -- (430,470) -- (427,474) -- (423,475) -- (422,476) -- (418,476) -- (413,481) -- (410,481) -- (407,484) -- (404,484) -- (401,487) -- (399,487) -- (396,489) -- (393,489) -- (387,494) -- (382,496) -- (380,498) -- (374,501) -- (372,501) -- (366,506) -- (363,506) -- (361,508) -- (354,510) -- (346,516) -- (342,517) -- (340,519) -- (338,519) -- (336,521) -- (331,523) -- (326,527) -- (318,529) -- (313,533) -- (309,534) -- (306,537) -- (304,536) -- (301,537) -- (296,542) -- (295,542) -- (289,547) -- (287,547) -- (282,551) -- (277,552) -- (274,556) -- (270,556) -- (264,561) -- (260,562) -- (257,564) -- (254,564) -- (248,569) -- (238,574) -- (234,578) -- (232,579) -- (229,579) -- (226,582) -- (224,582) -- (216,586) -- (211,590) -- (205,593) -- (203,593) -- (201,595) -- (197,595) -- (194,598) -- (193,598) -- (190,601) -- (188,601) -- (185,604) -- (183,604) -- (181,606) -- (178,606) -- (176,607) -- (166,615) -- (164,615) -- (160,617) -- (158,619) -- (155,619) -- (149,624) -- (147,624) -- (145,626) -- (143,626) -- (140,629) -- (135,630) -- (133,632) -- (128,634) -- (123,638) -- (120,639) -- (117,642) -- (112,643) -- (110,645) -- (108,645) -- (105,648) -- (103,648) -- (99,651) -- (97,651) -- (93,656) -- (93,657) -- (95,658) -- (131,658) -- (132,659) -- (137,659) -- (138,658) -- (140,659) -- (146,659) -- (147,658) -- (149,658) -- (150,659) -- (165,659) -- (166,658) -- (172,658) -- (173,659) -- (175,658) -- (188,658) -- (189,659) -- (226,659) -- (227,658) -- (228,659) -- (270,659) -- (271,658) -- (276,658) -- (277,659) -- (280,659) -- (281,658) -- (285,658) -- (286,657) -- (288,657) -- (289,658) -- (305,658) -- (306,659) -- (308,659) -- (311,657) -- (313,658) -- (315,658) -- (316,657) -- (321,657) -- (322,658) -- (345,658) -- (346,659) -- (348,658) -- (351,658) -- (352,657) -- (354,658) -- (359,658) -- (360,659) -- (362,659) -- (363,658) -- (377,658) -- (378,659) -- (381,659) -- (382,658) -- (405,658) -- (406,659) -- (408,659) -- (409,658) -- (416,658) -- (417,659) -- (422,659) -- (423,658) -- (428,658) -- (429,659) -- (447,659) -- (448,658) -- (450,659) -- (458,659) -- (459,658) -- (466,658) -- (467,659) -- (477,659) -- (478,658) -- (521,658) -- (522,659) -- (525,659) -- (526,658) -- (546,658) -- (547,659) -- (553,659) -- (554,658) -- (560,658) -- (561,659) -- (563,659) -- (564,658) -- (571,658) -- (572,659) -- (574,659) -- (575,658) -- (577,658) -- (578,659) -- (580,658) -- (582,659) -- (584,658) -- (599,658) -- (600,659) -- (602,659) -- (603,658) -- (644,658) -- (646,659) -- (648,658) -- (649,659) -- (651,659) -- (652,658) -- (659,658) -- (660,659) -- (664,659) -- (666,658) -- (667,659) -- (670,657) -- (685,657) -- (686,658) -- (688,658) -- (689,657) -- (696,657) -- (697,658) -- (702,658) -- (703,657) -- (705,657) -- (706,658) -- (720,658) -- (721,657) -- (723,657) -- (724,658) -- (728,658) -- (729,657) -- (732,657) -- (733,658) -- (736,658) -- (737,657) -- (749,657) -- (750,658) -- (755,658) -- (756,657) -- (763,657) -- (764,658) -- (775,658) -- (776,657) -- (779,657) -- (780,658) -- (794,658) -- (795,657) -- (798,657) -- (799,658) -- (821,658) -- (822,659) -- (824,658) -- (828,658) -- (830,659) -- (833,657) -- (835,658) -- (846,658) -- (847,659) -- (849,659) -- (850,658) -- (859,659) -- (862,657) -- (866,657) -- (867,658) -- (869,657) -- (884,657) -- (885,658) -- (891,658) -- (892,657) -- (899,657) -- (900,658) -- (902,657) -- (905,657) -- (906,658) -- (911,658) -- (912,657) -- (921,657) -- (922,658) -- (943,658) -- (944,659) -- (946,659) -- (947,658) -- (965,658) -- (966,657) -- (968,657) -- (969,658) -- (1011,658) -- (1012,659) -- (1014,658) -- (1035,658) -- (1036,659) -- (1038,658) -- (1041,658) -- (1042,659) -- (1047,658) -- (1048,659) -- (1050,659) -- (1052,657) -- (1058,657) -- (1059,658) -- (1060,657) -- (1064,657) -- (1065,658) -- (1080,658) -- (1081,657) -- (1083,657) -- (1087,659) -- (1091,659) -- (1094,657) -- (1105,657) -- (1106,658) -- (1117,657) -- (1121,659) -- (1123,658) -- (1161,658) -- (1162,659) -- (1164,658) -- (1170,658) -- (1171,659) -- (1173,658) -- (1185,658) -- (1186,659) -- (1197,659) -- (1198,660) -- (1200,659) -- (1203,659) -- (1204,658) -- (1208,658) -- (1209,659) -- (1220,659) -- (1221,660) -- (1223,659) -- (1241,659) -- (1242,658) -- (1254,658) -- (1255,659) -- (1257,658) -- (1259,659) -- (1263,659) -- (1264,658) -- (1272,658) -- (1273,659) -- (1283,659) -- (1284,658) -- (1286,658) -- (1287,659) -- (1289,658) -- (1299,658) -- (1300,659) -- (1302,659) -- (1303,658) -- (1306,658) -- (1307,659) -- (1337,659) -- (1338,658) -- (1340,658) -- (1341,659) -- (1350,659) -- (1351,658) -- (1353,659) -- (1365,659) -- (1366,658) -- (1371,658) -- (1372,659) -- (1374,659) -- (1375,658) -- (1377,658) -- (1378,659) -- (1397,659) -- (1398,658) -- (1404,658) -- (1407,660) -- (1412,660) -- (1413,659) -- (1418,659) -- (1419,660) -- (1428,660) -- (1429,659) -- (1437,659) -- (1438,660) -- (1443,660) -- (1444,659) -- (1466,659) -- (1467,660) -- (1470,660) -- (1471,659) -- (1477,659) -- (1478,660) -- (1480,660) -- (1481,659) -- (1486,659) -- (1487,660) -- (1491,660) -- (1492,659) -- (1494,659) -- (1495,660) -- (1521,660) -- (1522,659) -- (1530,659) -- (1531,660) -- (1539,660) -- (1540,661) -- (1542,661) -- cycle
(1518,653) -- (1517,655) -- (1501,655) -- (1500,654) -- (1498,654) -- (1497,655) -- (1480,655) -- (1479,654) -- (1474,654) -- (1472,655) -- (1471,654) -- (1454,654) -- (1453,653) -- (1433,653) -- (1432,654) -- (1427,654) -- (1426,655) -- (1424,655) -- (1423,654) -- (1419,654) -- (1418,655) -- (1413,655) -- (1412,654) -- (1410,654) -- (1409,655) -- (1407,655) -- (1406,654) -- (1396,654) -- (1395,655) -- (1392,655) -- (1391,654) -- (1381,654) -- (1380,653) -- (1378,654) -- (1371,654) -- (1370,653) -- (1368,653) -- (1367,654) -- (1360,654) -- (1358,652) -- (1356,653) -- (1352,653) -- (1351,654) -- (1349,653) -- (1347,654) -- (1341,654) -- (1340,653) -- (1334,653) -- (1333,654) -- (1329,654) -- (1328,653) -- (1326,653) -- (1325,654) -- (1323,653) -- (1321,653) -- (1320,654) -- (1318,653) -- (1316,653) -- (1315,654) -- (1313,654) -- (1312,653) -- (1308,653) -- (1307,654) -- (1303,654) -- (1302,653) -- (1294,653) -- (1292,654) -- (1291,653) -- (1264,653) -- (1263,652) -- (1242,652) -- (1241,653) -- (1237,653) -- (1236,654) -- (1234,654) -- (1233,653) -- (1201,653) -- (1200,652) -- (1197,652) -- (1196,653) -- (1194,653) -- (1193,652) -- (1149,652) -- (1148,651) -- (1145,651) -- (1144,652) -- (1142,651) -- (1140,652) -- (1116,652) -- (1115,651) -- (1113,651) -- (1112,652) -- (1104,652) -- (1103,651) -- (1101,652) -- (1069,652) -- (1068,651) -- (1066,651) -- (1065,652) -- (1061,652) -- (1059,651) -- (1058,652) -- (1040,652) -- (1039,653) -- (1034,653) -- (1033,652) -- (1023,652) -- (1022,653) -- (1020,653) -- (1019,652) -- (1016,652) -- (1015,653) -- (1013,653) -- (1012,652) -- (1009,652) -- (1008,653) -- (1004,653) -- (1003,652) -- (996,652) -- (995,653) -- (992,653) -- (990,652) -- (988,653) -- (987,652) -- (978,652) -- (977,651) -- (974,651) -- (973,652) -- (943,652) -- (942,651) -- (940,652) -- (900,652) -- (899,651) -- (894,651) -- (893,652) -- (890,652) -- (887,650) -- (884,650) -- (882,652) -- (860,652) -- (859,653) -- (856,653) -- (855,654) -- (850,654) -- (849,653) -- (840,653) -- (839,654) -- (831,654) -- (830,653) -- (740,653) -- (739,652) -- (737,652) -- (736,653) -- (717,653) -- (716,652) -- (709,652) -- (708,653) -- (705,653) -- (704,652) -- (676,652) -- (675,651) -- (673,652) -- (669,652) -- (668,653) -- (645,653) -- (644,652) -- (642,652) -- (641,653) -- (637,653) -- (636,652) -- (632,652) -- (631,653) -- (620,653) -- (619,652) -- (617,652) -- (616,653) -- (592,653) -- (591,652) -- (589,652) -- (588,653) -- (537,653) -- (536,654) -- (534,654) -- (533,653) -- (529,653) -- (528,654) -- (526,653) -- (516,653) -- (515,654) -- (512,654) -- (511,653) -- (507,653) -- (506,654) -- (504,653) -- (478,653) -- (477,654) -- (475,654) -- (474,653) -- (472,654) -- (466,654) -- (465,653) -- (463,654) -- (449,654) -- (448,653) -- (444,653) -- (443,654) -- (434,654) -- (433,653) -- (400,653) -- (399,654) -- (397,654) -- (396,653) -- (381,653) -- (380,652) -- (378,652) -- (377,653) -- (373,653) -- (372,652) -- (368,652) -- (367,653) -- (365,653) -- (362,651) -- (358,652) -- (357,653) -- (355,653) -- (354,652) -- (348,652) -- (347,653) -- (340,653) -- (339,652) -- (337,653) -- (335,652) -- (329,652) -- (328,653) -- (321,653) -- (320,652) -- (317,652) -- (316,653) -- (314,653) -- (313,652) -- (311,652) -- (310,653) -- (307,653) -- (306,654) -- (303,654) -- (300,652) -- (298,653) -- (286,653) -- (285,654) -- (272,654) -- (271,653) -- (265,653) -- (264,654) -- (185,654) -- (184,655) -- (166,655) -- (165,654) -- (161,654) -- (160,655) -- (146,655) -- (145,654) -- (143,654) -- (142,655) -- (124,655) -- (123,654) -- (121,654) -- (120,655) -- (114,655) -- (113,654) -- (113,652) -- (117,650) -- (125,649) -- (126,648) -- (131,648) -- (132,647) -- (135,647) -- (136,646) -- (141,646) -- (148,643) -- (153,643) -- (157,641) -- (161,641) -- (165,639) -- (173,638) -- (176,636) -- (179,636) -- (180,635) -- (183,635) -- (184,634) -- (191,634) -- (193,632) -- (199,632) -- (202,630) -- (205,630) -- (206,629) -- (211,629) -- (215,626) -- (223,626) -- (226,624) -- (231,624) -- (234,622) -- (238,622) -- (245,619) -- (250,619) -- (251,618) -- (258,617) -- (264,614) -- (269,614) -- (275,611) -- (282,610) -- (287,607) -- (293,607) -- (297,604) -- (306,603) -- (307,602) -- (310,602) -- (311,601) -- (316,601) -- (317,600) -- (323,599) -- (324,598) -- (327,598) -- (331,596) -- (337,596) -- (338,595) -- (340,595) -- (347,592) -- (355,591) -- (356,590) -- (360,590) -- (361,589) -- (364,589) -- (368,587) -- (373,587) -- (375,585) -- (382,585) -- (383,584) -- (387,584) -- (388,583) -- (391,583) -- (396,580) -- (400,580) -- (401,579) -- (404,579) -- (405,578) -- (415,577) -- (417,575) -- (422,575) -- (429,572) -- (436,572) -- (439,570) -- (447,569) -- (451,566) -- (454,566) -- (455,565) -- (462,565) -- (465,563) -- (469,563) -- (471,561) -- (475,561) -- (477,559) -- (479,559) -- (483,557) -- (487,557) -- (488,556) -- (495,555) -- (499,553) -- (507,552) -- (510,550) -- (514,550) -- (518,548) -- (521,548) -- (524,546) -- (532,545) -- (533,544) -- (536,544) -- (537,543) -- (542,543) -- (545,541) -- (551,541) -- (557,538) -- (567,537) -- (570,535) -- (573,535) -- (574,534) -- (585,533) -- (586,532) -- (588,532) -- (595,529) -- (600,529) -- (601,528) -- (604,528) -- (605,527) -- (611,527) -- (615,524) -- (618,524) -- (620,525) -- (624,523) -- (631,522) -- (635,519) -- (639,520) -- (642,518) -- (647,518) -- (652,515) -- (655,515) -- (656,514) -- (661,514) -- (662,513) -- (666,513) -- (670,510) -- (672,510) -- (673,509) -- (678,509) -- (679,508) -- (683,508) -- (684,507) -- (689,507) -- (694,504) -- (697,504) -- (698,503) -- (709,502) -- (716,499) -- (723,499) -- (726,496) -- (729,496) -- (730,495) -- (737,495) -- (738,494) -- (741,494) -- (745,492) -- (749,492) -- (750,491) -- (753,491) -- (754,490) -- (759,490) -- (762,488) -- (764,488) -- (770,485) -- (783,484) -- (786,482) -- (794,481) -- (795,480) -- (800,481) -- (801,480) -- (804,480) -- (805,479) -- (809,479) -- (812,476) -- (819,475) -- (823,473) -- (827,473) -- (835,469) -- (841,469) -- (842,468) -- (848,469) -- (851,471) -- (853,471) -- (857,473) -- (862,473) -- (863,474) -- (866,474) -- (869,476) -- (881,478) -- (885,481) -- (891,481) -- (892,482) -- (895,482) -- (899,485) -- (903,485) -- (906,487) -- (910,487) -- (911,488) -- (914,488) -- (915,489) -- (919,489) -- (922,491) -- (925,491) -- (929,493) -- (934,493) -- (935,494) -- (938,494) -- (942,496) -- (945,496) -- (947,498) -- (951,498) -- (952,499) -- (954,499) -- (955,500) -- (959,500) -- (966,503) -- (970,503) -- (974,505) -- (977,505) -- (978,506) -- (981,506) -- (982,507) -- (986,507) -- (994,511) -- (998,511) -- (1001,513) -- (1012,515) -- (1016,517) -- (1021,517) -- (1022,518) -- (1025,518) -- (1028,520) -- (1031,520) -- (1035,522) -- (1039,522) -- (1042,524) -- (1045,524) -- (1046,525) -- (1053,525) -- (1057,527) -- (1062,527) -- (1066,529) -- (1069,529) -- (1070,530) -- (1075,531) -- (1078,533) -- (1082,533) -- (1083,534) -- (1090,535) -- (1094,537) -- (1099,537) -- (1105,540) -- (1109,540) -- (1111,542) -- (1121,543) -- (1125,546) -- (1129,546) -- (1130,547) -- (1133,547) -- (1134,548) -- (1138,548) -- (1144,551) -- (1149,551) -- (1152,553) -- (1157,553) -- (1165,557) -- (1176,559) -- (1178,561) -- (1183,561) -- (1187,563) -- (1195,564) -- (1201,567) -- (1203,567) -- (1206,569) -- (1210,569) -- (1212,571) -- (1217,571) -- (1219,573) -- (1224,573) -- (1229,576) -- (1235,576) -- (1236,577) -- (1239,577) -- (1240,578) -- (1246,578) -- (1250,580) -- (1257,581) -- (1262,584) -- (1270,585) -- (1271,586) -- (1274,586) -- (1275,587) -- (1279,587) -- (1280,588) -- (1288,590) -- (1291,592) -- (1294,592) -- (1295,593) -- (1299,593) -- (1300,594) -- (1311,596) -- (1317,599) -- (1326,600) -- (1327,601) -- (1330,601) -- (1333,603) -- (1337,603) -- (1341,605) -- (1348,606) -- (1349,607) -- (1351,607) -- (1358,610) -- (1362,610) -- (1367,613) -- (1375,614) -- (1377,616) -- (1382,616) -- (1389,619) -- (1394,619) -- (1395,620) -- (1398,620) -- (1401,622) -- (1403,622) -- (1407,624) -- (1410,624) -- (1411,625) -- (1414,625) -- (1420,628) -- (1425,628) -- (1429,630) -- (1436,630) -- (1437,631) -- (1443,632) -- (1444,633) -- (1449,633) -- (1451,635) -- (1454,635) -- (1455,636) -- (1466,638) -- (1470,640) -- (1474,640) -- (1477,642) -- (1480,642) -- (1484,644) -- (1489,644) -- (1495,647) -- (1502,648) -- (1506,650) -- (1510,650) -- (1511,651) -- (1517,652) -- cycle
(846,260) -- (848,258) -- (850,258) -- (856,263) -- (861,263) -- (864,266) -- (875,271) -- (886,279) -- (890,281) -- (892,281) -- (902,288) -- (904,288) -- (909,292) -- (913,293) -- (918,298) -- (931,304) -- (934,307) -- (938,308) -- (940,310) -- (942,310) -- (952,318) -- (959,320) -- (962,323) -- (964,323) -- (977,332) -- (979,332) -- (982,335) -- (984,335) -- (989,340) -- (993,341) -- (997,343) -- (999,345) -- (1003,346) -- (1010,352) -- (1012,352) -- (1016,354) -- (1018,356) -- (1021,357) -- (1024,360) -- (1026,360) -- (1028,362) -- (1035,365) -- (1038,368) -- (1042,369) -- (1044,371) -- (1048,373) -- (1051,373) -- (1058,378) -- (1061,378) -- (1064,381) -- (1065,381) -- (1067,383) -- (1069,383) -- (1078,390) -- (1094,398) -- (1097,401) -- (1102,403) -- (1104,405) -- (1111,408) -- (1117,413) -- (1121,415) -- (1123,415) -- (1127,418) -- (1129,418) -- (1134,423) -- (1137,423) -- (1140,426) -- (1142,426) -- (1146,429) -- (1151,431) -- (1153,433) -- (1159,435) -- (1161,437) -- (1164,438) -- (1174,445) -- (1178,447) -- (1180,447) -- (1183,450) -- (1186,451) -- (1192,456) -- (1199,458) -- (1202,461) -- (1207,463) -- (1214,469) -- (1216,469) -- (1219,472) -- (1222,473) -- (1224,475) -- (1226,475) -- (1228,477) -- (1233,479) -- (1236,482) -- (1241,483) -- (1244,486) -- (1248,488) -- (1250,488) -- (1254,492) -- (1261,495) -- (1264,498) -- (1266,498) -- (1270,500) -- (1273,503) -- (1276,504) -- (1282,509) -- (1285,509) -- (1287,511) -- (1290,512) -- (1293,515) -- (1296,516) -- (1298,518) -- (1300,518) -- (1306,523) -- (1308,523) -- (1313,527) -- (1317,529) -- (1319,529) -- (1322,532) -- (1330,536) -- (1336,541) -- (1338,541) -- (1342,545) -- (1348,547) -- (1350,549) -- (1353,550) -- (1357,554) -- (1359,554) -- (1363,557) -- (1365,557) -- (1367,559) -- (1370,560) -- (1375,565) -- (1377,565) -- (1381,567) -- (1383,569) -- (1385,569) -- (1401,580) -- (1404,580) -- (1406,581) -- (1408,583) -- (1411,584) -- (1414,587) -- (1422,591) -- (1427,595) -- (1435,597) -- (1445,605) -- (1447,605) -- (1449,607) -- (1456,609) -- (1461,613) -- (1465,614) -- (1470,619) -- (1472,619) -- (1474,621) -- (1479,623) -- (1482,626) -- (1486,627) -- (1488,629) -- (1493,631) -- (1497,635) -- (1506,639) -- (1508,641) -- (1512,642) -- (1514,644) -- (1516,644) -- (1518,646) -- (1519,646) -- (1521,648) -- (1521,649) -- (1520,650) -- (1518,650) -- (1517,649) -- (1515,649) -- (1508,646) -- (1505,646) -- (1502,644) -- (1499,644) -- (1496,642) -- (1493,643) -- (1490,641) -- (1487,641) -- (1480,638) -- (1476,638) -- (1470,635) -- (1459,633) -- (1456,631) -- (1454,631) -- (1447,628) -- (1444,628) -- (1440,626) -- (1436,626) -- (1432,624) -- (1427,624) -- (1423,622) -- (1420,622) -- (1417,620) -- (1413,620) -- (1412,619) -- (1409,620) -- (1405,617) -- (1402,617) -- (1401,616) -- (1396,616) -- (1390,613) -- (1388,613) -- (1381,610) -- (1376,610) -- (1375,609) -- (1373,609) -- (1370,607) -- (1367,607) -- (1366,606) -- (1360,606) -- (1358,604) -- (1356,604) -- (1352,602) -- (1349,602) -- (1346,600) -- (1341,600) -- (1334,597) -- (1330,597) -- (1329,596) -- (1327,596) -- (1324,594) -- (1320,594) -- (1316,592) -- (1304,590) -- (1301,588) -- (1298,588) -- (1296,586) -- (1293,586) -- (1289,584) -- (1283,584) -- (1280,582) -- (1277,582) -- (1271,579) -- (1267,579) -- (1266,578) -- (1263,578) -- (1259,576) -- (1253,575) -- (1251,573) -- (1249,573) -- (1248,572) -- (1233,571) -- (1231,569) -- (1228,569) -- (1224,567) -- (1220,567) -- (1218,565) -- (1214,565) -- (1213,564) -- (1210,564) -- (1209,563) -- (1205,563) -- (1203,561) -- (1200,561) -- (1199,560) -- (1197,560) -- (1190,557) -- (1185,557) -- (1183,555) -- (1177,555) -- (1174,553) -- (1172,553) -- (1171,552) -- (1167,552) -- (1163,550) -- (1160,550) -- (1157,548) -- (1153,548) -- (1152,547) -- (1141,545) -- (1137,543) -- (1131,542) -- (1128,540) -- (1125,540) -- (1121,538) -- (1118,538) -- (1112,535) -- (1108,535) -- (1107,534) -- (1100,533) -- (1097,531) -- (1093,531) -- (1086,528) -- (1082,528) -- (1081,527) -- (1070,525) -- (1067,523) -- (1060,522) -- (1056,520) -- (1046,519) -- (1042,517) -- (1037,517) -- (1034,515) -- (1030,515) -- (1029,514) -- (1024,513) -- (1021,511) -- (1011,510) -- (1008,508) -- (1002,508) -- (998,505) -- (995,505) -- (994,504) -- (988,504) -- (986,502) -- (982,502) -- (979,500) -- (969,499) -- (957,494) -- (954,494) -- (953,493) -- (948,493) -- (947,492) -- (942,492) -- (940,490) -- (936,490) -- (935,489) -- (933,489) -- (930,487) -- (927,487) -- (926,486) -- (921,486) -- (915,483) -- (908,482) -- (904,480) -- (896,479) -- (893,477) -- (882,475) -- (879,473) -- (873,473) -- (869,471) -- (865,471) -- (864,470) -- (861,470) -- (860,469) -- (856,469) -- (855,468) -- (852,468) -- (851,467) -- (847,466) -- (845,464) -- (845,460) -- (846,459) -- (847,448) -- (848,447) -- (849,442) -- (847,440) -- (847,436) -- (846,435) -- (847,429) -- (848,428) -- (848,426) -- (847,425) -- (847,422) -- (848,421) -- (848,416) -- (846,414) -- (846,397) -- (847,396) -- (847,392) -- (845,390) -- (846,388) -- (846,376) -- (847,375) -- (846,373) -- (846,370) -- (847,369) -- (847,361) -- (846,360) -- (847,359) -- (846,358) -- (846,353) -- (847,352) -- (847,349) -- (848,348) -- (847,346) -- (847,338) -- (846,337) -- (846,329) -- (847,328) -- (846,326) -- (847,324) -- (847,322) -- (846,321) -- (846,308) -- (847,307) -- (847,278) -- (848,277) -- (847,275) -- (847,262) -- cycle
(843,260) -- (843,276) -- (844,277) -- (844,282) -- (843,284) -- (844,285) -- (844,287) -- (843,288) -- (843,300) -- (842,301) -- (842,306) -- (841,307) -- (842,309) -- (842,313) -- (840,315) -- (840,326) -- (841,327) -- (841,338) -- (840,339) -- (840,342) -- (841,344) -- (840,345) -- (840,347) -- (841,348) -- (841,353) -- (839,355) -- (839,358) -- (841,361) -- (841,372) -- (840,373) -- (840,377) -- (841,378) -- (840,379) -- (840,382) -- (841,384) -- (840,385) -- (840,387) -- (841,388) -- (841,392) -- (842,393) -- (842,394) -- (840,396) -- (840,399) -- (841,400) -- (841,402) -- (840,403) -- (840,415) -- (841,416) -- (841,421) -- (839,423) -- (839,424) -- (840,425) -- (840,438) -- (841,439) -- (841,442) -- (839,444) -- (840,446) -- (839,450) -- (840,451) -- (840,454) -- (841,455) -- (843,455) -- (844,456) -- (844,458) -- (842,460) -- (842,463) -- (840,465) -- (836,467) -- (829,467) -- (825,469) -- (821,469) -- (820,470) -- (817,470) -- (816,471) -- (810,470) -- (808,472) -- (804,472) -- (800,475) -- (794,475) -- (790,477) -- (781,478) -- (780,479) -- (778,479) -- (776,481) -- (763,482) -- (758,485) -- (748,486) -- (746,488) -- (742,488) -- (741,489) -- (732,490) -- (729,492) -- (725,492) -- (722,494) -- (707,495) -- (705,497) -- (701,497) -- (698,499) -- (692,499) -- (690,501) -- (683,501) -- (679,503) -- (674,503) -- (672,505) -- (670,505) -- (664,508) -- (658,508) -- (654,510) -- (651,510) -- (650,511) -- (645,511) -- (642,513) -- (637,513) -- (636,514) -- (633,514) -- (630,516) -- (625,516) -- (624,517) -- (621,517) -- (620,518) -- (613,518) -- (612,519) -- (608,520) -- (606,522) -- (598,523) -- (597,524) -- (593,524) -- (592,525) -- (585,526) -- (584,527) -- (582,527) -- (575,530) -- (569,530) -- (566,532) -- (558,532) -- (555,534) -- (551,534) -- (549,536) -- (547,536) -- (546,537) -- (541,537) -- (539,539) -- (534,539) -- (533,540) -- (530,540) -- (528,542) -- (520,542) -- (517,544) -- (509,545) -- (506,547) -- (497,548) -- (496,549) -- (493,549) -- (492,550) -- (487,550) -- (485,552) -- (482,552) -- (480,554) -- (473,555) -- (470,557) -- (466,557) -- (465,558) -- (462,558) -- (461,559) -- (457,559) -- (454,561) -- (449,561) -- (446,563) -- (442,563) -- (439,565) -- (436,564) -- (431,567) -- (428,567) -- (424,569) -- (419,569) -- (418,570) -- (412,570) -- (411,571) -- (408,571) -- (406,573) -- (402,573) -- (400,575) -- (396,575) -- (395,576) -- (392,576) -- (391,577) -- (388,577) -- (387,578) -- (384,578) -- (383,579) -- (375,579) -- (371,582) -- (368,582) -- (367,583) -- (361,583) -- (360,584) -- (357,584) -- (353,586) -- (350,586) -- (349,587) -- (345,587) -- (344,588) -- (337,589) -- (331,592) -- (326,592) -- (324,594) -- (321,594) -- (320,595) -- (317,595) -- (316,596) -- (306,597) -- (300,600) -- (294,600) -- (293,601) -- (290,601) -- (288,603) -- (286,603) -- (283,605) -- (277,606) -- (276,607) -- (271,607) -- (265,610) -- (262,610) -- (261,611) -- (254,612) -- (250,614) -- (247,614) -- (245,616) -- (239,616) -- (236,618) -- (230,618) -- (227,620) -- (223,620) -- (222,621) -- (218,620) -- (215,622) -- (212,622) -- (209,624) -- (206,624) -- (200,627) -- (196,627) -- (192,629) -- (185,630) -- (179,633) -- (176,633) -- (175,634) -- (168,634) -- (166,636) -- (164,636) -- (163,637) -- (160,637) -- (159,638) -- (156,638) -- (155,639) -- (152,639) -- (151,640) -- (146,640) -- (140,643) -- (135,643) -- (134,644) -- (128,645) -- (121,648) -- (119,648) -- (118,647) -- (118,645) -- (119,644) -- (123,643) -- (128,639) -- (131,639) -- (135,637) -- (140,633) -- (142,633) -- (146,630) -- (148,630) -- (154,627) -- (156,625) -- (162,623) -- (167,619) -- (173,618) -- (175,616) -- (180,614) -- (188,608) -- (190,608) -- (193,605) -- (197,604) -- (200,601) -- (203,601) -- (205,599) -- (208,598) -- (212,594) -- (214,594) -- (216,592) -- (218,592) -- (219,591) -- (226,590) -- (227,588) -- (233,583) -- (235,582) -- (238,582) -- (242,578) -- (247,577) -- (254,572) -- (257,572) -- (260,569) -- (264,568) -- (272,562) -- (275,561) -- (278,558) -- (285,556) -- (292,550) -- (294,550) -- (299,546) -- (303,546) -- (311,539) -- (314,539) -- (318,536) -- (322,535) -- (332,529) -- (334,529) -- (337,526) -- (346,522) -- (351,518) -- (354,518) -- (357,515) -- (360,515) -- (367,510) -- (371,509) -- (374,506) -- (377,506) -- (379,504) -- (383,503) -- (388,499) -- (396,495) -- (398,495) -- (402,491) -- (409,489) -- (413,485) -- (420,483) -- (422,481) -- (425,480) -- (429,477) -- (433,476) -- (435,474) -- (440,473) -- (442,471) -- (445,470) -- (452,465) -- (456,464) -- (459,461) -- (464,459) -- (466,457) -- (468,457) -- (472,454) -- (474,454) -- (478,450) -- (481,449) -- (483,447) -- (490,445) -- (495,441) -- (499,440) -- (505,437) -- (508,434) -- (510,434) -- (513,431) -- (516,431) -- (522,426) -- (525,426) -- (531,423) -- (533,421) -- (535,421) -- (541,418) -- (545,414) -- (547,414) -- (551,412) -- (554,409) -- (558,409) -- (561,406) -- (565,405) -- (567,403) -- (572,401) -- (575,398) -- (579,397) -- (582,394) -- (584,393) -- (587,393) -- (589,391) -- (593,390) -- (596,387) -- (602,384) -- (604,384) -- (606,382) -- (614,378) -- (617,378) -- (623,374) -- (626,374) -- (633,368) -- (635,368) -- (639,365) -- (649,361) -- (651,359) -- (655,359) -- (666,351) -- (672,348) -- (676,348) -- (678,347) -- (682,343) -- (686,342) -- (690,337) -- (700,332) -- (702,332) -- (704,330) -- (710,329) -- (714,325) -- (717,324) -- (719,322) -- (721,322) -- (725,320) -- (729,316) -- (732,316) -- (736,314) -- (738,312) -- (744,310) -- (746,308) -- (751,307) -- (755,303) -- (757,303) -- (763,300) -- (766,297) -- (780,290) -- (782,290) -- (789,285) -- (792,285) -- (795,282) -- (797,282) -- (801,280) -- (804,277) -- (808,277) -- (811,274) -- (814,273) -- (816,271) -- (820,270) -- (824,267) -- (827,267) -- (832,262) -- (836,261) -- (839,258) -- (841,258) -- cycle
(842,233) -- (841,234) -- (842,237) -- (841,238) -- (840,243) -- (842,245) -- (844,245) -- (846,243) -- (846,239) -- (845,238) -- (845,234) -- (844,233) -- cycle
(842,207) -- (840,209) -- (840,222) -- (842,224) -- (843,224) -- (845,222) -- (844,220) -- (845,218) -- (844,216) -- (845,214) -- (845,210) -- (844,208) -- cycle
(841,178) -- (840,179) -- (840,185) -- (838,188) -- (838,192) -- (839,193) -- (839,197) -- (841,199) -- (844,199) -- (845,198) -- (845,194) -- (846,193) -- (846,185) -- (845,184) -- (846,179) -- (845,178) -- cycle
(842,153) -- (840,156) -- (841,160) -- (839,164) -- (840,165) -- (840,168) -- (841,169) -- (844,169) -- (846,167) -- (847,155) -- (845,153) -- cycle
(845,122) -- (843,122) -- (841,124) -- (840,139) -- (842,141) -- (844,141) -- (846,139) -- (847,131) -- (848,130) -- (847,124) -- cycle
(843,88) -- (840,92) -- (840,96) -- (842,98) -- (842,101) -- (840,105) -- (841,106) -- (841,109) -- (842,110) -- (845,110) -- (847,108) -- (847,102) -- (848,101) -- (847,100) -- (847,97) -- (848,96) -- (848,91) -- (846,88) -- cycle
(842,66) -- (840,70) -- (840,72) -- (841,73) -- (840,77) -- (843,81) -- (845,81) -- (847,79) -- (847,72) -- (848,70) -- (846,66) -- cycle
(846,43) -- (844,43) -- (840,47) -- (841,56) -- (844,58) -- (845,58) -- (848,55) -- (848,51) -- (847,50) -- (848,46) -- cycle
(843,19) -- (841,22) -- (841,32) -- (842,33) -- (846,33) -- (848,31) -- (846,27) -- (847,25) -- (847,20) -- (846,19) -- cycle;
\node[inner sep=0pt,font=\fontsize{10}{10}\selectfont] at (884.5,212.5) {\scalebox{7.029}{$X$}};
\node[inner sep=0pt,font=\fontsize{10}{10}\selectfont] at (826.0,512.5) {\scalebox{5.841}{$B$}};
\node[inner sep=0pt,font=\fontsize{10}{10}\selectfont] at (45.0,678.0) {\scalebox{7.722}{$A$}};
\node[inner sep=0pt,font=\fontsize{10}{10}\selectfont] at (1581.5,667.5) {\scalebox{7.425}{$C$}};
\end{tikzpicture}%
}
\end{center}

Hier würde Ihnen die Rechnung den Punkt $X$ geben, der aber nur
passt, wenn Sie statt der Richtung $XB$ die entgegengesetzte
nehmen.

\begin{center}
\textsc{Gauss} an \textsc{Schumacher}. Göttingen, 23. April 1836.
\end{center}

\indent \mbox{. . . . . . .} Gleich nach Absendung meines letzten Briefes
dachte ich doch, dass der hinterher eingelegte Zettel\HistoricalNote{[*)]}{[*)}{Dieser Zettel fehlt. In seiner Antwort vom 26. April 1836 bemerkt
\textsc{Schumacher}, dass er in dem Briefe keine Einlage vorgefunden
habe.]} Sie wohl eher
verwirren, als das in dem Briefe angezogene ohne weiteres verständliche
Beispiel weiter erklären würde. Ich hatte nicht gehörig bedacht, dass
[\textls[120]{erstlich} das \textls[120]{allgemeine} auf dem Zettel ausgesprochene Theorem, um
vollkommen richtig aufgefasst zu werden, wohl noch einiger erläuternder
Bemerkungen bedürfen möchte. \textls[120]{Zweitens}, dass gerade der Fall des
gleichseitigen Dreiecks ohne solche erläuternde Bemerkungen
%
\SourceLeaf{318}
am leichtesten missverstanden werden kann. Drittens hätte auch hier ausdrücklich
in Beziehung auf das Beispiel eine literirte Figur beigefügt werden müssen. Und
viertens glaube ich sogar einen Schreibfehler begangen und anstatt
$60^0+60^0=120^0$ geschrieben zu haben $30^0+30^0=60^0$. Wollte ich nun diese
3 oder 4 Punkte alle und ganz aufklären, so würde dazu ein längerer Brief
nöthig sein, als ich heute schreiben kann. Einstweilen ignoriren Sie also das
Beispiel mit dem gleichseitigen Dreieck ganz (was ich übrigens, wenn es der
Mühe werth scheinen sollte, zu einer andern Zeit zu erklären gern erbötig
bin). Dagegen will ich heute zu weiterer Verständigung über das allgemeine
Theorem das Nöthige beifügen.

Vor allem aber muss ich, um jedes Missverständniss darüber auszuschliessen,
ausdrücklich bemerken, was die Data des \textsc{Pothenot}schen Problems nothwendig
enthalten müssen:

1) müssen die drei gegebenen [Punkte, wenigstens ihrer gegenseitigen \textls[120]{Lage}
nach, vollkommen bestimmt sein, so dass über das Rechts- oder Linksliegen
keine Zweideutigkeit ist. Es ist also z.\,B. nicht zureichend, die Grösse der
Winkel zu kennen, sondern auch ihre \textls[120]{Lage}; also z.\,B. nicht bloss $A=45^0$,
$B=45^0$, $C=90^0$, weil so das Dreieck

\begin{center}\begin{tabular}{c@{\quad}c@{\quad}c@{\quad}l}\resizebox{1.65cm}{!}{%
\begin{tikzpicture}[x=1pt,y=-1pt]
\path[use as bounding box] (0,0) rectangle (290,255);
\path[fill=black,draw=none,even odd rule]
(231,70) -- (227,67) -- (224,67) -- (222,68) -- (217,74) -- (216,74) -- (206,83) -- (205,83) -- (196,92) -- (195,92) -- (192,96) -- (191,96) -- (187,100) -- (184,101) -- (182,103) -- (182,105) -- (180,107) -- (179,107) -- (178,109) -- (176,111) -- (175,111) -- (160,127) -- (159,127) -- (153,133) -- (152,133) -- (142,143) -- (139,144) -- (138,146) -- (125,159) -- (124,159) -- (118,165) -- (116,166) -- (116,167) -- (110,174) -- (105,177) -- (89,193) -- (88,193) -- (82,199) -- (82,203) -- (84,205) -- (88,205) -- (89,206) -- (91,205) -- (95,205) -- (96,206) -- (100,206) -- (101,205) -- (116,205) -- (117,206) -- (118,205) -- (123,205) -- (124,204) -- (126,205) -- (141,205) -- (142,206) -- (149,206) -- (150,205) -- (154,205) -- (155,206) -- (166,205) -- (168,206) -- (169,205) -- (186,205) -- (187,204) -- (194,204) -- (195,205) -- (201,205) -- (202,206) -- (203,206) -- (205,204) -- (216,205) -- (217,206) -- (223,206) -- (225,205) -- (228,202) -- (228,191) -- (231,189) -- (231,187) -- (229,184) -- (229,181) -- (231,179) -- (231,177) -- (232,176) -- (232,174) -- (230,172) -- (230,166) -- (229,165) -- (229,154) -- (230,153) -- (230,138) -- (229,137) -- (229,132) -- (230,131) -- (230,119) -- (228,117) -- (230,113) -- (230,105) -- (229,103) -- (232,99) -- (229,93) -- (230,89) -- (228,86) -- (229,84) -- (229,78) -- (228,76) -- (230,74) -- cycle
(225,77) -- (225,79) -- (223,81) -- (223,83) -- (224,84) -- (223,86) -- (223,104) -- (224,105) -- (224,116) -- (223,117) -- (223,119) -- (224,120) -- (223,122) -- (224,124) -- (224,132) -- (223,133) -- (223,142) -- (224,143) -- (224,147) -- (223,148) -- (223,171) -- (222,172) -- (223,174) -- (223,177) -- (222,178) -- (223,180) -- (223,186) -- (225,189) -- (224,191) -- (225,198) -- (223,200) -- (208,200) -- (207,199) -- (205,200) -- (201,200) -- (200,199) -- (198,199) -- (197,200) -- (190,200) -- (189,199) -- (177,200) -- (176,199) -- (174,199) -- (173,200) -- (164,200) -- (163,201) -- (153,201) -- (152,200) -- (150,201) -- (148,200) -- (146,200) -- (145,201) -- (139,201) -- (138,200) -- (120,200) -- (119,199) -- (117,199) -- (116,200) -- (102,200) -- (101,199) -- (93,200) -- (92,199) -- (92,197) -- (96,193) -- (97,193) -- (118,172) -- (119,172) -- (134,158) -- (135,158) -- (152,141) -- (153,141) -- (159,135) -- (160,135) -- (162,133) -- (162,132) -- (167,127) -- (168,127) -- (174,120) -- (175,120) -- (183,112) -- (184,112) -- (189,107) -- (190,107) -- (192,105) -- (193,102) -- (196,99) -- (201,96) -- (210,86) -- (211,86) -- (218,79) -- (219,79) -- (222,75) -- (223,75) -- cycle;
\node[inner sep=0pt,font=\fontsize{10}{10}\selectfont] at (237.0,30.5) {\scalebox{5.643}{$A$}};
\node[inner sep=0pt,font=\fontsize{10}{10}\selectfont] at (34.0,212.5) {\scalebox{6.039}{$B$}};
\node[inner sep=0pt,font=\fontsize{10}{10}\selectfont] at (258.5,222.0) {\scalebox{6.138}{$C$}};
\end{tikzpicture}%
}
&nicht von&\resizebox{1.65cm}{!}{%
\begin{tikzpicture}[x=1pt,y=-1pt]
\path[use as bounding box] (0,0) rectangle (276,259);
\path[fill=black,draw=none,even odd rule]
(70,61) -- (67,65) -- (68,73) -- (70,76) -- (70,78) -- (69,79) -- (69,84) -- (68,85) -- (68,94) -- (70,97) -- (70,99) -- (69,100) -- (69,103) -- (70,105) -- (69,106) -- (69,116) -- (68,117) -- (68,125) -- (67,126) -- (67,131) -- (68,133) -- (66,135) -- (66,139) -- (68,142) -- (67,144) -- (67,149) -- (66,150) -- (66,152) -- (68,155) -- (68,160) -- (69,161) -- (69,165) -- (68,166) -- (69,170) -- (68,171) -- (68,174) -- (70,177) -- (70,187) -- (71,188) -- (70,189) -- (71,190) -- (71,193) -- (68,198) -- (69,201) -- (72,203) -- (76,202) -- (79,204) -- (82,204) -- (86,202) -- (87,203) -- (98,203) -- (99,204) -- (112,204) -- (114,206) -- (118,206) -- (119,205) -- (123,205) -- (124,204) -- (128,205) -- (129,204) -- (133,204) -- (134,203) -- (138,203) -- (139,204) -- (145,203) -- (148,205) -- (150,205) -- (151,204) -- (156,204) -- (157,203) -- (164,203) -- (165,204) -- (184,204) -- (187,206) -- (189,206) -- (190,205) -- (196,205) -- (197,204) -- (204,204) -- (206,202) -- (206,199) -- (205,197) -- (202,194) -- (202,193) -- (201,192) -- (199,192) -- (192,184) -- (192,183) -- (179,171) -- (179,169) -- (176,165) -- (176,163) -- (174,161) -- (174,160) -- (171,157) -- (168,157) -- (164,153) -- (164,152) -- (158,147) -- (158,146) -- (156,144) -- (156,142) -- (155,141) -- (153,141) -- (142,130) -- (142,129) -- (135,122) -- (135,121) -- (133,119) -- (130,118) -- (126,114) -- (124,110) -- (116,103) -- (114,99) -- (110,97) -- (105,92) -- (105,91) -- (99,85) -- (99,83) -- (89,74) -- (87,70) -- (86,70) -- (77,61) -- (75,60) -- (72,60) -- cycle
(76,71) -- (78,71) -- (80,73) -- (80,74) -- (87,82) -- (87,83) -- (94,90) -- (94,92) -- (96,93) -- (102,99) -- (103,101) -- (106,102) -- (108,104) -- (108,105) -- (112,109) -- (112,110) -- (118,115) -- (121,120) -- (136,135) -- (140,137) -- (143,140) -- (143,141) -- (159,157) -- (159,158) -- (165,164) -- (168,169) -- (171,172) -- (173,172) -- (181,180) -- (181,181) -- (195,195) -- (196,195) -- (198,197) -- (198,198) -- (196,200) -- (190,200) -- (189,199) -- (177,199) -- (176,198) -- (170,198) -- (169,199) -- (165,198) -- (163,199) -- (162,198) -- (152,198) -- (151,197) -- (133,197) -- (132,198) -- (104,198) -- (102,197) -- (100,198) -- (99,197) -- (77,197) -- (75,195) -- (76,189) -- (75,188) -- (74,183) -- (76,180) -- (76,176) -- (75,175) -- (74,168) -- (77,164) -- (76,158) -- (78,156) -- (78,152) -- (75,147) -- (75,136) -- (77,133) -- (77,130) -- (76,129) -- (76,124) -- (75,123) -- (76,121) -- (76,116) -- (75,115) -- (76,113) -- (76,111) -- (75,110) -- (75,108) -- (76,107) -- (76,96) -- (75,95) -- (75,74) -- (74,73) -- cycle;
\node[inner sep=0pt,font=\fontsize{10}{10}\selectfont] at (55.5,29.5) {\scalebox{5.841}{$A$}};
\node[inner sep=0pt,font=\fontsize{10}{10}\selectfont] at (30.5,217.5) {\scalebox{5.643}{$C$}};
\node[inner sep=0pt,font=\fontsize{10}{10}\selectfont] at (243.0,223.0) {\scalebox{6.732}{$B$}};
\end{tikzpicture}%
}
&unterschieden würde.\end{tabular}\end{center}

2) wird dieselbe vollständige Bestimmtheit für die Winkel an dem gesuchten
Punkte $X$ erfordert.

Es ist also z.\,B. nicht zureichend zu sagen, dass in $X$ der Winkel zwischen
$A$ und $B$ $40^0$, zwischen $B$ und $C$ $100^0$ sein soll, weil sonst

\begin{center}\resizebox{4.7cm}{!}{%
\begin{tikzpicture}[x=1pt,y=-1pt]
\path[use as bounding box] (0,0) rectangle (931,400);
\path[fill=black,draw=none,even odd rule]
(80,353) -- (78,351) -- (76,352) -- (66,352) -- (65,353) -- (65,356) -- (67,358) -- (76,358) -- (80,356) -- cycle
(163,351) -- (161,351) -- (160,350) -- (158,350) -- (157,351) -- (147,351) -- (145,353) -- (143,353) -- (141,351) -- (137,351) -- (136,352) -- (134,351) -- (125,351) -- (124,352) -- (124,355) -- (127,357) -- (133,357) -- (135,358) -- (136,357) -- (139,357) -- (141,356) -- (144,353) -- (148,357) -- (150,357) -- (151,358) -- (155,358) -- (158,356) -- (161,356) -- (163,354) -- cycle
(103,353) -- (103,355) -- (105,357) -- (109,357) -- (110,358) -- (112,357) -- (117,357) -- (119,354) -- (117,350) -- (110,350) -- (109,351) -- (105,351) -- cycle
(83,353) -- (83,355) -- (86,358) -- (88,358) -- (89,359) -- (95,358) -- (98,355) -- (98,353) -- (96,350) -- (89,350) -- cycle
(186,351) -- (186,354) -- (188,356) -- (197,356) -- (200,354) -- (200,351) -- (196,349) -- (192,349) -- (191,350) -- (187,350) -- cycle
(182,351) -- (178,349) -- (175,349) -- (174,350) -- (168,351) -- (167,352) -- (167,354) -- (171,357) -- (178,357) -- (182,355) -- cycle
(207,353) -- (208,355) -- (210,356) -- (216,362) -- (222,365) -- (227,365) -- (228,364) -- (228,362) -- (225,360) -- (223,360) -- (218,356) -- (220,354) -- (233,355) -- (234,356) -- (237,356) -- (239,354) -- (249,354) -- (250,353) -- (252,353) -- (253,354) -- (261,354) -- (262,355) -- (264,354) -- (273,354) -- (274,355) -- (278,355) -- (279,354) -- (294,354) -- (295,355) -- (302,355) -- (303,354) -- (306,354) -- (307,355) -- (320,355) -- (321,354) -- (324,354) -- (325,353) -- (347,353) -- (348,354) -- (351,354) -- (352,353) -- (354,353) -- (355,354) -- (363,354) -- (364,353) -- (366,353) -- (367,354) -- (385,354) -- (386,353) -- (401,353) -- (402,354) -- (404,353) -- (415,353) -- (416,354) -- (420,354) -- (421,353) -- (431,353) -- (432,354) -- (435,354) -- (436,353) -- (443,353) -- (444,352) -- (450,352) -- (451,353) -- (475,353) -- (476,352) -- (478,352) -- (479,353) -- (482,353) -- (483,354) -- (489,354) -- (492,352) -- (497,352) -- (498,353) -- (500,353) -- (504,351) -- (516,352) -- (517,351) -- (520,351) -- (521,349) -- (523,351) -- (524,351) -- (526,349) -- (525,347) -- (527,345) -- (532,342) -- (537,337) -- (541,335) -- (545,331) -- (548,330) -- (558,322) -- (564,319) -- (572,310) -- (573,310) -- (577,306) -- (579,306) -- (600,288) -- (601,288) -- (606,284) -- (609,283) -- (618,274) -- (621,273) -- (627,267) -- (631,266) -- (634,263) -- (636,263) -- (639,259) -- (640,259) -- (644,255) -- (645,255) -- (658,243) -- (666,238) -- (672,231) -- (674,231) -- (677,228) -- (678,228) -- (682,223) -- (684,223) -- (687,220) -- (688,220) -- (691,217) -- (696,214) -- (706,205) -- (708,205) -- (713,199) -- (715,199) -- (718,196) -- (720,196) -- (724,194) -- (739,178) -- (743,176) -- (756,163) -- (757,163) -- (759,165) -- (759,166) -- (755,170) -- (755,171) -- (757,173) -- (760,173) -- (761,172) -- (761,167) -- (767,156) -- (767,154) -- (765,152) -- (762,152) -- (756,155) -- (751,155) -- (750,156) -- (743,157) -- (742,158) -- (742,160) -- (743,161) -- (747,161) -- (749,163) -- (744,168) -- (739,171) -- (735,175) -- (731,177) -- (724,184) -- (720,186) -- (715,191) -- (712,192) -- (703,201) -- (698,203) -- (693,208) -- (687,211) -- (681,218) -- (680,218) -- (678,220) -- (676,220) -- (673,223) -- (670,224) -- (666,228) -- (663,229) -- (657,236) -- (656,236) -- (651,241) -- (650,241) -- (647,245) -- (639,250) -- (634,256) -- (632,256) -- (630,258) -- (629,258) -- (620,266) -- (613,270) -- (609,274) -- (606,275) -- (600,281) -- (597,282) -- (593,287) -- (590,288) -- (585,293) -- (582,294) -- (573,302) -- (569,303) -- (566,306) -- (566,308) -- (563,311) -- (559,313) -- (556,317) -- (555,317) -- (550,322) -- (546,324) -- (531,337) -- (526,340) -- (522,344) -- (518,344) -- (503,330) -- (502,330) -- (497,325) -- (496,325) -- (488,317) -- (484,316) -- (481,312) -- (478,311) -- (466,299) -- (463,298) -- (460,295) -- (459,295) -- (447,283) -- (446,283) -- (441,278) -- (440,278) -- (410,251) -- (405,248) -- (403,248) -- (400,245) -- (400,244) -- (390,235) -- (388,235) -- (382,229) -- (380,229) -- (372,221) -- (367,218) -- (360,211) -- (358,211) -- (351,203) -- (350,203) -- (340,194) -- (339,194) -- (324,180) -- (318,177) -- (305,164) -- (302,163) -- (299,160) -- (297,160) -- (293,156) -- (292,156) -- (290,154) -- (291,153) -- (298,153) -- (299,152) -- (299,151) -- (297,149) -- (294,149) -- (291,147) -- (287,147) -- (280,144) -- (274,144) -- (272,146) -- (272,149) -- (276,154) -- (276,156) -- (279,159) -- (280,163) -- (283,166) -- (285,166) -- (286,165) -- (286,162) -- (284,160) -- (284,158) -- (286,157) -- (294,166) -- (294,167) -- (296,169) -- (298,169) -- (301,172) -- (303,176) -- (304,176) -- (307,179) -- (310,179) -- (316,185) -- (319,186) -- (324,191) -- (325,191) -- (328,195) -- (331,196) -- (334,199) -- (336,199) -- (352,214) -- (356,216) -- (361,221) -- (366,224) -- (370,229) -- (371,229) -- (376,233) -- (379,234) -- (393,247) -- (398,250) -- (402,254) -- (404,254) -- (407,257) -- (407,258) -- (408,258) -- (413,263) -- (418,266) -- (425,273) -- (429,275) -- (442,288) -- (446,290) -- (448,292) -- (448,293) -- (457,301) -- (458,301) -- (463,306) -- (464,306) -- (468,310) -- (469,310) -- (474,315) -- (477,316) -- (480,319) -- (486,322) -- (500,336) -- (511,343) -- (513,345) -- (513,346) -- (512,347) -- (494,347) -- (493,346) -- (491,346) -- (485,349) -- (479,349) -- (478,348) -- (476,348) -- (475,349) -- (471,349) -- (470,348) -- (434,348) -- (433,347) -- (423,347) -- (421,348) -- (420,347) -- (415,347) -- (413,346) -- (412,347) -- (404,347) -- (403,348) -- (360,348) -- (359,347) -- (355,347) -- (354,348) -- (352,347) -- (350,347) -- (349,348) -- (330,348) -- (329,347) -- (323,347) -- (321,349) -- (318,349) -- (317,348) -- (308,348) -- (307,349) -- (302,349) -- (301,348) -- (298,348) -- (297,349) -- (293,349) -- (292,350) -- (290,349) -- (284,349) -- (283,348) -- (281,349) -- (274,349) -- (273,348) -- (268,348) -- (267,349) -- (262,349) -- (261,348) -- (248,348) -- (247,349) -- (235,349) -- (234,350) -- (230,349) -- (229,350) -- (223,350) -- (222,349) -- (222,347) -- (226,345) -- (227,343) -- (226,342) -- (222,342) -- (220,344) -- (218,344) -- (211,349) -- (208,350) -- (207,351) -- cycle
(254,133) -- (254,135) -- (258,137) -- (262,141) -- (268,141) -- (270,139) -- (270,137) -- (261,130) -- (256,130) -- cycle
(794,131) -- (791,128) -- (790,128) -- (788,130) -- (786,130) -- (780,137) -- (780,139) -- (779,140) -- (776,140) -- (773,143) -- (771,143) -- (768,146) -- (768,149) -- (771,151) -- (774,151) -- (781,144) -- (781,141) -- (782,140) -- (786,140) -- (787,139) -- (789,139) -- (794,134) -- cycle
(810,119) -- (809,118) -- (803,118) -- (796,124) -- (795,127) -- (797,129) -- (801,129) -- (803,128) -- (805,126) -- (809,124) -- (810,122) -- cycle
(240,120) -- (240,122) -- (248,129) -- (251,129) -- (254,127) -- (254,125) -- (245,116) -- (244,116) -- cycle
(226,105) -- (224,107) -- (224,110) -- (231,116) -- (233,117) -- (236,117) -- (238,115) -- (238,112) -- (234,108) -- (231,107) -- (229,105) -- cycle
(825,106) -- (821,103) -- (819,103) -- (812,110) -- (812,111) -- (810,113) -- (810,116) -- (816,116) -- (818,115) -- (825,109) -- cycle
(842,94) -- (838,91) -- (837,91) -- (835,93) -- (833,93) -- (827,100) -- (827,102) -- (831,104) -- (834,103) -- (838,99) -- (840,99) -- (842,97) -- cycle
(207,92) -- (208,96) -- (216,103) -- (219,103) -- (221,101) -- (221,98) -- (218,94) -- (215,93) -- (213,91) -- (208,91) -- cycle
(192,79) -- (192,83) -- (199,89) -- (201,89) -- (202,90) -- (205,90) -- (206,89) -- (206,85) -- (201,82) -- (198,79) -- cycle
(875,69) -- (870,65) -- (866,66) -- (865,67) -- (865,69) -- (859,75) -- (859,77) -- (857,79) -- (855,78) -- (852,78) -- (849,80) -- (847,84) -- (843,87) -- (843,90) -- (845,92) -- (847,92) -- (852,89) -- (854,89) -- (857,86) -- (858,84) -- (858,81) -- (860,79) -- (864,79) -- (868,76) -- (872,75) -- (874,73) -- cycle
(174,66) -- (175,69) -- (180,74) -- (181,74) -- (184,77) -- (190,78) -- (190,75) -- (191,74) -- (179,64) -- (176,64) -- cycle
(885,54) -- (882,54) -- (878,57) -- (878,59) -- (875,63) -- (876,65) -- (880,66) -- (888,58) -- (888,57) -- cycle;
\node[inner sep=0pt,font=\fontsize{10}{10}\selectfont] at (143.0,53.5) {\scalebox{6.039}{$B$}};
\node[inner sep=0pt,font=\fontsize{10}{10}\selectfont] at (905.5,27.0) {\scalebox{5.346}{$C$}};
\node[inner sep=0pt,font=\fontsize{10}{10}\selectfont] at (30.5,355.5) {\scalebox{5.643}{$A$}};
\node[inner sep=0pt,font=\fontsize{10}{10}\selectfont] at (544.0,372.0) {\scalebox{5.544}{$X$}};
\end{tikzpicture}%
}
\hspace{.8cm}\resizebox{3.4cm}{!}{%
\begin{tikzpicture}[x=1pt,y=-1pt]
\path[use as bounding box] (0,0) rectangle (629,583);
\path[fill=black,draw=none,even odd rule]
(58,525) -- (58,526) -- (61,529) -- (69,529) -- (72,527) -- (74,527) -- (78,530) -- (87,530) -- (91,527) -- (91,525) -- (88,523) -- (83,523) -- (82,522) -- (76,523) -- (74,525) -- (73,525) -- (70,522) -- (68,523) -- (60,523) -- cycle
(203,524) -- (201,522) -- (198,521) -- (197,522) -- (195,522) -- (193,524) -- (193,526) -- (197,528) -- (198,527) -- (201,527) -- (203,525) -- cycle
(188,524) -- (184,521) -- (182,522) -- (173,522) -- (172,523) -- (170,522) -- (156,522) -- (154,524) -- (154,526) -- (156,528) -- (163,528) -- (164,529) -- (167,529) -- (168,528) -- (178,528) -- (179,529) -- (181,529) -- (182,528) -- (185,528) -- (188,525) -- cycle
(97,525) -- (97,527) -- (99,529) -- (101,529) -- (102,530) -- (110,530) -- (111,529) -- (113,529) -- (115,527) -- (128,529) -- (129,528) -- (132,528) -- (133,527) -- (146,528) -- (149,526) -- (149,524) -- (147,522) -- (144,522) -- (143,521) -- (139,521) -- (136,523) -- (133,523) -- (132,522) -- (129,522) -- (128,521) -- (124,521) -- (123,522) -- (120,521) -- (118,523) -- (116,523) -- (115,524) -- (107,522) -- (106,523) -- (99,523) -- cycle
(256,322) -- (256,324) -- (258,326) -- (261,327) -- (267,332) -- (270,332) -- (272,330) -- (272,328) -- (265,322) -- (261,320) -- (259,320) -- cycle
(239,307) -- (239,309) -- (247,317) -- (249,317) -- (253,315) -- (252,312) -- (245,306) -- (242,305) -- cycle
(217,289) -- (217,291) -- (228,301) -- (231,301) -- (233,299) -- (233,297) -- (231,296) -- (223,288) -- (218,288) -- cycle
(195,272) -- (195,275) -- (198,276) -- (201,280) -- (204,281) -- (207,284) -- (209,285) -- (212,285) -- (214,283) -- (214,281) -- (211,278) -- (210,278) -- (206,274) -- (200,271) -- (198,269) -- cycle
(176,257) -- (176,259) -- (178,261) -- (178,263) -- (187,270) -- (192,270) -- (193,269) -- (193,266) -- (192,264) -- (188,260) -- (185,259) -- (180,255) -- (178,255) -- cycle
(158,238) -- (158,246) -- (159,246) -- (163,251) -- (167,253) -- (171,253) -- (172,252) -- (172,248) -- (166,243) -- (163,242) -- (159,238) -- cycle
(562,233) -- (560,231) -- (557,231) -- (552,236) -- (551,239) -- (546,245) -- (546,247) -- (547,248) -- (549,248) -- (553,245) -- (555,247) -- (555,252) -- (552,256) -- (552,259) -- (550,261) -- (550,264) -- (549,265) -- (549,272) -- (548,273) -- (548,275) -- (546,277) -- (546,280) -- (544,282) -- (544,284) -- (545,285) -- (541,288) -- (546,298) -- (541,302) -- (541,305) -- (543,308) -- (543,313) -- (542,314) -- (542,316) -- (540,318) -- (542,322) -- (541,324) -- (541,331) -- (540,332) -- (540,340) -- (539,341) -- (539,344) -- (538,345) -- (538,350) -- (537,351) -- (537,354) -- (535,358) -- (535,366) -- (534,367) -- (534,370) -- (533,371) -- (534,375) -- (532,379) -- (532,383) -- (529,387) -- (530,394) -- (529,395) -- (529,402) -- (527,406) -- (527,408) -- (528,409) -- (527,412) -- (528,414) -- (526,418) -- (526,423) -- (524,427) -- (523,440) -- (522,441) -- (522,445) -- (520,449) -- (519,459) -- (518,460) -- (518,463) -- (516,465) -- (515,482) -- (514,483) -- (514,486) -- (512,488) -- (514,492) -- (513,494) -- (512,502) -- (511,503) -- (511,506) -- (512,508) -- (510,511) -- (510,513) -- (508,515) -- (504,513) -- (488,498) -- (487,498) -- (484,495) -- (481,494) -- (478,491) -- (473,488) -- (464,479) -- (461,478) -- (453,470) -- (452,470) -- (445,464) -- (441,463) -- (438,459) -- (435,458) -- (428,451) -- (423,448) -- (412,438) -- (407,435) -- (388,418) -- (386,418) -- (378,410) -- (376,410) -- (373,407) -- (369,405) -- (363,399) -- (359,398) -- (350,389) -- (346,388) -- (340,382) -- (336,380) -- (333,376) -- (330,375) -- (317,364) -- (309,359) -- (303,352) -- (302,352) -- (293,344) -- (294,342) -- (298,342) -- (298,341) -- (294,340) -- (293,339) -- (287,338) -- (285,336) -- (283,336) -- (281,334) -- (277,333) -- (275,335) -- (274,339) -- (275,340) -- (277,340) -- (282,345) -- (282,348) -- (283,350) -- (287,354) -- (288,354) -- (290,352) -- (294,353) -- (303,361) -- (307,363) -- (312,368) -- (313,368) -- (319,375) -- (322,376) -- (328,383) -- (330,383) -- (335,388) -- (332,391) -- (333,393) -- (335,393) -- (336,392) -- (336,389) -- (338,388) -- (340,390) -- (340,393) -- (344,397) -- (346,397) -- (353,402) -- (357,403) -- (363,409) -- (366,410) -- (371,415) -- (376,417) -- (383,423) -- (387,425) -- (401,438) -- (402,438) -- (412,446) -- (413,446) -- (415,448) -- (415,449) -- (422,455) -- (427,457) -- (441,470) -- (445,471) -- (461,486) -- (464,487) -- (471,493) -- (473,493) -- (487,506) -- (492,508) -- (496,512) -- (499,513) -- (503,517) -- (501,519) -- (448,519) -- (447,518) -- (445,518) -- (444,519) -- (436,519) -- (435,520) -- (433,520) -- (432,519) -- (404,519) -- (403,520) -- (390,520) -- (389,519) -- (387,520) -- (385,519) -- (368,520) -- (367,519) -- (362,519) -- (361,520) -- (353,520) -- (352,521) -- (350,520) -- (346,520) -- (345,521) -- (338,521) -- (337,520) -- (333,520) -- (332,521) -- (329,521) -- (328,520) -- (290,520) -- (289,519) -- (280,519) -- (279,520) -- (262,520) -- (261,521) -- (259,521) -- (258,520) -- (255,520) -- (254,521) -- (252,520) -- (242,520) -- (241,521) -- (238,521) -- (237,520) -- (234,520) -- (230,515) -- (229,516) -- (224,517) -- (221,519) -- (210,520) -- (207,523) -- (207,526) -- (208,528) -- (210,530) -- (219,530) -- (222,532) -- (224,532) -- (226,534) -- (232,534) -- (233,532) -- (227,529) -- (225,527) -- (227,525) -- (233,525) -- (236,527) -- (238,527) -- (239,526) -- (245,526) -- (246,527) -- (256,527) -- (257,528) -- (259,527) -- (265,527) -- (266,528) -- (272,528) -- (273,527) -- (275,527) -- (276,528) -- (278,528) -- (280,526) -- (284,526) -- (285,527) -- (288,527) -- (289,526) -- (299,526) -- (300,527) -- (304,527) -- (305,526) -- (313,526) -- (314,527) -- (326,527) -- (327,526) -- (329,527) -- (338,527) -- (339,526) -- (346,526) -- (347,527) -- (349,527) -- (350,526) -- (361,526) -- (362,525) -- (369,525) -- (370,526) -- (373,526) -- (374,525) -- (380,525) -- (381,526) -- (391,526) -- (392,525) -- (428,525) -- (429,524) -- (431,525) -- (439,525) -- (440,524) -- (449,524) -- (450,525) -- (452,525) -- (453,524) -- (487,524) -- (488,523) -- (501,523) -- (502,522) -- (510,523) -- (510,521) -- (511,520) -- (513,520) -- (512,516) -- (514,512) -- (514,509) -- (516,506) -- (517,495) -- (518,494) -- (518,488) -- (519,487) -- (519,482) -- (520,481) -- (520,477) -- (522,473) -- (522,466) -- (524,463) -- (524,460) -- (525,459) -- (525,451) -- (527,448) -- (527,442) -- (528,441) -- (528,437) -- (530,433) -- (530,428) -- (531,427) -- (532,420) -- (534,417) -- (534,412) -- (535,410) -- (533,408) -- (533,407) -- (535,404) -- (535,396) -- (537,393) -- (537,390) -- (538,389) -- (537,387) -- (537,384) -- (538,383) -- (538,379) -- (539,378) -- (539,373) -- (541,370) -- (541,367) -- (542,366) -- (542,356) -- (544,354) -- (544,349) -- (545,348) -- (545,344) -- (546,343) -- (546,337) -- (547,336) -- (547,330) -- (548,329) -- (548,324) -- (547,323) -- (547,321) -- (549,318) -- (549,315) -- (550,314) -- (550,309) -- (551,308) -- (551,302) -- (553,299) -- (552,297) -- (552,295) -- (553,294) -- (553,289) -- (555,285) -- (555,282) -- (557,279) -- (557,274) -- (556,272) -- (558,269) -- (558,266) -- (560,264) -- (560,259) -- (561,258) -- (561,257) -- (559,255) -- (559,252) -- (561,250) -- (566,250) -- (567,249) -- (567,246) -- (564,243) -- (563,239) -- (562,238) -- cycle
(288,345) -- (290,347) -- (289,349) -- (288,349) -- (286,347) -- cycle
(563,211) -- (561,211) -- (559,213) -- (559,215) -- (557,218) -- (557,221) -- (560,224) -- (562,224) -- (564,222) -- (565,218) -- (566,217) -- (566,214) -- cycle
(565,190) -- (564,191) -- (562,191) -- (559,194) -- (558,200) -- (560,204) -- (565,204) -- (567,205) -- (570,202) -- (569,199) -- (570,198) -- (571,193) -- (569,191) -- cycle
(572,168) -- (568,168) -- (563,175) -- (563,178) -- (564,179) -- (563,184) -- (565,182) -- (571,182) -- (573,180) -- (573,174) -- (574,173) -- (574,170) -- cycle
(575,146) -- (572,146) -- (570,148) -- (570,151) -- (569,152) -- (569,155) -- (568,156) -- (568,157) -- (570,159) -- (574,159) -- (575,158) -- (575,155) -- (577,152) -- (577,150) -- cycle
(577,120) -- (575,121) -- (575,126) -- (574,127) -- (574,131) -- (573,132) -- (575,136) -- (577,136) -- (580,133) -- (580,131) -- (581,130) -- (581,123) -- cycle
(583,99) -- (581,99) -- (577,103) -- (577,105) -- (578,106) -- (577,114) -- (578,115) -- (583,115) -- (585,111) -- (585,104) -- (586,103) -- (586,101) -- cycle
(584,79) -- (583,80) -- (583,84) -- (580,88) -- (580,92) -- (582,94) -- (587,94) -- (590,89) -- (590,82) -- (587,79) -- cycle
(586,61) -- (585,62) -- (585,70) -- (587,72) -- (590,72) -- (591,71) -- (591,68) -- (593,66) -- (593,63) -- (591,61) -- cycle;
\node[inner sep=0pt,font=\fontsize{10}{10}\selectfont] at (590.5,32.0) {\scalebox{5.742}{$C$}};
\node[inner sep=0pt,font=\fontsize{10}{10}\selectfont] at (129.5,225.5) {\scalebox{5.445}{$A$}};
\node[inner sep=0pt,font=\fontsize{10}{10}\selectfont] at (29.0,525.0) {\scalebox{5.742}{$B$}};
\node[inner sep=0pt,font=\fontsize{10}{10}\selectfont] at (537.5,548.0) {\scalebox{6.534}{$X$}};
\end{tikzpicture}%
}
\end{center}
%
\SourceLeaf{319}
\begin{center}\resizebox{4.6cm}{!}{%
\begin{tikzpicture}[x=1pt,y=-1pt]
\path[use as bounding box] (0,0) rectangle (914,561);
\path[fill=black,draw=none,even odd rule]
(73,524) -- (73,525) -- (76,528) -- (90,528) -- (93,525) -- (92,523) -- (90,523) -- (89,522) -- (87,522) -- (86,523) -- (84,522) -- (77,522) -- cycle
(66,525) -- (62,522) -- (56,523) -- (56,528) -- (57,527) -- (61,528) -- (62,527) -- (64,527) -- cycle
(184,522) -- (184,524) -- (188,526) -- (194,526) -- (196,524) -- (196,522) -- (193,522) -- (192,521) -- (185,521) -- cycle
(156,523) -- (156,524) -- (158,526) -- (164,526) -- (165,527) -- (171,526) -- (173,524) -- (173,522) -- (172,521) -- (159,521) -- cycle
(132,525) -- (135,527) -- (139,527) -- (140,526) -- (145,526) -- (147,524) -- (147,523) -- (143,521) -- (137,521) -- (136,522) -- (134,522) -- (132,524) -- cycle
(104,524) -- (104,525) -- (108,528) -- (114,528) -- (115,527) -- (120,527) -- (121,526) -- (121,523) -- (120,522) -- (112,522) -- (110,521) -- cycle
(771,311) -- (769,309) -- (764,311) -- (757,317) -- (757,319) -- (758,320) -- (762,320) -- (765,317) -- (769,316) -- (770,315) -- cycle
(792,294) -- (791,293) -- (787,293) -- (785,294) -- (777,301) -- (777,303) -- (778,304) -- (781,304) -- (783,303) -- (786,300) -- (787,300) -- (792,295) -- cycle
(813,275) -- (812,274) -- (810,274) -- (803,278) -- (797,284) -- (797,286) -- (800,288) -- (805,286) -- (811,281) -- (811,280) -- (813,278) -- cycle
(835,258) -- (834,257) -- (830,257) -- (828,258) -- (819,266) -- (818,266) -- (817,268) -- (819,270) -- (823,270) -- (825,269) -- (835,260) -- cycle
(855,241) -- (850,240) -- (848,242) -- (847,242) -- (842,247) -- (842,250) -- (843,251) -- (847,251) -- (849,249) -- (850,249) -- (854,245) -- cycle
(875,225) -- (869,225) -- (867,226) -- (861,232) -- (861,233) -- (863,235) -- (869,233) -- (874,228) -- cycle
(755,327) -- (753,325) -- (744,325) -- (739,328) -- (728,329) -- (727,330) -- (727,332) -- (728,333) -- (732,333) -- (737,331) -- (739,333) -- (735,337) -- (734,337) -- (724,345) -- (721,346) -- (712,355) -- (711,355) -- (707,359) -- (703,361) -- (698,366) -- (693,369) -- (681,381) -- (680,381) -- (663,396) -- (658,399) -- (651,406) -- (646,409) -- (636,418) -- (633,419) -- (625,427) -- (624,427) -- (619,432) -- (618,432) -- (609,441) -- (607,441) -- (602,446) -- (597,449) -- (590,456) -- (585,459) -- (579,465) -- (578,465) -- (574,469) -- (569,472) -- (561,480) -- (560,480) -- (556,484) -- (553,485) -- (533,503) -- (532,503) -- (521,512) -- (519,512) -- (518,511) -- (518,508) -- (519,507) -- (519,503) -- (521,500) -- (521,494) -- (523,492) -- (523,485) -- (524,484) -- (524,480) -- (525,479) -- (526,474) -- (528,471) -- (528,462) -- (529,461) -- (529,456) -- (530,455) -- (530,452) -- (531,451) -- (532,438) -- (533,437) -- (532,434) -- (534,431) -- (534,427) -- (535,426) -- (535,419) -- (536,418) -- (536,415) -- (537,414) -- (537,409) -- (539,406) -- (539,398) -- (541,395) -- (542,384) -- (544,381) -- (544,374) -- (545,373) -- (545,369) -- (547,367) -- (547,363) -- (548,362) -- (547,357) -- (549,353) -- (549,348) -- (551,344) -- (553,329) -- (555,326) -- (554,320) -- (556,317) -- (556,309) -- (558,306) -- (558,304) -- (559,303) -- (559,299) -- (560,298) -- (560,291) -- (561,290) -- (560,289) -- (561,288) -- (560,287) -- (561,286) -- (561,282) -- (564,278) -- (563,276) -- (563,273) -- (564,272) -- (564,265) -- (565,264) -- (565,260) -- (566,259) -- (565,255) -- (568,251) -- (569,236) -- (571,234) -- (572,234) -- (573,235) -- (575,242) -- (578,244) -- (580,244) -- (582,242) -- (581,240) -- (581,237) -- (578,234) -- (577,230) -- (575,227) -- (574,218) -- (573,216) -- (571,214) -- (570,214) -- (566,218) -- (564,222) -- (560,226) -- (558,226) -- (549,236) -- (549,237) -- (551,239) -- (555,240) -- (557,239) -- (560,234) -- (563,231) -- (564,231) -- (566,233) -- (566,237) -- (564,240) -- (564,244) -- (563,245) -- (563,247) -- (560,251) -- (560,259) -- (558,261) -- (558,263) -- (557,264) -- (557,268) -- (556,269) -- (556,277) -- (555,278) -- (555,282) -- (554,283) -- (555,285) -- (555,288) -- (553,291) -- (553,296) -- (549,300) -- (551,302) -- (551,303) -- (548,306) -- (548,309) -- (550,311) -- (550,314) -- (549,315) -- (549,319) -- (548,320) -- (548,323) -- (547,324) -- (546,336) -- (545,337) -- (545,340) -- (544,341) -- (544,348) -- (543,349) -- (542,359) -- (541,360) -- (540,373) -- (539,374) -- (539,376) -- (537,378) -- (538,384) -- (536,387) -- (536,392) -- (535,393) -- (535,397) -- (534,398) -- (534,403) -- (533,404) -- (532,415) -- (531,416) -- (531,419) -- (530,420) -- (530,423) -- (529,424) -- (529,428) -- (528,429) -- (528,439) -- (527,440) -- (527,445) -- (526,446) -- (526,449) -- (525,450) -- (524,458) -- (521,463) -- (522,465) -- (522,468) -- (521,469) -- (521,476) -- (520,477) -- (520,481) -- (519,482) -- (519,485) -- (518,486) -- (517,501) -- (516,502) -- (516,506) -- (514,510) -- (514,516) -- (511,519) -- (501,519) -- (500,518) -- (498,518) -- (497,519) -- (489,519) -- (488,520) -- (483,520) -- (482,519) -- (472,519) -- (471,518) -- (469,518) -- (468,519) -- (464,519) -- (463,520) -- (453,520) -- (452,519) -- (448,519) -- (447,520) -- (437,520) -- (436,519) -- (430,519) -- (429,520) -- (427,519) -- (424,519) -- (423,520) -- (396,520) -- (395,519) -- (392,519) -- (391,520) -- (318,520) -- (317,521) -- (315,521) -- (314,520) -- (312,520) -- (311,521) -- (291,521) -- (290,520) -- (288,520) -- (287,521) -- (281,521) -- (280,520) -- (278,521) -- (271,521) -- (270,520) -- (267,520) -- (265,521) -- (264,520) -- (254,520) -- (253,521) -- (234,521) -- (233,522) -- (227,522) -- (224,519) -- (231,512) -- (230,510) -- (220,517) -- (217,517) -- (214,519) -- (212,519) -- (207,522) -- (207,524) -- (209,526) -- (211,527) -- (214,527) -- (216,529) -- (221,531) -- (224,534) -- (226,534) -- (227,533) -- (227,531) -- (224,528) -- (224,526) -- (225,525) -- (242,525) -- (243,526) -- (247,526) -- (248,525) -- (252,525) -- (253,526) -- (255,526) -- (256,525) -- (261,525) -- (262,526) -- (267,526) -- (268,525) -- (271,525) -- (272,526) -- (274,526) -- (275,525) -- (277,525) -- (278,526) -- (284,526) -- (285,525) -- (292,525) -- (293,526) -- (296,526) -- (297,525) -- (301,525) -- (302,526) -- (303,525) -- (328,525) -- (329,526) -- (331,525) -- (355,525) -- (357,526) -- (358,525) -- (376,525) -- (377,524) -- (385,524) -- (386,525) -- (388,525) -- (389,524) -- (407,525) -- (408,524) -- (417,524) -- (418,523) -- (420,523) -- (421,524) -- (438,524) -- (439,523) -- (441,523) -- (442,524) -- (453,524) -- (454,523) -- (460,523) -- (462,525) -- (464,524) -- (475,524) -- (476,523) -- (481,523) -- (482,524) -- (486,524) -- (487,523) -- (489,524) -- (502,524) -- (503,523) -- (505,524) -- (514,524) -- (516,523) -- (527,511) -- (533,508) -- (552,491) -- (554,491) -- (563,483) -- (566,482) -- (574,474) -- (576,474) -- (581,469) -- (586,466) -- (588,464) -- (588,463) -- (595,457) -- (597,457) -- (599,455) -- (603,453) -- (613,443) -- (616,442) -- (617,440) -- (629,429) -- (632,428) -- (639,421) -- (644,419) -- (650,413) -- (651,413) -- (662,403) -- (663,403) -- (684,383) -- (692,378) -- (697,373) -- (700,372) -- (710,362) -- (720,356) -- (724,351) -- (728,349) -- (732,344) -- (733,344) -- (737,340) -- (738,340) -- (743,335) -- (745,336) -- (745,338) -- (744,339) -- (744,342) -- (742,346) -- (742,350) -- (743,351) -- (745,351) -- (747,349) -- (747,345) -- (748,344) -- (749,338) -- (752,334) -- (752,332) -- (755,329) -- cycle
(579,171) -- (577,171) -- (575,173) -- (574,175) -- (574,180) -- (572,184) -- (572,193) -- (569,199) -- (569,202) -- (572,205) -- (574,205) -- (576,202) -- (577,191) -- (579,188) -- (579,185) -- (580,184) -- (580,180) -- (579,178) -- (581,175) -- (581,173) -- cycle
(583,142) -- (581,143) -- (581,145) -- (579,148) -- (579,152) -- (577,154) -- (577,158) -- (580,162) -- (581,162) -- (584,159) -- (584,155) -- (585,154) -- (585,151) -- (587,147) -- (587,145) -- cycle
(589,114) -- (587,114) -- (584,117) -- (585,120) -- (583,124) -- (583,128) -- (585,130) -- (588,130) -- (589,129) -- (589,126) -- (591,123) -- (591,116) -- cycle
(595,86) -- (592,86) -- (590,88) -- (590,91) -- (588,94) -- (588,98) -- (586,100) -- (586,101) -- (589,104) -- (592,104) -- (595,101) -- (595,98) -- (596,97) -- (596,92) -- (597,91) -- (596,87) -- cycle
(599,61) -- (597,61) -- (594,65) -- (594,68) -- (592,71) -- (592,74) -- (594,76) -- (598,76) -- (600,73) -- (600,66) -- (601,65) -- (601,63) -- cycle;
\node[inner sep=0pt,font=\fontsize{10}{10}\selectfont] at (599.5,29.5) {\scalebox{5.841}{$B$}};
\node[inner sep=0pt,font=\fontsize{10}{10}\selectfont] at (887.0,201.5) {\scalebox{5.841}{$A$}};
\node[inner sep=0pt,font=\fontsize{10}{10}\selectfont] at (28.0,526.0) {\scalebox{6.138}{$C$}};
\node[inner sep=0pt,font=\fontsize{10}{10}\selectfont] at (564.5,528.0) {\scalebox{6.534}{$X$}};
\end{tikzpicture}%
}
\hspace{.8cm}\resizebox{3.4cm}{!}{%
\begin{tikzpicture}[x=1pt,y=-1pt]
\path[use as bounding box] (0,0) rectangle (631,580);
\path[fill=black,draw=none,even odd rule]
(101,520) -- (101,521) -- (105,524) -- (113,524) -- (116,522) -- (116,521) -- (113,518) -- (103,518) -- cycle
(81,520) -- (83,524) -- (95,524) -- (98,522) -- (98,520) -- (96,518) -- (83,518) -- cycle
(60,522) -- (63,524) -- (68,524) -- (69,525) -- (71,525) -- (73,524) -- (76,521) -- (76,520) -- (74,518) -- (64,518) -- (62,519) -- (60,521) -- cycle
(167,519) -- (167,523) -- (171,526) -- (173,525) -- (177,525) -- (178,524) -- (183,524) -- (184,523) -- (186,523) -- (188,521) -- (189,521) -- (191,523) -- (194,523) -- (196,525) -- (198,524) -- (201,524) -- (202,523) -- (202,522) -- (198,518) -- (193,518) -- (189,521) -- (186,520) -- (182,517) -- (180,517) -- (179,518) -- (168,518) -- cycle
(143,520) -- (143,521) -- (145,523) -- (149,525) -- (152,525) -- (153,524) -- (157,525) -- (161,522) -- (161,519) -- (158,517) -- (152,517) -- (151,518) -- (147,518) -- cycle
(122,520) -- (122,522) -- (125,524) -- (129,524) -- (130,525) -- (132,524) -- (135,524) -- (138,522) -- (138,520) -- (136,518) -- (132,518) -- (131,517) -- (126,517) -- cycle
(354,250) -- (354,251) -- (358,255) -- (358,257) -- (361,261) -- (367,261) -- (369,259) -- (368,256) -- (365,253) -- (362,248) -- (359,247) -- cycle
(573,233) -- (571,231) -- (568,231) -- (566,234) -- (565,238) -- (554,250) -- (555,252) -- (557,252) -- (560,250) -- (564,249) -- (566,251) -- (566,254) -- (564,257) -- (564,263) -- (562,267) -- (562,271) -- (559,274) -- (560,276) -- (560,280) -- (559,281) -- (559,286) -- (558,287) -- (558,291) -- (557,292) -- (556,308) -- (554,311) -- (553,322) -- (550,328) -- (551,330) -- (551,336) -- (550,337) -- (550,341) -- (549,342) -- (549,345) -- (548,346) -- (548,349) -- (547,350) -- (548,356) -- (545,363) -- (544,375) -- (540,383) -- (540,393) -- (539,394) -- (540,400) -- (538,403) -- (537,416) -- (536,417) -- (536,420) -- (535,421) -- (535,424) -- (534,425) -- (534,431) -- (533,432) -- (533,441) -- (531,444) -- (531,449) -- (530,450) -- (531,453) -- (529,455) -- (529,461) -- (526,467) -- (527,471) -- (526,472) -- (525,480) -- (524,481) -- (524,485) -- (523,486) -- (523,490) -- (525,492) -- (525,495) -- (520,501) -- (520,507) -- (521,509) -- (518,513) -- (517,513) -- (515,511) -- (515,509) -- (512,506) -- (511,503) -- (509,501) -- (509,499) -- (504,493) -- (504,491) -- (503,490) -- (503,488) -- (501,486) -- (501,484) -- (496,478) -- (496,476) -- (491,469) -- (491,467) -- (485,460) -- (482,452) -- (477,447) -- (477,443) -- (473,439) -- (471,434) -- (467,429) -- (467,426) -- (465,424) -- (464,421) -- (461,418) -- (461,416) -- (458,413) -- (458,410) -- (456,408) -- (453,401) -- (447,395) -- (445,387) -- (441,384) -- (441,383) -- (439,381) -- (439,378) -- (436,375) -- (435,372) -- (428,363) -- (427,359) -- (421,353) -- (421,349) -- (419,347) -- (418,343) -- (414,338) -- (413,334) -- (411,332) -- (410,329) -- (405,324) -- (405,320) -- (398,312) -- (398,309) -- (395,305) -- (392,298) -- (388,293) -- (387,290) -- (383,287) -- (382,284) -- (380,282) -- (380,281) -- (382,279) -- (388,284) -- (391,284) -- (392,283) -- (392,281) -- (387,277) -- (384,277) -- (379,272) -- (378,272) -- (372,267) -- (370,267) -- (368,269) -- (367,275) -- (365,277) -- (365,280) -- (371,288) -- (370,290) -- (371,294) -- (373,296) -- (375,296) -- (376,295) -- (376,291) -- (375,290) -- (375,286) -- (374,285) -- (375,284) -- (376,284) -- (379,287) -- (380,290) -- (384,295) -- (384,300) -- (387,303) -- (388,306) -- (393,310) -- (392,311) -- (393,314) -- (395,316) -- (396,319) -- (401,324) -- (401,326) -- (405,331) -- (405,333) -- (410,337) -- (410,339) -- (412,341) -- (412,344) -- (415,347) -- (415,352) -- (419,356) -- (420,359) -- (424,364) -- (427,372) -- (430,375) -- (430,377) -- (436,384) -- (436,387) -- (442,395) -- (443,398) -- (446,401) -- (449,408) -- (452,412) -- (452,414) -- (454,416) -- (454,418) -- (459,423) -- (461,428) -- (464,431) -- (464,433) -- (466,437) -- (471,443) -- (472,448) -- (474,450) -- (475,453) -- (479,457) -- (478,460) -- (483,465) -- (484,468) -- (489,474) -- (490,477) -- (492,479) -- (492,481) -- (501,494) -- (501,496) -- (503,498) -- (505,503) -- (511,510) -- (511,512) -- (514,516) -- (514,517) -- (512,519) -- (426,519) -- (425,518) -- (418,518) -- (417,519) -- (413,519) -- (412,518) -- (409,518) -- (408,517) -- (406,517) -- (405,518) -- (403,517) -- (400,517) -- (399,518) -- (394,518) -- (393,517) -- (379,517) -- (378,518) -- (365,518) -- (364,517) -- (361,517) -- (360,518) -- (353,518) -- (352,517) -- (350,517) -- (349,518) -- (347,518) -- (346,517) -- (334,517) -- (332,518) -- (330,516) -- (328,517) -- (324,517) -- (323,516) -- (316,516) -- (315,517) -- (296,517) -- (295,518) -- (287,518) -- (286,519) -- (283,519) -- (282,518) -- (280,518) -- (279,519) -- (277,519) -- (276,518) -- (270,518) -- (269,519) -- (267,518) -- (260,518) -- (259,519) -- (230,519) -- (228,517) -- (228,516) -- (230,514) -- (236,511) -- (236,509) -- (235,508) -- (233,508) -- (232,509) -- (229,509) -- (219,516) -- (214,516) -- (212,517) -- (210,520) -- (210,522) -- (211,523) -- (213,523) -- (217,527) -- (221,529) -- (223,529) -- (226,532) -- (229,533) -- (232,537) -- (237,538) -- (239,540) -- (243,540) -- (244,538) -- (243,537) -- (241,537) -- (238,534) -- (235,533) -- (232,530) -- (226,527) -- (226,526) -- (229,523) -- (232,523) -- (233,522) -- (235,522) -- (239,524) -- (241,524) -- (242,523) -- (247,523) -- (248,524) -- (255,524) -- (256,523) -- (270,523) -- (271,524) -- (273,524) -- (274,523) -- (279,523) -- (280,524) -- (282,524) -- (283,523) -- (309,523) -- (310,524) -- (312,523) -- (321,523) -- (322,524) -- (326,524) -- (327,525) -- (329,524) -- (332,524) -- (333,523) -- (335,523) -- (336,524) -- (340,524) -- (341,523) -- (346,523) -- (347,524) -- (349,524) -- (350,523) -- (367,523) -- (368,524) -- (371,524) -- (372,523) -- (378,523) -- (379,524) -- (382,524) -- (383,523) -- (399,523) -- (400,524) -- (404,524) -- (405,523) -- (412,523) -- (413,524) -- (419,524) -- (420,523) -- (428,523) -- (429,524) -- (433,524) -- (434,523) -- (444,523) -- (445,524) -- (447,524) -- (448,523) -- (474,523) -- (475,524) -- (490,524) -- (491,523) -- (493,523) -- (494,524) -- (510,524) -- (511,522) -- (517,522) -- (519,520) -- (519,518) -- (524,514) -- (524,509) -- (525,508) -- (525,505) -- (526,504) -- (525,499) -- (526,498) -- (527,493) -- (529,490) -- (529,485) -- (530,484) -- (530,481) -- (532,479) -- (532,475) -- (534,473) -- (533,469) -- (534,468) -- (534,463) -- (536,460) -- (536,449) -- (535,448) -- (535,446) -- (536,445) -- (536,443) -- (538,441) -- (538,435) -- (540,432) -- (540,429) -- (541,428) -- (541,417) -- (543,414) -- (543,410) -- (544,409) -- (544,406) -- (545,405) -- (545,400) -- (547,397) -- (547,393) -- (548,392) -- (548,383) -- (549,382) -- (549,376) -- (547,374) -- (547,373) -- (550,371) -- (550,368) -- (551,367) -- (551,362) -- (552,361) -- (552,352) -- (554,350) -- (554,345) -- (553,344) -- (553,342) -- (555,340) -- (555,338) -- (556,337) -- (556,330) -- (559,324) -- (559,319) -- (558,318) -- (558,316) -- (559,315) -- (559,311) -- (560,310) -- (560,305) -- (562,302) -- (562,296) -- (564,294) -- (564,291) -- (565,290) -- (564,288) -- (565,286) -- (564,283) -- (566,281) -- (565,276) -- (566,275) -- (566,272) -- (569,267) -- (569,261) -- (567,259) -- (567,258) -- (569,256) -- (569,253) -- (571,251) -- (571,249) -- (572,248) -- (573,248) -- (578,254) -- (579,254) -- (581,252) -- (576,243) -- (576,240) -- (574,238) -- cycle
(224,520) -- (226,522) -- (226,523) -- (223,526) -- (221,526) -- (220,525) -- (220,523) -- (223,520) -- cycle
(346,226) -- (344,228) -- (344,231) -- (346,234) -- (346,236) -- (352,242) -- (359,241) -- (360,240) -- (357,237) -- (357,236) -- (353,233) -- (352,230) -- (348,226) -- cycle
(575,208) -- (574,208) -- (571,211) -- (571,213) -- (569,215) -- (569,221) -- (571,224) -- (577,223) -- (578,222) -- (577,218) -- (578,217) -- (578,213) -- (579,211) -- cycle
(335,207) -- (333,209) -- (334,215) -- (339,221) -- (343,221) -- (345,219) -- (345,215) -- (338,207) -- cycle
(579,186) -- (577,186) -- (575,188) -- (575,195) -- (574,196) -- (574,200) -- (575,201) -- (579,201) -- (581,199) -- (582,189) -- cycle
(321,186) -- (320,187) -- (321,195) -- (328,202) -- (331,201) -- (333,199) -- (333,197) -- (331,195) -- (330,192) -- (325,186) -- cycle
(308,163) -- (307,164) -- (307,166) -- (312,176) -- (316,180) -- (319,180) -- (321,178) -- (320,175) -- (317,172) -- (314,166) -- (311,163) -- cycle
(585,160) -- (581,160) -- (580,161) -- (580,164) -- (577,170) -- (578,172) -- (578,175) -- (579,176) -- (581,176) -- (585,173) -- (585,171) -- (586,170) -- (586,164) -- (585,163) -- (586,161) -- cycle
(586,139) -- (585,139) -- (583,141) -- (583,143) -- (581,146) -- (582,147) -- (582,151) -- (583,152) -- (587,152) -- (589,150) -- (589,144) -- (590,142) -- cycle
(289,137) -- (296,146) -- (296,149) -- (299,152) -- (301,157) -- (303,159) -- (306,159) -- (308,157) -- (308,155) -- (303,149) -- (302,145) -- (298,141) -- (296,137) -- cycle
(592,114) -- (589,115) -- (587,117) -- (587,121) -- (586,122) -- (586,129) -- (587,130) -- (591,130) -- (594,127) -- (594,122) -- (596,118) -- cycle
(596,87) -- (595,87) -- (593,89) -- (593,93) -- (591,97) -- (591,103) -- (592,104) -- (596,104) -- (598,101) -- (598,94) -- (599,93) -- (599,89) -- cycle
(603,62) -- (601,62) -- (599,64) -- (599,67) -- (596,70) -- (596,73) -- (592,75) -- (592,76) -- (595,79) -- (599,79) -- (603,76) -- (603,67) -- (604,66) -- (604,63) -- cycle;
\node[inner sep=0pt,font=\fontsize{10}{10}\selectfont] at (600.5,30.5) {\scalebox{6.039}{$B$}};
\node[inner sep=0pt,font=\fontsize{10}{10}\selectfont] at (278.0,107.5) {\scalebox{5.841}{$A$}};
\node[inner sep=0pt,font=\fontsize{10}{10}\selectfont] at (29.0,524.0) {\scalebox{6.336}{$C$}};
\node[inner sep=0pt,font=\fontsize{10}{10}\selectfont] at (544.0,551.5) {\scalebox{5.643}{$X$}};
\end{tikzpicture}%
}
\end{center}

die Fälle nicht gehörig unterschieden würden, oder, mit andern Worten, weil
man sonst nicht Ein Problem, sondern auf einmal vier verschiedene aufgeben
würde.

\textls[120]{Diess wohlverstanden}, sind nun alle bisherigen Behandlungen des Problems in
so fern unvollständig, als sie keine allgemeine Regel angeben, wonach man
vorgegebenen Datis gleich ansehen kann, ob sie zusammen bestehen können oder
nicht. Auch ohne bloss vom theoretischen Standpunkt auszugehen, scheint eine
Ergänzung dieses Mangels praktisch wünschenswerth, weil ja Data, die sich auf
einen wirklichen Fall in concreto beziehen sollen, wenn auch nicht durch
Neckerei, doch durch grobe Schreibfehler und Irrthum entstellt sein können,
deren Dasein sogleich in vielen Fällen erkennen zu können, doch auch dem
blossen Praktiker werth sein muss.

Allgemeine Regel verstehe ich so, dass man nicht für jeden speciellen Fall
eine eigene rohe Zeichnung machen soll. Denn wenn man das thun will, so wird
es jedesmal eben nicht schwer sein, über Möglichkeit oder Unmöglichkeit zu
urtheilen.

Das eben wird aber verlangt, dass Alle Fälle in Einer allgemeinen Regel
befasst werden sollen.

Sie schienen zu glauben, dass solche physische Unmöglichkeit oder
Unvereinbarkeit eben nur einzelne, vielleicht seltene, Fälle trifft, das ist
aber nicht so, sondern (abstrahirt von einer gewussten richtigen Beziehung
auf einen wirklichen Fall in concreto) gibt es dreimal so viel unmögliche als
mögliche Combinationen.

Das eben ist nun in dem allgemeinen Satz enthalten, welcher auf dem Zettel
steht. Nemlich:

\begin{quote}»Wenn die vorgeschriebenen Data möglich sind, so werden sie durch Abänderung
einer der drei Richtungen ($XA$, $XB$, $XC$) in die entgegengesetzte (um
$180^0$ verschieden) sogleich unmöglich«.\end{quote}
%
\SourceLeaf{320}
Derselbe Satz ist übrigens auch umgekehrt wahr, die unverträglichen Data
werden verträglich, wenn man Eine der drei Richtungen in die entgegengesetzte
abändert, nur ist hier die abzuändernde nicht mehr willkürlich.

Wenn also z.\,B. unter den obigen vier Zeichnungen die erste (für ein
bestimmtes gegebenes Dreieck $ABC$) möglich ist, nemlich

\begin{center}\resizebox{6.8cm}{!}{%
\begin{tikzpicture}[x=1pt,y=-1pt]
\path[use as bounding box] (0,0) rectangle (1198,529);
\path[fill=black,draw=none,even odd rule]
(84,497) -- (87,499) -- (97,499) -- (98,498) -- (118,498) -- (119,499) -- (132,499) -- (135,501) -- (139,500) -- (140,499) -- (149,499) -- (150,498) -- (152,498) -- (153,499) -- (158,499) -- (159,500) -- (167,500) -- (168,499) -- (172,499) -- (173,500) -- (180,500) -- (181,501) -- (183,501) -- (184,500) -- (201,500) -- (202,499) -- (204,500) -- (214,500) -- (215,499) -- (226,499) -- (227,500) -- (233,500) -- (234,499) -- (236,499) -- (237,500) -- (243,500) -- (244,499) -- (249,499) -- (250,500) -- (252,500) -- (253,499) -- (256,499) -- (257,500) -- (286,500) -- (287,501) -- (293,501) -- (294,500) -- (332,500) -- (333,501) -- (336,501) -- (337,500) -- (349,500) -- (350,501) -- (363,500) -- (364,501) -- (392,501) -- (393,500) -- (413,500) -- (414,499) -- (420,499) -- (421,500) -- (471,500) -- (472,501) -- (477,501) -- (478,500) -- (482,500) -- (483,501) -- (485,500) -- (492,500) -- (493,501) -- (505,501) -- (506,500) -- (509,500) -- (510,501) -- (516,501) -- (517,502) -- (525,502) -- (526,501) -- (534,501) -- (535,502) -- (537,502) -- (538,501) -- (572,501) -- (573,502) -- (579,502) -- (585,499) -- (587,499) -- (588,500) -- (594,500) -- (595,501) -- (597,500) -- (621,500) -- (622,501) -- (641,501) -- (642,502) -- (644,501) -- (649,501) -- (650,502) -- (652,502) -- (653,501) -- (656,501) -- (657,502) -- (660,502) -- (661,501) -- (667,501) -- (668,502) -- (672,501) -- (673,502) -- (678,502) -- (683,498) -- (689,495) -- (695,489) -- (698,488) -- (703,483) -- (708,480) -- (711,476) -- (719,471) -- (722,467) -- (723,467) -- (727,463) -- (728,463) -- (733,458) -- (734,458) -- (743,450) -- (747,448) -- (752,443) -- (755,442) -- (769,429) -- (771,429) -- (780,419) -- (781,419) -- (795,406) -- (796,406) -- (803,400) -- (804,400) -- (807,396) -- (808,396) -- (812,392) -- (813,392) -- (825,381) -- (830,378) -- (836,372) -- (841,369) -- (861,351) -- (869,346) -- (883,333) -- (888,330) -- (895,323) -- (898,322) -- (903,317) -- (906,316) -- (911,310) -- (919,305) -- (920,303) -- (923,300) -- (925,300) -- (930,294) -- (934,292) -- (941,285) -- (946,282) -- (951,277) -- (954,276) -- (969,262) -- (970,262) -- (974,258) -- (975,258) -- (985,249) -- (989,247) -- (997,239) -- (1000,238) -- (1006,232) -- (1007,232) -- (1010,228) -- (1011,228) -- (1016,223) -- (1017,223) -- (1021,219) -- (1022,219) -- (1026,215) -- (1027,215) -- (1036,206) -- (1041,203) -- (1045,199) -- (1049,197) -- (1071,177) -- (1072,177) -- (1075,174) -- (1077,174) -- (1084,166) -- (1089,164) -- (1095,157) -- (1096,157) -- (1100,153) -- (1106,150) -- (1113,143) -- (1116,142) -- (1131,129) -- (1133,129) -- (1135,127) -- (1136,127) -- (1139,123) -- (1145,123) -- (1148,120) -- (1148,118) -- (1145,118) -- (1143,116) -- (1140,116) -- (1139,117) -- (1137,116) -- (1132,120) -- (1129,121) -- (1123,127) -- (1120,128) -- (1114,134) -- (1113,134) -- (1109,138) -- (1105,140) -- (1090,154) -- (1087,155) -- (1084,159) -- (1083,159) -- (1079,163) -- (1076,164) -- (1072,169) -- (1069,170) -- (1066,173) -- (1065,173) -- (1062,177) -- (1057,180) -- (1049,188) -- (1047,188) -- (1027,206) -- (1024,207) -- (1018,213) -- (1014,215) -- (1007,222) -- (1006,222) -- (1003,225) -- (998,228) -- (998,229) -- (990,237) -- (989,237) -- (985,241) -- (980,244) -- (971,252) -- (967,254) -- (963,259) -- (959,261) -- (946,273) -- (945,273) -- (930,286) -- (929,286) -- (925,290) -- (924,290) -- (920,294) -- (919,294) -- (915,298) -- (914,298) -- (907,305) -- (899,310) -- (895,315) -- (892,316) -- (875,332) -- (874,332) -- (869,337) -- (864,340) -- (858,346) -- (857,346) -- (853,350) -- (852,350) -- (838,363) -- (837,363) -- (834,366) -- (833,366) -- (830,369) -- (825,372) -- (819,379) -- (816,380) -- (810,386) -- (805,389) -- (799,395) -- (798,395) -- (794,399) -- (793,399) -- (773,417) -- (772,417) -- (767,422) -- (763,424) -- (759,429) -- (757,429) -- (755,427) -- (756,426) -- (755,425) -- (751,429) -- (747,430) -- (747,433) -- (748,433) -- (750,431) -- (752,433) -- (750,435) -- (750,436) -- (740,446) -- (735,449) -- (718,464) -- (708,470) -- (703,476) -- (701,476) -- (697,480) -- (694,481) -- (681,493) -- (677,494) -- (673,490) -- (671,490) -- (640,461) -- (639,461) -- (635,457) -- (634,457) -- (628,451) -- (627,449) -- (625,449) -- (618,442) -- (617,442) -- (605,429) -- (594,422) -- (588,416) -- (580,411) -- (570,401) -- (569,401) -- (564,395) -- (561,394) -- (553,385) -- (550,384) -- (545,379) -- (544,379) -- (538,372) -- (537,372) -- (533,368) -- (532,368) -- (524,361) -- (523,361) -- (514,352) -- (513,352) -- (508,347) -- (507,347) -- (503,342) -- (502,342) -- (495,335) -- (494,335) -- (488,329) -- (487,329) -- (481,323) -- (480,323) -- (472,315) -- (471,315) -- (467,311) -- (466,311) -- (461,305) -- (458,304) -- (442,289) -- (438,287) -- (416,266) -- (415,266) -- (403,255) -- (402,255) -- (398,251) -- (397,251) -- (393,247) -- (393,246) -- (390,245) -- (381,236) -- (380,236) -- (377,233) -- (375,233) -- (371,229) -- (370,229) -- (361,220) -- (360,220) -- (354,214) -- (353,214) -- (351,212) -- (351,211) -- (349,210) -- (342,203) -- (341,203) -- (335,197) -- (334,197) -- (327,190) -- (325,190) -- (312,177) -- (311,177) -- (307,172) -- (306,172) -- (298,164) -- (297,164) -- (291,158) -- (290,156) -- (288,156) -- (280,149) -- (277,148) -- (275,146) -- (275,145) -- (264,135) -- (263,133) -- (260,132) -- (247,119) -- (246,119) -- (240,113) -- (239,113) -- (232,106) -- (228,105) -- (225,108) -- (225,110) -- (230,115) -- (231,115) -- (241,126) -- (243,126) -- (247,130) -- (250,131) -- (262,144) -- (263,144) -- (268,149) -- (272,150) -- (277,155) -- (278,155) -- (281,159) -- (284,159) -- (286,161) -- (286,163) -- (288,165) -- (289,165) -- (294,171) -- (297,172) -- (300,175) -- (301,175) -- (315,189) -- (318,190) -- (343,213) -- (348,216) -- (355,224) -- (360,227) -- (364,231) -- (367,232) -- (377,242) -- (378,242) -- (383,247) -- (384,247) -- (391,254) -- (396,256) -- (401,262) -- (406,265) -- (411,270) -- (411,271) -- (415,276) -- (417,276) -- (419,278) -- (420,278) -- (434,293) -- (439,295) -- (443,299) -- (448,302) -- (472,325) -- (473,325) -- (479,331) -- (480,331) -- (489,340) -- (490,340) -- (494,344) -- (495,344) -- (501,351) -- (502,351) -- (508,357) -- (512,359) -- (519,366) -- (520,366) -- (525,371) -- (526,371) -- (531,376) -- (532,376) -- (542,386) -- (543,388) -- (544,388) -- (548,392) -- (549,392) -- (555,398) -- (556,398) -- (561,403) -- (562,403) -- (567,408) -- (568,408) -- (573,413) -- (574,413) -- (582,421) -- (586,423) -- (588,423) -- (597,432) -- (600,433) -- (605,438) -- (606,438) -- (609,441) -- (615,445) -- (615,447) -- (613,449) -- (614,450) -- (616,450) -- (622,457) -- (624,455) -- (625,455) -- (658,485) -- (666,490) -- (669,493) -- (667,495) -- (640,495) -- (639,494) -- (637,494) -- (636,495) -- (630,495) -- (629,494) -- (627,495) -- (625,495) -- (624,494) -- (584,494) -- (583,495) -- (576,495) -- (575,494) -- (573,494) -- (572,495) -- (569,495) -- (568,494) -- (565,494) -- (563,495) -- (562,494) -- (501,494) -- (500,493) -- (498,493) -- (497,494) -- (477,494) -- (476,493) -- (458,493) -- (457,494) -- (438,494) -- (437,493) -- (434,493) -- (433,494) -- (425,494) -- (424,493) -- (394,493) -- (391,495) -- (390,494) -- (365,494) -- (364,495) -- (359,495) -- (358,494) -- (349,494) -- (348,495) -- (346,495) -- (345,494) -- (340,494) -- (339,495) -- (336,495) -- (335,494) -- (331,494) -- (330,495) -- (327,495) -- (326,494) -- (321,494) -- (320,495) -- (318,495) -- (317,494) -- (256,494) -- (255,493) -- (250,493) -- (249,494) -- (237,494) -- (236,493) -- (215,493) -- (214,492) -- (209,492) -- (207,493) -- (203,491) -- (199,494) -- (182,494) -- (181,493) -- (179,493) -- (178,494) -- (176,493) -- (169,493) -- (168,494) -- (161,494) -- (160,493) -- (158,493) -- (157,494) -- (154,494) -- (153,493) -- (147,494) -- (146,493) -- (140,493) -- (139,494) -- (135,494) -- (134,493) -- (108,493) -- (107,492) -- (104,492) -- (103,493) -- (97,492) -- (95,493) -- (94,492) -- (90,492) -- (89,491) -- (87,491) -- (84,494) -- cycle;
\node[inner sep=0pt,font=\fontsize{10}{10}\selectfont] at (667.5,27.5) {\scalebox{5.841}{$(1)$}};
\node[inner sep=0pt,font=\fontsize{10}{10}\selectfont] at (183.5,86.0) {\scalebox{5.94}{$B$}};
\node[inner sep=0pt,font=\fontsize{10}{10}\selectfont] at (1168.5,87.5) {\scalebox{6.039}{$C$}};
\node[inner sep=0pt,font=\fontsize{10}{10}\selectfont] at (38.5,495.5) {\scalebox{6.633}{$A$}};
\node[inner sep=0pt,font=\fontsize{10}{10}\selectfont] at (587.0,461.5) {\scalebox{4.851}{$40^0$}};
\node[inner sep=0pt,font=\fontsize{10}{10}\selectfont] at (692.0,443.5) {\scalebox{5.049}{$100^0$}};
\end{tikzpicture}%
}
\end{center}

möglich, so sind

\begin{center}\resizebox{3.1cm}{!}{%
\begin{tikzpicture}[x=1pt,y=-1pt]
\path[use as bounding box] (0,0) rectangle (497,525);
\path[fill=black,draw=none,even odd rule]
(400,59) -- (399,59) -- (397,61) -- (397,71) -- (398,72) -- (398,79) -- (400,83) -- (400,86) -- (401,87) -- (401,90) -- (402,91) -- (402,95) -- (404,98) -- (404,104) -- (405,105) -- (405,109) -- (407,113) -- (407,118) -- (408,119) -- (408,123) -- (409,124) -- (409,130) -- (413,138) -- (414,148) -- (416,151) -- (416,154) -- (417,155) -- (417,161) -- (419,165) -- (420,176) -- (423,179) -- (421,183) -- (424,188) -- (424,193) -- (425,194) -- (425,205) -- (426,206) -- (426,208) -- (428,210) -- (427,214) -- (428,215) -- (428,218) -- (431,223) -- (430,226) -- (432,228) -- (433,234) -- (434,235) -- (433,240) -- (434,241) -- (434,246) -- (436,250) -- (436,255) -- (438,258) -- (438,268) -- (440,271) -- (441,282) -- (444,286) -- (443,290) -- (444,291) -- (446,302) -- (448,304) -- (448,306) -- (447,307) -- (447,312) -- (448,313) -- (448,318) -- (449,319) -- (450,330) -- (452,333) -- (452,336) -- (453,337) -- (453,340) -- (454,341) -- (454,346) -- (456,350) -- (456,355) -- (457,356) -- (457,359) -- (455,360) -- (454,362) -- (455,364) -- (459,367) -- (461,378) -- (463,378) -- (464,379) -- (464,395) -- (466,399) -- (465,403) -- (466,404) -- (467,410) -- (469,412) -- (468,414) -- (469,416) -- (469,425) -- (470,426) -- (470,428) -- (472,430) -- (471,433) -- (473,437) -- (473,445) -- (474,446) -- (474,449) -- (475,450) -- (476,462) -- (478,466) -- (478,472) -- (479,473) -- (479,476) -- (481,479) -- (481,480) -- (479,482) -- (475,480) -- (472,476) -- (471,476) -- (462,468) -- (457,465) -- (449,457) -- (445,455) -- (438,448) -- (437,448) -- (434,445) -- (429,442) -- (417,431) -- (412,428) -- (408,424) -- (406,424) -- (399,417) -- (398,417) -- (393,412) -- (388,409) -- (381,402) -- (378,401) -- (376,399) -- (376,397) -- (375,396) -- (373,396) -- (370,393) -- (369,393) -- (362,386) -- (361,386) -- (357,382) -- (353,380) -- (345,372) -- (342,371) -- (325,356) -- (324,356) -- (320,352) -- (319,352) -- (313,346) -- (310,345) -- (306,341) -- (305,341) -- (301,337) -- (300,337) -- (297,333) -- (296,333) -- (294,331) -- (292,331) -- (286,326) -- (285,326) -- (274,316) -- (270,314) -- (265,309) -- (264,309) -- (261,306) -- (258,305) -- (253,300) -- (252,300) -- (242,290) -- (238,288) -- (234,284) -- (233,284) -- (230,281) -- (225,278) -- (216,270) -- (211,268) -- (208,265) -- (207,265) -- (204,261) -- (203,261) -- (194,253) -- (190,251) -- (179,240) -- (177,240) -- (168,232) -- (163,229) -- (156,222) -- (155,222) -- (146,214) -- (143,214) -- (141,216) -- (141,217) -- (147,225) -- (149,225) -- (151,227) -- (152,227) -- (161,235) -- (172,242) -- (177,247) -- (179,247) -- (181,249) -- (181,250) -- (190,258) -- (193,259) -- (203,268) -- (208,271) -- (213,276) -- (215,276) -- (223,283) -- (225,283) -- (236,293) -- (237,293) -- (246,302) -- (251,305) -- (263,316) -- (265,316) -- (281,330) -- (286,332) -- (289,335) -- (290,335) -- (297,342) -- (301,344) -- (324,364) -- (327,365) -- (332,370) -- (336,372) -- (343,378) -- (345,378) -- (354,388) -- (357,389) -- (364,396) -- (365,396) -- (366,398) -- (367,398) -- (374,404) -- (378,406) -- (384,412) -- (387,413) -- (393,419) -- (397,421) -- (403,427) -- (410,430) -- (421,440) -- (423,440) -- (433,450) -- (434,450) -- (439,455) -- (441,455) -- (444,458) -- (444,459) -- (446,460) -- (453,467) -- (454,467) -- (459,472) -- (462,473) -- (467,478) -- (471,480) -- (476,485) -- (476,487) -- (473,489) -- (471,488) -- (422,488) -- (421,487) -- (419,487) -- (418,488) -- (407,488) -- (406,489) -- (362,489) -- (361,488) -- (359,488) -- (358,489) -- (347,489) -- (346,488) -- (342,488) -- (341,489) -- (308,489) -- (307,490) -- (300,490) -- (299,489) -- (294,489) -- (293,490) -- (289,490) -- (288,489) -- (286,490) -- (282,490) -- (281,489) -- (267,489) -- (266,490) -- (260,490) -- (259,489) -- (257,490) -- (254,490) -- (253,489) -- (238,489) -- (237,490) -- (231,490) -- (230,489) -- (227,489) -- (226,490) -- (208,490) -- (207,491) -- (203,491) -- (202,490) -- (196,490) -- (195,491) -- (192,491) -- (191,490) -- (184,490) -- (183,491) -- (180,491) -- (179,490) -- (175,490) -- (174,491) -- (172,490) -- (164,490) -- (163,491) -- (159,491) -- (158,490) -- (155,490) -- (154,491) -- (149,491) -- (148,490) -- (145,490) -- (144,491) -- (142,490) -- (140,490) -- (139,491) -- (123,491) -- (120,489) -- (118,490) -- (95,490) -- (94,491) -- (91,491) -- (90,490) -- (56,490) -- (54,492) -- (54,495) -- (58,496) -- (59,495) -- (70,495) -- (71,494) -- (76,494) -- (77,495) -- (80,495) -- (81,494) -- (85,494) -- (86,495) -- (90,495) -- (91,494) -- (93,495) -- (95,495) -- (96,494) -- (101,494) -- (102,495) -- (113,495) -- (114,496) -- (120,496) -- (121,495) -- (131,495) -- (132,496) -- (134,496) -- (135,495) -- (138,495) -- (139,496) -- (144,496) -- (145,495) -- (156,495) -- (157,496) -- (162,496) -- (163,495) -- (167,495) -- (168,496) -- (171,496) -- (172,495) -- (174,495) -- (175,496) -- (212,496) -- (213,495) -- (222,495) -- (223,494) -- (225,494) -- (226,495) -- (235,495) -- (236,496) -- (242,496) -- (243,495) -- (256,495) -- (258,496) -- (262,494) -- (264,494) -- (265,495) -- (284,495) -- (285,494) -- (300,494) -- (301,495) -- (307,495) -- (308,494) -- (341,494) -- (342,495) -- (344,495) -- (345,494) -- (349,494) -- (350,495) -- (353,495) -- (354,494) -- (364,494) -- (365,493) -- (367,494) -- (377,494) -- (378,493) -- (380,493) -- (381,494) -- (387,494) -- (389,495) -- (390,494) -- (406,494) -- (407,493) -- (412,493) -- (413,492) -- (415,493) -- (452,493) -- (453,494) -- (456,494) -- (457,493) -- (470,493) -- (471,492) -- (474,493) -- (475,492) -- (485,492) -- (486,491) -- (486,486) -- (485,485) -- (484,479) -- (483,478) -- (482,467) -- (481,466) -- (481,462) -- (480,461) -- (480,457) -- (479,456) -- (479,448) -- (477,444) -- (477,435) -- (476,434) -- (476,431) -- (474,427) -- (474,420) -- (473,419) -- (473,415) -- (472,414) -- (472,409) -- (471,408) -- (470,402) -- (469,401) -- (469,393) -- (468,392) -- (468,388) -- (467,387) -- (467,383) -- (466,382) -- (466,378) -- (465,377) -- (465,373) -- (464,372) -- (463,362) -- (462,361) -- (462,358) -- (461,357) -- (461,350) -- (460,349) -- (460,346) -- (459,345) -- (459,341) -- (458,340) -- (457,331) -- (455,327) -- (455,323) -- (453,320) -- (453,311) -- (451,308) -- (450,296) -- (449,295) -- (449,292) -- (448,291) -- (448,288) -- (447,287) -- (447,279) -- (445,277) -- (446,272) -- (445,271) -- (444,265) -- (443,264) -- (443,261) -- (442,260) -- (442,254) -- (441,253) -- (441,247) -- (440,246) -- (440,242) -- (439,241) -- (439,236) -- (438,235) -- (437,228) -- (435,225) -- (435,222) -- (434,221) -- (434,216) -- (435,215) -- (435,213) -- (434,212) -- (434,209) -- (431,205) -- (431,200) -- (428,194) -- (429,190) -- (426,186) -- (427,185) -- (427,181) -- (426,180) -- (426,175) -- (424,173) -- (424,168) -- (423,167) -- (424,161) -- (421,157) -- (422,152) -- (420,150) -- (419,143) -- (417,139) -- (417,135) -- (418,134) -- (418,132) -- (415,129) -- (415,124) -- (412,117) -- (412,112) -- (411,111) -- (411,107) -- (410,106) -- (410,102) -- (409,101) -- (409,97) -- (408,96) -- (408,89) -- (406,87) -- (406,83) -- (404,81) -- (404,71) -- (403,70) -- (403,67) -- (402,66) -- (402,62) -- cycle;
\node[inner sep=0pt,font=\fontsize{10}{10}\selectfont] at (156.5,91.5) {\scalebox{5.247}{$(2)$}};
\node[inner sep=0pt,font=\fontsize{10}{10}\selectfont] at (400.5,27.0) {\scalebox{5.346}{$B$}};
\node[inner sep=0pt,font=\fontsize{10}{10}\selectfont] at (112.0,207.5) {\scalebox{6.039}{$A$}};
\node[inner sep=0pt,font=\fontsize{10}{10}\selectfont] at (27.0,492.5) {\scalebox{6.435}{$C$}};
\node[inner sep=0pt,font=\fontsize{10}{10}\selectfont] at (431.5,399.5) {\scalebox{4.059}{$40^0$}};
\node[inner sep=0pt,font=\fontsize{10}{10}\selectfont] at (395.0,460.0) {\scalebox{3.96}{$40^0$}};
\end{tikzpicture}%
}
\quad \resizebox{4.7cm}{!}{%
\begin{tikzpicture}[x=1pt,y=-1pt]
\path[use as bounding box] (0,0) rectangle (889,606);
\path[fill=black,draw=none,even odd rule]
(827,40) -- (826,39) -- (824,39) -- (815,47) -- (812,48) -- (809,51) -- (808,51) -- (801,58) -- (797,60) -- (794,63) -- (791,64) -- (778,76) -- (773,79) -- (769,83) -- (767,83) -- (763,87) -- (757,90) -- (750,97) -- (748,97) -- (739,106) -- (737,106) -- (727,115) -- (726,115) -- (723,119) -- (713,124) -- (701,134) -- (694,138) -- (677,153) -- (674,154) -- (671,157) -- (670,157) -- (666,162) -- (664,162) -- (662,164) -- (659,165) -- (653,171) -- (649,173) -- (642,180) -- (640,180) -- (635,185) -- (634,185) -- (632,187) -- (630,187) -- (607,206) -- (603,207) -- (593,216) -- (591,216) -- (586,222) -- (582,224) -- (577,229) -- (572,232) -- (565,239) -- (562,240) -- (554,247) -- (551,248) -- (547,252) -- (544,253) -- (538,259) -- (535,260) -- (530,265) -- (527,266) -- (523,270) -- (518,273) -- (512,279) -- (511,279) -- (508,282) -- (506,282) -- (500,288) -- (497,289) -- (495,291) -- (495,292) -- (485,301) -- (476,300) -- (475,301) -- (455,302) -- (454,301) -- (420,301) -- (419,300) -- (400,300) -- (399,301) -- (380,301) -- (378,302) -- (374,300) -- (371,302) -- (369,301) -- (333,301) -- (332,302) -- (330,301) -- (315,301) -- (314,302) -- (306,302) -- (305,301) -- (299,301) -- (298,302) -- (294,302) -- (293,301) -- (291,302) -- (281,302) -- (280,301) -- (274,301) -- (273,302) -- (271,301) -- (249,301) -- (248,302) -- (244,302) -- (243,301) -- (237,301) -- (235,300) -- (233,302) -- (212,302) -- (211,303) -- (180,303) -- (179,302) -- (174,302) -- (173,301) -- (171,302) -- (168,302) -- (167,303) -- (165,302) -- (162,302) -- (161,303) -- (155,303) -- (154,302) -- (152,303) -- (140,303) -- (139,302) -- (137,302) -- (136,303) -- (126,303) -- (125,302) -- (123,303) -- (120,303) -- (119,302) -- (109,302) -- (108,303) -- (106,302) -- (104,302) -- (103,303) -- (99,303) -- (98,302) -- (89,302) -- (88,303) -- (81,303) -- (80,302) -- (72,302) -- (69,304) -- (69,306) -- (70,307) -- (73,307) -- (74,308) -- (80,308) -- (81,307) -- (92,307) -- (94,306) -- (95,307) -- (99,307) -- (102,310) -- (106,307) -- (122,307) -- (126,309) -- (134,308) -- (135,307) -- (141,307) -- (142,308) -- (150,308) -- (151,307) -- (153,307) -- (154,308) -- (168,308) -- (169,309) -- (178,309) -- (179,310) -- (181,309) -- (193,309) -- (194,308) -- (207,308) -- (208,307) -- (216,307) -- (217,308) -- (228,307) -- (229,308) -- (231,308) -- (232,307) -- (267,307) -- (268,306) -- (270,307) -- (273,307) -- (274,308) -- (276,307) -- (318,307) -- (319,306) -- (321,307) -- (341,307) -- (342,306) -- (344,306) -- (345,307) -- (352,307) -- (353,306) -- (360,306) -- (361,307) -- (381,307) -- (382,306) -- (384,306) -- (385,307) -- (391,307) -- (392,308) -- (394,307) -- (399,308) -- (400,307) -- (403,307) -- (404,306) -- (407,307) -- (408,306) -- (426,306) -- (427,305) -- (429,306) -- (431,305) -- (443,305) -- (444,306) -- (454,306) -- (455,305) -- (468,305) -- (469,306) -- (471,306) -- (473,305) -- (474,306) -- (484,306) -- (485,307) -- (488,307) -- (491,311) -- (495,313) -- (500,318) -- (500,319) -- (507,325) -- (507,330) -- (508,330) -- (513,335) -- (524,338) -- (533,345) -- (534,345) -- (542,352) -- (549,356) -- (561,367) -- (564,368) -- (568,372) -- (571,373) -- (582,383) -- (584,383) -- (589,388) -- (594,391) -- (604,400) -- (607,401) -- (617,410) -- (618,410) -- (621,413) -- (622,413) -- (632,421) -- (633,421) -- (651,437) -- (654,438) -- (665,448) -- (666,448) -- (669,451) -- (670,451) -- (673,454) -- (678,457) -- (682,461) -- (683,461) -- (684,463) -- (688,465) -- (693,470) -- (700,474) -- (711,484) -- (717,487) -- (727,497) -- (733,500) -- (745,511) -- (748,512) -- (756,519) -- (760,521) -- (765,527) -- (768,528) -- (774,534) -- (775,534) -- (781,539) -- (784,540) -- (790,546) -- (791,546) -- (795,550) -- (799,552) -- (806,559) -- (809,560) -- (816,570) -- (820,570) -- (822,572) -- (827,568) -- (827,566) -- (825,564) -- (824,564) -- (820,560) -- (816,558) -- (804,547) -- (800,545) -- (796,541) -- (793,540) -- (787,534) -- (781,531) -- (774,525) -- (770,523) -- (763,516) -- (760,515) -- (756,511) -- (751,508) -- (746,503) -- (745,503) -- (736,495) -- (735,495) -- (729,489) -- (725,487) -- (721,483) -- (719,483) -- (717,481) -- (716,481) -- (714,479) -- (714,478) -- (707,472) -- (701,469) -- (694,462) -- (693,462) -- (690,459) -- (689,459) -- (686,456) -- (681,453) -- (658,434) -- (652,431) -- (645,424) -- (641,422) -- (635,416) -- (627,411) -- (619,403) -- (614,400) -- (610,396) -- (606,394) -- (601,389) -- (597,387) -- (591,381) -- (589,381) -- (586,378) -- (583,377) -- (578,372) -- (573,369) -- (564,361) -- (557,357) -- (550,350) -- (549,350) -- (543,345) -- (540,344) -- (533,337) -- (528,334) -- (517,324) -- (511,321) -- (506,316) -- (502,314) -- (497,309) -- (492,306) -- (490,304) -- (490,302) -- (493,299) -- (501,296) -- (506,290) -- (509,289) -- (512,286) -- (513,286) -- (518,280) -- (519,280) -- (522,277) -- (527,275) -- (535,268) -- (538,267) -- (552,254) -- (554,253) -- (557,253) -- (559,251) -- (560,251) -- (566,244) -- (569,243) -- (572,240) -- (574,240) -- (578,236) -- (579,236) -- (587,229) -- (588,229) -- (596,220) -- (598,220) -- (602,218) -- (606,213) -- (611,211) -- (621,202) -- (622,202) -- (626,198) -- (631,195) -- (636,190) -- (639,189) -- (645,183) -- (648,182) -- (657,174) -- (665,169) -- (675,160) -- (677,160) -- (679,158) -- (680,158) -- (685,153) -- (687,153) -- (701,141) -- (704,140) -- (718,127) -- (720,127) -- (724,123) -- (728,122) -- (734,115) -- (735,115) -- (749,103) -- (752,102) -- (757,96) -- (761,95) -- (765,91) -- (767,91) -- (776,82) -- (778,82) -- (785,75) -- (788,74) -- (795,67) -- (799,66) -- (808,58) -- (810,58) -- (812,56) -- (813,56) -- (819,50) -- (820,50) -- (827,44) -- cycle;
\node[inner sep=0pt,font=\fontsize{10}{10}\selectfont] at (405.5,73.0) {\scalebox{5.544}{$(3)$}};
\node[inner sep=0pt,font=\fontsize{10}{10}\selectfont] at (863.0,27.5) {\scalebox{5.445}{$C$}};
\node[inner sep=0pt,font=\fontsize{10}{10}\selectfont] at (29.5,304.5) {\scalebox{5.643}{$A$}};
\node[inner sep=0pt,font=\fontsize{10}{10}\selectfont] at (849.0,580.5) {\scalebox{5.049}{$B$}};
\node[inner sep=0pt,font=\fontsize{10}{10}\selectfont] at (456.0,261.0) {\scalebox{4.356}{$140^0$}};
\node[inner sep=0pt,font=\fontsize{10}{10}\selectfont] at (456.0,336.5) {\scalebox{4.059}{$140^0$}};
\node[inner sep=0pt,font=\fontsize{10}{10}\selectfont] at (569.5,300.5) {\scalebox{4.455}{$80^0$}};
\end{tikzpicture}%
}
\quad \resizebox{4.6cm}{!}{%
\begin{tikzpicture}[x=1pt,y=-1pt]
\path[use as bounding box] (0,0) rectangle (869,518);
\path[fill=black,draw=none,even odd rule]
(575,413) -- (568,419) -- (572,417) -- cycle
(829,216) -- (827,216) -- (825,214) -- (823,214) -- (814,222) -- (811,223) -- (808,226) -- (807,226) -- (798,235) -- (797,235) -- (792,240) -- (791,240) -- (785,246) -- (780,249) -- (767,261) -- (762,264) -- (755,271) -- (747,276) -- (740,283) -- (739,283) -- (733,289) -- (729,291) -- (722,298) -- (719,299) -- (713,305) -- (710,306) -- (702,314) -- (701,314) -- (694,320) -- (690,322) -- (684,329) -- (683,329) -- (679,333) -- (674,336) -- (670,340) -- (670,342) -- (668,344) -- (665,344) -- (663,345) -- (657,351) -- (656,351) -- (652,355) -- (647,358) -- (637,367) -- (633,369) -- (619,381) -- (615,383) -- (610,389) -- (607,390) -- (603,394) -- (599,396) -- (595,401) -- (592,402) -- (588,406) -- (585,405) -- (584,406) -- (584,409) -- (573,420) -- (569,422) -- (560,431) -- (556,433) -- (551,438) -- (548,439) -- (539,447) -- (535,448) -- (534,450) -- (528,456) -- (525,457) -- (521,461) -- (518,462) -- (506,474) -- (503,474) -- (502,473) -- (502,470) -- (503,469) -- (504,463) -- (505,462) -- (505,455) -- (506,454) -- (506,450) -- (508,447) -- (508,443) -- (509,442) -- (509,438) -- (510,437) -- (510,432) -- (512,428) -- (513,416) -- (515,413) -- (516,403) -- (517,401) -- (519,400) -- (519,396) -- (520,395) -- (520,392) -- (521,391) -- (521,386) -- (522,385) -- (522,379) -- (524,377) -- (524,373) -- (523,371) -- (526,365) -- (526,355) -- (528,353) -- (528,348) -- (530,346) -- (529,344) -- (529,341) -- (530,340) -- (530,335) -- (533,331) -- (533,326) -- (532,325) -- (533,323) -- (533,318) -- (534,317) -- (534,314) -- (535,313) -- (535,310) -- (536,309) -- (536,303) -- (538,299) -- (538,295) -- (539,294) -- (539,290) -- (540,289) -- (540,281) -- (542,277) -- (543,267) -- (544,266) -- (544,263) -- (546,259) -- (545,254) -- (547,252) -- (547,246) -- (549,243) -- (549,237) -- (550,236) -- (550,226) -- (552,223) -- (553,213) -- (554,212) -- (554,209) -- (555,208) -- (555,201) -- (554,199) -- (557,195) -- (557,189) -- (558,188) -- (558,183) -- (559,182) -- (559,177) -- (560,176) -- (559,169) -- (561,166) -- (561,162) -- (563,159) -- (564,151) -- (566,149) -- (566,146) -- (565,145) -- (566,140) -- (565,139) -- (565,137) -- (566,136) -- (567,130) -- (568,129) -- (568,120) -- (569,119) -- (569,116) -- (570,115) -- (571,109) -- (572,108) -- (572,103) -- (575,97) -- (575,91) -- (576,90) -- (577,77) -- (579,74) -- (579,70) -- (581,68) -- (579,64) -- (579,62) -- (576,60) -- (574,60) -- (572,62) -- (572,68) -- (569,74) -- (569,79) -- (568,80) -- (568,88) -- (566,92) -- (567,96) -- (565,98) -- (565,106) -- (564,107) -- (564,112) -- (563,113) -- (563,116) -- (562,117) -- (562,120) -- (561,121) -- (561,125) -- (562,126) -- (561,128) -- (561,131) -- (562,132) -- (562,134) -- (560,137) -- (560,144) -- (558,146) -- (557,162) -- (555,164) -- (554,174) -- (553,175) -- (553,183) -- (550,190) -- (549,204) -- (548,205) -- (548,208) -- (547,209) -- (547,213) -- (546,214) -- (546,217) -- (545,218) -- (545,224) -- (544,225) -- (544,233) -- (542,236) -- (542,239) -- (540,242) -- (540,253) -- (539,254) -- (539,258) -- (538,259) -- (538,264) -- (536,266) -- (537,269) -- (535,271) -- (534,285) -- (532,288) -- (532,295) -- (531,296) -- (531,300) -- (530,301) -- (530,308) -- (529,309) -- (529,313) -- (527,317) -- (527,323) -- (526,324) -- (527,326) -- (526,328) -- (526,334) -- (525,335) -- (525,338) -- (524,339) -- (524,342) -- (523,343) -- (523,348) -- (521,351) -- (521,359) -- (520,360) -- (520,364) -- (518,366) -- (518,372) -- (517,373) -- (516,381) -- (515,382) -- (515,391) -- (513,393) -- (513,399) -- (512,400) -- (512,404) -- (511,405) -- (510,414) -- (509,415) -- (509,418) -- (508,419) -- (508,426) -- (506,430) -- (506,437) -- (505,438) -- (504,446) -- (503,447) -- (503,450) -- (502,451) -- (502,458) -- (499,464) -- (499,472) -- (498,473) -- (498,476) -- (495,479) -- (457,479) -- (456,480) -- (450,480) -- (449,479) -- (447,480) -- (437,480) -- (436,479) -- (431,480) -- (430,479) -- (428,479) -- (427,480) -- (424,480) -- (423,481) -- (420,481) -- (419,480) -- (415,480) -- (414,481) -- (410,481) -- (409,480) -- (407,480) -- (406,481) -- (392,481) -- (391,480) -- (273,480) -- (272,481) -- (270,480) -- (261,480) -- (260,481) -- (254,481) -- (253,480) -- (245,480) -- (244,481) -- (229,481) -- (228,482) -- (208,482) -- (207,481) -- (201,482) -- (200,481) -- (195,481) -- (194,482) -- (166,482) -- (165,483) -- (163,482) -- (161,482) -- (160,483) -- (158,482) -- (145,482) -- (144,483) -- (118,483) -- (117,482) -- (115,482) -- (114,483) -- (102,483) -- (101,482) -- (90,482) -- (89,483) -- (87,483) -- (86,482) -- (72,482) -- (69,484) -- (69,486) -- (72,488) -- (77,488) -- (78,489) -- (80,489) -- (81,488) -- (83,489) -- (102,489) -- (103,488) -- (105,489) -- (112,489) -- (113,488) -- (124,488) -- (125,489) -- (128,489) -- (129,488) -- (147,488) -- (148,489) -- (150,488) -- (174,488) -- (176,489) -- (177,488) -- (208,488) -- (209,487) -- (212,487) -- (213,488) -- (232,488) -- (233,487) -- (236,487) -- (237,486) -- (263,486) -- (264,485) -- (268,485) -- (269,486) -- (305,486) -- (306,487) -- (309,487) -- (310,486) -- (353,486) -- (354,485) -- (356,485) -- (357,486) -- (381,486) -- (383,485) -- (384,486) -- (388,486) -- (390,487) -- (391,486) -- (423,486) -- (426,484) -- (428,485) -- (430,485) -- (431,484) -- (441,484) -- (442,485) -- (444,484) -- (449,484) -- (450,485) -- (453,485) -- (454,484) -- (456,485) -- (462,485) -- (463,484) -- (470,484) -- (471,485) -- (473,485) -- (474,484) -- (487,484) -- (488,485) -- (499,484) -- (508,477) -- (509,477) -- (514,472) -- (520,469) -- (527,462) -- (531,461) -- (536,456) -- (537,456) -- (542,450) -- (547,447) -- (553,441) -- (556,440) -- (559,437) -- (560,437) -- (566,431) -- (567,431) -- (572,426) -- (577,423) -- (582,418) -- (585,417) -- (591,411) -- (593,411) -- (598,406) -- (602,404) -- (606,400) -- (607,400) -- (611,396) -- (612,396) -- (629,380) -- (637,375) -- (643,369) -- (646,368) -- (656,359) -- (659,358) -- (667,350) -- (668,350) -- (671,346) -- (672,346) -- (676,342) -- (681,340) -- (688,333) -- (689,333) -- (693,329) -- (697,327) -- (704,320) -- (705,320) -- (708,317) -- (713,314) -- (732,297) -- (737,294) -- (752,281) -- (756,279) -- (771,265) -- (776,262) -- (784,255) -- (788,253) -- (795,246) -- (796,246) -- (801,241) -- (802,241) -- (810,232) -- (811,232) -- cycle;
\node[inner sep=0pt,font=\fontsize{10}{10}\selectfont] at (337.5,183.5) {\scalebox{5.049}{$(4)$}};
\node[inner sep=0pt,font=\fontsize{10}{10}\selectfont] at (586.5,27.0) {\scalebox{5.346}{$C$}};
\node[inner sep=0pt,font=\fontsize{10}{10}\selectfont] at (837.0,188.0) {\scalebox{5.544}{$A$}};
\node[inner sep=0pt,font=\fontsize{10}{10}\selectfont] at (29.5,487.5) {\scalebox{6.039}{$B$}};
\node[inner sep=0pt,font=\fontsize{10}{10}\selectfont] at (441.0,448.5) {\scalebox{4.851}{$100^0$}};
\node[inner sep=0pt,font=\fontsize{10}{10}\selectfont] at (545.5,422.0) {\scalebox{4.158}{$40^0$}};
\end{tikzpicture}%
}
\end{center}

alle unmöglich. Es entsteht hier

\[
\begin{array}{l}
\text{2 aus 1, indem ich für } C\\
\text{3 aus 1, } \text{,,} \text{ } \text{,,} \text{ } \text{,,} \text{ } B\\
\text{4 aus 1, } \text{,,} \text{ } \text{,,} \text{ } \text{,,} \text{ } A
\end{array}
\Big\}\ \text{die entgegengesetzte Richtung nehme.}
\]

Umgekehrt hingegen, wenn (1) unmöglich ist, so ist Eine der drei (2), (3),
(4) möglich, und eo ipso die beiden andern unmöglich.

Wenn die Data des Problems, gleichsam auf gut Glück, aufgegeben sind, so kann
man die Auflösung entweder geometrisch, oder halb geometrisch, halb
analytisch (Rechnung mit roher Zeichnung verbindend), oder rein analytisch
suchen. In den beiden ersten Fällen erkennt man etwanige Unmöglichkeit \textls[120]{im Laufe} des Geschäfts von selbst. Aber bei dem letzten Verfahren findet man \textls[120]{immer} ein Resultat; ob diess aber den vorgegebenen Zahlen entspreche
%
\SourceLeaf{321}
oder einer andern (unter den vier jedesmal zusammengehörigen Variationen),
erkennt man dann erst hinterher,\\\hspace*{2em}falls man nicht gleich anfangs mein
Criterium angewandt hat.

Ich setze das Criterium selbst nicht her, um Ihnen das Vergnügen, es selbst
zu finden, nicht zu rauben. Jedenfalls hoffe ich durch das Vorstehende
hinlänglich klar gemacht zu haben, \textls[120]{von was es sich handelt}. Übrigens ist wohl
unnöthig, zu bemerken, dass der ganze Gegenstand höchst elementarisch ist.
Nur ist um so mehr zu verwundern, dass man ihn so ganz übersehen hat.

\mbox{. . . . . . . . . .}

\Needspace{12\baselineskip}\begin{center}[2.]\end{center}

\begin{center}\textsc{Gauss} an \textsc{Gerling}. Göttingen, 24. October 1840.\end{center}

\indent \mbox{. . . . . . . .} Ihre kleine Schrift über die \textsc{Pothenot}sche Aufgabe habe ich
mit Vergnügen gelesen; sie wird allen, die die Aufgabe praktisch anwenden
wollen, sehr lehrreich und nützlich sein. Ich habe mich aber oft gewundert,
dass die vielen Schriftsteller, die sich damit beschäftigt haben, einen in
theoretischer Beziehung ganz wesentlichen Theil ganz übergangen haben,
nemlich die Angabe der mathematischen Bedingungen, unter denen die Aufgabe
\textls[120]{physisch möglich} ist, und die Darstellung dieser Bedingungen in einer
einfachen eleganten Form. Ich mache mich wohl auf folgende Art am
leichtesten verständlich.

Die Rechnungsvorschriften auf vorgegebene numerische Data angewandt geben
(wenn nicht die vier Punkte in Einem Kreise liegen, wo man $\frac{0}{0}$
erhält) immer Ein einziges bestimmtes Resultat. Aber man erkennt leicht,
dass man \textls[120]{einerlei} Resultat erhält, es möge aufgegeben sein, dass der zweite
Punkt von dem ersten um den Winkel $A$ rechts, oder dass er von demselben um
$180^0-A$ links liegt. Eben so beim dritten Punkt. Oder mit andern Worten:

Ist vorgeschrieben, der zweite Punkt solle um $A$ rechts vom ersten liegen,
und der dritte um $B$ rechts vom ersten, so findet man ein Resultat, welches
\textls[120]{dasselbe} ist, was man bei drei andern Datis finden würde, nemlich
%
\SourceLeaf{322}
\[
\begin{array}{rcll@{\qquad}rcll}
2 & \text{um} & 180^0+A & \text{rechts von }1 & 3 & \text{um} & B & \text{rechts von }1\\[2pt]
 & \text{,,} & A & \text{,,} & & \text{,,} & 180^0+B & \text{,,}\\[2pt]
 & \text{,,} & 180^0+A & \text{,,} & & \text{,,} & 180^0+B & \text{,,}\;.
\end{array}
\]

Das Resultat selbst aber quadrirt immer nur mit Einem dieser vier Systeme von Datis, und die andern drei sind \textls[120]{physisch unmöglich}.

Der metaphysische Grund dieser Erscheinung ist, dass man von der beobachteten Richtung weiter nichts benutzt, als dass die \textls[120]{geraden Linien} gewisse Winkel mit einander machen, indem man diese geraden Linien auf \textls[120]{beiden} Seiten in indefinitum sich erstrecken lässt, während das Fortschreiten des Lichts nur in Einem bestimmten Sinne geschieht, also die Fälle, wo der Gegenstand rückwärts läge, ausgeschlossen werden müssen.

Zur Vollständigkeit gehört also die Auflösung der Aufgabe, den Datis vor geführter Rechnung gleich anzusehen, ob sie vereinbar sind. Es lässt sich diess sehr zierlich abthun, ich will aber Ihnen nicht vorgreifen, es Sich selbst zu entwickeln. \mbox{. . . . . . .}



\centerline{\textsc{Gauss} an \textsc{Gerling}. Göttingen, 14. Januar 1842.}

\indent \mbox{. . . . . . . .} Das Criterium wegen der \textsc{Pothenot}schen Aufgabe theile ich Ihnen zwar gern mit, wünsche jedoch, dass Sie es für Sich behalten, und zwar aus dem Grunde, weil ich das Theorem, womit es zusammenhängt, selbst einmal bei schicklicher Gelegenheit zu behandeln mir vorbehalte, weniger wegen der Eleganz des Theorems an sich, als wegen der Eleganz, welche die Anwendung der complexen Grössen dabei darbietet; also namentlich bei einer Gelegenheit, wo ich mehr von dem Gebrauch der complexen Grössen sagen kann. Würde das Theorem vorher anderweit ins Publicum gebracht, so würde für mich nachher keine Lust mehr bleiben, selbst darauf zurückzukommen. Also jetzt zur Sache.

Man kann das Criterium auf vielfache Weise einkleiden; am zierlichsten wird es sein, es so zu thun, dass alle drei gegebenen Punkte symmetrisch dabei vorkommen. Ich bezeichne also die drei gegebenen Punkte mit 1, 2, 3;
%
\SourceLeaf{323} den verlangten mit 0. Ich bezeichne ferner die Neigung der Geraden 12 gegen eine feste $L$ schlechthin mit 12; diese Neigung wird in einem \textls[120]{bestimmten} Sinn als wachsend betrachtet von 0 bis $360^0$ (die Vorstellung zu fixiren denken Sie, $L$ sei der Meridian, und 12 das Azimuth von 2 aus 1 gesehen). Auf ähnliche Weise und auf dieselbe $L$ bezogen, sollen 21, 13, 31, 23, 32 verstanden werden, wo also sich von selbst versteht, dass $21-12=\pm180^0$ (oder allgemeiner einem ungeraden Vielfachen von $180^0$ gleich). Die Zeichen 01, 02, 03 sollen zwar ähnliche Bedeutung haben, jedoch mit dem Unterschiede, dass sie nicht auf $L$, sondern auf \textls[120]{eine andere ganz willkürliche} feste Linie bezogen werden.

Man könnte kurz sagen, 12, 13, 21, 23, 31, 32 sind die wahren Azimuthe der betreffenden Richtungen, hingegen 01, 02, 03 die noch nicht orientirten oder einer noch unbekannten Orientirungscorrection bedürftigen Azimuthe. Es ist unnöthig, zu bemerken, dass hier Azimuth und Meridian bloss zur Fixirung der Vorstellung dienen.

Das Criterium, ob die vorgeschriebenen 6 Grössen (wenn Sie 12 und 21 bloss für Eine zählen) physisch verträglich sind, ist nun dieses:
\[
\left.
\begin{array}{l}
\sin(12-13-02+03)\\
\sin(23-21-03+01)\\
\sin(31-32-01+02)
\end{array}
\right\}\ \text{müssen Einerlei Zeichen haben.}
\]

Das vorhin erwähnte, meines Wissens\HistoricalNote{*)}{*)}{Sollte ich darin mich irren, und Sie mir nachweisen können, dass das Theorem schon von sonst jemand aufgestellt sei, so werden Sie mich sehr verpflichten.} bisher unbekannte, in meinem Besitz seit fast einem halben Jahrhundert befindliche Theorem, die Vierecke betreffend, ist folgendes:
\[
\frac{\sin(12-13-02+03)}{(01).(23)}=\frac{\sin(23-21-03+01)}{(02).(13)}=\frac{\sin(31-32-01+02)}{(03).(12)}\,,
\]

wo die Factoren in den Nennern die durchweg positiv genommenen Abstände zwischen den betreffenden Punkten ausdrücken. Die \textls[120]{Ableitung} dieses schönen Theorems vermittelst des Gebrauchs der complexen Zahlen hat die höchste Einfachheit. Hier will ich sie bloss andeuten.



%
\SourceLeaf{324}
Wenn (0), (1), (2), (3) die complexen Zahlen sind, die die Lage der Punkte 0, 1, 2, 3 vorstellen (ich kann ohne Zweifel voraussetzen, dass Ihnen das Wesentliche der Theorie der complexen Zahlen gegenwärtig ist, bemerke aber zum Überflusse, dass $x+iy$, d.\,i. $x+y\sqrt{-1}$, die complexe Zahl ist, die die Lage des Punkts vorstellt, dessen Coordinaten $x$, $y$ sind); $A$, $B$, $C$ hingegen die Punkte, welche durch die complexen Zahlen
\[
(0)(1)+(2)(3),\qquad (0)(2)+(1)(3),\qquad (0)(3)+(1)(2)
\]
vorgestellt werden, so sind die Seiten des Dreiecks $ABC$ folgende:
\[
BC=(01).(23),\qquad AC=(02).(13),\qquad BC=(03).(12), \]
und die äussern Winkel des Dreiecks, oder die Unterschiede der Winkel $A$, $B$, $C$ von $180^0$, diejenigen, deren Sinus in den obigen Gleichungen vorkommen; die überaus kleine Mühe, diess unmittelbar selbst zu erkennen, will ich Ihnen überlassen. Man kann übrigens dasselbe Theorem fast eben so zierlich \textls[120]{geometrisch} beweisen. —

Will man das Theorem selbst, d.\,i. seinen Inhalt, geometrisch versinnlichen, so muss man die beiden Fälle besonders betrachten, wo die vier Punkte
\[
\begin{array}{c@{\quad}ccc@{\qquad}c@{\quad}ccc}
 & * & & * & & * & & *\\
\text{so} & & * & & \text{oder so} & & &\\
 & & * & & & * & & *
\end{array}
\]
liegen, d.\,i. wo einer derselben innerhalb des Dreiecks, welches die drei übrigen bilden, liegt, oder wo alle vier ein Viereck ohne einspringende Winkel bilden.

Im ersten Fall lässt sich das Theorem so enunciiren.

\Needspace{13\baselineskip}Theorem: In einem Dreieck, dessen Seiten $AD\times BC$, $BD\times AC$, $CD\times AB$ sind, sind die diesen gegenüberstehenden Winkel respective $\alpha-A$, $\beta-B$, $\gamma-C$.

\begin{center}\resizebox{5.0cm}{!}{%
\begin{tikzpicture}[x=1pt,y=-1pt]
\path[use as bounding box] (0,0) rectangle (895,483);
\path[fill=black,draw=none,even odd rule]
(812,419) -- (810,418) -- (803,411) -- (803,410) -- (798,406) -- (795,401) -- (784,390) -- (784,389) -- (779,385) -- (774,377) -- (761,364) -- (761,363) -- (755,358) -- (755,357) -- (719,320) -- (719,319) -- (705,306) -- (705,305) -- (700,300) -- (700,299) -- (693,293) -- (693,292) -- (654,252) -- (654,251) -- (650,247) -- (650,246) -- (648,244) -- (647,244) -- (644,241) -- (644,240) -- (635,231) -- (634,231) -- (632,229) -- (631,226) -- (601,195) -- (600,192) -- (597,191) -- (581,174) -- (579,174) -- (576,171) -- (571,163) -- (557,149) -- (557,148) -- (552,144) -- (552,143) -- (548,139) -- (548,138) -- (546,137) -- (537,128) -- (534,123) -- (521,110) -- (521,109) -- (517,105) -- (516,105) -- (510,98) -- (508,98) -- (506,96) -- (506,95) -- (502,91) -- (501,91) -- (495,84) -- (489,84) -- (487,86) -- (486,86) -- (484,88) -- (484,89) -- (480,93) -- (479,93) -- (474,99) -- (470,99) -- (468,98) -- (467,99) -- (467,101) -- (461,107) -- (460,107) -- (454,114) -- (451,115) -- (443,123) -- (441,123) -- (430,133) -- (429,133) -- (422,139) -- (419,140) -- (415,145) -- (412,146) -- (407,151) -- (404,152) -- (386,168) -- (384,168) -- (367,184) -- (363,186) -- (353,195) -- (349,197) -- (343,203) -- (340,204) -- (337,207) -- (333,209) -- (324,219) -- (321,220) -- (313,227) -- (308,230) -- (293,244) -- (291,244) -- (268,263) -- (257,270) -- (254,274) -- (253,274) -- (242,284) -- (241,284) -- (233,291) -- (232,291) -- (227,296) -- (224,297) -- (218,303) -- (215,304) -- (209,310) -- (204,313) -- (199,318) -- (197,318) -- (194,321) -- (193,321) -- (192,323) -- (189,326) -- (188,326) -- (183,332) -- (180,333) -- (175,338) -- (174,338) -- (171,341) -- (168,342) -- (162,348) -- (155,351) -- (148,358) -- (147,358) -- (143,362) -- (142,362) -- (139,366) -- (136,367) -- (131,372) -- (128,373) -- (114,386) -- (113,386) -- (110,389) -- (107,390) -- (95,401) -- (88,405) -- (79,414) -- (76,415) -- (74,417) -- (75,421) -- (78,424) -- (84,424) -- (85,423) -- (92,423) -- (93,422) -- (96,422) -- (97,423) -- (104,423) -- (105,424) -- (110,423) -- (111,424) -- (122,424) -- (123,425) -- (141,425) -- (142,424) -- (146,424) -- (147,423) -- (152,423) -- (153,424) -- (157,424) -- (158,425) -- (160,425) -- (161,424) -- (165,425) -- (167,423) -- (169,423) -- (170,422) -- (175,423) -- (177,425) -- (185,424) -- (186,423) -- (190,423) -- (192,425) -- (194,425) -- (198,422) -- (203,423) -- (206,425) -- (209,425) -- (212,423) -- (215,423) -- (216,424) -- (220,424) -- (221,423) -- (228,423) -- (229,424) -- (232,424) -- (233,423) -- (238,423) -- (239,424) -- (250,424) -- (251,423) -- (253,423) -- (254,424) -- (257,422) -- (263,422) -- (264,423) -- (269,423) -- (271,422) -- (272,423) -- (280,423) -- (281,424) -- (288,424) -- (289,423) -- (295,423) -- (296,424) -- (299,424) -- (300,423) -- (305,423) -- (306,424) -- (312,424) -- (313,423) -- (321,423) -- (322,424) -- (324,424) -- (325,423) -- (336,423) -- (337,424) -- (339,423) -- (345,423) -- (346,424) -- (348,424) -- (350,423) -- (351,424) -- (357,424) -- (358,425) -- (360,425) -- (361,424) -- (366,424) -- (367,425) -- (370,425) -- (372,423) -- (378,423) -- (379,422) -- (383,423) -- (384,424) -- (387,424) -- (388,423) -- (390,424) -- (394,424) -- (395,423) -- (402,423) -- (403,424) -- (407,424) -- (408,423) -- (446,423) -- (447,424) -- (451,424) -- (452,423) -- (454,424) -- (456,423) -- (458,424) -- (462,424) -- (463,423) -- (465,424) -- (475,424) -- (476,423) -- (480,423) -- (481,424) -- (485,424) -- (486,423) -- (496,423) -- (497,424) -- (500,424) -- (502,423) -- (503,424) -- (531,424) -- (532,425) -- (550,425) -- (551,424) -- (600,424) -- (601,423) -- (610,423) -- (611,422) -- (615,422) -- (616,423) -- (628,423) -- (629,422) -- (631,422) -- (632,423) -- (655,423) -- (657,422) -- (658,423) -- (675,423) -- (676,424) -- (678,424) -- (679,423) -- (684,423) -- (685,424) -- (694,424) -- (695,423) -- (697,424) -- (708,424) -- (709,423) -- (767,423) -- (768,424) -- (770,423) -- (783,423) -- (784,424) -- (786,423) -- (799,423) -- (800,424) -- (803,424) -- (804,423) -- (809,423) -- (812,420) -- cycle
(795,419) -- (794,420) -- (768,420) -- (767,419) -- (715,419) -- (714,420) -- (669,420) -- (668,419) -- (665,419) -- (664,420) -- (659,420) -- (658,419) -- (642,419) -- (640,420) -- (639,419) -- (634,419) -- (633,418) -- (627,418) -- (626,419) -- (624,418) -- (611,418) -- (610,419) -- (608,419) -- (607,418) -- (601,418) -- (598,420) -- (595,420) -- (594,419) -- (592,420) -- (574,420) -- (573,419) -- (563,419) -- (562,420) -- (560,419) -- (552,419) -- (551,420) -- (547,420) -- (546,419) -- (544,420) -- (541,420) -- (540,419) -- (535,419) -- (534,420) -- (532,419) -- (517,419) -- (516,420) -- (513,420) -- (512,419) -- (495,419) -- (494,418) -- (491,418) -- (490,419) -- (448,419) -- (447,418) -- (444,418) -- (443,419) -- (434,419) -- (433,418) -- (423,418) -- (422,419) -- (416,419) -- (415,418) -- (383,418) -- (382,419) -- (377,419) -- (376,418) -- (372,418) -- (370,420) -- (366,420) -- (365,419) -- (350,419) -- (348,420) -- (347,419) -- (287,419) -- (286,418) -- (284,419) -- (282,419) -- (281,418) -- (279,418) -- (278,419) -- (271,419) -- (270,418) -- (256,418) -- (255,419) -- (242,419) -- (241,418) -- (237,418) -- (236,419) -- (233,419) -- (232,418) -- (221,418) -- (220,419) -- (216,419) -- (215,418) -- (211,418) -- (210,419) -- (207,419) -- (206,418) -- (185,418) -- (184,417) -- (174,417) -- (173,418) -- (168,418) -- (167,417) -- (158,417) -- (157,418) -- (155,417) -- (145,417) -- (141,419) -- (139,419) -- (138,418) -- (135,418) -- (134,419) -- (130,419) -- (129,418) -- (114,418) -- (113,419) -- (111,419) -- (110,418) -- (103,418) -- (102,417) -- (102,415) -- (106,413) -- (115,412) -- (116,411) -- (119,411) -- (120,410) -- (124,410) -- (125,409) -- (131,409) -- (132,408) -- (141,407) -- (143,405) -- (147,405) -- (154,402) -- (159,403) -- (162,401) -- (171,400) -- (172,399) -- (175,399) -- (176,398) -- (179,398) -- (180,397) -- (191,396) -- (194,394) -- (200,394) -- (201,393) -- (205,393) -- (209,391) -- (212,392) -- (216,390) -- (222,390) -- (223,389) -- (227,389) -- (228,388) -- (233,388) -- (237,386) -- (242,386) -- (246,384) -- (255,384) -- (257,382) -- (265,381) -- (272,378) -- (280,378) -- (281,377) -- (286,377) -- (287,376) -- (292,376) -- (293,375) -- (300,375) -- (303,373) -- (308,373) -- (309,372) -- (312,372) -- (313,371) -- (316,371) -- (317,370) -- (320,370) -- (321,369) -- (325,369) -- (326,368) -- (329,368) -- (332,366) -- (340,366) -- (343,364) -- (348,364) -- (351,362) -- (356,362) -- (357,361) -- (360,361) -- (361,360) -- (368,360) -- (373,357) -- (379,357) -- (382,355) -- (388,355) -- (391,353) -- (396,353) -- (399,351) -- (406,351) -- (412,348) -- (420,348) -- (421,347) -- (427,346) -- (428,345) -- (437,345) -- (438,344) -- (441,344) -- (442,343) -- (453,342) -- (454,341) -- (461,340) -- (463,338) -- (470,338) -- (476,335) -- (484,334) -- (486,332) -- (492,332) -- (493,331) -- (496,331) -- (497,330) -- (502,330) -- (506,328) -- (511,328) -- (512,327) -- (516,327) -- (517,326) -- (528,326) -- (532,329) -- (541,331) -- (543,333) -- (545,333) -- (548,335) -- (551,335) -- (559,339) -- (566,340) -- (569,342) -- (571,342) -- (575,345) -- (586,347) -- (587,348) -- (592,349) -- (595,351) -- (601,352) -- (604,354) -- (609,354) -- (610,355) -- (612,355) -- (618,358) -- (626,359) -- (629,362) -- (632,362) -- (638,365) -- (641,365) -- (642,366) -- (645,366) -- (646,367) -- (650,367) -- (658,371) -- (669,374) -- (671,376) -- (676,376) -- (677,377) -- (680,377) -- (682,379) -- (686,379) -- (691,382) -- (693,382) -- (694,383) -- (698,383) -- (700,385) -- (704,387) -- (722,390) -- (724,392) -- (726,392) -- (729,394) -- (743,398) -- (748,401) -- (752,401) -- (757,404) -- (760,404) -- (764,406) -- (767,406) -- (774,410) -- (777,410) -- (778,411) -- (783,411) -- (789,415) -- (792,415) -- (795,417) -- cycle
(499,95) -- (500,95) -- (511,106) -- (511,107) -- (526,123) -- (526,125) -- (531,128) -- (534,131) -- (534,132) -- (544,141) -- (544,142) -- (559,157) -- (559,158) -- (566,164) -- (566,166) -- (568,168) -- (569,168) -- (583,183) -- (583,184) -- (593,193) -- (593,196) -- (596,199) -- (600,200) -- (602,202) -- (602,203) -- (604,205) -- (605,205) -- (613,214) -- (613,216) -- (615,218) -- (616,218) -- (622,224) -- (622,225) -- (624,226) -- (633,235) -- (633,236) -- (641,244) -- (641,245) -- (647,250) -- (647,251) -- (649,253) -- (651,253) -- (654,256) -- (655,259) -- (658,262) -- (658,263) -- (666,270) -- (666,271) -- (671,276) -- (672,279) -- (680,286) -- (680,287) -- (684,291) -- (685,294) -- (695,304) -- (698,305) -- (707,314) -- (710,319) -- (727,334) -- (727,335) -- (735,343) -- (735,344) -- (741,350) -- (741,351) -- (743,352) -- (746,355) -- (746,356) -- (748,357) -- (753,362) -- (753,363) -- (755,364) -- (765,374) -- (765,375) -- (771,381) -- (771,382) -- (775,386) -- (775,387) -- (782,394) -- (783,394) -- (793,405) -- (794,405) -- (796,407) -- (796,408) -- (800,412) -- (800,414) -- (799,415) -- (797,415) -- (794,413) -- (787,412) -- (777,407) -- (774,407) -- (770,404) -- (765,403) -- (762,401) -- (759,401) -- (756,399) -- (753,399) -- (751,397) -- (748,397) -- (745,395) -- (742,395) -- (740,393) -- (737,393) -- (732,390) -- (726,389) -- (718,385) -- (708,384) -- (704,381) -- (698,380) -- (697,379) -- (694,379) -- (689,376) -- (686,376) -- (677,372) -- (673,372) -- (669,369) -- (666,369) -- (663,367) -- (660,367) -- (659,366) -- (653,365) -- (650,363) -- (647,363) -- (637,358) -- (634,358) -- (633,357) -- (625,355) -- (622,353) -- (620,353) -- (613,350) -- (610,350) -- (607,348) -- (603,348) -- (601,346) -- (598,347) -- (595,345) -- (592,345) -- (587,342) -- (584,342) -- (583,341) -- (578,341) -- (577,340) -- (572,339) -- (569,337) -- (563,336) -- (558,333) -- (556,333) -- (549,330) -- (545,330) -- (544,329) -- (542,329) -- (541,328) -- (527,324) -- (525,322) -- (526,319) -- (525,318) -- (525,312) -- (523,308) -- (523,298) -- (522,297) -- (522,294) -- (521,293) -- (522,289) -- (520,286) -- (521,284) -- (521,278) -- (519,275) -- (520,273) -- (520,271) -- (519,270) -- (519,262) -- (518,261) -- (519,257) -- (518,256) -- (518,252) -- (517,251) -- (517,249) -- (518,248) -- (518,244) -- (516,242) -- (516,240) -- (515,239) -- (516,237) -- (516,233) -- (514,230) -- (514,226) -- (513,225) -- (515,217) -- (513,213) -- (511,211) -- (511,201) -- (510,200) -- (510,195) -- (508,192) -- (509,186) -- (507,184) -- (508,182) -- (508,179) -- (507,178) -- (507,171) -- (506,170) -- (506,161) -- (504,157) -- (504,155) -- (505,154) -- (504,152) -- (504,145) -- (503,144) -- (503,140) -- (502,139) -- (502,126) -- (501,125) -- (501,121) -- (500,120) -- (500,105) -- (499,104) -- (498,98) -- (497,97) -- cycle
(493,93) -- (493,96) -- (495,98) -- (495,104) -- (496,105) -- (496,108) -- (495,109) -- (495,113) -- (496,114) -- (496,124) -- (495,125) -- (495,129) -- (496,130) -- (496,132) -- (498,134) -- (498,137) -- (499,138) -- (498,144) -- (499,145) -- (499,150) -- (500,151) -- (500,158) -- (502,162) -- (502,165) -- (505,169) -- (505,170) -- (503,172) -- (503,177) -- (504,178) -- (504,185) -- (506,188) -- (506,192) -- (505,193) -- (505,198) -- (506,199) -- (506,202) -- (507,203) -- (506,207) -- (507,208) -- (507,212) -- (508,213) -- (508,216) -- (509,217) -- (509,222) -- (510,223) -- (511,243) -- (512,244) -- (512,248) -- (513,249) -- (513,258) -- (514,259) -- (514,266) -- (515,267) -- (516,280) -- (517,281) -- (517,291) -- (518,292) -- (519,298) -- (520,299) -- (520,307) -- (521,308) -- (520,312) -- (521,313) -- (521,317) -- (522,319) -- (519,322) -- (517,323) -- (513,323) -- (512,324) -- (501,325) -- (497,327) -- (493,327) -- (492,328) -- (485,329) -- (481,331) -- (475,331) -- (474,332) -- (471,332) -- (470,333) -- (466,333) -- (465,334) -- (462,334) -- (461,335) -- (458,335) -- (457,336) -- (454,336) -- (453,337) -- (442,338) -- (441,339) -- (437,339) -- (436,340) -- (430,340) -- (429,341) -- (422,342) -- (418,344) -- (410,344) -- (409,345) -- (398,346) -- (394,348) -- (389,348) -- (383,351) -- (374,352) -- (369,355) -- (364,355) -- (360,357) -- (353,357) -- (350,359) -- (344,359) -- (343,360) -- (340,360) -- (339,361) -- (324,363) -- (322,365) -- (309,366) -- (308,367) -- (305,367) -- (302,369) -- (296,369) -- (295,370) -- (293,370) -- (292,371) -- (288,371) -- (284,373) -- (270,374) -- (269,375) -- (266,375) -- (265,376) -- (256,377) -- (250,380) -- (241,381) -- (240,382) -- (237,382) -- (234,384) -- (221,385) -- (220,386) -- (217,386) -- (216,387) -- (204,388) -- (202,390) -- (195,390) -- (194,391) -- (191,391) -- (190,392) -- (183,392) -- (180,394) -- (176,394) -- (175,395) -- (169,396) -- (168,397) -- (159,397) -- (158,398) -- (155,398) -- (154,399) -- (151,399) -- (150,400) -- (145,400) -- (140,403) -- (136,403) -- (133,405) -- (130,404) -- (126,406) -- (123,406) -- (122,407) -- (119,407) -- (118,408) -- (111,408) -- (110,409) -- (107,409) -- (103,411) -- (98,411) -- (94,413) -- (90,413) -- (89,412) -- (89,410) -- (90,409) -- (94,408) -- (100,402) -- (101,402) -- (108,396) -- (112,394) -- (115,391) -- (118,390) -- (135,375) -- (137,375) -- (143,368) -- (146,367) -- (153,360) -- (160,356) -- (165,351) -- (169,349) -- (175,343) -- (182,339) -- (186,334) -- (189,333) -- (198,324) -- (205,320) -- (210,315) -- (214,313) -- (224,304) -- (225,304) -- (227,302) -- (229,302) -- (237,294) -- (240,294) -- (242,293) -- (244,289) -- (246,287) -- (252,284) -- (254,282) -- (255,279) -- (258,276) -- (261,275) -- (268,268) -- (271,267) -- (277,262) -- (278,262) -- (284,256) -- (288,254) -- (293,249) -- (296,248) -- (301,243) -- (309,238) -- (320,228) -- (324,226) -- (327,223) -- (329,223) -- (331,221) -- (332,221) -- (338,215) -- (340,211) -- (347,207) -- (354,200) -- (357,199) -- (361,195) -- (362,195) -- (367,190) -- (368,190) -- (371,186) -- (376,183) -- (385,174) -- (386,174) -- (396,166) -- (397,166) -- (402,161) -- (403,161) -- (406,158) -- (411,155) -- (415,151) -- (416,151) -- (420,147) -- (421,147) -- (426,142) -- (430,140) -- (439,132) -- (441,132) -- (447,126) -- (450,125) -- (455,120) -- (456,120) -- (460,115) -- (461,115) -- (465,111) -- (472,107) -- (475,103) -- (476,103) -- (480,99) -- (481,99) -- (489,90) -- (490,90) -- cycle;
\node[inner sep=0pt,font=\fontsize{10}{10}\selectfont] at (495.0,42.5) {\scalebox{7.425}{$C$}};
\node[inner sep=0pt,font=\fontsize{10}{10}\selectfont] at (36.0,424.0) {\scalebox{6.534}{$A$}};
\node[inner sep=0pt,font=\fontsize{10}{10}\selectfont] at (855.5,431.5) {\scalebox{6.633}{$B$}};
\node[inner sep=0pt,font=\fontsize{10}{10}\selectfont] at (490.0,362.0) {\scalebox{5.742}{$D$}};
\node[inner sep=0pt,font=\fontsize{10}{10}\selectfont] at (554.0,291.0) {\scalebox{5.346}{$\alpha$}};
\node[inner sep=0pt,font=\fontsize{10}{10}\selectfont] at (492.5,295.5) {\scalebox{6.831}{$\beta$}};
\node[inner sep=0pt,font=\fontsize{10}{10}\selectfont] at (542.5,362.0) {\scalebox{4.95}{$\gamma$}};
\end{tikzpicture}%
}
\end{center}

Für den zweiten Fall ist der Ausdruck ein ganz ähnlicher, den ich Kürze halber übergehe.

Um noch einmal auf das \textls[120]{Criterium} zurückzukommen, bemerke ich,
%
\SourceLeaf{325}
dass das obige Criterium in so fern nicht ganz allgemein ist, als es stillschweigend voraussetzt, dass keiner der betreffenden Sinus $=0$ ist. Zur Vervollständigung müsste noch beigefügt werden:

1) ist \textls[120]{Einer} der drei Sinus $=0$, z.\,B. der erste, aber keiner der beiden andern, so fällt 0 mit 1 zusammen, wo natürlich für die Richtung 01 jeder beliebige Werth genommen werden kann.

2) Sind \textls[120]{zwei} $=0$, so wird von selbst auch der dritte $=0$ (weil die Summe der drei Winkel identisch $=0$ ist). In diesem Fall ist die \textsc{Pothenot}sche Aufgabe entweder unbestimmt oder physisch unmöglich; das zweite, wenn die Cosinus der Winkel alle drei $=+1$ sind, das erste, wenn zwei Cosinus jeder $=-1$, der dritte $=+1$ ist (die Combinationen $+1,+1,-1$ und $-1,-1,-1$ sind offenbar unmöglich, weil die Summe aller drei Winkel identisch $=0$ ist).

Schliesslich bemerke ich noch,

1) dass das Criterium sich zierlich ganz allgemein so aussprechen lässt: die Data sind unverträglich, wenn die drei Kreispunkte, die auf der Peripherie eines Kreises von Einerlei Anfangspunkt, in Einerlei Sinn positiv gezählt, um
\[
(01)+(23),\qquad (02)+(31),\qquad (03)+(12)
\]
\begin{minipage}{\linewidth}\raggedright\noindent abstehen, \textls[120]{innerhalb} Eines Halbkreises liegen; oder: damit die Data verträglich sein sollen, müssen diese drei Punkte \textls[120]{wenigstens} Einen vollen Halb\-kreis umfassen. (Füllen Sie gerade einen Halbkreis,

\begin{center}\resizebox{2.4cm}{!}{%
\begin{tikzpicture}[x=1pt,y=-1pt]
\path[use as bounding box] (0,0) rectangle (330,340);
\path[fill=black,draw=none,even odd rule]
(164,18) -- (163,19) -- (163,21) -- (164,22) -- (164,26) -- (163,27) -- (146,27) -- (145,28) -- (141,28) -- (135,31) -- (129,31) -- (128,32) -- (125,32) -- (124,33) -- (120,33) -- (117,35) -- (109,37) -- (102,41) -- (97,42) -- (93,45) -- (90,46) -- (87,49) -- (83,49) -- (77,56) -- (74,57) -- (55,74) -- (55,76) -- (50,81) -- (50,82) -- (42,91) -- (41,94) -- (38,97) -- (36,102) -- (33,106) -- (33,108) -- (30,112) -- (30,114) -- (28,116) -- (28,119) -- (26,121) -- (26,123) -- (22,131) -- (22,134) -- (24,137) -- (21,140) -- (21,143) -- (20,144) -- (20,146) -- (21,147) -- (20,149) -- (20,153) -- (19,154) -- (20,158) -- (18,161) -- (18,166) -- (17,167) -- (17,177) -- (16,178) -- (16,180) -- (17,181) -- (17,191) -- (18,192) -- (18,197) -- (17,198) -- (17,200) -- (20,203) -- (20,206) -- (21,207) -- (21,212) -- (24,218) -- (24,221) -- (27,226) -- (27,229) -- (29,231) -- (29,234) -- (31,236) -- (31,238) -- (40,254) -- (44,258) -- (45,261) -- (51,267) -- (51,268) -- (55,272) -- (55,273) -- (59,277) -- (59,278) -- (75,292) -- (83,296) -- (89,301) -- (91,301) -- (93,303) -- (97,304) -- (101,308) -- (103,309) -- (106,309) -- (110,312) -- (113,312) -- (118,315) -- (121,315) -- (123,317) -- (129,317) -- (131,319) -- (137,319) -- (138,320) -- (144,320) -- (145,321) -- (149,321) -- (150,320) -- (153,320) -- (154,321) -- (168,321) -- (169,323) -- (164,329) -- (164,331) -- (165,332) -- (167,332) -- (174,324) -- (175,324) -- (179,328) -- (182,328) -- (183,326) -- (180,323) -- (180,322) -- (182,320) -- (186,320) -- (188,322) -- (189,322) -- (192,319) -- (198,319) -- (203,316) -- (211,315) -- (215,312) -- (217,312) -- (220,310) -- (222,310) -- (226,308) -- (229,308) -- (231,306) -- (234,306) -- (236,304) -- (238,304) -- (240,302) -- (244,301) -- (249,295) -- (253,294) -- (255,292) -- (259,290) -- (261,288) -- (261,287) -- (266,282) -- (270,280) -- (272,276) -- (283,266) -- (284,263) -- (291,255) -- (293,251) -- (293,249) -- (297,245) -- (297,242) -- (303,234) -- (304,230) -- (306,227) -- (306,223) -- (308,221) -- (310,217) -- (310,213) -- (312,209) -- (311,206) -- (312,205) -- (313,197) -- (314,196) -- (313,190) -- (314,189) -- (314,183) -- (315,182) -- (314,180) -- (314,168) -- (315,167) -- (314,166) -- (314,156) -- (313,155) -- (313,152) -- (312,151) -- (312,144) -- (310,140) -- (310,134) -- (308,132) -- (308,130) -- (306,127) -- (305,122) -- (302,118) -- (302,115) -- (299,110) -- (299,107) -- (289,92) -- (289,90) -- (291,88) -- (295,87) -- (295,85) -- (293,84) -- (289,84) -- (287,82) -- (287,80) -- (290,74) -- (290,73) -- (288,71) -- (287,71) -- (285,77) -- (282,80) -- (281,80) -- (269,68) -- (269,67) -- (260,59) -- (255,57) -- (245,49) -- (243,49) -- (240,46) -- (234,44) -- (229,40) -- (222,38) -- (220,36) -- (215,35) -- (212,33) -- (210,33) -- (209,32) -- (205,32) -- (202,30) -- (193,30) -- (192,29) -- (181,28) -- (179,26) -- (179,25) -- (182,22) -- (181,20) -- (178,21) -- (173,26) -- (170,26) -- (168,24) -- (168,22) -- (166,20) -- (166,19) -- cycle
(157,31) -- (163,31) -- (164,32) -- (164,34) -- (159,38) -- (159,40) -- (160,41) -- (162,41) -- (167,36) -- (169,36) -- (172,39) -- (176,40) -- (177,38) -- (174,35) -- (174,34) -- (176,32) -- (181,32) -- (182,33) -- (186,33) -- (187,34) -- (195,35) -- (196,36) -- (202,36) -- (204,38) -- (209,38) -- (212,40) -- (215,40) -- (217,42) -- (221,44) -- (224,44) -- (232,48) -- (234,50) -- (236,50) -- (242,56) -- (250,60) -- (253,63) -- (257,65) -- (262,70) -- (265,71) -- (269,75) -- (270,75) -- (274,79) -- (272,81) -- (270,81) -- (268,83) -- (268,84) -- (270,86) -- (275,86) -- (278,89) -- (277,92) -- (275,94) -- (275,97) -- (276,98) -- (279,98) -- (283,94) -- (284,94) -- (286,96) -- (286,98) -- (287,99) -- (287,103) -- (288,104) -- (290,104) -- (292,106) -- (292,109) -- (294,113) -- (296,115) -- (296,117) -- (299,122) -- (300,128) -- (302,130) -- (302,132) -- (304,134) -- (304,137) -- (306,141) -- (305,145) -- (307,147) -- (307,153) -- (308,154) -- (308,162) -- (309,163) -- (309,175) -- (311,177) -- (311,178) -- (309,180) -- (308,200) -- (307,201) -- (307,205) -- (305,209) -- (303,220) -- (300,224) -- (300,227) -- (298,229) -- (298,231) -- (296,233) -- (294,237) -- (294,239) -- (292,241) -- (290,245) -- (290,247) -- (285,252) -- (285,254) -- (282,257) -- (280,261) -- (273,268) -- (273,269) -- (269,273) -- (268,273) -- (257,285) -- (255,285) -- (251,289) -- (248,290) -- (243,294) -- (241,294) -- (239,296) -- (234,298) -- (231,301) -- (223,303) -- (221,305) -- (219,305) -- (218,306) -- (216,306) -- (209,309) -- (203,310) -- (199,313) -- (189,314) -- (188,315) -- (185,315) -- (183,313) -- (183,312) -- (186,309) -- (186,306) -- (185,305) -- (182,306) -- (180,310) -- (174,316) -- (169,312) -- (169,310) -- (166,306) -- (164,306) -- (163,307) -- (163,310) -- (166,313) -- (166,314) -- (164,316) -- (151,316) -- (150,315) -- (147,315) -- (144,313) -- (138,314) -- (137,313) -- (134,313) -- (130,311) -- (120,310) -- (114,306) -- (108,305) -- (104,303) -- (102,301) -- (100,301) -- (92,295) -- (88,294) -- (81,288) -- (79,288) -- (60,270) -- (53,259) -- (44,249) -- (43,247) -- (43,242) -- (39,238) -- (38,235) -- (35,232) -- (35,229) -- (34,228) -- (33,223) -- (31,221) -- (30,216) -- (28,214) -- (27,203) -- (26,202) -- (27,197) -- (25,195) -- (25,192) -- (24,191) -- (24,182) -- (23,181) -- (23,176) -- (26,173) -- (26,172) -- (24,170) -- (24,165) -- (25,164) -- (25,153) -- (26,152) -- (26,149) -- (27,148) -- (27,142) -- (29,140) -- (29,133) -- (31,131) -- (31,127) -- (32,126) -- (32,124) -- (35,119) -- (35,117) -- (39,109) -- (44,105) -- (44,103) -- (46,100) -- (46,97) -- (48,95) -- (49,95) -- (52,91) -- (54,90) -- (55,86) -- (61,81) -- (62,81) -- (66,77) -- (67,74) -- (69,72) -- (69,70) -- (71,68) -- (74,68) -- (76,67) -- (83,60) -- (87,58) -- (91,53) -- (93,53) -- (97,49) -- (99,50) -- (101,49) -- (105,49) -- (106,48) -- (106,46) -- (109,43) -- (114,43) -- (120,39) -- (123,39) -- (124,38) -- (127,38) -- (128,37) -- (132,37) -- (133,36) -- (138,36) -- (139,35) -- (142,35) -- (145,33) -- (150,33) -- (151,32) -- (156,32) -- cycle;
\end{tikzpicture}%
}
\end{center}

so ist das Problem unbestimmt.)
\end{minipage}

2) dass sich noch vielerlei zusetzen liesse, wodurch aber anstatt eines Briefes eine Abhandlung entstehen würde, die zu \textls[120]{ordnen} und zu schreiben ich vorerst keine Zeit habe. In der Hoffnung, dass Sie in den Andeutungen vielfachen Stoff zum Nachdenken finden, schliesse ich für heute mit der Bitte mich bald wieder mit einigen Zeilen zu erfreuen.


%
\SourceLeaf{326}

\begin{center}
[III.]

\textbf{DAS VIERECK.}

[1.]

[\emph{Handbuch, den astronomischen Wissenschaften gewidmet}, S.\,150.]
\end{center}

$p, p', p'', p'''$ die complexen Örter der Punkte $(0), (1), (2), (3)$. Die Betrachtung des Dreiecks, dessen Ecken $(4), (5), (6)$ durch
\[
pp'+p''p''', \qquad pp''+p'''p', \qquad pp'''+p'p'',
\]
dessen Seiten [65, 46, 54] ihrer Grösse und Richtung nach aus den Differenzen
\[
(p-p')(p''-p'''), \qquad (p-p'')(p'''-p'), \qquad (p-p''')(p'-p'')
\]
erkannt werden, also

\begin{center}
\begin{tabular}{rllcl}
65: & Grösse & $(01).(23)$ & Richtung & $01+23$\\
46: & ,, & $(02).(31)$ & ,, & $02+31$\\
54: & ,, & $(03).(12)$ & ,, & $03+12$,
\end{tabular}
\end{center}

gibt die Gleichungen (s.\,p.\,201 und p.\,269 [des Handbuchs])
\[
\begin{aligned}
\frac{\sin(02+31-03-12)}{(01).(23)} &= \frac{\sin(03+12-01-23)}{(02).(31)}\\
&= \frac{\sin(01+23-02-31)}{(03).(12)} = \Omega.
\end{aligned}
\]

Man findet ferner leicht (s.\,pag.\,284 [des Handbuchs]) aus der Entwickelung eines dieser Ausdrücke:
\[
\Omega = \frac{(01)\sin(03-02)+(02)\sin(01-03)+(03)\sin(02-01)}{(12).(13).(23)}.
\]

Bezeichnet man $\Omega$, in so fern es von obigen bestimmten Elementen abhängt, durch $\Omega^{0123}$; so ist, wenn $(1), (2), (3), (4)$ feste Punkte und $(0)$ einen beweglichen Punkt bedeuten, die Orientirung auch als veränderlich betrachtet, die Bedingungsgleichung zwischen den Differentialien:
\[
\left.
\begin{aligned}
&(01).(23).(24).(34).\Omega^{0234}\,d(01)\\
-\;&(02).(13).(14).(34).\Omega^{0134}\,d(02)\\
[+\;&(03).(12).(14).(24).\Omega^{0124}\,d(03)\\
-\;&(04).(12).(13).(23).\Omega^{0123}\,d(04)
\end{aligned}
\right\} = 0.
\]
%
\SourceLeaf{327}\Needspace{10\baselineskip}

\begin{center}
[2.]

[\emph{Handbuch}, S.\,284.]
\end{center}

Die Seite 150 [des Handbuchs] erwähnte Entwickelung ist folgende. Der erste Ausdruck wird
\[
= \frac{1}{(01)(23)}\{\sin(02-12)\cos(31-03)+\sin(31-03)\cos(02-12)\}.
\]

Es ist aber
\[
(12)\sin(02-12) = (01)\sin(01-02)
\]
\[
(13)\sin(31-03) = (01)\sin(01-03),
\]
also
\[
\Omega = \frac{1}{(12)(13)(23)}\{(13)\sin(01-02)\cos(31-03)+(12)\sin(01-03)\cos(02-12)\}.
\]

Ferner ist
\[
(13)\cos(31-03) = (01)\cos(01-03)-(03)
\]
\[
(12)\cos(02-12) = (02)-(01)\cos(01-02).
\]
Also
\[
\begin{aligned}
\Omega ={}& \frac{1}{(12)(13)(23)}\\
&\quad\left\{
\begin{aligned}
&(01)\big[\sin(01-02)\cos(01-03)-\sin(01-03)\cos(01-02)\big]\\
+\;&(02)\sin(01-03)-(03)\sin(01-02)
\end{aligned}
\right\}
\end{aligned}
\]
\[
= \frac{1}{(12)(13)(23)}\{(01)\sin(03-02)+(02)\sin(01-03)+(03)\sin(02-01)\}.
\]

\begin{center}
W. Z. B. W.

[3.]

[\emph{Handbuch}, S.\,201.]
\end{center}

Die Ableitung der oben S.\,150 [des Handbuchs] stehenden Sätze das Viereck betreffend kann einfacher so geschehen.

Die complexen Grössen $p^0, p', p'', p'''$ bezeichnen wie dort die vier Punkte, die Adjuncten dieser Grössen seien $P^0, P', P'', P'''$.
%
\SourceLeaf{328}

Man setze Kürze halber
\[
(p''-p^0)(p'''-p^0) = a
\]
\[
(p'''-p^0)(p'-p^0) = b
\]
\[
(p'-p^0)(p''-p^0) = c
\]
und die Adjuncten dieser Grössen $= A, B, C$.

\Needspace{13\baselineskip}Man setze
\begin{equation}
Q = aB+bC+cA-Ab-Bc-Ca,
\tag*{1)}
\end{equation}
also auch (identisch)
\begin{equation}
Q = (a-b)(A-C)-(A-B)(a-c)
\tag*{2)}
\end{equation}
\begin{equation}
Q = (b-c)(B-A)-(B-C)(b-a)
\tag*{3)}
\end{equation}
\begin{equation}
Q = (c-a)(C-B)-(C-A)(c-b).
\tag*{4)}
\end{equation}

Angewandt auf unsern Fall gibt
\begin{equation}
\begin{aligned}
Q = \;& (p'''-p^0)(P'''-P^0)\{(p''-p^0)(P'-P^0)-(P''-P^0)(p'-p^0)\}\\
+\;& (p'-p^0)(P'-P^0)\{(p'''-p^0)(P''-P^0)-(P'''-P^0)(p''-p^0)\}\\
+\;& (p''-p^0)(P''-P^0)\{(p'-p^0)(P'''-P^0)-(P'-P^0)(p'''-p^0)\},
\end{aligned}
\tag{1}
\end{equation}
\begin{equation}
Q=(p'''-p^0)(p''-p')(P''-P^0)(P'''-P') - (P'''-P^0)(P''-P')(p''-p^0)(p'''-p'),
\tag{2}
\end{equation}
\begin{equation}
Q=(p'-p^0)(p'''-p'')(P'''-P^0)(P'-P'') - (P'-P^0)(P'''-P'')(p'''-p^0)(p'-p''),
\tag{3}
\end{equation}
\begin{equation}
Q=(p''-p^0)(p'-p''')(P'-P^0)(P''-P''') - (P''-P^0)(P'-P''')(p'-p^0)(p''-p'''),
\tag{4}
\end{equation}

woraus alles leicht von selbst folgt.

Offenbar kann man unbeschadet der Allgemeinheit von Anfang an $p^0=0$ setzen, wodurch die Entwickelung noch viel kürzer wird.
%
\SourceLeaf{329}

\begin{center}
[4.]

[\emph{Handbuch}, S.\,203.]
\end{center}

Die Auflösung der \textsc{Pothenot}schen Aufgabe beruht auf den drei Gleichungen
\[
\frac{\alpha a+\beta b+\gamma c}{\alpha}\, . \, x' - \frac{\alpha' a'+\beta' b'+\gamma' c'}{\alpha'}\, . \, x - \frac{\beta b+\gamma c}{\alpha}\, . \, a' + \frac{\beta' b'+\gamma' c'}{\alpha'}\, . \, a = 0
\]
\[
\frac{\alpha a+\beta b+\gamma c}{\beta}\, . \, x' - \frac{\alpha' a'+\beta' b'+\gamma' c'}{\beta'}\, . \, x - \frac{\alpha a+\gamma c}{\beta}\, . \, b' + \frac{\alpha' a'+\gamma' c'}{\beta'}\, . \, b = 0
\]
\[
\frac{\alpha a+\beta b+\gamma c}{\gamma}\, . \, x' - \frac{\alpha' a'+\beta' b'+\gamma' c'}{\gamma'}\, . \, x - \frac{\alpha a+\beta b}{\gamma}\, . \, c' + \frac{\alpha' a'+\beta' b'}{\gamma'}\, . \, c = 0,
\]

wovon jede schon in den beiden andern enthalten ist. Also
\[
a\alpha+b\beta+c\gamma=\omega \quad [\text{und } a'\alpha'+b'\beta'+c'\gamma'=\omega']
\]
gesetzt:
\[
\left(\frac{\beta}{\beta'}-\frac{\gamma}{\gamma'}\right)\omega' x + (b'-c')\omega - \left(\frac{\beta b}{\beta'}-\frac{\gamma c}{\gamma'}\right)\omega' = 0
\]
\[
\left(\frac{\gamma}{\gamma'}-\frac{\alpha}{\alpha'}\right)\omega' x + (c'-a')\omega - \left(\frac{\gamma c}{\gamma'}-\frac{\alpha a}{\alpha'}\right)\omega' = 0
\]
\[
\left(\frac{\alpha}{\alpha'}-\frac{\beta}{\beta'}\right)\omega' x + (a'-b')\omega - \left(\frac{\alpha a}{\alpha'}-\frac{\beta b}{\beta'}\right)\omega' = 0.
\]

Noch zierlicher
\[
\begin{aligned}
&(\alpha\beta'-\beta\alpha')(\alpha' a'+\beta' b'+\gamma' c')x + a\alpha\beta'\gamma'(b'-c')\\
&\quad+ \beta b\alpha'\gamma'(c'-a') + \gamma c\beta'\alpha'(a'-b') = 0.
\end{aligned}
\]

\medskip

Sind $A, B, C$ drei complexe Zahlen, deren Winkel den beobachteten Azimuthen entsprechen, so kann man setzen
\[
\alpha = A(BC'-CB')
\]
\[
\beta = B(CA'-AC')
\]
\[
\gamma = C(AB'-BA').
\]

Die physische Möglichkeit beruht darauf, dass
\[
B'C(A-B)(A'-C') - BC'(A'-B')(A-C)
\]
mit den beiden andern symmetrisch daraus folgenden Grössen einerlei Zeichen haben [muss].
%
\SourceLeaf{330}\Needspace{28\baselineskip}

\begin{center}
[5.]

[\emph{Handbuch}, S.\ 269 bis 271.]
\end{center}

\begin{center}Zusatz zu p.\ 150 und 201 [des Handbuchs].\end{center}

In Fig.\ I sind die vier gegebenen Punkte und die zwischen ihnen gebildeten Verbindungslinien und Winkel mit ihren Bezeichnungen vorgestellt. Fig.\ II. entsteht, indem man die Winkel $\mathfrak{A}$, $\mathfrak{B}$, $\mathfrak{C}$ beibehält, aber anstatt der Radien $\alpha$, $\beta$, $\gamma$ die reciproken $\dfrac{\alpha\beta\gamma}{\alpha}$, $\dfrac{\alpha\beta\gamma}{\beta}$, $\dfrac{\alpha\beta\gamma}{\gamma}$ annimmt. Die übrigen Bezeichnungen erklären sich von selbst.

\begin{center}
\begin{minipage}[b]{5.5cm}\centering\resizebox{5.2cm}{!}{%
\begin{tikzpicture}[x=1pt,y=-1pt]
\path[use as bounding box] (0,0) rectangle (683,717);
\path[fill=black,draw=none,even odd rule]
(297,79) -- (294,82) -- (294,87) -- (291,91) -- (291,93) -- (289,96) -- (289,99) -- (286,103) -- (286,108) -- (285,109) -- (284,114) -- (281,117) -- (281,120) -- (276,130) -- (276,133) -- (273,136) -- (273,139) -- (271,142) -- (271,146) -- (270,147) -- (269,155) -- (268,156) -- (269,158) -- (267,161) -- (264,171) -- (260,174) -- (261,177) -- (256,185) -- (256,187) -- (253,193) -- (253,196) -- (251,198) -- (251,203) -- (249,206) -- (247,214) -- (245,217) -- (245,219) -- (242,226) -- (242,229) -- (240,231) -- (239,235) -- (236,239) -- (236,244) -- (234,246) -- (234,249) -- (231,253) -- (230,261) -- (228,265) -- (226,267) -- (224,276) -- (218,287) -- (218,290) -- (215,294) -- (214,300) -- (212,302) -- (212,307) -- (211,308) -- (211,311) -- (209,314) -- (209,317) -- (206,323) -- (205,329) -- (203,331) -- (203,335) -- (200,338) -- (199,343) -- (197,346) -- (197,349) -- (194,353) -- (193,359) -- (191,362) -- (190,368) -- (187,372) -- (186,378) -- (184,380) -- (184,383) -- (182,385) -- (182,388) -- (180,391) -- (179,396) -- (177,398) -- (177,402) -- (175,405) -- (175,408) -- (172,415) -- (171,421) -- (167,427) -- (167,431) -- (165,433) -- (165,435) -- (162,440) -- (161,445) -- (158,450) -- (157,457) -- (155,460) -- (155,463) -- (153,466) -- (152,472) -- (150,474) -- (148,478) -- (148,481) -- (146,484) -- (146,487) -- (143,492) -- (142,498) -- (140,501) -- (139,508) -- (137,510) -- (137,513) -- (133,520) -- (133,523) -- (129,531) -- (129,533) -- (126,540) -- (126,544) -- (121,553) -- (121,559) -- (120,560) -- (120,562) -- (118,565) -- (118,567) -- (116,569) -- (116,572) -- (114,575) -- (114,578) -- (112,581) -- (111,587) -- (109,589) -- (106,600) -- (103,604) -- (103,609) -- (101,612) -- (100,618) -- (98,621) -- (98,623) -- (96,626) -- (95,631) -- (93,634) -- (92,640) -- (90,642) -- (90,647) -- (86,652) -- (87,655) -- (89,657) -- (101,657) -- (102,658) -- (118,658) -- (119,659) -- (127,659) -- (128,658) -- (134,658) -- (135,657) -- (140,657) -- (141,658) -- (150,658) -- (151,657) -- (157,657) -- (158,658) -- (162,658) -- (163,657) -- (165,657) -- (166,658) -- (174,658) -- (175,657) -- (186,657) -- (187,658) -- (192,658) -- (193,657) -- (202,657) -- (203,658) -- (247,657) -- (248,658) -- (251,658) -- (253,657) -- (254,658) -- (265,658) -- (266,659) -- (277,659) -- (278,658) -- (285,659) -- (289,657) -- (291,658) -- (295,658) -- (296,657) -- (298,658) -- (303,658) -- (304,657) -- (315,657) -- (316,656) -- (321,656) -- (323,658) -- (324,658) -- (328,655) -- (337,655) -- (338,656) -- (341,655) -- (342,656) -- (351,656) -- (352,657) -- (354,656) -- (358,656) -- (359,657) -- (376,657) -- (377,656) -- (385,656) -- (386,657) -- (389,657) -- (390,656) -- (392,656) -- (393,657) -- (437,657) -- (438,656) -- (440,657) -- (452,657) -- (454,656) -- (455,657) -- (470,657) -- (471,658) -- (476,658) -- (477,657) -- (483,657) -- (484,656) -- (486,656) -- (487,657) -- (506,657) -- (507,656) -- (509,657) -- (512,657) -- (513,656) -- (518,656) -- (519,657) -- (530,657) -- (531,656) -- (537,656) -- (538,657) -- (547,657) -- (548,658) -- (552,658) -- (553,657) -- (563,658) -- (564,657) -- (576,657) -- (577,656) -- (583,656) -- (584,655) -- (597,655) -- (598,656) -- (600,655) -- (606,655) -- (606,651) -- (608,649) -- (609,649) -- (603,637) -- (598,630) -- (597,625) -- (595,623) -- (595,620) -- (592,617) -- (592,615) -- (589,611) -- (589,608) -- (587,606) -- (582,595) -- (576,587) -- (576,585) -- (571,578) -- (570,574) -- (568,572) -- (567,569) -- (564,566) -- (564,564) -- (561,560) -- (559,553) -- (554,548) -- (552,541) -- (547,537) -- (546,533) -- (544,530) -- (544,528) -- (542,526) -- (542,524) -- (538,517) -- (538,515) -- (535,511) -- (535,508) -- (528,500) -- (528,497) -- (525,494) -- (522,488) -- (519,485) -- (515,475) -- (512,472) -- (511,465) -- (506,459) -- (506,456) -- (505,454) -- (501,450) -- (498,444) -- (493,438) -- (493,433) -- (488,429) -- (488,427) -- (485,424) -- (485,420) -- (483,418) -- (481,413) -- (479,411) -- (479,409) -- (477,405) -- (475,403) -- (474,403) -- (470,398) -- (470,395) -- (469,394) -- (468,390) -- (462,383) -- (461,377) -- (456,370) -- (455,366) -- (453,364) -- (450,358) -- (450,356) -- (447,353) -- (447,351) -- (445,347) -- (441,342) -- (441,340) -- (439,337) -- (439,333) -- (436,330) -- (435,327) -- (430,320) -- (428,316) -- (428,314) -- (424,310) -- (424,308) -- (418,301) -- (416,293) -- (414,291) -- (414,288) -- (408,280) -- (408,277) -- (406,275) -- (405,271) -- (402,267) -- (402,265) -- (399,262) -- (396,255) -- (392,251) -- (392,248) -- (390,246) -- (388,242) -- (388,240) -- (381,230) -- (381,227) -- (379,223) -- (373,216) -- (373,213) -- (368,207) -- (367,202) -- (365,200) -- (363,195) -- (359,191) -- (359,186) -- (357,184) -- (354,177) -- (349,171) -- (349,166) -- (345,162) -- (344,159) -- (339,154) -- (337,147) -- (333,142) -- (333,139) -- (329,135) -- (329,131) -- (327,127) -- (324,124) -- (324,121) -- (319,116) -- (318,112) -- (316,110) -- (315,107) -- (312,104) -- (309,97) -- (305,92) -- (304,88) -- (301,85) -- (301,82) -- (298,79) -- cycle
(96,652) -- (96,650) -- (103,646) -- (107,642) -- (110,641) -- (113,638) -- (116,638) -- (118,636) -- (121,636) -- (134,624) -- (136,624) -- (148,614) -- (150,614) -- (153,611) -- (157,609) -- (161,605) -- (162,605) -- (167,601) -- (169,601) -- (171,599) -- (172,599) -- (179,592) -- (182,591) -- (196,580) -- (200,578) -- (211,568) -- (216,567) -- (221,564) -- (237,551) -- (238,551) -- (251,541) -- (252,541) -- (258,535) -- (263,532) -- (267,528) -- (270,527) -- (278,519) -- (282,517) -- (286,513) -- (288,513) -- (291,510) -- (294,509) -- (305,500) -- (312,502) -- (315,504) -- (317,504) -- (322,507) -- (324,507) -- (328,509) -- (330,511) -- (335,513) -- (339,516) -- (341,516) -- (347,521) -- (354,524) -- (356,526) -- (361,528) -- (363,530) -- (370,532) -- (372,534) -- (378,537) -- (380,537) -- (382,539) -- (384,539) -- (387,542) -- (392,543) -- (395,546) -- (399,547) -- (405,552) -- (410,552) -- (411,553) -- (413,553) -- (415,555) -- (415,557) -- (423,561) -- (428,561) -- (430,562) -- (437,568) -- (441,569) -- (444,572) -- (446,572) -- (454,576) -- (456,578) -- (458,578) -- (461,581) -- (465,582) -- (467,584) -- (469,584) -- (474,588) -- (478,588) -- (481,591) -- (483,591) -- (485,593) -- (492,595) -- (495,598) -- (499,600) -- (501,600) -- (504,602) -- (507,602) -- (510,605) -- (516,606) -- (520,610) -- (528,612) -- (534,617) -- (542,621) -- (547,622) -- (550,625) -- (560,630) -- (563,633) -- (566,634) -- (568,636) -- (571,636) -- (575,639) -- (579,639) -- (583,642) -- (587,644) -- (589,644) -- (591,646) -- (593,646) -- (595,648) -- (593,650) -- (579,650) -- (577,649) -- (576,650) -- (566,650) -- (565,651) -- (563,651) -- (560,649) -- (557,649) -- (555,651) -- (550,651) -- (549,650) -- (543,650) -- (542,651) -- (501,651) -- (500,650) -- (498,651) -- (492,651) -- (491,650) -- (477,650) -- (476,651) -- (455,651) -- (454,650) -- (445,650) -- (444,651) -- (442,651) -- (441,650) -- (439,650) -- (438,651) -- (433,651) -- (432,650) -- (427,650) -- (426,651) -- (401,651) -- (400,650) -- (398,651) -- (396,650) -- (377,650) -- (376,651) -- (367,651) -- (366,650) -- (357,650) -- (356,651) -- (351,651) -- (350,650) -- (346,650) -- (345,651) -- (340,651) -- (339,650) -- (334,650) -- (333,651) -- (330,651) -- (329,650) -- (327,650) -- (326,651) -- (287,651) -- (286,650) -- (278,650) -- (277,651) -- (271,651) -- (270,652) -- (267,652) -- (266,651) -- (262,651) -- (261,650) -- (259,651) -- (244,651) -- (243,650) -- (240,650) -- (239,651) -- (233,651) -- (232,652) -- (230,651) -- (228,651) -- (227,652) -- (189,652) -- (188,651) -- (182,651) -- (180,652) -- (178,651) -- (175,653) -- (146,653) -- (145,652) -- (128,652) -- (127,653) -- (97,653) -- cycle
(293,103) -- (298,107) -- (298,109) -- (297,110) -- (297,126) -- (296,127) -- (296,135) -- (297,136) -- (296,137) -- (296,145) -- (297,146) -- (297,151) -- (296,152) -- (296,155) -- (298,158) -- (298,167) -- (297,168) -- (297,175) -- (298,176) -- (298,181) -- (297,182) -- (298,184) -- (298,193) -- (297,194) -- (297,197) -- (296,198) -- (297,200) -- (297,213) -- (298,214) -- (298,217) -- (297,218) -- (298,220) -- (298,223) -- (297,224) -- (298,226) -- (298,231) -- (297,232) -- (297,237) -- (298,238) -- (298,242) -- (297,243) -- (297,255) -- (298,256) -- (297,276) -- (298,278) -- (297,279) -- (297,287) -- (296,288) -- (296,290) -- (297,291) -- (297,304) -- (298,305) -- (298,309) -- (297,310) -- (297,312) -- (299,314) -- (299,316) -- (298,317) -- (298,324) -- (299,325) -- (299,327) -- (298,328) -- (298,330) -- (299,331) -- (299,336) -- (298,337) -- (298,343) -- (299,345) -- (298,346) -- (298,350) -- (299,351) -- (299,390) -- (300,392) -- (298,394) -- (298,398) -- (299,399) -- (299,404) -- (300,405) -- (300,423) -- (299,424) -- (299,430) -- (300,431) -- (300,450) -- (301,451) -- (300,454) -- (301,455) -- (301,460) -- (302,461) -- (302,475) -- (301,476) -- (301,478) -- (302,479) -- (302,486) -- (303,487) -- (303,490) -- (302,491) -- (302,497) -- (299,500) -- (296,501) -- (292,505) -- (287,508) -- (282,513) -- (280,513) -- (278,515) -- (277,515) -- (273,520) -- (272,520) -- (267,524) -- (262,526) -- (257,531) -- (254,532) -- (248,537) -- (246,537) -- (242,541) -- (240,541) -- (238,543) -- (237,543) -- (233,547) -- (228,550) -- (220,558) -- (219,558) -- (217,560) -- (213,561) -- (208,566) -- (207,566) -- (200,572) -- (197,573) -- (185,582) -- (183,582) -- (171,593) -- (170,593) -- (166,597) -- (155,604) -- (146,612) -- (138,616) -- (132,621) -- (130,621) -- (124,627) -- (121,628) -- (117,632) -- (111,635) -- (105,641) -- (103,641) -- (97,646) -- (96,646) -- (95,645) -- (95,641) -- (102,630) -- (103,623) -- (105,621) -- (105,618) -- (108,614) -- (108,611) -- (109,610) -- (109,607) -- (110,606) -- (110,600) -- (112,596) -- (114,594) -- (117,588) -- (118,582) -- (120,578) -- (122,576) -- (123,570) -- (124,569) -- (124,566) -- (125,565) -- (125,562) -- (128,558) -- (128,552) -- (134,541) -- (135,536) -- (137,534) -- (137,531) -- (139,528) -- (139,526) -- (141,524) -- (142,516) -- (147,507) -- (148,499) -- (151,494) -- (152,489) -- (155,484) -- (155,481) -- (157,478) -- (157,475) -- (158,474) -- (158,470) -- (160,468) -- (160,465) -- (164,458) -- (164,454) -- (165,453) -- (165,451) -- (167,449) -- (167,445) -- (168,444) -- (168,442) -- (170,440) -- (170,436) -- (173,433) -- (173,431) -- (175,428) -- (175,425) -- (177,422) -- (177,418) -- (181,410) -- (181,406) -- (183,403) -- (183,399) -- (188,389) -- (188,386) -- (192,381) -- (192,377) -- (195,373) -- (195,368) -- (197,366) -- (197,364) -- (200,357) -- (200,354) -- (203,349) -- (204,341) -- (206,339) -- (209,332) -- (209,329) -- (211,327) -- (213,316) -- (217,312) -- (217,310) -- (219,308) -- (218,304) -- (221,299) -- (225,285) -- (228,280) -- (228,277) -- (231,273) -- (231,269) -- (234,265) -- (234,263) -- (236,260) -- (236,258) -- (237,257) -- (237,251) -- (240,246) -- (242,238) -- (246,231) -- (246,227) -- (249,224) -- (249,220) -- (251,218) -- (253,211) -- (255,208) -- (255,204) -- (257,200) -- (257,197) -- (260,193) -- (261,185) -- (264,181) -- (265,176) -- (267,174) -- (268,168) -- (270,165) -- (272,154) -- (275,149) -- (275,145) -- (278,141) -- (278,138) -- (281,134) -- (282,128) -- (284,126) -- (284,124) -- (286,121) -- (286,118) -- (288,116) -- (289,110) -- (291,108) -- (291,106) -- cycle
(302,95) -- (304,97) -- (304,100) -- (306,104) -- (309,107) -- (312,114) -- (315,117) -- (315,120) -- (319,125) -- (319,127) -- (321,129) -- (321,131) -- (324,134) -- (324,137) -- (327,140) -- (330,148) -- (338,156) -- (339,162) -- (343,167) -- (343,169) -- (345,171) -- (346,175) -- (348,177) -- (348,179) -- (352,186) -- (352,188) -- (358,200) -- (358,203) -- (364,209) -- (368,218) -- (370,220) -- (370,223) -- (375,229) -- (376,233) -- (380,239) -- (380,242) -- (383,245) -- (384,248) -- (388,253) -- (391,259) -- (392,263) -- (396,267) -- (396,269) -- (400,276) -- (400,278) -- (403,281) -- (404,284) -- (406,286) -- (408,294) -- (410,296) -- (412,301) -- (418,308) -- (420,315) -- (423,318) -- (428,330) -- (431,333) -- (433,337) -- (433,339) -- (437,344) -- (437,347) -- (447,360) -- (447,362) -- (449,365) -- (450,371) -- (454,375) -- (454,377) -- (458,383) -- (458,386) -- (461,389) -- (462,392) -- (464,394) -- (466,401) -- (470,405) -- (470,408) -- (471,410) -- (476,415) -- (478,419) -- (478,421) -- (481,425) -- (481,427) -- (483,429) -- (483,431) -- (485,435) -- (487,437) -- (487,439) -- (490,445) -- (495,451) -- (498,458) -- (501,462) -- (501,464) -- (505,469) -- (505,471) -- (508,476) -- (508,478) -- (510,480) -- (510,482) -- (514,487) -- (514,489) -- (516,491) -- (518,496) -- (520,498) -- (520,500) -- (527,509) -- (527,511) -- (529,513) -- (529,516) -- (535,523) -- (535,525) -- (537,528) -- (537,531) -- (540,534) -- (540,536) -- (542,538) -- (543,542) -- (545,544) -- (545,545) -- (547,547) -- (547,549) -- (551,554) -- (556,564) -- (557,568) -- (561,573) -- (562,576) -- (565,579) -- (565,581) -- (567,585) -- (570,588) -- (570,590) -- (574,598) -- (578,602) -- (579,606) -- (585,614) -- (585,616) -- (587,618) -- (587,621) -- (590,624) -- (590,625) -- (592,627) -- (592,629) -- (596,633) -- (598,638) -- (600,640) -- (602,644) -- (600,646) -- (599,646) -- (596,643) -- (593,643) -- (587,640) -- (585,638) -- (577,634) -- (573,633) -- (570,630) -- (568,630) -- (564,628) -- (562,626) -- (554,623) -- (552,621) -- (545,618) -- (543,616) -- (537,613) -- (535,613) -- (529,608) -- (527,607) -- (523,607) -- (517,604) -- (515,602) -- (511,601) -- (506,597) -- (502,596) -- (500,594) -- (498,594) -- (495,591) -- (484,586) -- (481,583) -- (474,582) -- (461,573) -- (454,571) -- (447,565) -- (443,565) -- (436,560) -- (434,560) -- (429,557) -- (422,555) -- (414,550) -- (412,550) -- (410,548) -- (405,546) -- (402,543) -- (398,541) -- (396,541) -- (391,537) -- (389,537) -- (384,533) -- (378,530) -- (374,530) -- (369,526) -- (365,525) -- (360,521) -- (357,520) -- (354,517) -- (351,517) -- (345,512) -- (343,512) -- (336,507) -- (330,504) -- (322,502) -- (320,500) -- (318,500) -- (310,496) -- (306,492) -- (306,488) -- (307,487) -- (307,479) -- (306,478) -- (307,470) -- (306,469) -- (306,466) -- (307,465) -- (307,459) -- (308,458) -- (308,456) -- (305,453) -- (306,452) -- (306,445) -- (305,444) -- (305,441) -- (306,440) -- (306,431) -- (305,430) -- (305,426) -- (306,425) -- (306,418) -- (307,417) -- (307,415) -- (305,413) -- (306,411) -- (306,405) -- (305,404) -- (305,401) -- (306,399) -- (305,398) -- (305,392) -- (304,391) -- (305,389) -- (305,384) -- (304,382) -- (306,380) -- (305,378) -- (305,374) -- (306,373) -- (305,372) -- (305,363) -- (304,362) -- (304,355) -- (303,354) -- (303,350) -- (305,348) -- (306,348) -- (306,339) -- (304,339) -- (303,338) -- (303,333) -- (302,332) -- (303,330) -- (303,320) -- (304,319) -- (304,317) -- (303,316) -- (303,310) -- (304,308) -- (302,304) -- (302,298) -- (304,295) -- (304,293) -- (303,292) -- (303,287) -- (304,286) -- (303,284) -- (303,273) -- (302,272) -- (303,270) -- (303,261) -- (302,259) -- (304,256) -- (302,253) -- (302,244) -- (301,243) -- (301,240) -- (302,239) -- (302,224) -- (303,223) -- (303,221) -- (301,219) -- (302,217) -- (302,210) -- (303,209) -- (303,207) -- (301,205) -- (301,202) -- (304,198) -- (304,194) -- (303,194) -- (301,192) -- (301,156) -- (300,155) -- (301,142) -- (300,141) -- (300,138) -- (299,137) -- (300,133) -- (301,132) -- (300,130) -- (300,123) -- (299,122) -- (301,118) -- (299,98) -- cycle;
\node[inner sep=0pt,font=\fontsize{10}{10}\selectfont] at (294,42) {\scalebox{5.8}{$A$}};
\node[inner sep=0pt,font=\fontsize{10}{10}\selectfont] at (644,653) {\scalebox{5.8}{$B$}};
\node[inner sep=0pt,font=\fontsize{10}{10}\selectfont] at (55,667) {\scalebox{5.8}{$C$}};
\node[inner sep=0pt,font=\fontsize{10}{10}\selectfont] at (335,685) {\scalebox{5.8}{$a$}};
\node[inner sep=0pt,font=\fontsize{10}{10}\selectfont] at (158,330) {\scalebox{5.2}{$b$}};
\node[inner sep=0pt,font=\fontsize{10}{10}\selectfont] at (478,347) {\scalebox{5.2}{$c$}};
\node[inner sep=0pt,font=\fontsize{10}{10}\selectfont] at (327,337) {\scalebox{5.2}{$\alpha$}};
\node[inner sep=0pt,font=\fontsize{10}{10}\selectfont] at (408,585) {\scalebox{4.8}{$\beta$}};
\node[inner sep=0pt,font=\fontsize{10}{10}\selectfont] at (222,596) {\scalebox{5.2}{$\gamma$}};
\node[inner sep=0pt,font=\fontsize{10}{10}\selectfont] at (266,464) {\scalebox{5.4}{$\mathcal{B}$}};
\node[inner sep=0pt,font=\fontsize{10}{10}\selectfont] at (339,465) {\scalebox{5.4}{$\mathcal{C}$}};
\node[inner sep=0pt,font=\fontsize{10}{10}\selectfont] at (300,534) {\scalebox{5.4}{$\mathcal{A}$}};
\node[inner sep=0pt,font=\fontsize{10}{10}\selectfont] at (277,179) {\scalebox{4.2}{$1$}};
\node[inner sep=0pt,font=\fontsize{10}{10}\selectfont] at (322,180) {\scalebox{4.2}{$2$}};
\node[inner sep=0pt,font=\fontsize{10}{10}\selectfont] at (538,591) {\scalebox{4.2}{$3$}};
\node[inner sep=0pt,font=\fontsize{10}{10}\selectfont] at (520,630) {\scalebox{4.2}{$4$}};
\node[inner sep=0pt,font=\fontsize{10}{10}\selectfont] at (155,637) {\scalebox{4.2}{$5$}};
\node[inner sep=0pt,font=\fontsize{10}{10}\selectfont] at (138,586) {\scalebox{4.2}{$6$}};
\end{tikzpicture}%
}
\\Fig.\ I.\end{minipage}\hspace{.3cm}%
\begin{minipage}[b]{6.6cm}\centering\resizebox{6.4cm}{!}{%
\begin{tikzpicture}[x=1pt,y=-1pt]
\path[use as bounding box] (0,0) rectangle (889,633);
\path[fill=black,draw=none,even odd rule]
(764,501) -- (762,499) -- (762,496) -- (758,492) -- (756,488) -- (750,482) -- (750,481) -- (746,476) -- (745,473) -- (736,465) -- (735,461) -- (730,455) -- (727,448) -- (724,445) -- (722,441) -- (717,436) -- (717,435) -- (711,430) -- (710,425) -- (694,405) -- (694,402) -- (684,390) -- (684,388) -- (679,383) -- (679,382) -- (669,370) -- (669,367) -- (663,361) -- (659,353) -- (654,348) -- (649,340) -- (644,336) -- (644,334) -- (641,331) -- (640,328) -- (636,324) -- (634,320) -- (629,315) -- (629,314) -- (623,307) -- (621,303) -- (613,293) -- (612,290) -- (607,285) -- (606,282) -- (601,277) -- (601,275) -- (597,271) -- (597,270) -- (592,264) -- (591,260) -- (587,255) -- (586,252) -- (581,248) -- (581,245) -- (569,232) -- (569,230) -- (564,225) -- (563,221) -- (559,217) -- (556,211) -- (547,202) -- (547,200) -- (540,193) -- (539,190) -- (531,182) -- (531,179) -- (517,162) -- (517,161) -- (515,159) -- (515,157) -- (514,157) -- (512,155) -- (513,154) -- (509,150) -- (504,142) -- (500,139) -- (499,136) -- (495,132) -- (494,129) -- (490,125) -- (490,122) -- (484,116) -- (483,113) -- (480,110) -- (478,105) -- (474,101) -- (474,99) -- (472,97) -- (471,97) -- (467,100) -- (467,103) -- (461,109) -- (458,116) -- (453,122) -- (453,124) -- (451,126) -- (451,128) -- (448,131) -- (447,134) -- (442,141) -- (442,143) -- (438,148) -- (436,152) -- (436,154) -- (429,162) -- (428,165) -- (423,171) -- (423,173) -- (416,182) -- (415,185) -- (411,190) -- (410,193) -- (407,195) -- (407,197) -- (405,199) -- (404,201) -- (404,204) -- (402,206) -- (401,209) -- (395,216) -- (394,219) -- (386,229) -- (383,235) -- (378,240) -- (379,243) -- (373,250) -- (373,251) -- (371,253) -- (371,255) -- (368,258) -- (366,262) -- (362,266) -- (361,269) -- (357,274) -- (357,276) -- (355,280) -- (349,287) -- (344,299) -- (335,309) -- (334,313) -- (332,315) -- (320,334) -- (314,340) -- (313,346) -- (307,352) -- (302,362) -- (297,367) -- (297,369) -- (291,376) -- (288,382) -- (284,386) -- (284,388) -- (282,390) -- (280,395) -- (278,397) -- (274,404) -- (270,408) -- (266,415) -- (262,419) -- (261,422) -- (258,425) -- (254,433) -- (247,441) -- (247,443) -- (245,447) -- (241,452) -- (240,456) -- (234,462) -- (234,465) -- (230,468) -- (229,471) -- (217,486) -- (217,488) -- (214,491) -- (213,495) -- (208,502) -- (208,504) -- (205,507) -- (201,514) -- (191,526) -- (191,528) -- (185,535) -- (181,543) -- (178,546) -- (178,548) -- (175,551) -- (171,559) -- (168,562) -- (168,564) -- (164,568) -- (162,573) -- (159,576) -- (158,580) -- (150,589) -- (150,592) -- (152,594) -- (155,594) -- (161,591) -- (167,591) -- (169,592) -- (173,590) -- (198,589) -- (199,588) -- (208,587) -- (212,585) -- (231,584) -- (233,582) -- (249,582) -- (250,581) -- (253,581) -- (254,580) -- (261,580) -- (262,579) -- (268,579) -- (269,578) -- (277,578) -- (279,576) -- (283,577) -- (284,576) -- (303,575) -- (304,574) -- (307,574) -- (308,573) -- (313,573) -- (314,572) -- (322,572) -- (323,571) -- (326,571) -- (327,570) -- (334,570) -- (335,569) -- (341,569) -- (342,568) -- (351,568) -- (352,567) -- (357,567) -- (359,565) -- (362,565) -- (363,564) -- (374,564) -- (375,563) -- (382,563) -- (383,562) -- (388,562) -- (389,561) -- (398,561) -- (401,559) -- (407,559) -- (408,558) -- (412,558) -- (413,557) -- (417,557) -- (418,556) -- (430,555) -- (431,554) -- (437,554) -- (438,553) -- (444,553) -- (445,552) -- (451,552) -- (452,551) -- (458,551) -- (459,550) -- (465,550) -- (466,549) -- (477,549) -- (478,548) -- (486,548) -- (487,547) -- (493,547) -- (494,546) -- (503,546) -- (507,544) -- (513,545) -- (514,544) -- (535,543) -- (536,542) -- (540,542) -- (541,541) -- (545,541) -- (546,540) -- (551,540) -- (557,537) -- (570,537) -- (574,535) -- (579,535) -- (583,533) -- (585,534) -- (595,534) -- (596,533) -- (599,533) -- (601,531) -- (605,532) -- (607,530) -- (613,530) -- (614,529) -- (617,529) -- (619,530) -- (620,529) -- (624,529) -- (625,528) -- (635,528) -- (636,527) -- (640,527) -- (641,526) -- (655,525) -- (658,523) -- (665,523) -- (668,521) -- (679,521) -- (682,519) -- (684,520) -- (688,520) -- (689,519) -- (691,520) -- (694,520) -- (695,519) -- (702,519) -- (706,516) -- (717,516) -- (718,515) -- (730,514) -- (731,513) -- (735,513) -- (736,512) -- (752,511) -- (757,508) -- (763,508) -- (763,504) -- (764,503) -- cycle
(755,504) -- (753,506) -- (749,506) -- (748,507) -- (744,507) -- (743,508) -- (735,508) -- (734,509) -- (724,509) -- (723,510) -- (716,510) -- (715,511) -- (712,511) -- (711,512) -- (706,512) -- (705,513) -- (700,513) -- (699,514) -- (694,514) -- (693,515) -- (685,515) -- (684,516) -- (679,516) -- (678,517) -- (667,517) -- (666,518) -- (656,518) -- (653,520) -- (647,520) -- (646,521) -- (641,521) -- (640,522) -- (630,522) -- (629,523) -- (624,523) -- (623,524) -- (619,523) -- (615,525) -- (610,525) -- (606,527) -- (599,527) -- (598,528) -- (583,528) -- (580,530) -- (577,530) -- (576,531) -- (571,531) -- (569,530) -- (565,532) -- (557,532) -- (556,533) -- (553,533) -- (552,534) -- (545,534) -- (544,535) -- (532,535) -- (531,536) -- (528,536) -- (527,537) -- (518,537) -- (517,538) -- (511,538) -- (510,539) -- (507,539) -- (504,541) -- (493,541) -- (492,542) -- (489,542) -- (488,543) -- (484,543) -- (483,542) -- (481,543) -- (475,543) -- (474,544) -- (470,544) -- (469,545) -- (458,546) -- (456,548) -- (454,548) -- (453,549) -- (440,549) -- (439,550) -- (434,550) -- (433,551) -- (414,552) -- (413,553) -- (408,553) -- (407,554) -- (398,554) -- (397,555) -- (391,555) -- (390,556) -- (384,556) -- (383,557) -- (362,558) -- (361,559) -- (355,560) -- (352,562) -- (339,562) -- (338,563) -- (331,563) -- (330,564) -- (325,564) -- (324,565) -- (320,564) -- (319,565) -- (308,566) -- (304,568) -- (296,569) -- (295,570) -- (288,570) -- (287,571) -- (279,571) -- (278,572) -- (267,572) -- (266,573) -- (261,573) -- (260,574) -- (254,575) -- (253,576) -- (250,576) -- (249,577) -- (232,578) -- (231,579) -- (228,578) -- (224,580) -- (214,580) -- (213,581) -- (210,580) -- (207,582) -- (204,582) -- (203,583) -- (193,583) -- (192,584) -- (178,585) -- (177,586) -- (170,586) -- (169,587) -- (164,587) -- (163,586) -- (163,584) -- (164,583) -- (167,582) -- (172,578) -- (175,577) -- (181,572) -- (187,569) -- (190,566) -- (198,562) -- (216,547) -- (222,544) -- (227,539) -- (228,539) -- (230,537) -- (232,537) -- (245,527) -- (252,524) -- (266,512) -- (268,512) -- (272,510) -- (279,504) -- (287,500) -- (299,489) -- (300,489) -- (303,486) -- (308,483) -- (312,478) -- (319,474) -- (322,471) -- (325,470) -- (328,467) -- (332,465) -- (335,462) -- (338,461) -- (341,458) -- (346,456) -- (349,453) -- (355,450) -- (359,446) -- (370,439) -- (374,435) -- (377,434) -- (381,430) -- (385,429) -- (391,423) -- (396,421) -- (409,409) -- (413,407) -- (419,402) -- (423,401) -- (425,399) -- (426,399) -- (431,394) -- (434,393) -- (444,384) -- (452,380) -- (458,374) -- (462,372) -- (471,365) -- (473,365) -- (476,362) -- (486,366) -- (488,368) -- (493,370) -- (495,372) -- (497,372) -- (501,376) -- (503,376) -- (508,379) -- (512,380) -- (521,387) -- (523,387) -- (531,392) -- (534,392) -- (538,394) -- (543,398) -- (547,399) -- (549,401) -- (553,403) -- (556,402) -- (558,404) -- (560,404) -- (564,408) -- (567,409) -- (569,411) -- (571,411) -- (574,413) -- (577,413) -- (582,417) -- (586,418) -- (589,421) -- (595,423) -- (598,426) -- (602,426) -- (605,428) -- (607,428) -- (613,431) -- (616,434) -- (620,435) -- (624,438) -- (627,438) -- (629,439) -- (631,441) -- (645,448) -- (647,450) -- (650,450) -- (653,453) -- (658,454) -- (664,458) -- (671,460) -- (673,462) -- (677,463) -- (681,466) -- (684,466) -- (686,468) -- (691,470) -- (693,472) -- (696,472) -- (699,475) -- (704,478) -- (709,479) -- (711,481) -- (714,481) -- (722,487) -- (724,487) -- (731,491) -- (733,491) -- (738,495) -- (742,497) -- (744,497) -- (746,499) -- (748,499) -- (750,501) -- (752,501) -- cycle
(466,114) -- (468,116) -- (468,121) -- (470,123) -- (470,127) -- (469,128) -- (469,140) -- (470,141) -- (470,145) -- (471,146) -- (471,151) -- (470,152) -- (471,154) -- (471,157) -- (470,158) -- (469,165) -- (470,166) -- (470,168) -- (472,170) -- (472,171) -- (470,173) -- (470,178) -- (469,180) -- (470,181) -- (470,190) -- (471,191) -- (471,206) -- (472,207) -- (471,209) -- (471,220) -- (470,221) -- (471,225) -- (470,226) -- (470,228) -- (471,229) -- (471,244) -- (472,245) -- (472,253) -- (471,254) -- (471,261) -- (472,262) -- (472,271) -- (471,272) -- (471,274) -- (472,275) -- (472,283) -- (471,284) -- (471,289) -- (473,291) -- (473,296) -- (472,297) -- (472,301) -- (473,302) -- (473,307) -- (472,308) -- (472,310) -- (473,311) -- (473,323) -- (472,324) -- (473,336) -- (472,337) -- (472,339) -- (474,342) -- (474,344) -- (473,345) -- (473,358) -- (467,364) -- (460,368) -- (452,375) -- (449,376) -- (446,379) -- (438,384) -- (420,398) -- (417,399) -- (406,407) -- (400,410) -- (391,418) -- (389,418) -- (387,420) -- (386,420) -- (381,425) -- (379,425) -- (371,431) -- (367,432) -- (362,438) -- (359,439) -- (346,451) -- (342,453) -- (340,453) -- (338,455) -- (337,455) -- (334,458) -- (329,461) -- (324,466) -- (318,469) -- (314,473) -- (309,475) -- (304,480) -- (300,482) -- (293,488) -- (291,488) -- (285,493) -- (278,497) -- (274,501) -- (266,505) -- (248,520) -- (244,521) -- (238,527) -- (236,527) -- (232,531) -- (229,532) -- (225,536) -- (222,537) -- (217,542) -- (216,542) -- (202,554) -- (193,559) -- (186,565) -- (179,569) -- (171,576) -- (167,578) -- (165,578) -- (164,577) -- (167,574) -- (171,567) -- (175,563) -- (175,561) -- (177,559) -- (177,558) -- (184,550) -- (185,546) -- (188,543) -- (188,540) -- (195,532) -- (197,528) -- (201,524) -- (201,522) -- (205,518) -- (205,515) -- (210,511) -- (209,508) -- (219,497) -- (219,496) -- (221,494) -- (221,492) -- (227,485) -- (230,478) -- (236,472) -- (237,469) -- (242,463) -- (243,461) -- (243,458) -- (246,455) -- (249,455) -- (251,453) -- (251,449) -- (254,446) -- (256,442) -- (256,439) -- (260,435) -- (260,434) -- (264,429) -- (265,426) -- (271,419) -- (272,416) -- (277,409) -- (282,399) -- (289,391) -- (289,389) -- (292,386) -- (293,382) -- (300,374) -- (305,364) -- (311,358) -- (311,355) -- (312,353) -- (315,350) -- (317,350) -- (318,349) -- (318,346) -- (323,339) -- (324,335) -- (328,331) -- (331,325) -- (336,320) -- (338,316) -- (338,314) -- (343,309) -- (347,302) -- (352,297) -- (352,294) -- (354,292) -- (354,290) -- (356,288) -- (356,286) -- (358,284) -- (360,280) -- (364,276) -- (364,274) -- (371,265) -- (372,262) -- (375,258) -- (375,256) -- (380,250) -- (381,247) -- (384,244) -- (385,239) -- (392,231) -- (393,229) -- (393,226) -- (397,223) -- (400,217) -- (402,215) -- (404,211) -- (404,209) -- (408,206) -- (409,203) -- (411,201) -- (414,195) -- (415,190) -- (419,186) -- (420,183) -- (424,178) -- (425,175) -- (428,172) -- (434,162) -- (437,159) -- (437,157) -- (439,155) -- (442,149) -- (447,143) -- (447,140) -- (451,135) -- (455,127) -- (460,121) -- (461,118) -- (465,114) -- cycle
(473,111) -- (475,111) -- (477,113) -- (477,115) -- (484,122) -- (485,125) -- (487,127) -- (488,130) -- (491,133) -- (492,136) -- (500,144) -- (506,154) -- (514,163) -- (514,165) -- (517,168) -- (518,171) -- (523,177) -- (525,182) -- (535,193) -- (536,196) -- (540,200) -- (541,204) -- (547,211) -- (548,214) -- (553,219) -- (554,222) -- (559,228) -- (561,233) -- (573,246) -- (575,251) -- (578,254) -- (581,260) -- (588,268) -- (591,274) -- (597,280) -- (598,284) -- (602,288) -- (602,289) -- (603,289) -- (604,291) -- (610,296) -- (610,298) -- (616,305) -- (617,308) -- (622,314) -- (623,317) -- (629,324) -- (630,327) -- (634,331) -- (634,332) -- (636,334) -- (636,336) -- (642,341) -- (642,343) -- (645,346) -- (645,348) -- (654,358) -- (659,366) -- (665,372) -- (665,374) -- (668,378) -- (669,381) -- (671,382) -- (674,385) -- (674,386) -- (679,390) -- (679,392) -- (689,405) -- (690,408) -- (694,412) -- (695,415) -- (705,427) -- (705,428) -- (707,430) -- (707,432) -- (714,439) -- (715,442) -- (718,445) -- (718,447) -- (724,453) -- (726,458) -- (735,468) -- (735,472) -- (741,478) -- (742,481) -- (755,495) -- (756,497) -- (754,499) -- (750,497) -- (748,495) -- (741,492) -- (739,490) -- (737,490) -- (735,488) -- (733,488) -- (729,486) -- (727,484) -- (716,479) -- (711,475) -- (703,473) -- (698,469) -- (694,467) -- (692,467) -- (682,461) -- (678,460) -- (673,456) -- (667,453) -- (665,453) -- (663,451) -- (660,450) -- (658,448) -- (652,445) -- (647,444) -- (642,440) -- (636,439) -- (631,435) -- (629,435) -- (621,431) -- (619,429) -- (617,429) -- (610,424) -- (597,419) -- (593,415) -- (591,415) -- (588,412) -- (583,411) -- (581,409) -- (579,409) -- (575,406) -- (571,405) -- (569,403) -- (567,403) -- (565,401) -- (559,399) -- (556,396) -- (553,396) -- (548,393) -- (545,393) -- (543,391) -- (535,388) -- (532,385) -- (530,385) -- (524,382) -- (518,377) -- (516,376) -- (513,376) -- (505,371) -- (502,371) -- (496,368) -- (494,366) -- (492,366) -- (490,364) -- (481,360) -- (478,357) -- (477,355) -- (477,343) -- (478,342) -- (477,341) -- (477,332) -- (476,331) -- (477,329) -- (477,320) -- (478,319) -- (478,314) -- (477,313) -- (477,310) -- (478,309) -- (477,307) -- (477,298) -- (478,297) -- (478,293) -- (476,290) -- (476,283) -- (478,282) -- (478,279) -- (476,278) -- (476,248) -- (477,247) -- (477,242) -- (476,241) -- (476,232) -- (474,230) -- (477,224) -- (477,220) -- (475,216) -- (475,209) -- (476,208) -- (476,206) -- (475,205) -- (475,201) -- (476,200) -- (476,195) -- (475,194) -- (475,191) -- (474,190) -- (474,183) -- (473,182) -- (473,180) -- (474,179) -- (474,154) -- (475,153) -- (475,151) -- (474,150) -- (474,141) -- (473,140) -- (473,138) -- (474,137) -- (474,133) -- (472,130) -- (472,128) -- (473,127) -- (473,117) -- (472,116) -- (472,112) -- cycle;
\node[inner sep=0pt,font=\fontsize{10}{10}\selectfont] at (466,50) {\scalebox{5.6}{$\mathcal{A}-A$}};
\node[inner sep=0pt,font=\fontsize{10}{10}\selectfont] at (821,514) {\scalebox{5.6}{\scalebox{.72}[1]{$\mathcal{B}-B$}}};
\node[inner sep=0pt,font=\fontsize{10}{10}\selectfont] at (66,600) {\scalebox{5.6}{$\mathcal{C}-C$}};
\node[inner sep=0pt,font=\fontsize{10}{10}\selectfont] at (468,569) {\scalebox{5.6}{$a\alpha$}};
\node[inner sep=0pt,font=\fontsize{10}{10}\selectfont] at (288,302) {\scalebox{5.6}{$b\beta$}};
\node[inner sep=0pt,font=\fontsize{10}{10}\selectfont] at (650,287) {\scalebox{5.6}{$c\gamma$}};
\node[inner sep=0pt,font=\fontsize{10}{10}\selectfont] at (517,263) {\scalebox{4.8}{$\beta\gamma$}};
\node[inner sep=0pt,font=\fontsize{10}{10}\selectfont] at (355,480) {\scalebox{5.2}{$\alpha\beta$}};
\node[inner sep=0pt,font=\fontsize{10}{10}\selectfont] at (581,456) {\scalebox{5.2}{$\alpha\gamma$}};
\node[inner sep=0pt,font=\fontsize{10}{10}\selectfont] at (433,335) {\scalebox{5.4}{$\mathcal{B}$}};
\node[inner sep=0pt,font=\fontsize{10}{10}\selectfont] at (508,337) {\scalebox{5.4}{$\mathcal{C}$}};
\node[inner sep=0pt,font=\fontsize{10}{10}\selectfont] at (480,397) {\scalebox{5.4}{$\mathcal{A}$}};
\node[inner sep=0pt,font=\fontsize{10}{10}\selectfont] at (446,177) {\scalebox{4.2}{$6$}};
\node[inner sep=0pt,font=\fontsize{10}{10}\selectfont] at (496,177) {\scalebox{4.2}{$3$}};
\node[inner sep=0pt,font=\fontsize{10}{10}\selectfont] at (692,441) {\scalebox{4.2}{$2$}};
\node[inner sep=0pt,font=\fontsize{10}{10}\selectfont] at (684,492) {\scalebox{4.2}{$5$}};
\node[inner sep=0pt,font=\fontsize{10}{10}\selectfont] at (215,521) {\scalebox{4.2}{$1$}};
\node[inner sep=0pt,font=\fontsize{10}{10}\selectfont] at (236,561) {\scalebox{4.2}{$4$}};
\end{tikzpicture}%
}
\\Fig.\ II.\end{minipage}
\end{center}

Man hat

\begin{equation}
1+2+3+4+5+6=180^0
\tag{1}
\end{equation}

\begin{equation}
\sin 1 \, . \, \sin 3 \, . \, \sin 5 = \sin 2 \, . \, \sin 4 \, . \, \sin 6.
\tag{2}
\end{equation}

Aus (1) folgt

\[
\sin(2+3+5) = \sin(1+4+6)
\]

oder

\[
\sin 1 \, . \, \sin 2 \left\{\begin{array}{l} \sin 2 \cos 3 \cos 5 \\ + \sin 3 \cos 2 \cos 5 \\ + \sin 5 \cos 2 \cos 3 \\ - \sin 2 \sin 3 \sin 5 \end{array}\right\}
= \sin 1 \, . \, \sin 2 \left\{\begin{array}{l} \sin 1 \cos 4 \cos 6 \\ + \sin 4 \cos 1 \cos 6 \\ + \sin 6 \cos 1 \cos 4 \\ - \sin 1 \sin 4 \sin 6 \end{array}\right\}
\]

oder indem man $\sin 1 \, . \, \sin 3 \, . \, \sin 5 = \sin 2 \, . \, \sin 4 \, . \, \sin 6$ auf beiden Seiten hinzu addirt:

\[
\sin 1 \left\{\begin{array}{l} \sin 2^2 \cos 3 \cos 5 \\ + \sin 2 \, \sin 3 \cos 2 \cos 5 \\ + \sin 2 \, \sin 5 \cos 2 \cos 3 \\ + \cos 2^2 \sin 3 \sin 5 \end{array}\right\}
= \sin 2 \left\{\begin{array}{l} \sin 1^2 \cos 4 \cos 5 \\ + \sin 1 \, \sin 4 \cos 1 \cos 6 \\ + \sin 1 \, \sin 6 \cos 1 \cos 4 \\ + \cos 1^2 \sin 4 \sin 6 \end{array}\right\}
\]

%
\SourceLeaf{331}

oder

\begin{equation}
\sin 1 \, . \, \sin(2+3)\sin(2+5) = \sin 2 \, . \, \sin(1+4)\sin(1+6).
\tag{3}
\end{equation}

Auf gleiche Weise oder auch bloss durch Permutation

\begin{equation}
\sin 1 \, . \, \sin(4+3)\sin(4+5) = \sin 4 \, . \, \sin(1+2)\sin(1+6)
\tag{4}
\end{equation}

\begin{equation}
\sin 1 \, . \, \sin(6+3)\sin(6+5) = \sin 6 \, . \, \sin(1+2)\sin(1+4)
\tag{5}
\end{equation}

\begin{equation}
\sin 3 \, . \, \sin(2+1)\sin(2+5) = \sin 2 \, . \, \sin(3+4)\sin(3+6)
\tag{6}
\end{equation}

\begin{equation}
\sin 3 \, . \, \sin(4+1)\sin(4+5) = \sin 4 \, . \, \sin(3+2)\sin(3+6)
\tag{7}
\end{equation}

\begin{equation}
\sin 3 \, . \, \sin(6+1)\sin(6+5) = \sin 6 \, . \, \sin(3+2)\sin(3+4)
\tag{8}
\end{equation}

\begin{equation}
\sin 5 \, . \, \sin(2+1)\sin(2+3) = \sin 2 \, . \, \sin(5+4)\sin(5+6)
\tag{9}
\end{equation}

\begin{equation}
\sin 5 \, . \, \sin(4+1)\sin(4+3) = \sin 4 \, . \, \sin(5+2)\sin(5+6)
\tag{10}
\end{equation}

\begin{equation}
\sin 5 \, . \, \sin(6+1)\sin(6+3) = \sin 6 \, . \, \sin(5+2)\sin(5+4).
\tag{11}
\end{equation}

\Needspace{5\baselineskip}Diese neun Gleichungen können auch so geschrieben werden:

\begin{equation}
\sin 1 \, . \, \sin\mathfrak{C} \, \sin(\mathfrak{B}-B) = \sin 2 \, . \, \sin(\mathfrak{C}-C)\sin\mathfrak{B}
\tag{3}
\end{equation}

\begin{equation}
\sin 1 \, . \, \sin B \, \sin\mathfrak{A} = \sin 4 \, . \, \sin A \, \sin\mathfrak{B}
\tag{4}
\end{equation}

\begin{equation}
\sin 1 \, . \, \sin(\mathfrak{A}-A)\sin C = \sin 6 \, . \, \sin A \, \sin(\mathfrak{C}-C)
\tag{5}
\end{equation}

\begin{equation}
\sin 3 \, . \, \sin A \, \sin(\mathfrak{B}-B) = \sin 2 \, . \, \sin B \, \sin(\mathfrak{A}-A)
\tag{6}
\end{equation}

\begin{equation}
\sin 3 \, . \, \sin(\mathfrak{C}-C)\sin\mathfrak{A} = \sin 4 \, . \, \sin\mathfrak{C} \, \sin(\mathfrak{A}-A)
\tag{7}
\end{equation}

\begin{equation}
\sin 3 \, . \, \sin\mathfrak{B} \, \sin C = \sin 6 \, . \, \sin\mathfrak{C} \, \sin B
\tag{8}
\end{equation}

\begin{equation}
\sin 5 \, . \, \sin A \, \sin\mathfrak{C} = \sin 2 \, . \, \sin\mathfrak{A} \, \sin C
\tag{9}
\end{equation}

\begin{equation}
\sin 5 \, . \, \sin(\mathfrak{C}-C)\sin B = \sin 4 \, . \, \sin(\mathfrak{B}-B)\sin C
\tag{10}
\end{equation}

\begin{equation}
\sin 5 \, . \, \sin\mathfrak{B} \, \sin(\mathfrak{A}-A) = \sin 6 \, . \, \sin(\mathfrak{B}-B)\sin\mathfrak{A},
\tag{11}
\end{equation}

indem

\[
A=1+2, \qquad B=3+4, \qquad C=5+6,
\]
\[
180^0-\mathfrak{A}=4+5, \qquad 180^0-\mathfrak{B}=6+1, \qquad 180^0-\mathfrak{C}=2+3,
\]
\[
\mathfrak{A}-A=3+6, \qquad \mathfrak{B}-B=2+5, \qquad \mathfrak{C}-C=1+4
\]
ist.

\Needspace{9\baselineskip}Die Ableitung kann auch aus Betrachtung der Figur geschehen, nemlich
\begin{center}(1) und (2) nach Gefallen aus I oder II\end{center}
%
\SourceLeaf{332}

\[
(4), (8), (9) \text{ aus I}
\]
\[
(3), (7), (11) \text{ aus II}
\]
\[
(5), (6), (10) \text{ aus I und II verbunden.}
\]

Jede der Gleichungen (3)—(11) kann zur Auflösung des \textsc{Pothenot}schen Problems dienen, da die zwei Unbekannten, die in jeder Gleichung vorkommen, eine bekannte Summe haben.

Aus Figur II ergeben sich die Viereckgleichungen

\[
\frac{\sin(\mathfrak{A}-A)}{a\alpha} = \frac{\sin(\mathfrak{B}-B)}{b\beta} = \frac{\sin(\mathfrak{C}-C)}{c\gamma} = \frac{rr-\rho\rho}{2\alpha\beta\gamma r}
\]

von selbst, wenn $r$ Halbmesser des Kreises durch $A$, $B$, $C$, $\rho$ Distanz des vierten Punkts vom Mittelpunkte dieses Kreises; und die Gleichheit der Zeichen von

\[
\sin(\mathfrak{A}-A), \qquad \sin(\mathfrak{B}-B), \qquad \sin(\mathfrak{C}-C)
\]

ist das Criterium der physischen Möglichkeit der Aufgabe.

Die geometrische Bedeutung der Formeln (2)—(11) kann sehr zierlich an die fünf Kreise geknüpft werden, die resp.\ durch je drei Punkte von Fig.\ I oder von Fig.\ II beschrieben werden und deren Durchmesser $p$, $q$, $r$, $s$, $t$ sein mögen. Es ist dann

\[
\sin A = \frac{a}{p}, \qquad \sin B = \frac{b}{p}, \qquad \sin C = \frac{c}{p},
\]
\[
\sin\mathfrak{A} = \frac{a}{q}, \qquad \sin 4 = \frac{\gamma}{q}, \qquad \sin 5 = \frac{\beta}{q},
\]
\[
\sin\mathfrak{B} = \frac{b}{r}, \qquad \sin 6 = \frac{\alpha}{r}, \qquad \sin 1 = \frac{\gamma}{r},
\]
\[
\sin\mathfrak{C} = \frac{c}{s}, \qquad \sin 2 = \frac{\beta}{s}, \qquad \sin 3 = \frac{\alpha}{s},
\]
\[
\sin(\mathfrak{A}-A) = \frac{a\alpha}{t}, \qquad \sin(\mathfrak{B}-B) = \frac{b\beta}{t}, \qquad \sin(\mathfrak{C}-C) = \frac{c\gamma}{t}.
\]

Durch Substitution dieser Ausdrücke für die 15 Sinus in den 10 Gleichungen (2)—(11) werden diese zu identischen.

\Needspace{7\baselineskip}Ganz auf dieselbe Weise, wie aus Gleichung (2) die übrigen (3)—(11) identisch folgen, lassen sich auch aus jeder von diesen alle übrigen 9 ableiten. Man braucht z.\,B. nur die Gleichung (3) in die Form zu setzen
%
\SourceLeaf{333}

\[
\sin(-1)\sin(2+3)\sin(2+5) = \sin(-2)\sin(1+4)\sin(1+6),
\]

wo die Grössen

\[
-1, \qquad 2+3, \qquad 2+5; \qquad -2, \qquad 1+4, \qquad 1+6
\]

denselben Bedingungen entsprechen, welche für die Grössen

\[
+1, \qquad 3, \qquad 5; \qquad +2, \qquad 4, \qquad 6
\]

durch die Gleichungen (1), (2) ausgedrückt waren, daher auch alle Folgerungen sich übertragen lassen.

\begin{center}
[6.]

[\emph{Handbuch}, S.\ 278 bis 281.]

\textbf{NOCH ZUR BEURTHEILUNG DER PHYSISCHEN MÖGLICHKEIT DER DATEN IN POTHENOTS AUFGABE.}
\end{center}

$a$, $b$, $c$, $A$, $B$, $C$, $\alpha$, $\beta$, $\gamma$, $\mathfrak{A}$, $\mathfrak{B}$, $\mathfrak{C}$ haben dieselbe Bedeutung wie oben. Nemlich

\begin{center}\begin{tabular}{@{}lp{9.8cm}@{}}
$a,b,c$ & Seiten;\\
$A,B,C$ & Winkel des aus den gegebenen Punkten (1), (2), (3) gebildeten Dreiecks, dessen Radius $=r$;\\
$\alpha,\beta,\gamma,\rho$ & Entfernungen des gesuchten Punkts (0) von $A,B,C$ und dem Mittelpunkt jenes Kreises;\\
$\mathfrak{A}$ & Winkel von $02$ nach $03$, eben so gemessen wie $A$ von $12$ nach $13$ (d.\,i. beide von Rechts nach Links oder beide umgekehrt), ebenso $\mathfrak{B},\mathfrak{C}$.
\end{tabular}\end{center}

Es ist dann, wie schon p.\ 270 [des Handbuchs] bemerkt ist:

\[
\frac{\sin(\mathfrak{A}-A)}{a\alpha} = \frac{\sin(\mathfrak{B}-B)}{b\beta} = \frac{\sin(\mathfrak{C}-C)}{c\gamma} = \frac{rr-\rho\rho}{2\alpha\beta\gamma r},
\]

woraus folgt, dass, je nachdem (0) innerhalb oder ausserhalb jenes Kreises liegt,

\[
\sin(\mathfrak{A}-A), \qquad \sin(\mathfrak{B}-B), \qquad \sin(\mathfrak{C}-C)
\]

alle positiv oder negativ werden.
%
\SourceLeaf{334}

\Needspace{4\baselineskip}Hieraus ist klar:

I. Haben diese drei Sinus nicht alle gleiche Zeichen, so ist die Aufgabe physisch unmöglich.

In diesem Falle wird also der Schnitt der beiden Kreise, die auf $b$ und $c$ mit den Winkeln $\mathfrak{B}, \mathfrak{C}$ beschrieben sind, nicht in den vorgeschriebenen Abschnitten liegen, und folglich in keinem von beiden, wenn
\[ \sin(\mathfrak{B}-B), \qquad \sin(\mathfrak{C}-C) \]
unter sich gleiche Zeichen haben, aber
\[ \sin(\mathfrak{A}-A) \]
das entgegengesetzte. Man wird also anstatt $\mathfrak{B}$, $\mathfrak{B}+180^0$ und anstatt $\mathfrak{C}$, $\mathfrak{C}+180^0$ an dem Schnittpunkte vorfinden.

II. Jene Bedingung ist aber zur Möglichkeit zureichend. Diess wird so bewiesen:

1) Sind die drei Sinus positiv, so fällt der Schnitt der Kreise auf $b$ und $c$ innerhalb des Kreises durch 1, 2, 3.

\textls[120]{Beweis.} Gesetzt, er fiele ausserhalb, so müsste daselbst erscheinen $\mathfrak{B}' = \mathfrak{B}+180^0$ anstatt $\mathfrak{B}$ und $\mathfrak{C}' = \mathfrak{C}+180^0$ anstatt $\mathfrak{C}$ (damit $\sin(\mathfrak{B}'-B)$ negativ werde und eben so $\sin(\mathfrak{C}'-C)$). Es bleibt aber daselbst der Winkel auf $a$
\[ =360^0-\mathfrak{B}'-\mathfrak{C}'=360^0-\mathfrak{B}-\mathfrak{C}=\mathfrak{A}. \]

2) Eben so wird für den Fall, wo die drei Sinus negativ sind, die Absurdität der Voraussetzung, der Schnitt zweier Hülfskreise falle innerhalb des Hauptkreises, bewiesen.

Die Pointe ist, dass $\mathfrak{A}, \mathfrak{B}, \mathfrak{C}$ nothwendig die Summe $=360^0$ geben müssen (d. i. durch $360^0$ divisibel). Die vorgegebenen Werthe haben diese Eigenschaft von selbst. Sind solche unverträglich, so müssen \textls[120]{zwei} abgeändert werden (um $180^0$), damit jene Eigenschaft conservirt bleibe.

Sind $b, c$ die beiden Seiten, auf denen die Hülfskreise errichtet werden, so werden die Bedingungen sich so aussprechen lassen:
\[ \sin(\mathfrak{B}-B), \quad \sin(\mathfrak{C}-C) \quad \text{und} \quad \sin(\mathfrak{B}-B+\mathfrak{C}-C) \]
sollen gleiche Zeichen haben.
%
\SourceLeaf{335}

Sind Werthe von $\mathfrak{B}, \mathfrak{C}$ vorgegeben, die diesen Bedingungen, oder einer von ihnen, nicht genügen, so steht fest, dass diese Aufgabe eine physische Unmöglichkeit implicirt.

Umgekehrt, genügen die aufgegebenen Werthe den Bedingungen, so ist die Aufgabe möglich, d. i. die beiden Hülfskreise treffen sich in einem Punkte, der für beide Kreise im geeigneten Abschnitte liegt. Sind nemlich 1) die Zeichen jener Sinus $+\;+\;+$, so kann jener Durchschnitt nicht ausserhalb des Hauptkreises liegen, wo die ersten beiden Sinus negativ, der dritte positiv sein müsste; die beiden ersten Bedingungen würden nemlich erfordern, dass $\mathfrak{B}+180^0$, $\mathfrak{C}+180^0$ an die Stelle von $\mathfrak{B}, \mathfrak{C}$ treten, was aber wieder der dritten Bedingung widerspricht. Also liegt jener Durchschnitt innerhalb des Hauptkreises, wo eben die geeigneten Abschnitte liegen.

Auf ähnliche Weise wird 2) für die Zeichen $-\;-\;+$ bewiesen, dass der Durchschnitt nur ausserhalb liegen kann, wo dann die geeigneten Abschnitte liegen.

\begin{center}
\begin{tabular}{c@{\hspace{2em}}cc}
Sind die Zeichen & \multicolumn{2}{c}{so liegt der Durchschnitt} \\
 & im geeigneten Abschnitt & im ungeeigneten \\
$+\;+\;+$ & Kreis auf $a$ & [Kreis auf] $b, c$ \\
$+\;-\;+$ & $b$ & $a, c$ \\
$+\;-\;-$ & $c$ & $a, b$ \\
$-\;-\;-$ & $a$ & $b, c$ \\
$-\;+\;+$ & $c$ & $a, b$ \\
$-\;+\;-$ & $b$ & $a, c$ \\
\end{tabular}
\end{center}

Die elementarische Construction für das \textsc{Pothenot}sche Problem besteht darin, dass man Punkte in der durch (0) gehenden gegen 10 normalen Geraden bestimmt. Von einem solchen sind die Coordinaten polar relativ gegen (0):
\[ \text{Distanz } \frac{(12)}{\sin(02-01)}, \qquad \text{Richtung } 12+90^0-(02-01). \]

Eine andere neue Constructionsart ist die, dass man zunächst nicht den Punkt (0) sucht, sondern einen Punkt $(*)$, der in der Geraden 10 liegt und in der Distanz $(1*)=\frac{1}{(10)}$. Dieser Punkt wird durch Schnitte gerader Linien bestimmt, eine weniger, als von (0) aus Objecte geschnitten sind. Eine dieser
%
\SourceLeaf{336}
geraden Linien geht durch [den] Punkt, dessen Polarcoordinaten $\frac{1}{(12)}$, 12 sind (nemlich Distanz $\frac{1}{(12)}$, Richtung 12) und macht dort mit (12) den Winkel $02-01$ oder hat Azimuth $21-(02-01)$.

Am elegantesten wird die Construction, wenn man $(12)\times(13)$ anstatt 1 in die Zähler der Brüche setzt.

\medskip

\Needspace{13\baselineskip}\begin{center}[7.]\end{center}

Gegebene Punkte 1, 2, 3; Winkel an denselben $A, B, C$; Seiten $23=a$, $31=b$, $12=c$; $\mathfrak{A}, \mathfrak{B}, \mathfrak{C}$ vorgeschriebene Winkel an einem vierten Punkte.

$\beta$ Kreis, bestehend aus den Abschnitten

$\beta'$, wo Sehne $b$ unter Winkel $\mathfrak{B}$,

$\beta''$, wo diese Sehne unter Winkel $\mathfrak{B}\pm180^0$ erscheint.

$\gamma$ Kreis, bestehend aus den Abschnitten

$\gamma'$, wo Sehne $c$ unter Winkel $\mathfrak{C}$,

$\gamma''$, wo Sehne $c$ unter Winkel $\mathfrak{C}\pm180^0$ erscheint.

Die Kreise schneiden sich in $A$ und $D$.

In $D$ seien die durch $b$ und $c$ subtendirten Winkel $\mathfrak{B}^*, \mathfrak{C}^*$; die Sinus von $\mathfrak{B}^*-B$, $\mathfrak{C}^*-C$, $\mathfrak{B}^*+\mathfrak{C}^*-B-C$ werden alle drei positiv sein, wenn I) $D$ innerhalb des Kreises durch 1, 2, 3 liegt, alle drei hingegen negativ, wenn II) $D$ ausserhalb dieses Kreises liegt.

\Needspace{9\baselineskip}In Beziehung auf die Lage von $D$ in den Kreisen $\beta$ und $\gamma$ sind vier Fälle möglich. Es liegt $D$

1) in $\beta'$ und $\gamma'$, wo dann $\mathfrak{B}^*=\mathfrak{B}$, $\mathfrak{C}^*=\mathfrak{C}$

2) in $\beta''$ und $\gamma'$, wo $\mathfrak{B}^*=\mathfrak{B}\pm180^0$, $\mathfrak{C}^*=\mathfrak{C}$

3) in $\beta'$ und $\gamma''$, wo $\mathfrak{B}^*=\mathfrak{B}$, $\mathfrak{C}^*=\mathfrak{C}\pm180^0$

4) in $\beta''$ und $\gamma''$, wo $\mathfrak{B}^*=\mathfrak{B}\pm180^0$, $\mathfrak{C}^*=\mathfrak{C}\pm180^0$.

Im Falle 1) ist die Aufgabe physisch möglich, in den drei andern unmöglich.
%
\SourceLeaf{337}

Die Zeichen von
\[ \sin(\mathfrak{B}-B), \qquad \sin(\mathfrak{C}-C), \qquad \sin(\mathfrak{B}+\mathfrak{C}-B-C) \]
werden in den acht möglichen Combinationen diese sein

\begin{center}
\begin{tabular}{cc|cc}
I 1 & $+\;+\;+$ & II 1 & $-\;-\;-$ \\
I 2 & $-\;+\;-$ & II 2 & $+\;-\;+$ \\
I 3 & $+\;-\;-$ & II 3 & $-\;+\;+$ \\
I 4 & $-\;-\;+$ & II 4 & $+\;+\;-$ \\
\end{tabular}
\end{center}

Zur physischen Möglichkeit der Aufgabe ist also erforderlich und zureichend, dass die drei Sinus gleiche Zeichen haben.

Hat der erste Sinus ein von den beiden andern abweichendes Zeichen, so (Fall 3) gibt die Construction einen Punkt, der im Kreise $\gamma$ im ungeeigneten Segmente liegt.

Hat der zweite Sinus ein von den beiden andern verschiedenes Zeichen, so [Fall 2] liegt der durch die Construction gefundene Punkt im ungeeigneten Segmente des Kreises $\beta$.

\Needspace{6\baselineskip}Weicht hingegen der dritte Sinus rücksichtlich des Zeichens von den beiden andern ab, so [Fall 4] liegt der durch die Construction gefundene Punkt in beiden Kreisen $\beta$ und $\gamma$ in ungeeigneten Segmenten.

\hfill 1852 Juli 29.

\medskip

\begin{center}
\textbf{BEMERKUNGEN.}
\end{center}

Die meisten der vorstehenden Notizen sind in einem Handbuche enthalten, nur [7] befindet sich auf einem einzelnen Blatte. Was die Zeit der Abfassung betrifft, so stammt die Notiz [1] aus dem Jahre 1832, während [5] und [6] in das Jahr 1849, [7] in das Jahr 1852 zu setzen ist.

Vom December 1839 bis Ostern 1840 hat \textsc{Gauss} eine Vorlesung über die \emph{Theorie der imaginären Grössen} gehalten, die E. \textsc{Heine} gehört und ausgearbeitet hat. Dieser Ausarbeitung sind die folgenden auf das Viereck bezüglichen Ausführungen entnommen, die zur Erläuterung der unter [III] zusammengestellten Notizen dienen mögen. \hfill \textsc{Stäckel}.
%
\SourceLeaf{338}

\begin{center}
[\emph{Vorlesung über die Theorie der imaginären Grössen.}]

Vom Viereck.

[1.]
\end{center}

Hat man ein Viereck (1), (2), (3), (4), das durch die complexen Zahlen $p'$, $p''$, $p'''$,
$p^{\mathrm{IV}}$ bestimmt ist, so lässt sich eine allgemeine Formel vermittelst der Lateralen
bestimmen, indem man auf ein Dreieck (I), (II), (III) recurrirt, das daraus auf die alleinig
symmetrische Art entstanden ist, wo also die Eckpunkte folgendermassen festgelegt sind:
\[
\begin{array}{rll}
\text{(I)} & \text{durch} & p'p'' + p'''p^{\mathrm{IV}} \\
\text{(II)} & \text{,,} & p'p''' + p'' p^{\mathrm{IV}} \\
\text{(III)} & \text{,,} & p'p^{\mathrm{IV}} + p'' p'''.
\end{array}
\]
Ist z.\,B. die Complexe, die
\[
\begin{array}{rlll}
\text{(1)} & \text{festlegt} & : & 2 + \sqrt{-1} = p', \\
\text{die (2)} & \text{,,} & : & 3 \hspace{2.8em} = p'', \\
\text{,, (3)} & \text{,,} & : & 1 + 2\sqrt{-1} = p''', \\
\text{,, (4)} & \text{,,} & : & 1 + 3\sqrt{-1} = p^{\mathrm{IV}},
\end{array}
\]
so wird
\[
\begin{array}{rlll}
\text{(I)} & \text{bestimmt durch} & : & 1 + 8\sqrt{-1} \\
\text{(II)} & \text{,, ,,} & : & 3 + 14\sqrt{-1} \\
\text{(III)} & \text{,, ,,} & : & 2 + 13\sqrt{-1}.
\end{array}
\]

Beziehe ich (II) auf (I), so erhalte ich, indem ich die Complexen, die zu diesen Ecken gehören,
von einander abziehe:
\[
\text{(II)} - \text{(I)} = p'(p''' - p'') + p^{\mathrm{IV}}(p'' - p'''),
\]
und ebenso:
\[
\text{(III)} - \text{(II)} = p'(p^{\mathrm{IV}} - p''') + p'' (p''' - p^{\mathrm{IV}})
\]
\[
\text{(I)} - \text{(III)} = p'(p'' - p^{\mathrm{IV}}) + p'''(p^{\mathrm{IV}} - p'')
\]
oder
\[
\begin{array}{rl}
\text{(II)} - \text{(I)} & = (p' - p^{\mathrm{IV}})(p''' - p'') \\
\text{(III)} - \text{(II)} & = (p' - p'')(p^{\mathrm{IV}} - p''') \\
\text{(I)} - \text{(III)} & = (p' - p''')(p'' - p^{\mathrm{IV}}).
\end{array}
\]

Beziehe ich die Winkel des Vierecks auf eine feste Axe, und bezeichne z.\,B. den Winkel, den (4,1)
mit dieser bildet, mit [4,1], so ist:
\[
\begin{array}{rl}
-p^{\mathrm{IV}} + p' & = (4,1)\{\cos[4,1] + i\sin[4,1]\} \\
p''' - p'' & = (3,2)\{\cos[3,2] + i\sin[3,2]\}
\end{array}
\]
\[
\text{etc.}
\]
Man hat daher:
\[
\begin{array}{rl}
\text{(I, II)} & = (4,1) . (3,2) \\
\text{(II, III)} & = (2,1) . (4,3) \\
\text{(III, I)} & = (3,1) . (4,2).
\end{array}
\]
%
\SourceLeaf{339}

Ferner, bezeichnen wir z.\,B. den Winkel, welchen (I, II) mit der festen Axe bildet, mit [I, II],
so haben wir:
\[
\begin{array}{rl}
\text{[I, II]} & = [4,1] + [3,2] \\
\text{[II, III]} & = [2,1] + [4,3] \\
\text{[III, I]} & = [3,1] + [4,2].
\end{array}
\]
Folglich haben wir: Es ist der Winkel:
\[
\begin{array}{rl}
\text{II, I, III} & = [4,1] + [3,2] - [3,1] - [4,2] \\
\text{III, II, I} & = [2,1] + [4,3] - [4,1] - [3,2] \\
\text{I, III, II} & = [3,1] + [4,2] - [2,1] - [4,3].
\end{array}
\]

Wendet man also den Satz an, dass sich in einem Dreieck die Seiten verhalten, wie die Sinus der
gegenüberliegenden Winkel, so hat man für das ursprüngliche Viereck:
\[
\begin{aligned}
&\frac{\sin\{[4,1]+[3,2]-[3,1]-4,2]\}}{(2,1)(4,3)}\\
&=\frac{\sin\{[2,1]+[4,3]-[4,1]-[3,2]\}}{(3,1)(4,2)}\\
&=\frac{\sin\{[3,1]+[4,2]-[2,1]-[4,3]\}}{(4,1)(3,2)}.
\end{aligned}
\]

\Needspace{27\baselineskip}\begin{center}
[2.]
\end{center}

Unser Satz lässt sich ebenso einfach geometrisch beweisen, weshalb wir den letztern Beweis
gleichfalls hierher setzen, obgleich er eigentlich nicht in diesen Vortrag gehört.

Wir nehmen zuerst den Fall an, es läge der Punkt 4 innerhalb des Dreiecks 1, 2, 3. Zieht man dann
(4,1), (4,2), (4,3) und verlängert die Richtungen dieser Linien bis nach I, II, III, so dass man
habe:

\begin{center}\resizebox{12.8cm}{!}{%
\begin{tikzpicture}[x=1pt,y=-1pt]
\path[use as bounding box] (0,0) rectangle (1606,989);
\path[fill=black,draw=none,even odd rule]
(1452,949) -- (1452,950) -- (1454,952) -- (1460,952) -- (1461,953) -- (1464,953) -- (1465,954) -- (1471,954) -- (1475,956) -- (1480,956) -- (1481,955) -- (1481,953) -- (1478,951) -- (1472,951) -- (1471,950) -- (1468,950) -- (1464,948) -- (1453,948) -- cycle
(1420,943) -- (1420,944) -- (1422,946) -- (1426,946) -- (1429,948) -- (1434,948) -- (1435,949) -- (1443,949) -- (1445,947) -- (1444,945) -- (1442,944) -- (1439,944) -- (1435,942) -- (1421,942) -- cycle
(1385,937) -- (1385,938) -- (1388,941) -- (1401,942) -- (1402,943) -- (1409,943) -- (1410,942) -- (1410,941) -- (1408,939) -- (1406,939) -- (1405,938) -- (1398,938) -- (1395,936) -- (1386,936) -- cycle
(1353,931) -- (1353,934) -- (1356,936) -- (1363,935) -- (1366,937) -- (1370,937) -- (1371,938) -- (1378,938) -- (1379,937) -- (1379,935) -- (1377,933) -- (1357,933) -- (1356,932) -- (1356,930) -- (1354,930) -- cycle
(1315,925) -- (1315,928) -- (1318,930) -- (1325,930) -- (1326,931) -- (1334,932) -- (1335,933) -- (1340,933) -- (1342,931) -- (1342,928) -- (1339,926) -- (1332,926) -- (1331,925) -- (1327,925) -- (1326,924) -- (1316,924) -- cycle
(1276,919) -- (1278,923) -- (1283,923) -- (1287,925) -- (1295,925) -- (1296,926) -- (1302,926) -- (1303,925) -- (1303,923) -- (1301,921) -- (1297,921) -- (1295,919) -- (1293,919) -- (1292,920) -- (1290,920) -- (1286,918) -- (1277,918) -- cycle
(1466,917) -- (1467,921) -- (1474,927) -- (1475,927) -- (1482,935) -- (1487,935) -- (1488,934) -- (1486,930) -- (1484,929) -- (1480,925) -- (1480,924) -- (1473,918) -- (1472,916) -- (1468,915) -- cycle
(1238,912) -- (1238,914) -- (1239,915) -- (1244,915) -- (1250,918) -- (1254,917) -- (1257,919) -- (1264,919) -- (1266,917) -- (1266,915) -- (1265,914) -- (1258,914) -- (1257,913) -- (1251,912) -- (1250,911) -- (1239,911) -- cycle
(1205,907) -- (1205,909) -- (1209,911) -- (1215,911) -- (1216,912) -- (1220,912) -- (1221,913) -- (1226,913) -- (1228,911) -- (1227,909) -- (1225,909) -- (1222,907) -- (1215,907) -- (1214,906) -- (1206,906) -- cycle
(1164,900) -- (1164,903) -- (1165,904) -- (1177,904) -- (1180,906) -- (1182,906) -- (1183,907) -- (1188,907) -- (1189,908) -- (1191,908) -- (1194,905) -- (1193,903) -- (1187,902) -- (1186,901) -- (1176,901) -- (1174,899) -- (1165,899) -- cycle
(1446,897) -- (1446,899) -- (1449,902) -- (1449,903) -- (1459,912) -- (1462,912) -- (1464,910) -- (1464,909) -- (1458,903) -- (1458,902) -- (1452,897) -- (1450,896) -- (1447,896) -- cycle
(1126,894) -- (1126,896) -- (1128,898) -- (1133,898) -- (1134,899) -- (1138,898) -- (1141,900) -- (1146,900) -- (1147,901) -- (1155,901) -- (1156,900) -- (1156,898) -- (1155,897) -- (1150,897) -- (1144,894) -- (1142,895) -- (1138,895) -- (1132,892) -- (1130,892) -- cycle
(1088,888) -- (1088,889) -- (1091,891) -- (1094,891) -- (1101,894) -- (1114,894) -- (1115,893) -- (1115,890) -- (1112,888) -- (1102,888) -- (1101,887) -- (1096,887) -- (1095,886) -- (1090,886) -- cycle
(1050,881) -- (1051,883) -- (1054,884) -- (1056,886) -- (1059,886) -- (1060,887) -- (1074,888) -- (1076,886) -- (1076,884) -- (1075,883) -- (1071,883) -- (1070,882) -- (1057,882) -- (1055,880) -- (1051,880) -- cycle
(1018,876) -- (1018,878) -- (1020,880) -- (1025,880) -- (1028,882) -- (1033,882) -- (1034,883) -- (1037,882) -- (1039,883) -- (1041,881) -- (1041,879) -- (1037,877) -- (1033,878) -- (1032,877) -- (1029,877) -- (1025,875) -- (1023,875) -- (1022,876) -- (1019,875) -- cycle
(1425,873) -- (1423,875) -- (1423,877) -- (1432,887) -- (1433,887) -- (1436,890) -- (1440,891) -- (1442,889) -- (1442,887) -- (1438,883) -- (1437,883) -- (1429,874) -- cycle
(982,870) -- (982,871) -- (987,875) -- (989,874) -- (991,875) -- (994,875) -- (996,874) -- (997,875) -- (999,875) -- (1001,877) -- (1010,877) -- (1011,876) -- (1011,875) -- (1008,873) -- (1004,874) -- (1002,872) -- (996,872) -- (992,870) -- (986,870) -- (985,869) -- (983,869) -- cycle
(949,864) -- (951,868) -- (954,868) -- (955,869) -- (961,869) -- (963,871) -- (970,871) -- (972,869) -- (971,867) -- (968,867) -- (966,865) -- (960,865) -- (958,863) -- (950,863) -- cycle
(914,859) -- (914,861) -- (917,863) -- (921,863) -- (923,865) -- (939,866) -- (941,864) -- (940,862) -- (936,862) -- (933,860) -- (927,860) -- (926,859) -- (922,860) -- (918,858) -- (915,858) -- cycle
(878,854) -- (878,855) -- (882,860) -- (884,860) -- (885,859) -- (893,859) -- (894,860) -- (899,860) -- (900,859) -- (902,859) -- (904,857) -- (903,855) -- (901,855) -- (900,854) -- (894,854) -- (893,853) -- (883,852) -- (882,851) -- (881,851) -- cycle
(842,847) -- (842,849) -- (844,851) -- (846,851) -- (852,854) -- (865,855) -- (867,853) -- (867,851) -- (863,848) -- (855,848) -- (854,849) -- (852,849) -- (849,847) -- (846,847) -- (845,846) -- (843,846) -- cycle
(1382,846) -- (1386,849) -- (1387,848) -- (1387,846) -- (1389,844) -- (1390,844) -- (1393,847) -- (1394,849) -- (1394,852) -- (1402,856) -- (1409,865) -- (1412,866) -- (1414,864) -- (1415,862) -- (1414,860) -- (1412,860) -- (1409,856) -- (1392,845) -- (1389,842) -- (1387,841) -- cycle
(809,843) -- (809,845) -- (813,848) -- (818,848) -- (819,849) -- (830,849) -- (832,847) -- (832,846) -- (829,843) -- (816,842) -- (815,841) -- (813,841) -- cycle
(774,838) -- (774,840) -- (775,841) -- (778,841) -- (779,842) -- (785,842) -- (787,844) -- (800,844) -- (802,842) -- (802,841) -- (800,839) -- (791,839) -- (790,838) -- (786,838) -- (783,836) -- (781,837) -- (775,837) -- cycle
(738,832) -- (738,834) -- (739,835) -- (741,835) -- (742,836) -- (748,836) -- (751,838) -- (757,838) -- (758,839) -- (761,839) -- (765,837) -- (765,835) -- (763,833) -- (761,833) -- (760,832) -- (753,832) -- (752,833) -- (751,833) -- (749,831) -- (739,831) -- cycle
(701,826) -- (701,828) -- (702,829) -- (704,829) -- (705,830) -- (710,830) -- (714,832) -- (718,831) -- (722,833) -- (725,833) -- (728,831) -- (728,829) -- (727,828) -- (723,828) -- (721,826) -- (719,827) -- (717,826) -- (710,826) -- (708,824) -- (703,824) -- cycle
(1366,820) -- (1366,823) -- (1380,836) -- (1388,836) -- (1385,833) -- (1385,832) -- (1383,831) -- (1371,819) -- (1367,819) -- cycle
(668,820) -- (668,822) -- (671,824) -- (674,824) -- (675,825) -- (681,824) -- (684,827) -- (690,827) -- (692,825) -- (692,824) -- (690,822) -- (688,822) -- (684,820) -- (674,820) -- (671,818) -- (670,818) -- cycle
(635,815) -- (635,817) -- (637,819) -- (641,819) -- (642,820) -- (646,819) -- (649,821) -- (653,821) -- (654,820) -- (658,820) -- (659,819) -- (659,817) -- (655,814) -- (652,814) -- (651,815) -- (648,815) -- (646,816) -- (642,813) -- (639,813) -- cycle
(603,809) -- (603,811) -- (605,813) -- (608,813) -- (609,814) -- (618,814) -- (621,816) -- (623,815) -- (626,815) -- (628,813) -- (627,811) -- (622,811) -- (619,809) -- (613,809) -- (609,807) -- (605,807) -- cycle
(565,803) -- (565,804) -- (567,806) -- (572,807) -- (574,809) -- (580,808) -- (581,809) -- (585,809) -- (586,810) -- (592,810) -- (594,808) -- (594,807) -- (592,805) -- (590,805) -- (589,804) -- (582,804) -- (581,803) -- (578,803) -- (577,802) -- (574,802) -- (573,801) -- (567,801) -- cycle
(535,798) -- (535,800) -- (536,801) -- (538,801) -- (539,802) -- (543,802) -- (547,804) -- (551,804) -- (552,805) -- (554,805) -- (557,803) -- (557,800) -- (554,798) -- (547,798) -- (544,796) -- (538,796) -- cycle
(1343,794) -- (1341,796) -- (1342,799) -- (1355,812) -- (1359,813) -- (1361,811) -- (1361,809) -- (1357,805) -- (1357,804) -- (1355,803) -- (1350,797) -- (1345,794) -- cycle
(505,794) -- (505,795) -- (507,797) -- (510,797) -- (511,798) -- (528,799) -- (530,797) -- (530,796) -- (527,794) -- (524,794) -- (520,792) -- (514,792) -- (513,791) -- (507,792) -- cycle
(470,787) -- (470,790) -- (471,791) -- (473,791) -- (475,793) -- (484,793) -- (485,794) -- (492,794) -- (493,795) -- (495,795) -- (498,793) -- (498,792) -- (496,790) -- (492,789) -- (491,788) -- (487,788) -- (486,787) -- (482,787) -- (481,786) -- (476,786) -- (475,785) -- (471,786) -- cycle
(436,782) -- (436,784) -- (438,786) -- (443,786) -- (444,787) -- (454,788) -- (457,790) -- (459,790) -- (462,788) -- (462,785) -- (460,783) -- (450,782) -- (449,781) -- (444,781) -- (441,779) -- (438,780) -- cycle
(398,776) -- (398,777) -- (401,780) -- (413,781) -- (414,782) -- (418,782) -- (419,783) -- (423,783) -- (424,782) -- (424,779) -- (422,777) -- (418,777) -- (415,775) -- (400,774) -- cycle
(1321,773) -- (1319,775) -- (1320,779) -- (1323,782) -- (1328,785) -- (1331,789) -- (1332,789) -- (1335,792) -- (1337,792) -- (1339,790) -- (1339,787) -- (1328,776) -- (1327,776) -- (1324,773) -- cycle
(363,769) -- (363,771) -- (365,773) -- (373,774) -- (374,775) -- (377,775) -- (378,776) -- (381,776) -- (382,777) -- (385,777) -- (388,774) -- (388,772) -- (386,770) -- (383,770) -- (382,769) -- (375,769) -- (371,767) -- (366,767) -- cycle
(329,764) -- (329,765) -- (331,767) -- (339,768) -- (340,769) -- (343,769) -- (344,770) -- (350,770) -- (351,771) -- (353,771) -- (355,769) -- (355,767) -- (351,764) -- (345,764) -- (344,763) -- (334,762) -- (333,761) -- (332,761) -- cycle
(295,758) -- (295,759) -- (298,762) -- (303,762) -- (305,764) -- (315,764) -- (316,765) -- (321,765) -- (322,764) -- (322,761) -- (320,759) -- (310,758) -- (306,756) -- (297,756) -- cycle
(256,751) -- (256,754) -- (258,756) -- (263,756) -- (266,758) -- (274,758) -- (275,759) -- (280,759) -- (283,757) -- (283,756) -- (279,752) -- (269,752) -- (265,750) -- (257,750) -- cycle
(1296,748) -- (1295,749) -- (1295,752) -- (1296,754) -- (1306,764) -- (1306,765) -- (1308,767) -- (1310,768) -- (1314,768) -- (1316,766) -- (1316,765) -- (1300,749) -- cycle
(216,746) -- (216,748) -- (217,749) -- (221,750) -- (222,751) -- (227,751) -- (228,752) -- (236,752) -- (237,753) -- (243,753) -- (244,752) -- (244,749) -- (241,746) -- (239,747) -- (235,747) -- (234,746) -- (226,745) -- (222,743) -- (219,743) -- cycle
(179,740) -- (179,742) -- (182,745) -- (190,745) -- (191,746) -- (194,746) -- (195,747) -- (201,747) -- (202,748) -- (205,748) -- (207,746) -- (207,743) -- (203,740) -- (197,740) -- (196,739) -- (190,739) -- (189,738) -- (181,738) -- cycle
(144,734) -- (144,735) -- (148,739) -- (150,740) -- (153,740) -- (154,741) -- (164,741) -- (165,742) -- (168,742) -- (172,740) -- (172,737) -- (170,735) -- (168,735) -- (167,734) -- (162,734) -- (160,735) -- (159,734) -- (154,734) -- (152,732) -- (147,732) -- cycle
(111,728) -- (111,730) -- (113,732) -- (115,732) -- (116,733) -- (120,733) -- (121,734) -- (124,734) -- (125,735) -- (134,735) -- (136,733) -- (136,731) -- (132,728) -- (123,728) -- (120,726) -- (113,726) -- cycle
(1273,725) -- (1271,727) -- (1271,729) -- (1280,738) -- (1280,739) -- (1282,741) -- (1287,744) -- (1290,741) -- (1290,740) -- (1276,726) -- cycle
(1249,703) -- (1248,704) -- (1249,707) -- (1257,715) -- (1257,716) -- (1260,719) -- (1264,719) -- (1265,718) -- (1264,714) -- (1262,713) -- (1259,710) -- (1259,709) -- (1252,703) -- cycle
(112,684) -- (107,687) -- (99,697) -- (99,698) -- (101,700) -- (105,699) -- (110,693) -- (110,692) -- (114,688) -- cycle
(1223,680) -- (1223,682) -- (1232,691) -- (1233,693) -- (1234,693) -- (1239,697) -- (1242,697) -- (1243,696) -- (1243,694) -- (1241,692) -- (1240,692) -- (1238,690) -- (1238,689) -- (1235,686) -- (1234,686) -- (1228,679) -- (1224,679) -- cycle
(133,658) -- (131,658) -- (119,671) -- (119,672) -- (122,675) -- (124,675) -- (129,670) -- (131,666) -- (134,663) -- (135,660) -- cycle
(1198,653) -- (1197,654) -- (1198,657) -- (1211,670) -- (1215,671) -- (1216,670) -- (1216,667) -- (1214,666) -- (1210,662) -- (1210,661) -- (1204,656) -- (1203,654) -- (1201,653) -- cycle
(1177,632) -- (1176,633) -- (1176,636) -- (1187,647) -- (1189,648) -- (1192,648) -- (1193,647) -- (1193,644) -- (1181,633) -- cycle
(171,611) -- (168,612) -- (165,615) -- (165,616) -- (154,626) -- (151,631) -- (151,634) -- (146,638) -- (145,641) -- (140,646) -- (140,649) -- (141,650) -- (143,650) -- (152,640) -- (154,641) -- (158,635) -- (158,631) -- (160,629) -- (162,630) -- (162,629) -- (164,627) -- (164,624) -- (170,619) -- (172,619) -- (172,612) -- cycle
(1148,606) -- (1148,608) -- (1154,614) -- (1154,615) -- (1156,616) -- (1158,618) -- (1158,619) -- (1165,625) -- (1169,625) -- (1170,624) -- (1170,622) -- (1164,617) -- (1163,617) -- (1153,605) -- (1149,605) -- cycle
(190,589) -- (187,588) -- (182,593) -- (182,594) -- (176,601) -- (176,604) -- (177,605) -- (181,605) -- (183,603) -- (183,602) -- (187,599) -- (188,596) -- (191,593) -- (191,590) -- cycle
(1123,582) -- (1124,585) -- (1132,594) -- (1133,594) -- (1138,598) -- (1141,598) -- (1142,597) -- (1142,596) -- (1138,592) -- (1138,591) -- (1129,582) -- (1127,581) -- (1124,581) -- cycle
(209,564) -- (206,564) -- (203,566) -- (203,568) -- (194,578) -- (194,580) -- (196,582) -- (199,582) -- (207,573) -- (210,568) -- (210,565) -- cycle
(1098,555) -- (1097,556) -- (1097,558) -- (1099,560) -- (1101,564) -- (1112,574) -- (1116,574) -- (1117,573) -- (1117,571) -- (1113,569) -- (1111,567) -- (1111,565) -- (1107,561) -- (1106,561) -- (1101,556) -- cycle
(225,544) -- (222,543) -- (219,546) -- (217,550) -- (213,554) -- (213,558) -- (214,559) -- (218,559) -- (224,552) -- (227,547) -- (227,546) -- cycle
(1071,532) -- (1071,534) -- (1073,535) -- (1075,537) -- (1075,538) -- (1084,547) -- (1086,548) -- (1090,548) -- (1091,547) -- (1090,544) -- (1087,541) -- (1084,540) -- (1075,530) -- (1073,530) -- cycle
(243,519) -- (241,519) -- (234,527) -- (234,528) -- (228,534) -- (228,536) -- (230,538) -- (234,538) -- (244,526) -- (245,524) -- (245,521) -- cycle
(1044,505) -- (1044,508) -- (1047,511) -- (1048,511) -- (1057,521) -- (1058,521) -- (1061,524) -- (1065,522) -- (1065,521) -- (1062,518) -- (1061,518) -- (1059,516) -- (1059,515) -- (1047,504) -- (1045,504) -- cycle
(263,494) -- (258,494) -- (257,495) -- (256,499) -- (249,507) -- (249,509) -- (251,511) -- (254,511) -- (259,507) -- (260,504) -- (264,499) -- (264,495) -- cycle
(1021,484) -- (1022,487) -- (1024,489) -- (1025,489) -- (1035,499) -- (1039,499) -- (1040,498) -- (1040,496) -- (1035,493) -- (1026,483) -- (1022,483) -- cycle
(283,468) -- (280,468) -- (276,470) -- (275,474) -- (268,482) -- (268,484) -- (270,486) -- (274,486) -- (277,482) -- (279,482) -- (280,481) -- (279,479) -- (284,473) -- (284,469) -- cycle
(996,457) -- (995,458) -- (996,462) -- (1002,468) -- (1003,468) -- (1006,472) -- (1007,472) -- (1011,476) -- (1012,476) -- (1014,474) -- (1014,472) -- (998,457) -- cycle
(302,442) -- (300,442) -- (294,448) -- (292,452) -- (287,458) -- (287,460) -- (288,461) -- (291,462) -- (297,456) -- (299,452) -- (304,448) -- (304,444) -- cycle
(970,432) -- (969,433) -- (969,436) -- (974,441) -- (975,441) -- (982,449) -- (987,451) -- (989,449) -- (989,448) -- (987,446) -- (986,446) -- (979,438) -- (977,438) -- (971,432) -- cycle
(324,415) -- (321,415) -- (311,426) -- (311,427) -- (307,432) -- (307,433) -- (309,435) -- (313,435) -- (317,431) -- (317,429) -- (324,422) -- (324,420) -- (325,419) -- (325,416) -- cycle
(945,411) -- (945,413) -- (956,424) -- (958,425) -- (962,425) -- (963,424) -- (963,422) -- (961,421) -- (951,411) -- (949,410) -- (946,410) -- cycle
(344,390) -- (341,390) -- (332,398) -- (329,403) -- (329,405) -- (331,407) -- (334,407) -- (344,396) -- (345,394) -- (345,391) -- cycle
(919,386) -- (919,388) -- (921,390) -- (921,391) -- (929,399) -- (930,399) -- (934,403) -- (938,403) -- (939,402) -- (939,400) -- (934,395) -- (933,395) -- (925,386) -- (923,385) -- (920,385) -- cycle
(362,366) -- (360,366) -- (357,369) -- (357,370) -- (353,373) -- (353,375) -- (348,379) -- (348,381) -- (350,383) -- (353,383) -- (358,377) -- (358,376) -- (362,372) -- (363,370) -- (363,367) -- cycle
(892,361) -- (893,364) -- (897,368) -- (897,369) -- (899,370) -- (906,377) -- (908,378) -- (912,378) -- (913,377) -- (911,371) -- (909,369) -- (903,366) -- (898,360) -- (893,360) -- cycle
(379,344) -- (376,344) -- (366,356) -- (366,357) -- (368,359) -- (371,359) -- (373,358) -- (378,352) -- (381,347) -- (381,346) -- cycle
(863,334) -- (863,335) -- (883,354) -- (887,354) -- (888,353) -- (888,351) -- (880,343) -- (878,343) -- (876,341) -- (876,339) -- (874,338) -- (869,332) -- (865,332) -- cycle
(397,321) -- (395,321) -- (389,327) -- (388,330) -- (383,335) -- (383,338) -- (384,339) -- (388,338) -- (392,333) -- (392,332) -- (399,325) -- (399,323) -- cycle
(840,308) -- (839,309) -- (839,313) -- (842,316) -- (845,317) -- (847,319) -- (848,322) -- (851,325) -- (857,326) -- (858,325) -- (858,322) -- (856,320) -- (856,319) -- (855,318) -- (852,318) -- (850,316) -- (850,314) -- (843,308) -- cycle
(416,297) -- (414,297) -- (407,305) -- (407,306) -- (401,312) -- (401,314) -- (402,315) -- (405,316) -- (414,307) -- (414,306) -- (417,303) -- (417,301) -- (418,300) -- (418,299) -- cycle
(814,282) -- (813,283) -- (813,287) -- (828,302) -- (833,302) -- (834,301) -- (834,299) -- (830,295) -- (829,295) -- (820,286) -- (820,285) -- (818,283) -- cycle
(434,275) -- (431,275) -- (423,284) -- (423,285) -- (420,288) -- (420,291) -- (421,292) -- (426,292) -- (427,291) -- (427,289) -- (436,279) -- (436,277) -- cycle
(787,257) -- (786,258) -- (786,261) -- (791,266) -- (792,266) -- (800,275) -- (802,276) -- (805,276) -- (806,275) -- (806,272) -- (796,262) -- (795,262) -- (791,257) -- cycle
(454,249) -- (452,249) -- (445,256) -- (443,260) -- (437,266) -- (437,267) -- (439,269) -- (442,269) -- (447,265) -- (447,263) -- (455,253) -- (455,250) -- cycle
(764,234) -- (762,236) -- (762,239) -- (764,241) -- (765,241) -- (775,251) -- (776,251) -- (778,253) -- (780,253) -- (782,251) -- (781,248) -- (769,236) -- (768,236) -- (766,234) -- cycle
(474,224) -- (472,224) -- (459,239) -- (459,240) -- (461,242) -- (464,242) -- (473,231) -- (473,230) -- (475,228) -- (475,225) -- cycle
(739,211) -- (738,213) -- (741,218) -- (745,222) -- (746,222) -- (751,228) -- (755,230) -- (757,230) -- (759,228) -- (759,227) -- (742,210) -- cycle
(493,198) -- (491,198) -- (487,204) -- (478,214) -- (478,216) -- (479,217) -- (483,217) -- (494,205) -- (495,203) -- (495,200) -- cycle
(712,184) -- (711,185) -- (711,188) -- (725,202) -- (730,204) -- (732,202) -- (732,199) -- (716,184) -- cycle
(515,174) -- (511,174) -- (508,176) -- (504,183) -- (501,186) -- (501,188) -- (502,189) -- (505,189) -- (515,179) -- (516,177) -- (516,175) -- cycle
(689,161) -- (688,162) -- (688,165) -- (702,180) -- (706,181) -- (708,179) -- (708,176) -- (706,175) -- (703,172) -- (703,171) -- (701,170) -- (692,161) -- cycle
(531,151) -- (529,151) -- (518,163) -- (518,165) -- (520,167) -- (523,167) -- (525,166) -- (532,158) -- (533,156) -- (533,153) -- cycle
(667,140) -- (666,141) -- (666,143) -- (668,147) -- (670,148) -- (679,157) -- (683,157) -- (684,156) -- (684,153) -- (671,140) -- cycle
(549,129) -- (546,129) -- (533,143) -- (534,147) -- (536,149) -- (539,149) -- (541,150) -- (546,145) -- (546,138) -- (550,133) -- (550,130) -- cycle
(555,118) -- (551,122) -- (551,125) -- (552,126) -- (554,126) -- (558,121) -- (558,120) -- (557,120) -- cycle
(641,114) -- (640,115) -- (641,119) -- (643,121) -- (643,122) -- (650,128) -- (659,128) -- (660,129) -- (660,135) -- (662,133) -- (662,131) -- (659,128) -- (658,128) -- (656,126) -- (656,125) -- (654,124) -- (644,113) -- cycle
(631,101) -- (631,109) -- (633,111) -- (637,111) -- (638,110) -- (638,108) -- (636,106) -- (635,106) -- cycle
(619,95) -- (620,97) -- (626,97) -- (626,96) -- (624,94) -- (620,94) -- cycle
(582,87) -- (579,88) -- (569,100) -- (569,102) -- (571,103) -- (573,105) -- (573,106) -- (578,100) -- (576,98) -- (582,92) -- (584,91) -- (584,89) -- cycle
(77,722) -- (77,723) -- (81,726) -- (84,726) -- (85,727) -- (90,727) -- (92,729) -- (96,729) -- (97,730) -- (101,730) -- (103,729) -- (105,727) -- (105,726) -- (102,723) -- (92,723) -- (91,722) -- (91,720) -- (93,718) -- (96,718) -- (98,716) -- (104,716) -- (107,712) -- (108,712) -- (111,709) -- (114,708) -- (116,706) -- (121,706) -- (134,694) -- (137,693) -- (140,690) -- (142,690) -- (149,684) -- (152,683) -- (160,676) -- (163,676) -- (165,675) -- (171,670) -- (173,670) -- (180,665) -- (182,665) -- (187,662) -- (193,662) -- (197,660) -- (207,651) -- (207,649) -- (209,647) -- (213,646) -- (222,638) -- (224,638) -- (231,634) -- (233,634) -- (243,625) -- (246,625) -- (260,614) -- (265,612) -- (269,608) -- (271,608) -- (273,606) -- (276,606) -- (281,601) -- (289,597) -- (290,595) -- (296,590) -- (299,590) -- (301,589) -- (305,584) -- (313,580) -- (316,577) -- (319,576) -- (325,571) -- (331,568) -- (333,566) -- (337,565) -- (342,560) -- (346,558) -- (353,552) -- (356,551) -- (359,548) -- (364,546) -- (376,537) -- (382,537) -- (383,536) -- (386,536) -- (387,535) -- (392,535) -- (393,534) -- (401,533) -- (402,532) -- (408,532) -- (414,529) -- (417,529) -- (418,528) -- (424,528) -- (425,527) -- (431,526) -- (435,524) -- (440,524) -- (443,522) -- (446,523) -- (450,520) -- (459,520) -- (462,518) -- (465,519) -- (466,518) -- (469,518) -- (472,516) -- (475,516) -- (479,514) -- (485,514) -- (488,512) -- (494,511) -- (495,510) -- (499,510) -- (500,509) -- (507,509) -- (515,505) -- (519,506) -- (520,505) -- (523,505) -- (527,503) -- (533,503) -- (535,501) -- (541,501) -- (544,499) -- (554,498) -- (557,496) -- (565,495) -- (566,494) -- (569,494) -- (570,493) -- (576,492) -- (577,491) -- (582,491) -- (585,489) -- (591,489) -- (594,487) -- (597,487) -- (598,486) -- (601,486) -- (602,485) -- (607,485) -- (608,484) -- (614,483) -- (615,482) -- (618,482) -- (619,481) -- (623,481) -- (628,484) -- (630,484) -- (633,487) -- (635,487) -- (638,490) -- (642,491) -- (644,493) -- (646,493) -- (648,495) -- (652,497) -- (654,497) -- (657,500) -- (662,503) -- (664,503) -- (667,506) -- (672,507) -- (677,511) -- (679,511) -- (681,513) -- (682,513) -- (685,516) -- (689,517) -- (691,519) -- (696,521) -- (701,525) -- (708,527) -- (712,531) -- (716,532) -- (718,534) -- (720,534) -- (724,536) -- (727,536) -- (730,540) -- (734,541) -- (737,544) -- (741,546) -- (744,546) -- (747,549) -- (757,554) -- (760,557) -- (763,558) -- (765,560) -- (767,560) -- (769,562) -- (772,562) -- (774,564) -- (778,565) -- (783,569) -- (787,570) -- (789,572) -- (796,575) -- (798,577) -- (800,577) -- (804,580) -- (806,580) -- (808,582) -- (815,585) -- (817,587) -- (819,587) -- (824,591) -- (831,593) -- (834,596) -- (840,599) -- (844,600) -- (852,606) -- (857,606) -- (861,611) -- (863,611) -- (867,613) -- (869,615) -- (874,617) -- (876,619) -- (878,619) -- (880,621) -- (884,623) -- (886,623) -- (893,628) -- (897,630) -- (899,630) -- (901,632) -- (910,636) -- (912,638) -- (916,640) -- (919,640) -- (924,645) -- (926,645) -- (930,647) -- (932,649) -- (934,649) -- (937,652) -- (939,652) -- (941,654) -- (954,660) -- (956,662) -- (960,664) -- (962,664) -- (966,668) -- (969,668) -- (971,670) -- (973,670) -- (975,672) -- (984,676) -- (990,681) -- (993,681) -- (996,684) -- (998,684) -- (1006,688) -- (1009,691) -- (1015,693) -- (1018,696) -- (1024,698) -- (1026,700) -- (1033,703) -- (1035,705) -- (1037,706) -- (1040,706) -- (1043,708) -- (1045,708) -- (1047,710) -- (1047,711) -- (1051,714) -- (1058,716) -- (1068,723) -- (1070,723) -- (1076,728) -- (1080,729) -- (1084,732) -- (1087,732) -- (1091,736) -- (1098,738) -- (1100,740) -- (1110,746) -- (1113,746) -- (1118,750) -- (1120,750) -- (1126,753) -- (1128,755) -- (1131,756) -- (1134,759) -- (1136,759) -- (1141,763) -- (1145,764) -- (1147,766) -- (1154,769) -- (1156,771) -- (1162,772) -- (1165,776) -- (1168,776) -- (1172,780) -- (1174,780) -- (1181,783) -- (1185,787) -- (1189,788) -- (1193,792) -- (1196,792) -- (1199,795) -- (1201,795) -- (1208,798) -- (1210,800) -- (1226,808) -- (1228,810) -- (1231,811) -- (1238,816) -- (1242,818) -- (1244,818) -- (1247,821) -- (1254,823) -- (1256,825) -- (1259,826) -- (1262,829) -- (1271,833) -- (1273,835) -- (1277,836) -- (1279,838) -- (1284,840) -- (1287,843) -- (1292,845) -- (1294,847) -- (1300,850) -- (1302,850) -- (1307,854) -- (1310,854) -- (1312,856) -- (1316,857) -- (1320,861) -- (1330,866) -- (1332,868) -- (1336,870) -- (1338,870) -- (1342,872) -- (1344,874) -- (1346,874) -- (1360,881) -- (1364,882) -- (1366,884) -- (1369,885) -- (1372,888) -- (1374,888) -- (1376,890) -- (1379,890) -- (1383,892) -- (1385,894) -- (1394,898) -- (1395,899) -- (1395,901) -- (1396,901) -- (1398,903) -- (1400,904) -- (1404,904) -- (1409,908) -- (1411,908) -- (1413,910) -- (1417,911) -- (1423,915) -- (1425,915) -- (1428,918) -- (1431,919) -- (1433,921) -- (1435,921) -- (1437,923) -- (1439,923) -- (1441,925) -- (1445,926) -- (1448,929) -- (1453,931) -- (1456,934) -- (1461,935) -- (1465,937) -- (1470,941) -- (1474,942) -- (1478,945) -- (1480,945) -- (1482,947) -- (1488,949) -- (1489,950) -- (1489,952) -- (1485,954) -- (1485,956) -- (1487,958) -- (1494,958) -- (1498,960) -- (1500,959) -- (1500,946) -- (1497,942) -- (1494,942) -- (1493,943) -- (1494,946) -- (1492,947) -- (1487,943) -- (1485,943) -- (1478,938) -- (1474,937) -- (1460,930) -- (1457,927) -- (1451,924) -- (1449,922) -- (1446,922) -- (1443,919) -- (1439,918) -- (1433,913) -- (1431,913) -- (1427,910) -- (1424,910) -- (1420,907) -- (1417,906) -- (1415,904) -- (1411,902) -- (1409,902) -- (1407,900) -- (1403,898) -- (1401,898) -- (1396,895) -- (1394,895) -- (1390,891) -- (1388,891) -- (1383,888) -- (1379,887) -- (1378,886) -- (1379,884) -- (1384,884) -- (1375,880) -- (1373,878) -- (1372,880) -- (1368,880) -- (1361,876) -- (1357,872) -- (1350,870) -- (1345,866) -- (1341,865) -- (1334,861) -- (1332,861) -- (1327,857) -- (1318,853) -- (1316,851) -- (1314,851) -- (1310,848) -- (1303,845) -- (1301,843) -- (1295,841) -- (1292,838) -- (1290,838) -- (1288,836) -- (1285,836) -- (1274,828) -- (1271,828) -- (1270,827) -- (1265,826) -- (1259,820) -- (1255,819) -- (1254,818) -- (1251,818) -- (1244,812) -- (1242,812) -- (1238,809) -- (1234,808) -- (1232,806) -- (1227,804) -- (1225,802) -- (1218,800) -- (1211,794) -- (1205,791) -- (1202,791) -- (1200,790) -- (1196,786) -- (1194,786) -- (1187,781) -- (1183,780) -- (1180,777) -- (1178,776) -- (1174,776) -- (1166,770) -- (1158,767) -- (1154,763) -- (1152,762) -- (1148,762) -- (1141,756) -- (1137,755) -- (1133,752) -- (1128,751) -- (1126,749) -- (1125,749) -- (1121,744) -- (1116,744) -- (1113,742) -- (1111,742) -- (1107,738) -- (1102,735) -- (1100,735) -- (1097,732) -- (1091,730) -- (1086,726) -- (1080,725) -- (1077,722) -- (1073,721) -- (1069,717) -- (1065,716) -- (1061,713) -- (1057,712) -- (1053,709) -- (1051,709) -- (1044,703) -- (1033,699) -- (1028,694) -- (1024,692) -- (1021,692) -- (1015,689) -- (1012,686) -- (1009,686) -- (1005,682) -- (1003,682) -- (999,680) -- (997,678) -- (994,677) -- (989,673) -- (987,673) -- (981,670) -- (970,662) -- (965,662) -- (961,658) -- (957,657) -- (954,654) -- (952,654) -- (948,652) -- (943,648) -- (940,648) -- (936,646) -- (931,641) -- (927,640) -- (923,637) -- (921,637) -- (918,634) -- (912,631) -- (910,631) -- (906,628) -- (903,628) -- (897,623) -- (895,623) -- (892,620) -- (890,620) -- (884,616) -- (881,616) -- (874,610) -- (869,609) -- (866,606) -- (864,606) -- (860,603) -- (856,603) -- (852,599) -- (848,598) -- (845,596) -- (843,596) -- (840,593) -- (836,592) -- (833,589) -- (831,589) -- (828,586) -- (823,585) -- (819,582) -- (816,582) -- (811,577) -- (808,577) -- (806,575) -- (804,575) -- (801,572) -- (796,570) -- (791,566) -- (789,566) -- (786,564) -- (784,564) -- (782,562) -- (778,561) -- (775,558) -- (772,557) -- (767,553) -- (760,551) -- (753,546) -- (750,546) -- (747,543) -- (743,542) -- (739,538) -- (737,538) -- (731,535) -- (725,530) -- (723,530) -- (719,527) -- (716,527) -- (709,522) -- (707,522) -- (704,519) -- (701,519) -- (698,516) -- (695,515) -- (693,513) -- (690,512) -- (681,506) -- (677,506) -- (673,502) -- (670,501) -- (668,499) -- (664,498) -- (660,496) -- (658,494) -- (655,494) -- (648,489) -- (644,487) -- (642,487) -- (640,485) -- (638,485) -- (631,480) -- (626,478) -- (623,475) -- (623,473) -- (622,472) -- (621,447) -- (620,446) -- (620,437) -- (619,436) -- (620,432) -- (619,431) -- (619,425) -- (618,424) -- (618,420) -- (619,419) -- (619,417) -- (618,416) -- (618,411) -- (617,410) -- (617,404) -- (615,401) -- (616,397) -- (614,394) -- (614,391) -- (615,390) -- (615,387) -- (614,386) -- (614,373) -- (613,372) -- (613,365) -- (612,364) -- (612,356) -- (611,355) -- (611,352) -- (610,351) -- (611,348) -- (609,344) -- (610,335) -- (609,334) -- (609,322) -- (608,321) -- (608,313) -- (606,310) -- (606,294) -- (605,293) -- (606,289) -- (605,288) -- (604,275) -- (603,274) -- (604,272) -- (604,270) -- (603,269) -- (604,264) -- (602,262) -- (603,253) -- (601,250) -- (601,238) -- (599,235) -- (599,219) -- (598,218) -- (598,209) -- (597,208) -- (597,205) -- (596,204) -- (597,196) -- (595,194) -- (595,184) -- (594,183) -- (595,177) -- (594,176) -- (594,173) -- (593,172) -- (593,169) -- (594,167) -- (593,166) -- (593,161) -- (592,160) -- (591,146) -- (589,144) -- (589,139) -- (590,138) -- (590,133) -- (591,132) -- (591,125) -- (592,124) -- (592,119) -- (593,118) -- (593,112) -- (594,111) -- (594,107) -- (595,106) -- (596,95) -- (599,89) -- (599,81) -- (600,80) -- (602,80) -- (609,86) -- (612,86) -- (613,85) -- (613,83) -- (608,80) -- (600,72) -- (600,70) -- (599,69) -- (596,69) -- (588,77) -- (588,79) -- (589,80) -- (592,80) -- (594,78) -- (596,80) -- (593,88) -- (593,100) -- (590,103) -- (590,106) -- (589,107) -- (589,112) -- (590,113) -- (590,115) -- (589,116) -- (589,118) -- (587,119) -- (587,124) -- (588,126) -- (587,127) -- (587,133) -- (585,135) -- (585,138) -- (584,139) -- (584,144) -- (583,145) -- (581,144) -- (579,145) -- (578,148) -- (576,150) -- (576,152) -- (574,154) -- (573,157) -- (570,160) -- (570,162) -- (568,164) -- (565,171) -- (563,173) -- (560,179) -- (560,181) -- (557,184) -- (554,191) -- (549,198) -- (548,203) -- (544,207) -- (544,209) -- (539,216) -- (539,218) -- (536,221) -- (533,229) -- (529,234) -- (528,239) -- (524,244) -- (524,246) -- (522,248) -- (520,252) -- (520,254) -- (517,257) -- (517,259) -- (515,263) -- (511,268) -- (508,275) -- (504,280) -- (503,284) -- (500,290) -- (498,292) -- (496,297) -- (494,299) -- (492,305) -- (490,307) -- (488,312) -- (486,314) -- (486,316) -- (484,320) -- (482,322) -- (482,324) -- (480,326) -- (473,340) -- (471,342) -- (471,344) -- (468,350) -- (466,352) -- (465,356) -- (463,358) -- (461,364) -- (458,367) -- (452,381) -- (448,386) -- (447,390) -- (444,394) -- (444,396) -- (441,399) -- (441,402) -- (437,407) -- (435,411) -- (435,413) -- (433,415) -- (433,417) -- (429,422) -- (424,432) -- (424,434) -- (421,437) -- (421,439) -- (415,449) -- (415,452) -- (411,456) -- (409,461) -- (407,463) -- (406,468) -- (404,470) -- (402,475) -- (400,477) -- (400,479) -- (398,481) -- (397,485) -- (392,492) -- (387,503) -- (383,508) -- (383,510) -- (381,512) -- (380,516) -- (378,520) -- (376,522) -- (376,524) -- (374,528) -- (372,530) -- (372,532) -- (364,540) -- (360,541) -- (357,544) -- (354,545) -- (350,549) -- (342,553) -- (338,557) -- (334,558) -- (328,563) -- (324,564) -- (315,571) -- (313,571) -- (306,577) -- (305,577) -- (300,581) -- (297,582) -- (287,590) -- (283,592) -- (281,592) -- (276,598) -- (272,599) -- (270,601) -- (260,607) -- (257,610) -- (254,610) -- (242,620) -- (240,620) -- (235,623) -- (229,629) -- (223,630) -- (217,636) -- (215,636) -- (213,638) -- (206,642) -- (203,645) -- (199,646) -- (186,656) -- (177,661) -- (172,666) -- (169,666) -- (166,669) -- (160,672) -- (156,676) -- (154,676) -- (149,680) -- (145,681) -- (143,683) -- (138,685) -- (136,687) -- (129,691) -- (126,694) -- (121,696) -- (112,703) -- (108,704) -- (100,710) -- (98,710) -- (96,712) -- (94,712) -- (92,710) -- (90,710) -- (80,720) -- (79,720) -- cycle
(615,476) -- (615,477) -- (613,479) -- (605,480) -- (601,483) -- (594,483) -- (590,485) -- (587,484) -- (585,486) -- (583,486) -- (581,484) -- (580,485) -- (577,485) -- (577,486) -- (575,488) -- (573,488) -- (572,489) -- (568,489) -- (567,490) -- (565,490) -- (563,488) -- (558,492) -- (554,492) -- (550,494) -- (543,494) -- (542,495) -- (539,495) -- (533,498) -- (529,498) -- (525,500) -- (517,500) -- (511,503) -- (504,503) -- (501,505) -- (495,506) -- (494,507) -- (483,508) -- (482,509) -- (479,509) -- (475,511) -- (469,511) -- (463,514) -- (462,514) -- (459,511) -- (458,512) -- (455,512) -- (453,513) -- (454,513) -- (456,515) -- (456,516) -- (454,518) -- (453,517) -- (448,517) -- (447,518) -- (439,519) -- (438,520) -- (435,520) -- (434,521) -- (431,521) -- (430,522) -- (425,522) -- (424,523) -- (415,524) -- (413,526) -- (402,527) -- (401,528) -- (399,528) -- (398,529) -- (394,529) -- (390,531) -- (388,529) -- (395,523) -- (397,523) -- (401,520) -- (406,518) -- (408,516) -- (411,515) -- (414,512) -- (424,507) -- (427,504) -- (428,504) -- (430,506) -- (435,502) -- (438,501) -- (443,497) -- (446,496) -- (448,494) -- (446,493) -- (446,492) -- (448,490) -- (450,490) -- (456,485) -- (460,484) -- (463,481) -- (466,480) -- (471,476) -- (475,475) -- (480,470) -- (484,468) -- (491,462) -- (494,462) -- (498,458) -- (500,458) -- (504,456) -- (509,451) -- (515,448) -- (518,445) -- (524,442) -- (527,439) -- (531,438) -- (534,435) -- (542,436) -- (549,441) -- (551,441) -- (554,444) -- (556,444) -- (566,450) -- (569,450) -- (570,451) -- (570,453) -- (569,454) -- (571,455) -- (573,453) -- (576,453) -- (578,454) -- (585,460) -- (590,462) -- (592,464) -- (598,467) -- (600,467) -- (602,469) -- (613,474) -- cycle
(586,158) -- (588,158) -- (589,159) -- (589,168) -- (590,169) -- (590,184) -- (591,185) -- (591,189) -- (592,190) -- (592,196) -- (591,197) -- (591,199) -- (592,200) -- (592,205) -- (593,206) -- (593,213) -- (594,214) -- (594,217) -- (595,218) -- (594,220) -- (592,220) -- (591,218) -- (591,223) -- (595,227) -- (595,233) -- (596,234) -- (596,239) -- (597,240) -- (597,243) -- (596,244) -- (596,250) -- (597,251) -- (597,270) -- (598,271) -- (598,274) -- (599,275) -- (599,279) -- (598,280) -- (598,282) -- (601,288) -- (600,290) -- (600,299) -- (601,300) -- (601,306) -- (602,307) -- (602,314) -- (601,315) -- (601,317) -- (602,318) -- (602,324) -- (603,325) -- (603,330) -- (604,331) -- (604,334) -- (605,335) -- (605,341) -- (606,342) -- (605,344) -- (605,347) -- (606,348) -- (607,367) -- (608,368) -- (608,374) -- (609,375) -- (609,384) -- (610,385) -- (609,387) -- (609,390) -- (611,393) -- (610,395) -- (610,398) -- (611,399) -- (611,402) -- (612,403) -- (612,406) -- (613,407) -- (612,409) -- (612,414) -- (613,415) -- (613,420) -- (615,422) -- (615,427) -- (614,429) -- (616,433) -- (616,435) -- (615,436) -- (615,438) -- (616,439) -- (616,444) -- (617,445) -- (617,450) -- (616,451) -- (616,453) -- (618,457) -- (618,466) -- (619,467) -- (619,472) -- (618,473) -- (612,472) -- (603,465) -- (592,461) -- (590,459) -- (583,456) -- (581,454) -- (576,452) -- (574,450) -- (564,445) -- (563,443) -- (561,442) -- (558,442) -- (555,438) -- (552,438) -- (550,436) -- (548,436) -- (546,434) -- (541,432) -- (538,429) -- (539,428) -- (539,425) -- (540,424) -- (540,422) -- (542,418) -- (542,414) -- (543,413) -- (543,409) -- (542,408) -- (542,406) -- (544,404) -- (544,401) -- (545,400) -- (544,398) -- (545,396) -- (545,392) -- (547,390) -- (548,376) -- (549,375) -- (549,371) -- (551,367) -- (551,365) -- (550,364) -- (550,360) -- (551,359) -- (552,353) -- (553,352) -- (553,346) -- (554,345) -- (554,342) -- (555,341) -- (556,326) -- (558,323) -- (559,311) -- (561,309) -- (561,306) -- (564,302) -- (564,300) -- (562,298) -- (562,295) -- (563,294) -- (563,285) -- (565,282) -- (565,273) -- (567,270) -- (568,257) -- (570,254) -- (570,250) -- (571,249) -- (571,244) -- (572,243) -- (572,234) -- (573,233) -- (573,227) -- (574,226) -- (574,219) -- (577,214) -- (577,205) -- (578,204) -- (578,200) -- (579,199) -- (579,194) -- (578,192) -- (580,190) -- (580,186) -- (582,183) -- (582,177) -- (583,176) -- (583,171) -- (584,170) -- (584,162) -- cycle
(583,150) -- (584,152) -- (583,153) -- (583,157) -- (581,160) -- (580,171) -- (579,172) -- (579,174) -- (578,175) -- (578,179) -- (577,180) -- (578,183) -- (577,184) -- (576,196) -- (575,197) -- (575,200) -- (574,201) -- (574,205) -- (572,207) -- (571,207) -- (570,209) -- (570,213) -- (569,214) -- (569,219) -- (568,220) -- (568,222) -- (569,222) -- (571,224) -- (570,226) -- (570,232) -- (569,233) -- (569,238) -- (568,239) -- (567,250) -- (564,257) -- (564,261) -- (563,262) -- (563,265) -- (562,266) -- (562,273) -- (561,274) -- (560,285) -- (559,286) -- (559,289) -- (558,290) -- (558,295) -- (557,296) -- (557,300) -- (556,301) -- (556,304) -- (555,305) -- (555,311) -- (554,312) -- (553,322) -- (552,323) -- (552,330) -- (551,331) -- (551,336) -- (550,337) -- (550,340) -- (549,341) -- (550,344) -- (549,345) -- (549,347) -- (547,349) -- (548,351) -- (548,354) -- (547,355) -- (547,361) -- (545,365) -- (545,372) -- (544,373) -- (544,376) -- (543,377) -- (543,385) -- (542,386) -- (542,389) -- (541,390) -- (541,393) -- (542,394) -- (541,400) -- (540,401) -- (541,406) -- (539,408) -- (539,414) -- (537,418) -- (537,421) -- (538,423) -- (536,426) -- (536,430) -- (530,433) -- (527,436) -- (524,437) -- (522,439) -- (520,439) -- (517,442) -- (515,442) -- (513,444) -- (510,445) -- (507,448) -- (504,449) -- (495,456) -- (488,459) -- (485,462) -- (479,465) -- (470,472) -- (465,474) -- (458,480) -- (455,481) -- (452,484) -- (449,485) -- (445,489) -- (439,492) -- (434,497) -- (432,497) -- (430,499) -- (425,501) -- (419,506) -- (415,507) -- (413,509) -- (410,510) -- (406,514) -- (396,519) -- (394,521) -- (390,523) -- (386,527) -- (380,530) -- (378,530) -- (377,529) -- (377,527) -- (381,523) -- (384,517) -- (386,515) -- (386,513) -- (388,509) -- (391,506) -- (391,504) -- (393,502) -- (394,498) -- (397,495) -- (398,491) -- (402,483) -- (405,480) -- (406,475) -- (411,469) -- (411,466) -- (414,460) -- (417,457) -- (418,454) -- (420,452) -- (422,447) -- (425,443) -- (425,441) -- (427,439) -- (429,434) -- (431,432) -- (431,430) -- (434,427) -- (434,424) -- (437,421) -- (440,413) -- (443,410) -- (443,407) -- (446,404) -- (447,400) -- (450,397) -- (450,395) -- (455,386) -- (455,384) -- (458,381) -- (460,377) -- (460,375) -- (465,369) -- (468,363) -- (468,361) -- (471,358) -- (471,356) -- (473,352) -- (475,350) -- (480,339) -- (483,336) -- (484,330) -- (486,328) -- (488,323) -- (490,321) -- (490,319) -- (492,315) -- (494,313) -- (494,311) -- (497,308) -- (497,306) -- (502,297) -- (502,295) -- (506,291) -- (506,288) -- (508,286) -- (508,284) -- (512,279) -- (514,274) -- (517,271) -- (517,269) -- (519,267) -- (520,263) -- (522,261) -- (522,259) -- (524,257) -- (528,249) -- (528,247) -- (531,244) -- (533,238) -- (537,234) -- (537,231) -- (539,229) -- (541,224) -- (543,222) -- (546,216) -- (546,214) -- (549,211) -- (550,211) -- (551,209) -- (551,205) -- (554,201) -- (556,196) -- (558,194) -- (558,192) -- (561,186) -- (565,182) -- (565,180) -- (568,174) -- (571,171) -- (571,168) -- (573,166) -- (574,162) -- (579,155) -- (580,151) -- (581,150) -- cycle;
\node[inner sep=0pt,font=\fontsize{10}{10}\selectfont] at (45,739) {\scalebox{6.7}{$I$}};
\node[inner sep=0pt,font=\fontsize{10}{10}\selectfont] at (619,36) {\scalebox{6.7}{$II$}};
\node[inner sep=0pt,font=\fontsize{10}{10}\selectfont] at (1547,951) {\scalebox{6.7}{$III$}};
\node[inner sep=0pt,font=\fontsize{10}{10}\selectfont] at (340,518) {\scalebox{5.2}{$1$}};
\node[inner sep=0pt,font=\fontsize{10}{10}\selectfont] at (635,152) {\scalebox{5.2}{$2$}};
\node[inner sep=0pt,font=\fontsize{10}{10}\selectfont] at (657,463) {\scalebox{5.2}{$3$}};
\node[inner sep=0pt,font=\fontsize{10}{10}\selectfont] at (515,417) {\scalebox{5.2}{$4$}};
\node[inner sep=0pt,font=\fontsize{10}{10}\selectfont] at (426,476) {\scalebox{5.5}{$\beta$}};
\node[inner sep=0pt,font=\fontsize{10}{10}\selectfont] at (454,500) {\scalebox{3.5}{$\alpha$}};
\node[inner sep=0pt,font=\fontsize{10}{10}\selectfont] at (547,251) {\scalebox{3.6}{$\gamma$}};
\node[inner sep=0pt,font=\fontsize{10}{10}\selectfont] at (581,261) {\scalebox{3.6}{$\delta$}};
\node[inner sep=0pt,font=\fontsize{10}{10}\selectfont] at (590,438) {\scalebox{5.5}{$\varepsilon$}};
\node[inner sep=0pt,font=\fontsize{10}{10}\selectfont] at (567,468) {\scalebox{3.7}{$\zeta$}};
\node[inner sep=0pt,font=\fontsize{10}{10}\selectfont] at (155,647) {\scalebox{4.2}{$\delta$}};
\node[inner sep=0pt,font=\fontsize{10}{10}\selectfont] at (161,705) {\scalebox{5.5}{$\zeta$}};
\node[inner sep=0pt,font=\fontsize{10}{10}\selectfont] at (575,121) {\scalebox{4.0}{$\beta$}};
\node[inner sep=0pt,font=\fontsize{10}{10}\selectfont] at (610,122) {\scalebox{5.5}{$\varepsilon$}};
\node[inner sep=0pt,font=\fontsize{10}{10}\selectfont] at (1387,870) {\scalebox{4.0}{$\delta$}};
\node[inner sep=0pt,font=\fontsize{10}{10}\selectfont] at (1380,912) {\scalebox{5.5}{$\alpha$}};
\end{tikzpicture}%
}
\end{center}

\[
\begin{array}{rl}
(4,1) : (4,3) & = (4,\text{III}) : (4,\text{I}) \\
(4,1) : (4,2) & = (4,\text{II}) : (4,\text{I}),
\end{array}
\]
dann folgt schon von selbst:
\[
(4,3) : (4,2) = (4,\text{II}) : (4,\text{III}).
\]

Man hat demnach
\[
\triangle\ (4,1,3) \sim \triangle\ (4,\text{III},\text{I}),
\]
%
\SourceLeaf{340}
also:
\[
\begin{array}{rl}
(4,1) : (4,\text{III}) & = (3,1) : (\text{I}, \text{III}) \\
(4,1) : (4,\text{II}) & = (2,1) : (\text{I}, \text{II}) \\
(4,3) : (4,\text{II}) & = (2,3) : (\text{III}, \text{II}).
\end{array}
\]
Die beiden letzten Gleichungen geben durch Division:
\begin{equation}
(4,1) : (4,3) = \frac{(2,1)}{(2,3)} : \frac{(\text{I}, \text{II})}{(\text{III}, \text{II})}
= \frac{(2,1)}{(3,2)} : \frac{(\text{I}, \text{II})}{(\text{II}, \text{III})}.
\tag{1}
\end{equation}
Fügt man zum System der drei Gleichungen noch hinzu:
\[
(4,2) : (4,\text{III}) = (3,2) : (\text{II}, \text{III}),
\]
so hat man
\begin{equation}
(4,1) : (4,2) = \frac{(3,1)}{(3,2)} : \frac{(\text{I}, \text{III})}{(\text{II}, \text{III})}
\tag{2}
\end{equation}
und ebenso:
\begin{equation}
(4,2) : (4,3) = \frac{(2,1)}{(3,1)} : \frac{(\text{I}, \text{II})}{(\text{I}, \text{III})}.
\tag{3}
\end{equation}
Aus (1) folgt:
\[
\frac{(\text{I}, \text{II})}{(\text{II}, \text{III})} = \frac{(2,1) . (4,3)}{(4,1) . (3,2)},
\]
aus (2) folgt:
\[
\frac{(\text{I}, \text{III})}{(\text{II}, \text{III})} = \frac{(3,1) . (4,2)}{(4,1) . (3,2)}.
\]
Demnach ist:
\[
(\text{I}, \text{II}) : (\text{II}, \text{III}) : (\text{III}, \text{I}) = (2,1) . (4,3) :
(4,1) . (3,2) : (4,2) . (3,1).
\]
Bezeichnen wir die Winkel wie vorher, und setzen:
\[
\sphericalangle\ \text{I}, \text{II}, \text{III} = [\text{I}, \text{III}],
\]
so ist:
\[
(\text{I}, \text{II}) : (\text{II}, \text{III}) : (\text{III}, \text{I}) = \sin[\text{I}, \text{II}] :
\sin[\text{II}, \text{III}] : \sin[\text{III}, \text{I}].
\]
Nun ist:
\[
\begin{array}{rl}
{[\text{I}, \text{II}]} & = [\text{I}, 4] + [4, \text{II}] = \alpha + \delta \\
{[\text{II}, \text{III}]} & = [\text{II}, 4] + [4, \text{III}] = \gamma + \zeta \\
{[\text{III}, \text{I}]} & = [\text{III}, 4] + [4, \text{I}] = \beta + \varepsilon,
\end{array}
\]
also:
\[
\frac{\sin(\alpha+\delta)}{(2,1) . (4,3)} = \frac{\sin(\gamma+\zeta)}{(4,1) . (3,2)} =
\frac{\sin(\beta+\varepsilon)}{(4,2) . (3,1)},
\]
welches mit der frühern Gleichung übereinstimmt.

Ähnlich stellt sich die Aufgabe, liegt 4 ausserhalb des Dreiecks 1, 2, 3. Man nimmt dann beliebig
Eine der vier Ecken, z.\,B. wieder 4, zieht (4,1), (4,2), (4,3) und bestimmt ähnlich die Punkte
I, II, III und hierauf das Dreieck (I), (II), (III), indem man wieder die Relationen zwischen
Seiten und Winkeln aufsucht.
%
\SourceLeaf{341}

\begin{center}
[3.]
\end{center}

Schliesslich wollen wir zeigen, wie man von dem Beweis in [1] auf den in [2] kommen kann, indem
wirklich der Beweis zuerst durch Laterale gefunden ist, und erst später durch die geometrische
Betrachtung.

Untersuchen wir, wie sich der Beweis in [1] stellt, wenn ein Punkt des Vierecks, z.\,B. (4), mit
dem Anfangspunkt der Zählung coincidirt. Man hat dann $p^{\mathrm{IV}} = 0$, also sind die
Coordinaten
\[
\begin{array}{rl}
\text{von I:} & p'p'' \\
\text{von II:} & p'p''' \\
\text{von III:} & p''p'''.
\end{array}
\]
Also sind nach [2] die Seiten:
\[
\begin{array}{rl}
(\text{I}, \text{II}) & = (4,1) . (4,2) \\
(\text{II}, \text{III}) & = (4,1) . (4,3) \\
(\text{III}, \text{I}) & = (4,2) . (4,3).
\end{array}
\]
Folglich verhalten sich die Seiten des Dreiecks
\[
(\text{I}, \text{II}) : (\text{II}, \text{III}) : (\text{III}, \text{I}) = (4,1) . (4,2) :
(4,1) . (4,3) : (4,2) . (4,3),
\]
wie wir sie in [2] durch Construction bestimmten.
%
\clearpage
\clearpage\SourceLeaf{342}
\begin{center}
\textbf{[DER KREIS.]}\par\medskip
[I.]\par\medskip
\textbf{COMPLEXE ZAHLEN}\par
\textbf{AUF DEN RAUM ANGEWANDT.}
\end{center}

$\alpha$, $\beta$, $\gamma$, $\delta$ u.\,s.\,w. cyklische Grössen, zugleich Bezeichnung der Punkte in der Kreisperipherie, mit Radius 1.

\Needspace{10\baselineskip}1. Inhalt des Dreiecks $\alpha$, $\beta$, $\gamma$
\[
=\frac{(\beta-\alpha)(\gamma-\beta)(\alpha-\gamma)}{4\alpha\beta\gamma}\cdot i
\qquad \text{des gleichseitigen: }\frac{\BareRoot{27}}{4}.
\]
\[
\begin{aligned}
&=\frac{\alpha\alpha\beta+\beta\beta\gamma+\gamma\gamma\alpha-\alpha\beta\beta-\beta\gamma\gamma-\gamma\alpha\alpha}{4\alpha\beta\gamma}\cdot i\\
&=\tfrac{1}{4}i\cdot\left\{\frac{\beta}{\alpha}+\frac{\gamma}{\beta}+\frac{\alpha}{\gamma}-\frac{\alpha}{\beta}-\frac{\beta}{\gamma}-\frac{\gamma}{\alpha}\right\}.
\end{aligned}
\]

\Needspace{9\baselineskip}2. Inhalt des Dreiecks, welches den Kreis in den Punkten $\alpha$, $\beta$, $\gamma$ berührt, und dessen Winkelpunkte durch
\[
\frac{2\beta\gamma}{\beta+\gamma},\qquad \frac{2\gamma\alpha}{\gamma+\alpha},\qquad \frac{2\alpha\beta}{\alpha+\beta}
\]
dargestellt werden;
\[
\text{Inhalt: }\frac{(\gamma-\beta)(\alpha-\gamma)(\beta-\alpha)}{(\gamma+\beta)(\alpha+\gamma)(\beta+\alpha)}\cdot i,
\qquad\text{des gleichseitigen: }\BareRoot{27}.
\]
%
\SourceLeaf{343}

Das Centrum des darum beschriebenen Kreises

\[
= \frac{2\alpha\beta\gamma(\alpha+\beta+\gamma)}{(\beta+\gamma)(\gamma+\alpha)(\alpha+\beta)}.
\]

Product des Halbmessers in den Modulus von $(\beta+\gamma)(\gamma+\alpha)(\alpha+\beta)$

\[
= 2.
\]

Clisen von $\dfrac{\beta\gamma}{\alpha}$, $\dfrac{\alpha\gamma}{\beta}$, $\dfrac{\alpha\beta}{\gamma}$ entsprechen den drei Halbmessern dieses Kreises.

\medskip
\noindent 3. Vier Zahlen $a$, $b$, $c$, $d$ sind einander harmonisch zugeordnet, wenn das
Product der Summen jedes Paars doppelt so gross ist wie die Summe der
Producte jedes Paars, oder wenn

\[
(a+b)(c+d) = 2ab+2cd.
\]

Sind die Zahlen reell, so stellen sie vier Punkte in gerader Linie vor, die
hiernach harmonisch getheilt ist.

Mit jener Gleichung ist verbunden

\[
(c-a)(d-b)+(d-a)(c-b) = 0.
\]

\Needspace{9\baselineskip}Die Berechnung jeder aus den drei \"ubrigen geschieht durch die Formeln

\[
\begin{aligned}
a &= \frac{bc+bd-2cd}{-d-c+2b} = \frac{b(c+d)-2cd}{2b-(c+d)} = \frac{2bb-2cd}{2b-(c+d)}-b \\
  &= +b-\frac{2(c-b)(d-b)}{2b-c-d} \\
  &= +c-\frac{(d-c)(c-b)}{2b-c-d}.
\end{aligned}
\]

Sind $a$, $b$, $c$, $d$ complexe Zahlen, so liegen die dadurch vorgestellten Punkte
in einem Kreise und die Producte der Seiten $ac.bd$ und $ad.bc$ sind gleich.
%
\SourceLeaf{344}

\Needspace{7\baselineskip}\begin{center}
[II.]

\textbf{[KREIS DURCH VIER PUNKTE.]}

[1.]
\end{center}

Der Kreis durch die Punkte, die den complexen Gr\"ossen $a$, $b$, $c$ entsprechen, deren Adjuncten $A$, $B$, $C$, hat zum Mittelpunkte das Correlat von

\[
\begin{aligned}
&\frac{Aa(b-c)+Bb(c-a)+Cc(a-b)}{A(b-c)+B(c-a)+C(a-b)} \\
&= \frac{(Aa-Bb)(a-c)-(a-b)(Aa-Cc)}{(A-B)(a-c)-(a-b)(A-C)};
\end{aligned}
\]

einfacher, wenn $a=A=0$:

\[
\frac{(B-C)bc}{Bc-Cb}.
\]

\Needspace{7\baselineskip}\begin{center}
[2.]
\end{center}

Die Bedingungsgleichung, dass vier Punkte $0$, $a$, $b$, $c$, deren Adjuncten
$0$, $A$, $B$, $C$, in Einem Kreise liegen, ist

\[
bcAB+caBC+abCA-abBC-bcCA-caAB = 0
\]
oder einfacher

\[
\frac{1}{aC}+\frac{1}{bA}+\frac{1}{cB}-\frac{1}{cA}-\frac{1}{aB}-\frac{1}{bC} = 0
\]
oder auch

\[
c(a-b)B(A-C) = b(a-c)C(A-B)
\]
etc.

\Needspace{7\baselineskip}\begin{center}
[3.]
\end{center}

Die Punkte, welche dreien jener Punkte in Beziehung auf den vierten
reciprok sind, liegen in Einer geraden Linie.

Dieser Satz ist unter dem allgemeinern begriffen:

\begin{quote}
»Alle Punkte, welche in Beziehung auf einen gegebenen Punkt den
Punkten eines Kreises reciprok sind, liegen in Einem Kreise«.
\end{quote}
%
\SourceLeaf{345}

\Needspace{8\baselineskip}\begin{center}
[III.]

\textbf{GEOMETRISCHE RELATIONEN DURCH COMPLEXE ZAHLEN.}

I.
\end{center}

$M$, $M'$ zwei Punkte, deren complexe Zahlen $m$, $m'$.

$P$ ein unbestimmter Punkt, wof\"ur Azimuthe $PM'-PM$ einen gegebenen
Unterschied $= \varphi$ bilden.

Der geometrische Ort von $P$ ist ein Kreis, dessen

\[
\text{Mittelpunkt } \frac{m\,e^{i\varphi}-m'\,e^{-i\varphi}}{2i\sin\varphi},
\]

\[
\text{Halbmesser } = \text{Modulus von } \frac{m'-m}{2\sin\varphi}.
\]

\Needspace{8\baselineskip}\begin{center}
II.
\end{center}

$M$, $M'$, $m$, $m'$ wie vorhin.

$P$ ein unbestimmter Punkt, dessen Distanzen von $M$, $M'$ in einem gegebenen Verh\"altnisse $k:k'$ stehen.

Der geometrische Ort von $P$ ist ein Kreis, dessen

\[
\text{Mittelpunkt } \frac{m.k'k'-m'.kk}{k'k'-kk} = c,
\]

\[
\text{Halbmesser } = \text{Modulus von } \frac{kk'}{k'k'-kk}\cdot(m'-m).
\]



\Needspace{7\baselineskip}\begin{center}
\textbf{BEMERKUNGEN.}
\end{center}

Die drei vorstehenden Notizen r\"uhren vermuthlich aus der Zeit um 1840
her; [I] steht auf einem einzelnen Blatte, [II] und [III] sind in einem Handbuche
eingetragen. Auf den Kreis bezieht sich auch eine bereits Bd.\,IV.\,S.\,397 abgedruckte
»handschriftliche Notiz«. \hfill \textsc{Stäckel}.
%
\clearpage
\clearpage\SourceLeaf{346}

\begin{center}
\textbf{[DIE KEGELSCHNITTE.]}
\end{center}

\begin{center}
[I.]

\textbf{LEHRSATZ.}
\end{center}

Sind $t$, $t'$ die complexen Werthe zweier conjugirter Halbaxen einer
Ellipse und $e$ der complexe Werth der Excentricität (Abstand eines
Brennpunkts), so ist
\[
tt + t't' = ee.
\]

Die Punkte einer Ellipse sind durch die Formel gegeben:
\[
g e^{ix} + h e^{-ix},
\]
wo $x$ alle Werthe von $x = 0$ bis $x = 2\pi$ durchlaufen muss.

Lehrsatz. Es sind dann zwei conjugirte Halbaxen
\[
g e^{ix} + h e^{-ix} \quad \text{und} \quad g i e^{ix} - h i e^{-ix}.
\]
Also
\[
4gh = ee.
\]
%
\SourceLeaf{347}

\begin{center}
[II.]

\textbf{LEHRSÄTZE UND AUFGABEN}

\textbf{DIE COMPLEXEN ZAHLEN BETREFFEND.}
\end{center}

\Needspace{7\baselineskip}1. $\dfrac{a}{b-\rho}$ stellt, wenn $a$, $b$ bestimmte Zahlen, $\rho$ eine
unbestimmte cyklische Zahl bedeutet, einen Kreis vor,
\[
\dfrac{ab}{bb-\beta\beta}
\]
den Mittelpunkt, wenn $\beta$ der in $b$ enthaltene cyklische Factor,
\[
\dfrac{a}{bb-\beta\beta}
\]
(oder der Modulus dieser Grösse) den Halbmesser des Kreises.

\Needspace{7\baselineskip}2. (vor 1.) Ist wie oben $\beta$ der in $b$ enthaltene cyklische Factor, so
haben
\[
b + \beta\rho \quad \text{und} \quad b + \dfrac{\beta}{\rho}
\]
einerlei Modulus, imgleichen
\[
b + \beta\rho \quad \text{und} \quad \dfrac{\rho\beta\beta}{b\rho+\beta}
\]
einerlei Clise.

\Needspace{7\baselineskip}3. $\rho - \beta$ und $i\sqrt{\rho\beta}$ haben gleiche Clise.

\Needspace{7\baselineskip}4. $\Omega = \dfrac{a\rho+b}{1-\rho}$ stellt eine gerade Linie vor. Man hat
nemlich
\[
2\Omega = \dfrac{(a+b)(1+\rho) - (a-b)(1-\rho)}{1-\rho}
\]
\[
= (a+b)t - (a-b),
\]
wenn
\[
t = \dfrac{1+\rho}{1-\rho}
\]
gesetzt wird, wodurch also $t$ eine unbestimmte imaginäre Zahl wird. Es ist
ferner
\[
\dfrac{2\Omega}{a+b} = t - \dfrac{a-b}{a+b}.
\]

Derjenige Werth dieser Grösse, dessen Modulus ein Minimum ist, ist offenbar
nichts anderes als der reelle Theil von $\dfrac{a-b}{a+b}$.

Ist demnach $k$ der reelle Theil von $\dfrac{a-b}{a+b}$, so ist der
gesuchte Werth von $\Omega$
\[
= \tfrac{1}{2} k (a+b).
\]
%
\SourceLeaf{348}

\begin{center}
[III.]

\textbf{ZUR THEORIE DER KEGELSCHNITTE.}
\end{center}

Zunächst beziehen sich die Untersuchungen auf die Ellipse, sie lassen sich
aber mit geringen Modificationen auf die Hyperbel übertragen.

\Needspace{8\baselineskip}\begin{center}
[1.]
\end{center}

Es seien $m$, $n$ zunächst zwei reelle Zahlen, $\rho$ eine Grösse von der
Form $\cos\varphi + i\sin\varphi$, so dass $\rho$ alle seine Werthe
durchläuft, während $\varphi$ von $0^0$ bis $360^0$ wächst.

Es sei ferner
\[
r = m\rho + \dfrac{n}{\rho},
\]
und $R$ der Punkt, der die complexe Grösse $r$ repräsentirt; $R$ ist also
ein unbestimmter Punkt in einer Ellipse, deren Hauptaxen $m+n$, $m-n$ in
der Linie der reellen und imaginären Zahlen liegen, während ihr
Mittelpunkt mit dem Nullpunkte zusammenfällt.

Den vier Werthen von $\rho$: $\alpha$, $\beta$, $\gamma$, $\delta$
entsprechen die vier Werthe von $r$: $a$, $b$, $c$, $d$ oder die vier
Punkte $A$, $B$, $C$, $D$.

\Needspace{7\baselineskip}Endlich setze man
\[
\dfrac{-\alpha\beta\gamma m^3 + (\alpha+\beta+\gamma) m m n
  - \delta(\alpha\beta+\alpha\gamma+\beta\gamma) m n n + \delta n^3}
  {(1-\alpha\beta\gamma\delta) m n} = t
\]
und nenne den Punkt, auf welchen sich die complexe Zahl $t$ bezieht, $T$.

Man sieht leicht, dass, wenn von den vier Punkten $A$, $B$, $C$, $D$ drei
constant sind, und der vierte den Kegelschnitt durchläuft, $T$ eine gerade
Linie durchlaufen werde.

\Needspace{9\baselineskip}Durch eine leichte Rechnung findet man
\[
t - a = -\dfrac{(\alpha\beta m - n)(\alpha\gamma m - n)(m - \alpha\delta n)}
  {\alpha(1-\alpha\beta\gamma\delta) m n}
\]
\[
= -\dfrac{(a-b)(a-c)}{a-d} \cdot
  \dfrac{\alpha\beta\gamma(\alpha-\delta)\left(mm -
  \left(\alpha\delta+\dfrac{1}{\alpha\delta}\right) m n + n n\right)}
  {(\alpha-\beta)(\alpha-\gamma)(1-\alpha\beta\gamma\delta) m n}.
\]
%
\SourceLeaf{349}

Hieraus folgt sogleich von selbst, dass
\[
(t-a)(d-a) \quad \text{und} \quad (b-a)(c-a)
\]
gleiche Clisen haben oder dass die Winkel
\[
\text{zwischen } AT, AC \quad \text{und} \quad \text{zwischen } AB, AD
\]
gleich sind.

\Needspace{8\baselineskip}Es lässt sich auch $t$ in folgende elegante Form setzen
\[
t = \dfrac{\left(m-\dfrac{n}{\beta\gamma}\right)
  \left(m-\dfrac{n}{\alpha\gamma}\right)
  \left(m-\dfrac{n}{\alpha\beta}\right)}
  {\left(\delta - \dfrac{1}{\alpha\beta\gamma}\right) m n}
  + \left(\left(\dfrac{1}{\alpha}+\dfrac{1}{\beta}+\dfrac{1}{\gamma}\right) m
  - \dfrac{1}{\alpha\beta\gamma} n\right) \dfrac{n}{m}
\]
oder auch in folgende
\[
t = \dfrac{(\beta\gamma m - n)(m - \beta\delta n)(m-\gamma\delta n)}
  {\alpha\beta\gamma\delta\left(\beta\gamma\delta - \dfrac{1}{\alpha}\right) m n}
  + \dfrac{\beta\gamma m m - m n + (\beta+\gamma)\delta n n}{\beta\gamma\delta n}.
\]

\Needspace{10\baselineskip}\begin{center}
[2.]
\end{center}

Das Centrum eines durch drei Punkte der Ellipse $A$, $B$, $C$ gehenden
Kreises ist
\[
\dfrac{\left\{\alpha\beta\gamma +
  \left(\dfrac{1}{\alpha}+\dfrac{1}{\beta}+\dfrac{1}{\gamma}\right)\right\} m
  - \left\{(\alpha+\beta+\gamma) + \dfrac{1}{\alpha\beta\gamma}\right\} n}
  {\dfrac{m}{n} - \dfrac{n}{m}},
\]
der Halbmesser aequal dem Modulus von \[
\dfrac{(m-\beta\gamma n)(m-\alpha\gamma n)(m-\alpha\beta n)}{mm - nn}
\]
oder (was dasselbe ist) von
\[
\dfrac{\left(m-\dfrac{n}{\beta\gamma}\right)\left(m-\dfrac{n}{\alpha\gamma}\right)
  \left(m-\dfrac{n}{\alpha\beta}\right)}{mm-nn}.
\]

Das Centrum des erwähnten Kreises kann auch so ausgedrückt werden, wenn zur
Abkürzung
\[
\beta\gamma = \cos 2\mathfrak{A} + i\sin 2\mathfrak{A}
\]
\[
\alpha\gamma = \cos 2\mathfrak{B} + i\sin 2\mathfrak{B}
\]
\[
\alpha\beta = \cos 2\mathfrak{C} + i\sin 2\mathfrak{C}
\]%
%
\SourceLeaf{350}%
geschrieben, oder die Clisen von $\beta\gamma$, $\alpha\gamma$, $\alpha\beta$ mit $2\mathfrak{A}$, $2\mathfrak{B}$, $2\mathfrak{C}$ bezeichnet werden:
\[
\frac{4\cos\mathfrak{A}\cos\mathfrak{B}\cos\mathfrak{C}}{\dfrac{1}{n}+\dfrac{1}{m}} - \frac{4\sin\mathfrak{A}\sin\mathfrak{B}\sin\mathfrak{C}}{\dfrac{1}{n}-\dfrac{1}{m}}\, i.
\]
Es sind hierbei aber die Halbkreise, in welchen $\mathfrak{A}$, $\mathfrak{B}$, $\mathfrak{C}$ genommen werden, nur für zwei willkürlich, z.\,B. für $\mathfrak{A}$, $\mathfrak{B}$; für den dritten muss werden
\[
\gamma = \cos(\mathfrak{A}+\mathfrak{B}-\mathfrak{C}) + i\sin(\mathfrak{A}+\mathfrak{B}-\mathfrak{C}).
\]
\Needspace{12\baselineskip}Mit andern Worten, hätte man gemacht
\[
a = \cos a^{*} + i\sin a^{*}
\]
\[
\beta = \cos \beta^{*} + i\sin \beta^{*}
\]
\[
\gamma = \cos \gamma^{*} + i\sin \gamma^{*},
\]
so würde man zu setzen haben
\[
2\mathfrak{A} = \beta^{*}+\gamma^{*}
\]
\[
2\mathfrak{B} = a^{*}+\gamma^{*}
\]
\[
2\mathfrak{C} = a^{*}+\beta^{*}.
\]
\Needspace{6\baselineskip}Der Halbmesser des Kreises wird
\begingroup\small
\[
= \frac{\BareRoot{((mm+nn-2mn\cos 2\mathfrak{A})(mm+nn-2mn\cos 2\mathfrak{B})(mm+nn-2mn\cos 2\mathfrak{C}))}}{mm-nn}.
\]
\endgroup

\Needspace{11\baselineskip}\begin{center}[3.]\end{center}

Bei der \textsc{Newton}schen Construction hat die Differenz zweier Werthe von $t$ die Clise wie
\[
\frac{\dfrac{\beta\gamma}{\beta-\gamma}(b-c)\times\dfrac{\beta-\delta}{b-d}\times\dfrac{\gamma-\delta}{c-d}}{\beta\gamma\delta}
\left\{\frac{1}{\alpha\beta\gamma\delta-1}-\frac{1}{\alpha'\beta\gamma\delta-1}\right\}
\]
oder, weil $\beta-\delta$ die Clise [wie] $i\sqrt{\beta\delta}$ u.\,s.\,w.,
\[
\frac{(\beta-\delta)(\gamma-\delta)}{\delta(\beta-\gamma)}
\left\{\frac{(\alpha-\alpha')\beta\gamma\delta}{(\alpha\beta\gamma\delta-1).(\alpha'\beta\gamma\delta-1)}\right\}
\frac{b-c}{(b-d)(c-d)}
\]
oder, weil Clise von $\alpha-\alpha'$ wie $i\sqrt{\alpha\alpha'}$ u.\,s.\,w., wie
\[
\frac{b-c}{(b-d)(c-d)}.
\]
%
\SourceLeaf{351}
In meiner (der ersten) Construction, wo $\delta$ unbestimmt, ist die Clise wie
\[
(a-b)(a-c)(b-c).
\]
Nennt man $e$ den vierten Schnitt[punkt] der Ellipse mit dem durch $b$, $c$, $d$ gelegten Kreise und setzt
\[
e = \varepsilon m + \frac{n}{\varepsilon},
\]
wodurch
\[
\beta\gamma\delta\varepsilon = 1
\]
wird, so ist, in \textsc{Newton}s Construction, die Clise der Differenz zweier Werthe von $t$ wie
\[
(c-b)(b-e)(c-e).
\]
\Needspace{10\baselineskip}Auch kann obige Formel ($a=\varepsilon\rho$) besser so geschrieben werden:
\[
t = -\frac{(m-\delta\beta n)(m-\delta\gamma n)(m-\delta\varepsilon n)\rho}{\delta(1-\rho)mn} + \frac{(\beta+\gamma)mm-\dfrac{1}{\varepsilon}mn+\delta nn}{m}
\]
\[
= -\frac{(b-c)(b-e)(c-e)}{\delta(\beta-\gamma)(\beta-\varepsilon)(\gamma-\varepsilon)}\cdot\frac{\rho}{(1-\rho)mn} + \frac{(\beta+\gamma)mm-\dfrac{1}{\varepsilon}mn+\delta nn}{m}.
\]

\Needspace{10\baselineskip}\begin{center}\textbf{BEMERKUNGEN.}\end{center}

Die Notiz [III], die auf einem einzelnen Blatte steht, gibt den Inhalt von Untersuchungen wieder, über die \textsc{Gauss} in dem S.\,295 bis 296 dieses Bandes abgedruckten Briefe an \textsc{Schumacher} vom 12. Mai 1843 berichtet hat; auf diesen möge zur Erläuterung des Textes verwiesen werden. Die Notizen [I] und [II], die mit [III] in engem Zusammenhange stehen, finden sich in einem Handbuche und auf einem Zettel; [I] stammt wahrscheinlich bereits aus dem Jahre 1831, [II] wohl aus der Zeit um 1840.
\hfill \textsc{Stäckel}.
%
\clearpage
\SourceLeaf{352}
\begin{center}\large\bfseries [PROJECTION DES WÜRFELS.]\end{center}

\Needspace{8\baselineskip}\begin{center}[1.]\end{center}

Sind die complexen Werthe der orthographischen Projectionen von drei gleich langen und unter einander senkrechten Geraden
\[
a,\ b,\ c,
\]
so ist
\[
aa+bb+cc=0,
\]
z.\,B.
\[
a=1+8i,\qquad b=8+i,\qquad c=4-4i.
\]
Allgemein kann man setzen
\[
a=pp+qq,\qquad b=2ipq,\qquad c=ipp-iqq.
\]

\Needspace{9\baselineskip}\begin{center}[2.]\end{center}

Die zierlichste Form der Auflösung der allgemeinen Gleichung
\[
pp+qq+rr=0
\]
ist, $a$ und $b$ beliebige complexe Zahlen bedeutend,
\[
p=(a-b)(b-aii)
\]
\[
q=(b-bi)(aii-ai)
\]
\[
r=(bi-a)(ai-b).
\]

\Needspace{9\baselineskip}\begin{center}\bfseries BEMERKUNGEN.\end{center}

Ein Theil dieser beiden Notizen, die sich in einem Handbuche befinden und aus dem Jahre 1831 stammen, ist bereits in Bd.\,II, S.\,309 abgedruckt worden.

\indent Vom December 1839 bis Ostern 1840 hat \textsc{Gauss} eine Vorlesung über die \emph{Theorie der imaginären Grössen} gehalten, die E.\,\textsc{Heine} gehört und ausgearbeitet hat, wie bereits S.\,330 angeführt. Dieser Ausarbeitung sind die folgenden auf die Projection des Würfels bezüglichen Ausführungen entnommen.
\par\noindent\hfill\textsc{Stäckel}.\par
%
\clearpage\SourceLeaf{353}

\begin{center}[\emph{Vorlesung über die Theorie der imaginären Grössen.}]\end{center}

\begin{center}\bfseries Projection des Würfels.\end{center}

Will man einen Würfel auf eine Ebene projiciren, und zwar durch Senkrechte von den Ecken auf die Ebene, so muss die Projection die Eigenschaft haben, dass die Summe der Quadrate der Zahlen, die die Projection je dreier auf einander stossender Kanten anzeigen, $=0$ sei.

Um diess zu zeigen, bezeichnen wir mit $0,1,7,3$ die Ecken der einen, mit $2,4,5,6$ die der damit parallelen Fläche. Ist $(0,1)$, die Seite des Würfels, $=r$, so haben wir, dass an der Ecke $0$ die drei Kanten zusammenstossen, deren jede $=r$ ist:

\begin{center}\resizebox{4.6cm}{!}{%
\begin{tikzpicture}[x=1pt,y=-1pt]
\path[use as bounding box] (0,0) rectangle (767,728);
\path[fill=black,draw=none,even odd rule]
(118,628) -- (117,627) -- (115,627) -- (110,632) -- (105,635) -- (102,640) -- (103,641) -- (106,641) -- (108,639) -- (109,639) -- (118,631) -- cycle
(140,604) -- (138,604) -- (135,608) -- (134,608) -- (126,615) -- (126,616) -- (124,618) -- (124,619) -- (126,621) -- (130,621) -- (135,616) -- (135,615) -- (140,611) -- (140,610) -- (142,608) -- (142,606) -- cycle
(163,586) -- (162,585) -- (159,585) -- (145,598) -- (145,601) -- (146,602) -- (150,602) -- (163,589) -- cycle
(186,566) -- (185,565) -- (183,565) -- (179,569) -- (174,572) -- (170,576) -- (170,577) -- (172,579) -- (174,578) -- (177,578) -- (186,570) -- cycle
(208,547) -- (206,545) -- (205,545) -- (203,547) -- (200,548) -- (197,551) -- (197,553) -- (191,558) -- (191,560) -- (192,561) -- (195,561) -- (197,559) -- (198,559) -- (208,550) -- cycle
(230,527) -- (228,525) -- (227,525) -- (225,527) -- (224,527) -- (215,536) -- (215,537) -- (216,538) -- (220,538) -- (224,536) -- (229,532) -- (230,530) -- cycle
(251,510) -- (249,508) -- (246,508) -- (235,518) -- (235,520) -- (237,522) -- (241,521) -- (250,513) -- cycle
(270,492) -- (269,491) -- (266,491) -- (264,492) -- (256,499) -- (257,502) -- (258,503) -- (262,503) -- (270,495) -- cycle
(647,485) -- (647,487) -- (648,488) -- (654,488) -- (655,487) -- (664,487) -- (665,486) -- (665,484) -- (664,483) -- (655,483) -- (654,484) -- (648,484) -- cycle
(618,485) -- (619,487) -- (624,487) -- (625,488) -- (632,488) -- (635,486) -- (634,484) -- (625,484) -- (624,483) -- (622,483) -- cycle
(589,484) -- (590,486) -- (595,486) -- (596,487) -- (608,487) -- (609,486) -- (609,485) -- (607,483) -- (590,483) -- cycle
(557,483) -- (558,485) -- (560,485) -- (561,486) -- (578,486) -- (579,485) -- (579,484) -- (577,482) -- (568,482) -- (567,481) -- (565,482) -- (560,481) -- cycle
(525,482) -- (525,484) -- (527,486) -- (529,486) -- (530,487) -- (532,486) -- (540,486) -- (541,485) -- (546,485) -- (547,484) -- (547,482) -- (546,481) -- (534,481) -- (533,482) -- (531,482) -- (530,481) -- (526,481) -- cycle
(464,481) -- (464,483) -- (465,484) -- (468,484) -- (469,485) -- (479,485) -- (480,484) -- (483,485) -- (486,483) -- (485,481) -- (482,481) -- (479,479) -- (470,479) -- (469,480) -- (465,480) -- cycle
(404,483) -- (405,484) -- (413,484) -- (414,485) -- (419,485) -- (420,484) -- (424,484) -- (426,481) -- (425,480) -- (420,480) -- (419,479) -- (406,479) -- (404,481) -- cycle
(434,480) -- (434,481) -- (436,483) -- (438,483) -- (439,484) -- (454,484) -- (456,482) -- (456,481) -- (454,479) -- (443,479) -- (441,478) -- (440,479) -- (435,479) -- cycle
(369,481) -- (369,482) -- (372,484) -- (379,484) -- (380,485) -- (382,484) -- (385,484) -- (387,485) -- (390,482) -- (390,481) -- (387,479) -- (381,479) -- (380,478) -- (374,478) -- cycle
(338,480) -- (338,481) -- (340,483) -- (360,483) -- (361,482) -- (361,480) -- (358,478) -- (340,478) -- cycle
(309,480) -- (310,482) -- (315,482) -- (316,481) -- (319,481) -- (320,482) -- (326,482) -- (327,483) -- (329,483) -- (333,481) -- (333,479) -- (332,478) -- (327,478) -- (324,476) -- (312,477) -- cycle
(301,479) -- (297,477) -- (293,477) -- (292,478) -- (290,478) -- (288,476) -- (288,466) -- (287,465) -- (285,465) -- (284,466) -- (284,469) -- (283,470) -- (284,472) -- (284,477) -- (282,479) -- (278,479) -- (275,482) -- (274,486) -- (275,487) -- (279,487) -- (285,483) -- (295,483) -- (296,482) -- (300,482) -- (301,481) -- cycle
(288,436) -- (286,436) -- (284,438) -- (284,441) -- (285,442) -- (284,444) -- (284,454) -- (285,455) -- (287,455) -- (289,453) -- (289,451) -- (288,450) -- (288,448) -- (289,447) -- (289,437) -- cycle
(289,408) -- (287,408) -- (284,411) -- (284,413) -- (285,414) -- (284,424) -- (286,426) -- (288,426) -- (290,423) -- (290,409) -- cycle
(286,381) -- (284,384) -- (285,386) -- (285,397) -- (287,399) -- (289,399) -- (290,398) -- (291,386) -- (290,385) -- (290,383) -- (288,381) -- cycle
(288,355) -- (287,355) -- (285,358) -- (285,370) -- (284,371) -- (284,373) -- (285,374) -- (289,374) -- (291,371) -- (290,370) -- (290,363) -- (291,362) -- (291,358) -- cycle
(289,330) -- (285,332) -- (285,339) -- (286,341) -- (284,345) -- (286,347) -- (289,347) -- (294,340) -- (291,337) -- (291,334) -- cycle
(287,304) -- (285,308) -- (285,321) -- (287,323) -- (290,323) -- (292,319) -- (292,316) -- (291,315) -- (291,312) -- (292,310) -- (291,309) -- (290,305) -- (289,304) -- cycle
(288,280) -- (286,282) -- (286,298) -- (287,299) -- (290,299) -- (292,295) -- (292,288) -- (291,287) -- (292,283) -- (289,280) -- cycle
(288,257) -- (286,259) -- (286,262) -- (285,263) -- (285,270) -- (286,271) -- (286,274) -- (288,276) -- (290,276) -- (292,274) -- (292,272) -- (290,269) -- (291,268) -- (291,266) -- (290,265) -- (290,258) -- cycle
(717,69) -- (716,68) -- (700,68) -- (699,67) -- (696,68) -- (695,67) -- (690,67) -- (689,66) -- (687,66) -- (686,67) -- (660,67) -- (659,66) -- (656,66) -- (655,67) -- (626,67) -- (625,66) -- (621,66) -- (619,67) -- (618,66) -- (611,66) -- (610,65) -- (608,65) -- (607,66) -- (605,66) -- (604,65) -- (565,65) -- (564,64) -- (562,65) -- (560,65) -- (559,64) -- (535,64) -- (534,65) -- (524,65) -- (523,64) -- (518,64) -- (516,65) -- (515,64) -- (486,64) -- (485,63) -- (481,63) -- (480,64) -- (469,64) -- (468,63) -- (457,63) -- (456,62) -- (453,62) -- (452,63) -- (448,63) -- (447,62) -- (439,63) -- (438,62) -- (426,62) -- (425,61) -- (423,62) -- (405,62) -- (403,63) -- (402,62) -- (368,62) -- (367,61) -- (358,61) -- (357,62) -- (355,62) -- (354,61) -- (309,60) -- (308,61) -- (303,61) -- (302,62) -- (291,62) -- (282,70) -- (278,72) -- (272,78) -- (269,79) -- (264,84) -- (261,85) -- (244,101) -- (242,101) -- (221,119) -- (217,121) -- (215,123) -- (215,125) -- (213,127) -- (211,127) -- (207,129) -- (188,147) -- (187,147) -- (182,152) -- (179,153) -- (174,159) -- (173,159) -- (158,173) -- (155,174) -- (151,178) -- (146,181) -- (140,187) -- (135,190) -- (129,196) -- (124,199) -- (117,206) -- (116,206) -- (112,210) -- (111,210) -- (104,217) -- (101,218) -- (98,221) -- (98,222) -- (90,229) -- (89,229) -- (82,235) -- (81,235) -- (76,241) -- (76,244) -- (77,245) -- (79,245) -- (81,243) -- (90,243) -- (91,244) -- (92,243) -- (93,244) -- (96,244) -- (97,245) -- (103,245) -- (104,244) -- (113,244) -- (114,245) -- (138,245) -- (139,246) -- (141,246) -- (142,245) -- (145,245) -- (146,246) -- (152,246) -- (153,245) -- (154,246) -- (177,246) -- (178,247) -- (180,246) -- (182,247) -- (194,247) -- (195,246) -- (215,246) -- (216,247) -- (218,246) -- (221,246) -- (222,247) -- (225,247) -- (226,246) -- (234,246) -- (235,247) -- (237,246) -- (240,246) -- (241,247) -- (255,247) -- (257,246) -- (258,247) -- (285,247) -- (288,250) -- (293,247) -- (301,247) -- (302,248) -- (304,248) -- (305,247) -- (307,248) -- (309,248) -- (310,247) -- (318,247) -- (319,248) -- (323,248) -- (324,247) -- (341,248) -- (342,247) -- (345,247) -- (346,248) -- (356,248) -- (357,249) -- (384,249) -- (385,248) -- (387,248) -- (388,249) -- (392,249) -- (393,250) -- (420,250) -- (421,251) -- (423,251) -- (426,249) -- (435,249) -- (436,250) -- (438,249) -- (471,249) -- (472,250) -- (481,249) -- (482,250) -- (487,250) -- (488,251) -- (491,251) -- (492,250) -- (497,250) -- (499,251) -- (500,250) -- (505,250) -- (506,249) -- (508,249) -- (509,250) -- (513,250) -- (515,252) -- (515,269) -- (516,270) -- (516,275) -- (515,276) -- (516,278) -- (516,282) -- (514,285) -- (514,289) -- (515,290) -- (515,293) -- (514,294) -- (515,296) -- (515,302) -- (514,304) -- (516,306) -- (516,307) -- (514,309) -- (515,312) -- (514,313) -- (514,316) -- (513,317) -- (514,319) -- (514,328) -- (515,329) -- (514,331) -- (515,335) -- (513,338) -- (513,340) -- (514,341) -- (514,343) -- (513,344) -- (513,348) -- (514,349) -- (513,351) -- (513,353) -- (514,354) -- (513,358) -- (514,359) -- (515,363) -- (514,364) -- (514,374) -- (513,375) -- (513,384) -- (514,385) -- (514,396) -- (513,397) -- (513,402) -- (514,403) -- (514,405) -- (513,406) -- (513,412) -- (514,413) -- (514,417) -- (513,418) -- (514,420) -- (513,422) -- (513,436) -- (514,437) -- (514,439) -- (513,440) -- (513,461) -- (512,462) -- (512,467) -- (513,468) -- (513,477) -- (514,479) -- (512,481) -- (505,481) -- (504,480) -- (502,480) -- (501,481) -- (497,481) -- (496,483) -- (498,485) -- (511,485) -- (512,486) -- (512,499) -- (513,500) -- (513,504) -- (512,505) -- (513,506) -- (513,513) -- (512,514) -- (513,516) -- (513,519) -- (512,520) -- (512,530) -- (513,531) -- (513,533) -- (511,536) -- (512,538) -- (512,557) -- (511,558) -- (511,572) -- (512,574) -- (511,575) -- (511,585) -- (510,586) -- (510,590) -- (511,591) -- (511,598) -- (510,599) -- (511,600) -- (511,604) -- (510,605) -- (510,612) -- (512,614) -- (511,625) -- (513,627) -- (514,627) -- (517,625) -- (517,621) -- (516,620) -- (517,618) -- (517,614) -- (516,613) -- (516,593) -- (517,592) -- (517,564) -- (518,563) -- (517,562) -- (518,561) -- (517,553) -- (518,552) -- (518,515) -- (519,514) -- (518,504) -- (519,503) -- (519,500) -- (517,498) -- (517,492) -- (516,491) -- (516,489) -- (517,488) -- (517,474) -- (518,473) -- (518,463) -- (519,462) -- (519,455) -- (518,454) -- (519,452) -- (519,439) -- (520,438) -- (520,432) -- (519,431) -- (519,428) -- (520,427) -- (520,415) -- (521,414) -- (520,412) -- (521,407) -- (519,405) -- (520,401) -- (521,400) -- (521,393) -- (520,392) -- (521,390) -- (520,388) -- (521,385) -- (520,384) -- (520,369) -- (519,368) -- (519,365) -- (520,364) -- (520,348) -- (521,347) -- (521,345) -- (520,344) -- (520,341) -- (521,340) -- (521,338) -- (520,337) -- (520,319) -- (521,318) -- (521,313) -- (522,312) -- (521,310) -- (521,306) -- (522,305) -- (522,303) -- (520,300) -- (520,298) -- (521,297) -- (521,293) -- (520,292) -- (521,291) -- (521,287) -- (520,285) -- (522,283) -- (522,279) -- (519,275) -- (519,272) -- (520,271) -- (519,269) -- (519,262) -- (520,261) -- (520,248) -- (522,245) -- (522,239) -- (523,238) -- (530,238) -- (544,225) -- (545,225) -- (554,216) -- (555,216) -- (559,212) -- (560,212) -- (564,208) -- (565,208) -- (576,197) -- (577,197) -- (583,191) -- (584,191) -- (596,180) -- (597,180) -- (608,169) -- (609,169) -- (613,165) -- (614,165) -- (617,161) -- (622,159) -- (626,155) -- (626,153) -- (636,144) -- (640,142) -- (646,136) -- (647,136) -- (647,135) -- (653,130) -- (653,128) -- (657,124) -- (661,123) -- (667,116) -- (668,116) -- (677,107) -- (680,106) -- (687,98) -- (689,98) -- (695,91) -- (696,91) -- (701,86) -- (702,86) -- (711,77) -- (713,76) -- (715,78) -- (715,79) -- (714,80) -- (714,88) -- (713,89) -- (713,96) -- (714,97) -- (713,99) -- (713,102) -- (714,103) -- (714,109) -- (713,110) -- (713,114) -- (712,115) -- (712,133) -- (710,136) -- (710,139) -- (712,142) -- (712,146) -- (711,147) -- (711,153) -- (712,154) -- (712,166) -- (711,167) -- (711,169) -- (712,170) -- (712,178) -- (711,179) -- (711,181) -- (712,182) -- (712,185) -- (711,186) -- (711,192) -- (710,193) -- (710,198) -- (711,199) -- (711,204) -- (710,205) -- (710,209) -- (711,210) -- (710,225) -- (711,227) -- (710,228) -- (710,239) -- (708,242) -- (708,244) -- (709,245) -- (709,267) -- (710,268) -- (709,272) -- (710,273) -- (711,279) -- (710,280) -- (710,290) -- (709,291) -- (709,297) -- (708,298) -- (708,302) -- (709,303) -- (709,305) -- (708,306) -- (709,310) -- (708,312) -- (709,313) -- (709,316) -- (708,317) -- (708,328) -- (707,329) -- (708,331) -- (708,335) -- (710,339) -- (707,342) -- (708,344) -- (708,347) -- (706,349) -- (708,353) -- (708,358) -- (707,359) -- (707,362) -- (706,363) -- (706,366) -- (707,367) -- (706,368) -- (707,369) -- (707,373) -- (706,374) -- (706,376) -- (708,378) -- (706,382) -- (707,384) -- (707,387) -- (706,389) -- (708,391) -- (708,393) -- (706,396) -- (706,400) -- (707,401) -- (707,403) -- (706,404) -- (706,406) -- (707,407) -- (707,416) -- (706,417) -- (706,425) -- (707,427) -- (706,428) -- (706,433) -- (707,435) -- (706,436) -- (706,450) -- (705,451) -- (706,453) -- (706,458) -- (705,459) -- (706,461) -- (706,464) -- (705,465) -- (705,478) -- (706,479) -- (706,482) -- (703,485) -- (701,485) -- (700,486) -- (698,485) -- (695,485) -- (694,484) -- (675,484) -- (674,486) -- (676,488) -- (692,488) -- (693,487) -- (699,487) -- (701,489) -- (687,502) -- (686,502) -- (680,509) -- (679,509) -- (673,515) -- (672,515) -- (669,518) -- (669,519) -- (666,522) -- (665,522) -- (650,538) -- (645,541) -- (641,546) -- (640,546) -- (635,551) -- (634,551) -- (627,559) -- (626,559) -- (612,572) -- (611,572) -- (606,577) -- (605,577) -- (596,586) -- (595,586) -- (584,596) -- (583,596) -- (581,598) -- (581,599) -- (579,601) -- (578,601) -- (569,611) -- (568,611) -- (554,625) -- (553,625) -- (548,630) -- (547,630) -- (539,639) -- (538,639) -- (518,658) -- (517,658) -- (515,656) -- (515,647) -- (516,646) -- (516,631) -- (515,630) -- (512,630) -- (511,631) -- (512,635) -- (511,636) -- (511,640) -- (510,641) -- (510,646) -- (511,647) -- (511,663) -- (510,664) -- (503,664) -- (501,663) -- (500,664) -- (494,664) -- (492,665) -- (488,663) -- (477,663) -- (476,664) -- (474,664) -- (473,663) -- (437,663) -- (436,662) -- (431,662) -- (430,663) -- (427,663) -- (426,664) -- (402,664) -- (401,663) -- (376,663) -- (375,662) -- (370,662) -- (368,663) -- (367,662) -- (315,662) -- (314,661) -- (312,661) -- (311,662) -- (276,662) -- (275,663) -- (273,662) -- (254,662) -- (252,663) -- (251,662) -- (232,662) -- (231,661) -- (221,661) -- (220,662) -- (218,662) -- (217,661) -- (215,662) -- (211,662) -- (210,661) -- (206,661) -- (205,662) -- (203,662) -- (202,661) -- (197,661) -- (196,662) -- (193,662) -- (192,663) -- (188,663) -- (187,662) -- (172,662) -- (170,663) -- (169,662) -- (166,662) -- (165,661) -- (163,662) -- (151,662) -- (150,661) -- (148,662) -- (146,661) -- (132,661) -- (131,662) -- (123,662) -- (122,661) -- (91,661) -- (90,662) -- (88,662) -- (86,661) -- (84,659) -- (91,653) -- (94,652) -- (95,651) -- (95,649) -- (94,648) -- (90,649) -- (83,656) -- (82,656) -- (80,654) -- (81,652) -- (81,649) -- (82,648) -- (82,643) -- (81,642) -- (81,635) -- (82,634) -- (82,617) -- (81,616) -- (81,610) -- (82,609) -- (81,605) -- (82,604) -- (82,602) -- (81,601) -- (81,594) -- (82,593) -- (82,587) -- (81,585) -- (82,584) -- (82,577) -- (83,576) -- (82,574) -- (82,546) -- (83,545) -- (83,542) -- (82,541) -- (83,540) -- (81,537) -- (82,535) -- (82,532) -- (81,531) -- (82,528) -- (81,527) -- (81,525) -- (83,522) -- (83,519) -- (82,518) -- (82,515) -- (81,514) -- (82,512) -- (82,505) -- (83,504) -- (83,503) -- (81,501) -- (82,499) -- (81,495) -- (83,492) -- (83,491) -- (81,489) -- (81,485) -- (82,484) -- (82,468) -- (81,467) -- (81,465) -- (82,464) -- (82,447) -- (80,445) -- (80,444) -- (82,442) -- (82,436) -- (81,435) -- (81,433) -- (82,432) -- (82,419) -- (83,418) -- (83,415) -- (82,414) -- (83,408) -- (81,403) -- (83,401) -- (83,392) -- (81,390) -- (82,388) -- (82,381) -- (81,380) -- (81,378) -- (82,377) -- (82,375) -- (81,374) -- (81,372) -- (82,371) -- (82,367) -- (81,366) -- (82,364) -- (82,357) -- (81,356) -- (81,354) -- (83,350) -- (82,349) -- (82,343) -- (81,342) -- (81,339) -- (84,335) -- (83,333) -- (83,329) -- (82,328) -- (83,316) -- (82,315) -- (82,312) -- (81,311) -- (82,309) -- (82,305) -- (83,304) -- (83,300) -- (84,299) -- (83,297) -- (83,283) -- (82,282) -- (83,280) -- (83,272) -- (82,271) -- (82,264) -- (81,263) -- (82,248) -- (80,247) -- (77,250) -- (77,252) -- (76,253) -- (76,255) -- (77,256) -- (76,258) -- (76,264) -- (77,265) -- (77,270) -- (78,271) -- (78,274) -- (76,278) -- (76,281) -- (77,282) -- (77,287) -- (76,288) -- (76,294) -- (77,295) -- (76,297) -- (76,315) -- (77,316) -- (77,318) -- (76,319) -- (76,326) -- (77,327) -- (77,329) -- (76,330) -- (76,337) -- (75,339) -- (77,343) -- (75,345) -- (75,352) -- (76,353) -- (75,363) -- (76,365) -- (74,367) -- (74,371) -- (75,372) -- (75,381) -- (74,382) -- (74,386) -- (76,388) -- (75,390) -- (75,397) -- (76,398) -- (76,400) -- (75,401) -- (75,405) -- (76,406) -- (76,415) -- (75,416) -- (75,420) -- (74,421) -- (75,422) -- (75,427) -- (78,430) -- (76,432) -- (76,434) -- (74,437) -- (74,440) -- (75,441) -- (75,445) -- (74,446) -- (74,449) -- (75,450) -- (74,452) -- (75,458) -- (74,459) -- (74,469) -- (73,470) -- (73,472) -- (75,474) -- (75,479) -- (73,483) -- (73,486) -- (75,488) -- (74,490) -- (74,493) -- (75,494) -- (75,497) -- (74,499) -- (75,500) -- (75,509) -- (76,510) -- (76,517) -- (75,518) -- (75,520) -- (77,522) -- (77,524) -- (75,527) -- (75,528) -- (77,530) -- (77,532) -- (76,533) -- (77,536) -- (75,538) -- (75,541) -- (74,542) -- (74,556) -- (75,557) -- (75,562) -- (74,563) -- (74,568) -- (75,569) -- (75,574) -- (74,575) -- (74,594) -- (75,596) -- (73,600) -- (73,602) -- (74,603) -- (73,613) -- (75,617) -- (74,619) -- (74,623) -- (73,624) -- (73,626) -- (75,628) -- (75,637) -- (74,638) -- (74,641) -- (75,642) -- (74,651) -- (75,652) -- (75,657) -- (76,658) -- (76,660) -- (75,661) -- (75,665) -- (78,667) -- (81,667) -- (82,666) -- (90,666) -- (91,667) -- (134,667) -- (135,668) -- (137,667) -- (140,667) -- (141,668) -- (150,668) -- (151,667) -- (168,668) -- (171,666) -- (173,666) -- (174,667) -- (180,667) -- (181,668) -- (185,668) -- (186,667) -- (190,667) -- (191,666) -- (195,668) -- (197,667) -- (221,667) -- (222,666) -- (226,666) -- (227,667) -- (281,667) -- (282,668) -- (302,669) -- (303,670) -- (306,670) -- (309,668) -- (318,668) -- (321,671) -- (323,669) -- (326,669) -- (328,670) -- (329,669) -- (333,669) -- (334,668) -- (336,668) -- (337,669) -- (340,669) -- (341,668) -- (353,668) -- (356,670) -- (358,670) -- (359,669) -- (381,669) -- (382,668) -- (397,668) -- (398,669) -- (426,669) -- (427,668) -- (444,668) -- (445,669) -- (447,669) -- (448,668) -- (471,668) -- (472,669) -- (477,669) -- (478,668) -- (480,668) -- (481,669) -- (490,669) -- (491,668) -- (509,668) -- (510,669) -- (513,669) -- (516,667) -- (517,664) -- (519,662) -- (524,659) -- (536,648) -- (537,648) -- (544,640) -- (547,639) -- (549,637) -- (549,636) -- (560,625) -- (563,624) -- (564,622) -- (566,620) -- (567,620) -- (568,618) -- (570,616) -- (571,616) -- (573,614) -- (573,613) -- (583,603) -- (588,600) -- (594,594) -- (597,593) -- (599,591) -- (599,588) -- (601,586) -- (604,585) -- (613,576) -- (614,576) -- (618,571) -- (619,571) -- (629,561) -- (630,561) -- (638,553) -- (639,553) -- (649,543) -- (649,542) -- (651,540) -- (654,539) -- (655,537) -- (665,528) -- (667,528) -- (668,526) -- (674,520) -- (675,520) -- (678,516) -- (682,514) -- (697,499) -- (697,498) -- (709,487) -- (710,485) -- (711,461) -- (712,460) -- (712,458) -- (711,457) -- (712,440) -- (711,438) -- (712,437) -- (712,424) -- (713,423) -- (713,421) -- (712,420) -- (712,418) -- (713,417) -- (713,404) -- (712,403) -- (712,401) -- (713,400) -- (713,392) -- (714,391) -- (713,381) -- (712,380) -- (713,378) -- (712,376) -- (712,365) -- (713,364) -- (713,361) -- (712,360) -- (713,359) -- (712,358) -- (713,357) -- (713,337) -- (714,336) -- (713,334) -- (713,331) -- (714,330) -- (714,328) -- (713,327) -- (713,323) -- (714,322) -- (714,311) -- (715,310) -- (714,308) -- (714,306) -- (715,305) -- (715,302) -- (714,301) -- (714,298) -- (715,297) -- (714,295) -- (715,293) -- (714,287) -- (715,286) -- (715,282) -- (716,281) -- (715,279) -- (716,275) -- (715,274) -- (715,271) -- (713,267) -- (715,263) -- (715,251) -- (716,250) -- (716,248) -- (715,247) -- (715,240) -- (716,239) -- (716,234) -- (715,233) -- (715,231) -- (716,230) -- (716,221) -- (717,220) -- (717,218) -- (716,217) -- (716,210) -- (717,209) -- (716,208) -- (717,207) -- (716,200) -- (717,199) -- (717,192) -- (718,191) -- (716,187) -- (717,181) -- (718,180) -- (717,178) -- (717,174) -- (718,173) -- (718,171) -- (717,170) -- (717,158) -- (718,157) -- (718,155) -- (716,152) -- (716,148) -- (717,147) -- (717,134) -- (718,133) -- (718,114) -- (719,113) -- (718,111) -- (718,79) -- (717,78) -- (718,72) -- (717,71) -- cycle
(86,238) -- (90,234) -- (93,233) -- (96,229) -- (101,227) -- (113,216) -- (118,213) -- (124,207) -- (128,205) -- (134,199) -- (138,197) -- (143,191) -- (148,189) -- (160,177) -- (165,175) -- (173,167) -- (176,166) -- (182,160) -- (186,159) -- (188,155) -- (199,144) -- (202,143) -- (209,136) -- (210,136) -- (221,126) -- (222,126) -- (237,113) -- (238,113) -- (251,100) -- (252,100) -- (257,95) -- (261,93) -- (270,85) -- (271,85) -- (275,81) -- (276,81) -- (288,69) -- (289,69) -- (291,67) -- (294,67) -- (295,66) -- (297,66) -- (298,67) -- (308,67) -- (311,65) -- (312,66) -- (334,66) -- (335,67) -- (337,67) -- (338,66) -- (341,66) -- (342,67) -- (349,67) -- (350,66) -- (352,66) -- (353,67) -- (371,67) -- (372,68) -- (385,68) -- (386,69) -- (408,69) -- (409,70) -- (412,70) -- (413,69) -- (423,69) -- (424,68) -- (427,68) -- (428,69) -- (458,69) -- (459,70) -- (462,70) -- (463,69) -- (465,69) -- (466,70) -- (471,70) -- (472,69) -- (474,70) -- (476,70) -- (477,69) -- (479,69) -- (480,70) -- (493,70) -- (495,69) -- (496,70) -- (516,70) -- (517,71) -- (520,71) -- (521,70) -- (527,70) -- (528,71) -- (530,71) -- (531,70) -- (535,70) -- (537,71) -- (538,70) -- (541,70) -- (543,69) -- (544,70) -- (570,70) -- (571,71) -- (573,70) -- (575,71) -- (611,71) -- (612,72) -- (616,72) -- (617,71) -- (621,71) -- (622,72) -- (653,72) -- (654,71) -- (669,71) -- (670,72) -- (672,71) -- (676,71) -- (677,72) -- (687,72) -- (688,71) -- (690,71) -- (691,72) -- (702,72) -- (703,71) -- (706,71) -- (709,73) -- (709,75) -- (708,76) -- (707,76) -- (703,80) -- (702,80) -- (698,84) -- (697,84) -- (678,101) -- (677,101) -- (668,109) -- (667,109) -- (641,133) -- (636,136) -- (627,145) -- (626,145) -- (617,155) -- (616,155) -- (609,162) -- (604,165) -- (585,183) -- (582,184) -- (578,188) -- (577,188) -- (571,195) -- (567,197) -- (561,203) -- (560,203) -- (553,211) -- (552,211) -- (542,221) -- (541,221) -- (536,226) -- (535,226) -- (527,234) -- (526,234) -- (517,243) -- (516,243) -- (514,245) -- (512,245) -- (511,244) -- (493,244) -- (492,243) -- (490,244) -- (485,244) -- (484,243) -- (480,243) -- (479,244) -- (455,244) -- (454,243) -- (422,243) -- (421,244) -- (418,244) -- (416,243) -- (414,244) -- (413,243) -- (404,243) -- (402,244) -- (401,243) -- (386,243) -- (385,242) -- (378,242) -- (377,243) -- (371,243) -- (370,242) -- (344,242) -- (343,243) -- (339,243) -- (338,242) -- (331,242) -- (330,241) -- (324,241) -- (323,242) -- (321,241) -- (315,241) -- (314,242) -- (310,241) -- (307,243) -- (305,243) -- (304,242) -- (292,243) -- (290,241) -- (290,237) -- (291,236) -- (291,232) -- (290,231) -- (288,231) -- (287,232) -- (287,241) -- (285,243) -- (275,243) -- (274,242) -- (250,242) -- (249,241) -- (242,241) -- (241,242) -- (239,242) -- (238,241) -- (226,241) -- (225,240) -- (205,240) -- (204,239) -- (202,239) -- (201,240) -- (194,240) -- (193,241) -- (177,241) -- (176,240) -- (148,240) -- (147,239) -- (145,240) -- (140,240) -- (139,239) -- (90,239) -- (88,240) -- cycle
(291,206) -- (289,206) -- (288,207) -- (288,210) -- (287,211) -- (287,213) -- (288,214) -- (287,216) -- (287,222) -- (288,223) -- (290,223) -- (292,219) -- (292,207) -- cycle
(289,177) -- (288,178) -- (288,185) -- (287,186) -- (287,196) -- (288,197) -- (291,197) -- (293,195) -- (293,180) -- (290,177) -- cycle
(289,149) -- (287,153) -- (288,159) -- (289,160) -- (289,161) -- (287,163) -- (287,167) -- (288,168) -- (292,168) -- (294,166) -- (294,162) -- (293,161) -- (293,158) -- (294,156) -- (293,155) -- (293,152) -- (291,149) -- cycle
(290,117) -- (289,118) -- (289,123) -- (288,124) -- (288,130) -- (287,131) -- (288,136) -- (287,138) -- (289,140) -- (292,140) -- (293,139) -- (293,129) -- (294,128) -- (293,122) -- (294,121) -- (294,119) -- (292,117) -- cycle
(291,93) -- (289,95) -- (290,99) -- (288,102) -- (288,106) -- (290,109) -- (293,109) -- (295,107) -- (294,105) -- (294,96) -- (293,95) -- (293,93) -- cycle
(294,71) -- (293,71) -- (290,74) -- (290,80) -- (291,81) -- (290,83) -- (290,88) -- (292,90) -- (293,90) -- (295,88) -- (295,85) -- (294,84) -- (295,72) -- cycle;
\node[inner sep=0pt,font=\fontsize{10}{10}\selectfont] at (282,43) {\scalebox{4.4}{$5$}};
\node[inner sep=0pt,font=\fontsize{10}{10}\selectfont] at (738,54) {\scalebox{4.6}{$6$}};
\node[inner sep=0pt,font=\fontsize{10}{10}\selectfont] at (39,246) {\scalebox{4.2}{$4$}};
\node[inner sep=0pt,font=\fontsize{10}{10}\selectfont] at (541,262) {\scalebox{4.6}{$2$}};
\node[inner sep=0pt,font=\fontsize{10}{10}\selectfont] at (306,445) {\scalebox{4.2}{$7$}};
\node[inner sep=0pt,font=\fontsize{10}{10}\selectfont] at (720,516) {\scalebox{4.2}{$1$}};
\node[inner sep=0pt,font=\fontsize{10}{10}\selectfont] at (51,677) {\scalebox{4.4}{$3$}};
\node[inner sep=0pt,font=\fontsize{10}{10}\selectfont] at (518,692) {\scalebox{4.0}{$0$}};
\end{tikzpicture}%
}
\end{center}

\[
(0,1),\qquad (0,2),\qquad (0,3).
\]

Legen wir zwei rechtwinklige Axen $AB$ und $AC$, die Axe der $x$ und der $y$, durch den Punkt $0$ oder $A$ der Projectionsebene, so ist, wenn die Complexen, die die Projectionen von $1,2,3$ ausdrücken, $A,B,C$ sind:

\begin{center}\begin{minipage}{5.2cm}\resizebox{5.2cm}{!}{%
\begin{tikzpicture}[x=1pt,y=-1pt]
\path[use as bounding box] (0,0) rectangle (878,733);
\path[fill=black,draw=none,even odd rule]
(418,703) -- (416,704) -- (416,710) -- (414,712) -- (415,716) -- (416,717) -- (418,717) -- (419,716) -- (419,709) -- (418,708) -- (419,704) -- cycle
(418,692) -- (416,693) -- (416,695) -- (415,696) -- (415,699) -- (416,700) -- (417,700) -- (420,697) -- (420,695) -- cycle
(417,650) -- (415,652) -- (415,655) -- (416,656) -- (416,666) -- (415,667) -- (416,669) -- (416,674) -- (415,676) -- (417,678) -- (419,678) -- (420,677) -- (420,671) -- (421,670) -- (421,668) -- (420,667) -- (420,658) -- (419,657) -- (420,656) -- (421,652) -- (419,650) -- cycle
(418,619) -- (417,619) -- (414,623) -- (414,624) -- (416,626) -- (414,630) -- (414,632) -- (416,634) -- (416,635) -- (414,637) -- (414,639) -- (416,641) -- (418,641) -- (419,640) -- (419,625) -- (420,623) -- cycle
(416,589) -- (414,591) -- (414,596) -- (417,598) -- (417,600) -- (414,604) -- (414,610) -- (415,611) -- (418,611) -- (419,610) -- (419,607) -- (418,606) -- (419,604) -- (419,595) -- (418,594) -- (419,590) -- (418,589) -- cycle
(416,560) -- (415,561) -- (415,566) -- (414,567) -- (414,573) -- (416,575) -- (418,575) -- (419,574) -- (419,566) -- (418,565) -- (418,561) -- cycle
(415,523) -- (414,524) -- (413,531) -- (414,532) -- (414,535) -- (416,538) -- (416,542) -- (415,544) -- (417,545) -- (418,544) -- (418,540) -- (417,539) -- (418,537) -- (418,532) -- (417,531) -- (417,529) -- (418,528) -- (418,524) -- (417,523) -- cycle
(416,495) -- (414,497) -- (414,503) -- (415,505) -- (413,507) -- (413,513) -- (415,515) -- (416,515) -- (418,513) -- (418,507) -- (417,506) -- (417,496) -- cycle
(488,466) -- (478,466) -- (477,467) -- (462,467) -- (462,469) -- (465,469) -- (466,470) -- (472,470) -- (474,469) -- (476,470) -- (480,468) -- (488,468) -- cycle
(351,467) -- (351,469) -- (352,470) -- (355,470) -- (356,471) -- (364,471) -- (365,470) -- (370,470) -- (371,469) -- (374,469) -- (375,468) -- (374,466) -- (365,466) -- (364,467) -- (360,467) -- (359,466) -- (352,466) -- cycle
(192,467) -- (192,469) -- (193,470) -- (196,470) -- (197,471) -- (202,471) -- (203,470) -- (205,471) -- (211,471) -- (212,470) -- (220,470) -- (221,469) -- (221,467) -- (220,466) -- (193,466) -- cycle
(148,467) -- (148,468) -- (152,471) -- (154,471) -- (155,470) -- (163,470) -- (164,469) -- (166,470) -- (175,470) -- (177,468) -- (176,466) -- (165,466) -- (164,467) -- (155,467) -- (154,466) -- (149,466) -- cycle
(112,467) -- (112,470) -- (113,471) -- (118,471) -- (119,470) -- (124,470) -- (125,471) -- (133,471) -- (135,469) -- (134,467) -- (125,467) -- (124,466) -- (122,467) -- cycle
(505,467) -- (506,469) -- (509,469) -- (510,470) -- (521,470) -- (522,469) -- (529,469) -- (530,468) -- (530,467) -- (528,465) -- (507,465) -- cycle
(307,467) -- (307,468) -- (310,470) -- (316,470) -- (317,471) -- (321,471) -- (322,470) -- (339,470) -- (340,469) -- (340,468) -- (338,466) -- (325,466) -- (324,465) -- (322,466) -- (318,466) -- (317,465) -- (310,465) -- cycle
(270,467) -- (270,468) -- (272,470) -- (295,470) -- (296,469) -- (299,469) -- (300,468) -- (299,466) -- (296,466) -- (295,465) -- (273,465) -- cycle
(260,467) -- (258,465) -- (250,465) -- (249,466) -- (236,466) -- (235,467) -- (233,467) -- (232,469) -- (236,471) -- (238,470) -- (258,470) -- (260,468) -- cycle
(73,468) -- (74,470) -- (78,470) -- (79,471) -- (83,471) -- (84,470) -- (86,471) -- (98,471) -- (101,469) -- (100,467) -- (98,467) -- (97,466) -- (89,466) -- (88,465) -- (79,465) -- (78,466) -- (75,466) -- cycle
(34,467) -- (34,469) -- (36,470) -- (41,470) -- (42,471) -- (56,471) -- (57,470) -- (59,470) -- (61,468) -- (59,466) -- (56,466) -- (55,465) -- (45,465) -- (44,466) -- (37,466) -- (36,467) -- cycle
(775,466) -- (775,467) -- (778,469) -- (798,469) -- (798,464) -- (795,464) -- (794,465) -- (792,464) -- (778,464) -- cycle
(736,466) -- (738,468) -- (761,468) -- (763,466) -- (763,465) -- (762,464) -- (738,464) -- cycle
(657,466) -- (657,467) -- (659,469) -- (682,469) -- (684,467) -- (683,465) -- (681,465) -- (680,464) -- (659,464) -- cycle
(569,469) -- (569,466) -- (568,465) -- (564,465) -- (563,464) -- (561,464) -- (560,465) -- (544,465) -- (543,466) -- (544,468) -- (548,470) -- (556,470) -- (557,469) -- (568,470) -- cycle
(724,465) -- (723,464) -- (709,463) -- (705,465) -- (703,464) -- (701,464) -- (700,465) -- (694,465) -- (693,466) -- (693,468) -- (694,469) -- (722,468) -- (724,466) -- cycle
(616,465) -- (616,467) -- (617,468) -- (622,468) -- (623,469) -- (644,469) -- (645,468) -- (645,465) -- (643,464) -- (628,464) -- (627,463) -- (621,463) -- (620,464) -- (617,464) -- cycle
(605,466) -- (601,463) -- (596,463) -- (595,464) -- (590,464) -- (589,465) -- (587,464) -- (580,464) -- (577,466) -- (577,468) -- (580,470) -- (588,470) -- (589,469) -- (594,470) -- (595,469) -- (603,469) -- (605,467) -- cycle
(416,437) -- (414,438) -- (414,444) -- (413,445) -- (414,447) -- (414,456) -- (415,457) -- (417,455) -- (417,452) -- (416,451) -- (416,443) -- (417,442) -- (417,438) -- cycle
(413,408) -- (411,411) -- (411,415) -- (413,417) -- (413,425) -- (414,426) -- (416,426) -- (417,425) -- (417,414) -- (416,413) -- (416,409) -- (415,408) -- cycle
(412,378) -- (411,379) -- (411,385) -- (412,386) -- (412,391) -- (411,393) -- (413,396) -- (416,396) -- (417,395) -- (416,393) -- (416,384) -- (415,383) -- (416,379) -- (415,378) -- cycle
(413,349) -- (411,351) -- (411,357) -- (413,359) -- (411,367) -- (413,369) -- (415,369) -- (418,366) -- (418,365) -- (416,363) -- (416,361) -- (417,360) -- (417,358) -- (416,357) -- (416,350) -- (415,349) -- cycle
(415,316) -- (411,318) -- (412,320) -- (411,322) -- (412,324) -- (412,329) -- (411,330) -- (411,332) -- (412,333) -- (412,337) -- (413,338) -- (415,338) -- (417,336) -- (417,331) -- (416,330) -- (416,328) -- (418,326) -- (417,319) -- cycle
(411,286) -- (410,287) -- (411,289) -- (411,291) -- (410,292) -- (410,296) -- (411,297) -- (411,301) -- (412,302) -- (412,308) -- (414,310) -- (415,310) -- (417,308) -- (417,302) -- (416,301) -- (416,298) -- (415,297) -- (416,295) -- (416,287) -- (415,286) -- cycle
(412,257) -- (411,258) -- (411,276) -- (412,277) -- (415,277) -- (416,276) -- (416,269) -- (417,268) -- (417,265) -- (416,264) -- (417,260) -- (415,257) -- cycle
(412,228) -- (411,229) -- (411,238) -- (412,240) -- (409,245) -- (412,249) -- (418,249) -- (419,247) -- (417,247) -- (416,246) -- (416,234) -- (417,233) -- (417,229) -- (416,228) -- cycle
(413,200) -- (411,202) -- (412,204) -- (412,206) -- (411,207) -- (411,216) -- (413,218) -- (413,221) -- (417,221) -- (417,206) -- (415,203) -- (416,201) -- (415,200) -- cycle
(415,164) -- (413,164) -- (412,165) -- (412,168) -- (410,171) -- (410,172) -- (412,174) -- (410,181) -- (412,183) -- (412,186) -- (413,187) -- (415,186) -- (414,184) -- (416,182) -- (416,181) -- (414,179) -- (414,177) -- (416,175) -- (416,165) -- cycle
(416,136) -- (414,136) -- (411,139) -- (411,146) -- (413,148) -- (413,150) -- (411,152) -- (411,155) -- (412,156) -- (415,156) -- (417,154) -- (416,152) -- (416,148) -- (417,147) -- (417,142) -- (416,141) -- (416,139) -- (417,137) -- cycle
(415,109) -- (414,109) -- (412,111) -- (412,119) -- (413,120) -- (413,122) -- (411,126) -- (412,128) -- (416,128) -- (417,127) -- (417,124) -- (416,123) -- (417,111) -- cycle
(772,309) -- (769,307) -- (767,307) -- (753,315) -- (748,316) -- (743,320) -- (738,321) -- (733,324) -- (730,324) -- (726,327) -- (719,328) -- (714,332) -- (712,332) -- (707,336) -- (701,338) -- (698,341) -- (694,341) -- (687,345) -- (682,346) -- (677,350) -- (673,350) -- (671,352) -- (667,353) -- (665,355) -- (662,355) -- (659,357) -- (655,358) -- (653,360) -- (650,360) -- (640,366) -- (638,366) -- (635,368) -- (632,368) -- (629,371) -- (627,371) -- (624,373) -- (618,374) -- (616,376) -- (610,379) -- (607,379) -- (605,381) -- (599,384) -- (597,384) -- (592,388) -- (590,388) -- (587,390) -- (582,391) -- (580,393) -- (576,394) -- (574,396) -- (571,396) -- (568,399) -- (565,399) -- (563,401) -- (561,401) -- (553,405) -- (551,405) -- (548,408) -- (540,410) -- (535,414) -- (529,415) -- (523,419) -- (520,419) -- (516,422) -- (514,422) -- (511,424) -- (508,424) -- (503,427) -- (501,427) -- (496,431) -- (493,431) -- (489,434) -- (482,436) -- (479,439) -- (474,442) -- (471,442) -- (467,445) -- (464,445) -- (460,448) -- (458,448) -- (449,453) -- (447,453) -- (438,458) -- (436,458) -- (434,460) -- (432,460) -- (424,464) -- (422,462) -- (422,459) -- (425,454) -- (426,449) -- (428,447) -- (428,444) -- (430,441) -- (430,437) -- (433,433) -- (433,429) -- (435,427) -- (435,424) -- (438,418) -- (437,416) -- (438,412) -- (442,405) -- (442,402) -- (443,401) -- (443,398) -- (445,395) -- (445,393) -- (447,391) -- (447,388) -- (448,387) -- (449,381) -- (451,378) -- (451,375) -- (455,370) -- (455,367) -- (457,364) -- (458,356) -- (461,352) -- (461,349) -- (464,343) -- (464,339) -- (466,337) -- (467,335) -- (467,332) -- (469,330) -- (471,319) -- (473,317) -- (473,314) -- (476,309) -- (476,306) -- (478,303) -- (478,298) -- (480,295) -- (480,292) -- (482,289) -- (483,283) -- (487,276) -- (487,272) -- (491,264) -- (491,261) -- (494,257) -- (493,253) -- (496,250) -- (496,247) -- (498,244) -- (498,241) -- (501,237) -- (501,235) -- (503,232) -- (503,228) -- (504,227) -- (504,225) -- (508,218) -- (508,215) -- (512,207) -- (512,202) -- (514,199) -- (514,197) -- (517,193) -- (517,190) -- (519,186) -- (518,182) -- (521,178) -- (521,175) -- (523,172) -- (522,168) -- (526,164) -- (526,162) -- (527,161) -- (527,157) -- (531,149) -- (532,143) -- (536,138) -- (535,136) -- (535,133) -- (538,129) -- (538,126) -- (540,123) -- (541,117) -- (544,112) -- (545,105) -- (547,102) -- (547,99) -- (546,98) -- (544,98) -- (543,99) -- (543,103) -- (539,108) -- (538,116) -- (536,118) -- (536,122) -- (533,127) -- (532,132) -- (530,135) -- (530,138) -- (531,139) -- (528,142) -- (525,148) -- (526,152) -- (523,155) -- (523,158) -- (521,160) -- (521,163) -- (522,165) -- (517,170) -- (517,176) -- (516,177) -- (516,180) -- (512,184) -- (512,189) -- (510,192) -- (510,196) -- (508,198) -- (506,206) -- (504,209) -- (504,211) -- (502,215) -- (502,218) -- (498,227) -- (498,231) -- (495,235) -- (495,238) -- (493,240) -- (493,243) -- (491,247) -- (490,253) -- (484,265) -- (484,267) -- (482,271) -- (482,275) -- (481,276) -- (481,278) -- (479,280) -- (479,284) -- (478,285) -- (477,290) -- (474,294) -- (474,298) -- (472,300) -- (473,303) -- (470,308) -- (470,312) -- (468,314) -- (468,316) -- (466,319) -- (465,324) -- (460,334) -- (462,338) -- (459,341) -- (456,347) -- (456,350) -- (453,354) -- (453,358) -- (451,361) -- (450,368) -- (448,371) -- (448,375) -- (446,378) -- (445,383) -- (443,386) -- (443,390) -- (439,398) -- (439,403) -- (437,405) -- (435,409) -- (435,412) -- (433,414) -- (433,420) -- (432,421) -- (430,429) -- (428,432) -- (427,438) -- (424,443) -- (424,446) -- (422,449) -- (422,452) -- (420,454) -- (420,458) -- (419,459) -- (419,461) -- (416,464) -- (412,464) -- (409,462) -- (407,462) -- (405,460) -- (402,460) -- (390,454) -- (387,454) -- (385,452) -- (383,452) -- (382,451) -- (378,451) -- (377,450) -- (375,450) -- (373,448) -- (371,448) -- (367,446) -- (363,446) -- (360,444) -- (355,443) -- (352,441) -- (347,441) -- (345,439) -- (343,439) -- (338,436) -- (334,436) -- (331,434) -- (328,434) -- (321,431) -- (318,431) -- (315,429) -- (313,429) -- (309,427) -- (306,427) -- (304,425) -- (301,425) -- (298,423) -- (295,423) -- (292,421) -- (289,421) -- (286,419) -- (280,418) -- (279,417) -- (265,413) -- (263,411) -- (260,411) -- (254,408) -- (247,408) -- (245,406) -- (243,406) -- (240,404) -- (235,403) -- (232,401) -- (226,400) -- (223,398) -- (216,398) -- (214,396) -- (212,396) -- (209,394) -- (206,394) -- (201,391) -- (185,387) -- (180,384) -- (178,384) -- (177,383) -- (175,383) -- (174,382) -- (172,382) -- (171,381) -- (169,381) -- (168,380) -- (154,376) -- (152,374) -- (148,374) -- (144,371) -- (140,371) -- (136,369) -- (131,369) -- (124,365) -- (121,365) -- (119,363) -- (113,363) -- (110,360) -- (104,359) -- (94,354) -- (88,353) -- (87,352) -- (82,351) -- (79,349) -- (77,349) -- (74,347) -- (68,346) -- (63,343) -- (58,343) -- (56,345) -- (56,346) -- (58,348) -- (65,350) -- (70,353) -- (72,353) -- (73,354) -- (77,354) -- (78,355) -- (86,357) -- (89,359) -- (95,360) -- (100,363) -- (103,363) -- (104,364) -- (109,364) -- (113,367) -- (121,369) -- (124,371) -- (128,371) -- (130,373) -- (137,373) -- (138,374) -- (145,375) -- (149,378) -- (152,378) -- (155,380) -- (157,380) -- (163,383) -- (171,385) -- (174,387) -- (181,388) -- (184,390) -- (188,390) -- (190,392) -- (192,392) -- (195,394) -- (200,395) -- (202,397) -- (205,397) -- (206,398) -- (212,399) -- (222,404) -- (225,404) -- (228,406) -- (231,406) -- (234,408) -- (238,408) -- (241,411) -- (245,411) -- (248,413) -- (251,412) -- (257,415) -- (260,415) -- (261,416) -- (272,419) -- (275,421) -- (278,421) -- (284,424) -- (288,424) -- (292,427) -- (295,427) -- (303,431) -- (307,431) -- (308,432) -- (313,433) -- (320,437) -- (324,437) -- (325,438) -- (330,439) -- (333,441) -- (336,441) -- (340,444) -- (344,444) -- (354,449) -- (357,449) -- (359,451) -- (371,452) -- (375,454) -- (377,456) -- (381,456) -- (385,458) -- (389,458) -- (391,460) -- (394,460) -- (398,462) -- (404,463) -- (405,464) -- (404,466) -- (387,466) -- (386,468) -- (387,469) -- (407,469) -- (410,466) -- (411,466) -- (415,470) -- (415,475) -- (414,476) -- (414,484) -- (415,485) -- (417,485) -- (418,484) -- (418,482) -- (416,478) -- (418,475) -- (418,472) -- (421,469) -- (448,469) -- (448,467) -- (446,466) -- (443,466) -- (442,467) -- (431,467) -- (429,465) -- (431,463) -- (433,463) -- (436,461) -- (441,460) -- (444,458) -- (447,458) -- (449,456) -- (451,456) -- (456,453) -- (463,451) -- (465,449) -- (467,449) -- (474,445) -- (478,444) -- (481,441) -- (486,440) -- (490,437) -- (493,437) -- (498,434) -- (500,434) -- (503,431) -- (509,430) -- (514,426) -- (516,426) -- (524,422) -- (529,421) -- (531,419) -- (533,419) -- (537,416) -- (549,411) -- (551,409) -- (557,408) -- (559,406) -- (565,403) -- (567,403) -- (569,401) -- (571,401) -- (573,399) -- (575,399) -- (581,396) -- (584,396) -- (586,394) -- (592,392) -- (595,389) -- (597,389) -- (599,387) -- (604,386) -- (607,384) -- (610,384) -- (612,382) -- (614,382) -- (616,380) -- (620,378) -- (627,376) -- (631,373) -- (637,372) -- (641,370) -- (643,368) -- (647,368) -- (648,367) -- (652,366) -- (654,364) -- (656,364) -- (659,362) -- (662,362) -- (668,358) -- (673,357) -- (678,353) -- (681,353) -- (685,351) -- (687,349) -- (692,348) -- (694,346) -- (697,346) -- (701,344) -- (703,342) -- (706,342) -- (713,337) -- (716,337) -- (718,335) -- (721,334) -- (723,332) -- (725,332) -- (727,330) -- (733,329) -- (735,327) -- (739,326) -- (740,325) -- (743,325) -- (747,322) -- (749,322) -- (751,320) -- (753,320) -- (755,318) -- (759,316) -- (765,315) -- (767,313) -- (770,312) -- (772,310) -- cycle
(416,84) -- (414,84) -- (413,85) -- (413,90) -- (411,94) -- (411,99) -- (414,101) -- (415,101) -- (417,99) -- (417,85) -- cycle;
\node[inner sep=0pt,font=\fontsize{10}{10}\selectfont] at (419,49) {\scalebox{5.6}{$C$}};
\node[inner sep=0pt,font=\fontsize{10}{10}\selectfont] at (556,77) {\scalebox{4.4}{$2$}};
\node[inner sep=0pt,font=\fontsize{10}{10}\selectfont] at (787,297) {\scalebox{4.0}{$1$}};
\node[inner sep=0pt,font=\fontsize{10}{10}\selectfont] at (34,342) {\scalebox{4.5}{$3$}};
\node[inner sep=0pt,font=\fontsize{10}{10}\selectfont] at (832,464) {\scalebox{5.6}{$B$}};
\node[inner sep=0pt,font=\fontsize{10}{10}\selectfont] at (381,499) {\scalebox{4.0}{$0$}};
\node[inner sep=0pt,font=\fontsize{10}{10}\selectfont] at (466,511) {\scalebox{5.6}{$A$}};
\end{tikzpicture}%
}
\end{minipage}\hspace{.3cm}\begin{minipage}{7.2cm}

\[
r(\cos(1,0,x)+i\cos(1,0,y)) = A
\]
\[
r(\cos(2,0,x)+i\cos(2,0,y)) = B
\]
\[
r(\cos(3,0,x)+i\cos(3,0,y)) = C.
\]
\end{minipage}\end{center}
Ferner haben wir die Bedingungsgleichungen
\[
\cos(1,0,x)^2+\cos(2,0,x)^2+\cos(3,0,x)^2 = 1
\]
\[
\cos(1,0,y)^2+\cos(2,0,y)^2+\cos(3,0,y)^2 = 1
\]
und
\[
\begin{gathered}\cos(1,0,x)\cos(1,0,y)+\cos(2,0,x)\cos(2,0,y)\\{}+\cos(3,0,x)\cos(3,0,y)\\{}=\cos(x,0,y)=\cos 90^0=0.\end{gathered}
\]
Hieraus folgt, es sei wirklich
\[
A^2+B^2+C^2=0.
\]
Ein Gleiches folgt für die übrigen Ecken.

Es kann gewünscht werden, $A$, $B$, $C$ so zu wählen, dass die reellen Theile ganze Zahlen sind. Um diess thun zu können, kommen wir auf die reellen Zahlen zurück, und erinnern, dass, wenn $a^2+b^2=c^2$ sein soll, und $a,b,c$ ganze Zahlen, $p$ und $q$ auch ganze Zahlen sind, man haben müsse
\[
a^2=(p^2-q^2)^2,\qquad b^2=4p^2q^2,\qquad c^2=(p^2+q^2)^2,
\]
also
\[
a=p^2-q^2,\qquad b=2pq,\qquad c=p^2+q^2.
\]
Z.\,B., ist $p=2$, $q=1$, so ist $a=3$, $b=4$, $c=5$ und
\[
5^2=4^2+3^2.
\]
%
\SourceLeaf{354}

Soll nun
\[
A^2 + B^2 = -C^2
\]
sein, so werden wir haben:
\[\begin{aligned}
A &= (m+m'i)^2 - (n+n'i)^2 \\
B &= 2(m+m'i)(n+n'i) \\
C &= (m+m'i)^2 i + (n+n'i)^2 i.
\end{aligned}\]

Führt man diess aus, so ist:
\[\begin{aligned}
(m+m'i)^2 &= m^2 - m'^2 + 2mm'i \\
(n+n'i)^2 &= n^2 - n'^2 + 2nn'i \\
(m+m'i)(n+n'i) &= mn - m'n' + i(m'n+n'm),
\end{aligned}\]
also:
\[\begin{aligned}
A &= m^2 - m'^2 - n^2 + n'^2 + 2i(mm' - nn') \\
B &= 2(mn - m'n') + 2i(m'n+n'm) \\
C &= -2(mm'+nn') + i(m^2+n^2-m'^2-n'^2).
\end{aligned}\]

Umgekehrt kann man auch von diesen Zahlen auf die reellen zurückkommen.

Ist hier $m = 2$, $n = 1$, $m' = 1$, $n' = 1$, so erhält man:
\[
\begin{array}{lcl}
A = & 3+2i & \\
B = & 2+6i & \qquad A^2+B^2+C^2 = \left\{\begin{array}{r} 5+12i \\ -32+24i \\ +27-36i \end{array}\right\} = 0. \\
C = & -6+3i &
\end{array}
\]

\textls{Zusatz.} Sind die complexen Zahlen für die Ecken 1, 2, 3 gefunden, so hat man auch die für
die andern 4. Es sind nemlich die für:
\begin{center}\begin{minipage}{3.3cm}\[
\begin{array}{cc}
0 & 0 \\
1 & A \\
2 & B \\
3 & C \\
6 & A+B \\
7 & A+C \\
5 & A+B+C \\
4 & B+C
\end{array}
\]

\end{minipage}\hspace{.4cm}\begin{minipage}{7.1cm}\resizebox{7.1cm}{!}{%
\begin{tikzpicture}[x=1pt,y=-1pt]
\path[use as bounding box] (0,0) rectangle (1311,1217);
\path[fill=black,draw=none,even odd rule]
(124,876) -- (122,874) -- (119,874) -- (117,876) -- (114,877) -- (109,882) -- (107,882) -- (103,886) -- (103,888) -- (104,889) -- (106,889) -- (108,887) -- (110,887) -- (120,880) -- (123,879) -- (124,878) -- cycle
(149,859) -- (146,856) -- (143,856) -- (140,859) -- (129,866) -- (127,868) -- (127,870) -- (129,872) -- (131,872) -- (136,868) -- (140,866) -- (143,863) -- (147,862) -- (149,860) -- cycle
(172,841) -- (171,840) -- (167,840) -- (165,841) -- (162,844) -- (154,849) -- (151,852) -- (151,853) -- (153,855) -- (156,855) -- (161,851) -- (164,850) -- (169,845) -- (170,845) -- (172,843) -- cycle
(603,835) -- (603,837) -- (604,838) -- (608,839) -- (612,843) -- (615,843) -- (617,845) -- (623,845) -- (624,844) -- (624,842) -- (623,841) -- (622,841) -- (620,839) -- (618,839) -- (616,837) -- (612,835) -- (610,835) -- (608,833) -- (605,833) -- cycle
(196,823) -- (195,822) -- (191,823) -- (189,825) -- (185,827) -- (181,832) -- (179,832) -- (175,836) -- (175,838) -- (179,840) -- (186,832) -- (187,832) -- (189,830) -- (191,830) -- (196,825) -- cycle
(574,820) -- (576,824) -- (579,824) -- (585,829) -- (588,829) -- (592,831) -- (596,831) -- (597,830) -- (597,829) -- (595,827) -- (594,827) -- (589,823) -- (587,823) -- (580,819) -- (575,819) -- cycle
(218,808) -- (216,806) -- (214,806) -- (211,809) -- (207,809) -- (206,811) -- (199,818) -- (199,820) -- (200,821) -- (203,821) -- (205,819) -- (207,819) -- (210,815) -- (215,813) -- (218,810) -- cycle
(542,804) -- (542,807) -- (545,809) -- (549,809) -- (551,811) -- (556,813) -- (559,816) -- (561,816) -- (562,817) -- (568,817) -- (569,816) -- (569,814) -- (568,813) -- (552,805) -- (550,803) -- (543,803) -- cycle
(243,792) -- (240,789) -- (238,789) -- (230,795) -- (228,795) -- (226,797) -- (225,797) -- (222,800) -- (222,803) -- (224,805) -- (230,803) -- (235,798) -- (240,796) -- (243,793) -- cycle
(509,789) -- (509,792) -- (510,794) -- (515,797) -- (518,796) -- (523,800) -- (531,801) -- (532,802) -- (536,801) -- (536,799) -- (535,799) -- (532,796) -- (526,793) -- (523,793) -- (520,790) -- (518,790) -- (516,788) -- (511,787) -- cycle
(267,774) -- (266,773) -- (261,773) -- (258,776) -- (256,776) -- (249,782) -- (248,782) -- (245,785) -- (245,787) -- (247,789) -- (250,789) -- (253,786) -- (255,786) -- (262,780) -- (266,779) -- (267,778) -- cycle
(478,772) -- (478,774) -- (483,779) -- (485,779) -- (491,782) -- (494,785) -- (498,785) -- (500,786) -- (503,784) -- (503,782) -- (502,781) -- (498,780) -- (496,778) -- (494,778) -- (487,773) -- (485,773) -- (481,771) -- (479,771) -- cycle
(445,758) -- (445,760) -- (449,765) -- (460,767) -- (462,769) -- (467,769) -- (469,767) -- (469,765) -- (468,764) -- (467,764) -- (462,760) -- (460,760) -- (456,757) -- (453,757) -- (452,756) -- (447,756) -- cycle
(293,756) -- (289,754) -- (276,764) -- (275,764) -- (271,768) -- (271,770) -- (272,771) -- (278,770) -- (287,762) -- (293,758) -- cycle
(413,741) -- (413,744) -- (415,746) -- (416,746) -- (418,748) -- (421,748) -- (425,750) -- (432,755) -- (434,755) -- (437,752) -- (437,750) -- (436,749) -- (432,748) -- (429,745) -- (427,745) -- (418,740) -- (414,740) -- cycle
(315,741) -- (313,739) -- (308,741) -- (306,743) -- (304,743) -- (295,751) -- (295,753) -- (296,754) -- (302,753) -- (310,746) -- (312,746) -- (315,743) -- cycle
(337,726) -- (336,725) -- (333,725) -- (327,730) -- (323,731) -- (319,735) -- (319,737) -- (320,738) -- (324,738) -- (330,732) -- (334,731) -- (337,728) -- cycle
(380,725) -- (380,727) -- (382,729) -- (387,731) -- (389,733) -- (391,733) -- (395,735) -- (398,735) -- (399,736) -- (401,736) -- (402,735) -- (402,733) -- (398,730) -- (396,730) -- (385,724) -- (381,724) -- cycle
(360,684) -- (359,684) -- (355,688) -- (355,694) -- (353,698) -- (351,700) -- (351,703) -- (349,707) -- (349,710) -- (351,713) -- (349,715) -- (342,718) -- (340,720) -- (340,723) -- (341,724) -- (344,724) -- (348,720) -- (353,718) -- (355,716) -- (358,716) -- (366,722) -- (369,722) -- (373,720) -- (372,718) -- (369,718) -- (365,716) -- (363,714) -- (358,714) -- (352,709) -- (352,708) -- (357,702) -- (357,700) -- (360,694) -- (360,691) -- (362,688) -- (362,686) -- cycle
(371,654) -- (369,654) -- (366,656) -- (366,658) -- (367,659) -- (364,662) -- (364,667) -- (359,674) -- (359,677) -- (360,678) -- (364,679) -- (367,676) -- (368,667) -- (370,666) -- (370,661) -- (372,658) -- (372,655) -- cycle
(380,627) -- (377,627) -- (374,630) -- (373,639) -- (369,644) -- (369,647) -- (370,648) -- (375,648) -- (378,643) -- (378,640) -- (381,634) -- (381,628) -- cycle
(390,595) -- (387,595) -- (385,597) -- (385,600) -- (383,603) -- (381,611) -- (379,613) -- (379,618) -- (380,619) -- (383,619) -- (385,617) -- (386,617) -- (387,616) -- (387,614) -- (390,610) -- (390,606) -- (389,605) -- (391,601) -- (391,596) -- cycle
(401,563) -- (397,563) -- (396,564) -- (395,570) -- (392,575) -- (392,578) -- (390,580) -- (390,585) -- (391,586) -- (394,586) -- (397,583) -- (397,578) -- (402,566) -- (402,564) -- cycle
(409,531) -- (407,531) -- (405,533) -- (405,538) -- (402,542) -- (402,545) -- (401,546) -- (400,552) -- (399,553) -- (399,554) -- (401,556) -- (403,556) -- (406,553) -- (406,551) -- (407,550) -- (407,548) -- (408,547) -- (408,545) -- (411,538) -- (411,533) -- cycle
(418,505) -- (415,505) -- (414,506) -- (414,511) -- (411,515) -- (411,517) -- (408,523) -- (408,524) -- (410,526) -- (413,526) -- (415,523) -- (415,520) -- (416,519) -- (417,513) -- (419,509) -- (419,506) -- cycle
(428,475) -- (425,475) -- (423,477) -- (423,481) -- (421,483) -- (420,488) -- (418,491) -- (418,493) -- (420,495) -- (422,495) -- (425,492) -- (425,489) -- (427,486) -- (428,480) -- (430,477) -- cycle
(437,447) -- (435,447) -- (432,450) -- (432,453) -- (430,456) -- (430,458) -- (428,460) -- (427,466) -- (429,468) -- (431,468) -- (434,464) -- (434,461) -- (435,460) -- (436,454) -- (437,453) -- (437,450) -- (438,449) -- cycle
(626,23) -- (623,29) -- (624,30) -- (623,31) -- (624,32) -- (624,36) -- (625,37) -- (624,39) -- (624,44) -- (625,45) -- (624,47) -- (625,49) -- (625,52) -- (624,53) -- (624,57) -- (625,58) -- (625,61) -- (624,62) -- (624,91) -- (625,92) -- (625,97) -- (624,98) -- (624,102) -- (625,103) -- (625,111) -- (624,112) -- (625,114) -- (625,116) -- (624,117) -- (625,119) -- (624,121) -- (624,128) -- (625,129) -- (625,131) -- (624,132) -- (625,134) -- (624,136) -- (624,151) -- (625,152) -- (625,158) -- (624,159) -- (624,180) -- (625,181) -- (625,183) -- (620,187) -- (620,189) -- (623,189) -- (625,191) -- (625,195) -- (623,197) -- (622,197) -- (620,195) -- (618,195) -- (613,191) -- (611,191) -- (609,189) -- (604,187) -- (597,182) -- (595,182) -- (584,176) -- (582,176) -- (580,174) -- (578,174) -- (570,170) -- (568,168) -- (560,164) -- (556,163) -- (552,159) -- (550,159) -- (545,156) -- (540,155) -- (537,153) -- (535,153) -- (531,151) -- (527,152) -- (523,155) -- (520,156) -- (512,163) -- (505,166) -- (483,184) -- (478,186) -- (473,191) -- (470,192) -- (467,195) -- (461,198) -- (458,201) -- (457,201) -- (452,207) -- (447,209) -- (440,216) -- (438,216) -- (435,219) -- (431,221) -- (427,225) -- (424,226) -- (414,235) -- (413,235) -- (411,237) -- (409,237) -- (401,244) -- (395,247) -- (383,257) -- (380,258) -- (377,261) -- (375,261) -- (355,279) -- (349,282) -- (344,287) -- (343,287) -- (339,291) -- (337,291) -- (329,298) -- (321,302) -- (305,315) -- (303,315) -- (302,314) -- (302,313) -- (304,312) -- (305,310) -- (300,312) -- (299,313) -- (301,316) -- (300,317) -- (300,319) -- (298,321) -- (291,324) -- (288,327) -- (285,328) -- (277,336) -- (272,339) -- (262,348) -- (259,349) -- (254,356) -- (254,360) -- (251,365) -- (251,371) -- (249,373) -- (247,377) -- (247,379) -- (245,382) -- (244,387) -- (241,392) -- (241,395) -- (239,397) -- (239,400) -- (237,404) -- (237,407) -- (235,410) -- (235,414) -- (233,417) -- (233,419) -- (231,423) -- (231,426) -- (228,430) -- (228,437) -- (226,439) -- (226,441) -- (223,446) -- (224,449) -- (222,452) -- (222,454) -- (220,458) -- (220,461) -- (216,470) -- (216,474) -- (214,477) -- (214,480) -- (213,481) -- (213,486) -- (211,489) -- (212,492) -- (209,495) -- (208,502) -- (210,504) -- (212,504) -- (213,503) -- (213,500) -- (212,498) -- (216,495) -- (217,491) -- (219,489) -- (219,486) -- (222,481) -- (221,478) -- (224,473) -- (225,469) -- (227,467) -- (226,465) -- (227,463) -- (227,459) -- (230,452) -- (232,450) -- (231,446) -- (232,445) -- (232,442) -- (235,436) -- (236,428) -- (240,421) -- (240,418) -- (242,415) -- (241,412) -- (244,409) -- (244,406) -- (247,399) -- (247,396) -- (250,392) -- (250,388) -- (252,385) -- (252,382) -- (254,379) -- (254,376) -- (256,373) -- (256,370) -- (259,365) -- (259,361) -- (261,358) -- (265,358) -- (267,359) -- (275,367) -- (278,366) -- (281,369) -- (286,369) -- (292,372) -- (297,373) -- (300,376) -- (306,379) -- (308,381) -- (315,383) -- (317,385) -- (325,388) -- (334,394) -- (337,394) -- (339,396) -- (343,397) -- (349,402) -- (353,402) -- (356,405) -- (362,408) -- (367,409) -- (370,411) -- (373,411) -- (377,414) -- (379,414) -- (383,417) -- (385,417) -- (389,419) -- (391,421) -- (397,423) -- (399,425) -- (404,426) -- (408,429) -- (420,434) -- (422,436) -- (425,436) -- (427,438) -- (431,440) -- (436,441) -- (441,445) -- (443,445) -- (451,449) -- (453,451) -- (457,451) -- (461,455) -- (465,455) -- (470,459) -- (473,459) -- (475,461) -- (478,462) -- (481,465) -- (483,465) -- (490,469) -- (493,469) -- (496,472) -- (504,476) -- (506,476) -- (508,478) -- (514,481) -- (517,481) -- (521,483) -- (524,486) -- (526,486) -- (528,488) -- (530,488) -- (540,494) -- (545,495) -- (552,499) -- (554,499) -- (558,501) -- (561,504) -- (569,508) -- (573,508) -- (578,512) -- (580,512) -- (584,514) -- (587,517) -- (595,519) -- (600,523) -- (610,528) -- (615,529) -- (619,532) -- (622,532) -- (626,535) -- (625,539) -- (627,542) -- (627,545) -- (625,548) -- (625,551) -- (626,552) -- (626,556) -- (627,558) -- (625,560) -- (625,564) -- (626,565) -- (626,574) -- (625,575) -- (626,589) -- (625,590) -- (625,592) -- (627,594) -- (627,595) -- (626,596) -- (626,599) -- (625,600) -- (625,604) -- (624,606) -- (625,607) -- (625,610) -- (627,612) -- (627,613) -- (625,615) -- (625,617) -- (626,618) -- (626,625) -- (625,627) -- (627,630) -- (626,632) -- (626,642) -- (627,643) -- (627,650) -- (625,653) -- (625,654) -- (627,656) -- (627,660) -- (625,662) -- (625,665) -- (626,666) -- (626,688) -- (625,689) -- (626,691) -- (626,698) -- (625,699) -- (626,703) -- (625,705) -- (626,706) -- (626,714) -- (627,715) -- (627,729) -- (625,732) -- (625,734) -- (627,737) -- (626,739) -- (626,742) -- (627,743) -- (626,767) -- (627,769) -- (626,770) -- (626,793) -- (625,795) -- (627,798) -- (627,801) -- (628,802) -- (627,804) -- (627,829) -- (628,830) -- (627,832) -- (627,839) -- (628,840) -- (628,847) -- (625,850) -- (625,851) -- (627,853) -- (626,855) -- (626,863) -- (625,864) -- (626,866) -- (626,868) -- (625,869) -- (626,871) -- (625,877) -- (627,881) -- (626,883) -- (627,896) -- (625,900) -- (626,901) -- (626,906) -- (627,907) -- (626,909) -- (626,921) -- (627,922) -- (627,924) -- (626,925) -- (626,929) -- (628,931) -- (627,937) -- (629,939) -- (629,941) -- (626,943) -- (626,953) -- (625,954) -- (625,972) -- (628,975) -- (628,976) -- (626,978) -- (627,980) -- (627,985) -- (626,986) -- (626,988) -- (628,991) -- (628,993) -- (627,994) -- (627,1002) -- (626,1003) -- (627,1005) -- (626,1007) -- (627,1011) -- (626,1012) -- (626,1024) -- (625,1025) -- (625,1027) -- (626,1028) -- (626,1032) -- (625,1034) -- (627,1036) -- (627,1039) -- (626,1040) -- (626,1062) -- (627,1063) -- (627,1069) -- (628,1070) -- (627,1072) -- (627,1077) -- (628,1078) -- (628,1080) -- (627,1081) -- (627,1088) -- (628,1090) -- (627,1091) -- (627,1100) -- (626,1101) -- (626,1103) -- (627,1104) -- (627,1115) -- (626,1116) -- (626,1123) -- (627,1124) -- (627,1127) -- (626,1129) -- (627,1130) -- (627,1150) -- (629,1152) -- (629,1155) -- (627,1157) -- (624,1157) -- (620,1155) -- (615,1151) -- (612,1151) -- (605,1147) -- (603,1147) -- (601,1145) -- (591,1140) -- (589,1140) -- (587,1138) -- (584,1137) -- (580,1134) -- (577,1134) -- (573,1131) -- (566,1129) -- (563,1126) -- (559,1125) -- (557,1123) -- (555,1123) -- (553,1121) -- (551,1121) -- (550,1120) -- (546,1120) -- (541,1116) -- (536,1115) -- (533,1112) -- (526,1110) -- (524,1108) -- (514,1103) -- (509,1102) -- (507,1100) -- (505,1100) -- (503,1098) -- (501,1098) -- (496,1095) -- (494,1095) -- (490,1092) -- (486,1090) -- (484,1090) -- (477,1085) -- (474,1085) -- (472,1084) -- (469,1081) -- (467,1081) -- (464,1079) -- (457,1078) -- (455,1076) -- (443,1070) -- (441,1070) -- (437,1066) -- (434,1066) -- (432,1064) -- (429,1064) -- (415,1057) -- (413,1057) -- (411,1055) -- (404,1053) -- (401,1050) -- (397,1049) -- (395,1047) -- (392,1047) -- (390,1045) -- (385,1043) -- (383,1041) -- (380,1041) -- (374,1037) -- (370,1037) -- (366,1035) -- (364,1033) -- (362,1033) -- (358,1031) -- (354,1028) -- (349,1027) -- (344,1023) -- (341,1023) -- (339,1021) -- (333,1018) -- (331,1018) -- (329,1016) -- (325,1014) -- (323,1014) -- (319,1012) -- (315,1009) -- (311,1009) -- (306,1005) -- (302,1004) -- (298,1001) -- (296,1001) -- (289,996) -- (287,996) -- (286,995) -- (282,995) -- (274,991) -- (271,988) -- (269,988) -- (267,986) -- (264,986) -- (261,984) -- (257,983) -- (250,978) -- (245,977) -- (240,974) -- (238,974) -- (236,972) -- (230,971) -- (227,968) -- (223,967) -- (220,964) -- (214,963) -- (212,961) -- (210,961) -- (208,959) -- (206,959) -- (202,957) -- (199,954) -- (193,954) -- (191,952) -- (181,947) -- (174,945) -- (172,941) -- (167,942) -- (165,940) -- (167,938) -- (167,937) -- (163,936) -- (157,932) -- (155,932) -- (151,934) -- (146,930) -- (142,928) -- (137,927) -- (133,924) -- (127,921) -- (125,921) -- (118,917) -- (116,917) -- (110,913) -- (104,912) -- (103,911) -- (100,911) -- (96,908) -- (92,906) -- (90,906) -- (87,903) -- (89,901) -- (91,901) -- (99,895) -- (98,893) -- (94,893) -- (87,899) -- (85,897) -- (87,895) -- (88,887) -- (92,879) -- (91,875) -- (94,872) -- (94,867) -- (95,866) -- (95,864) -- (97,862) -- (97,859) -- (99,856) -- (99,851) -- (101,849) -- (101,846) -- (103,844) -- (103,841) -- (104,840) -- (105,834) -- (107,831) -- (108,824) -- (111,819) -- (111,817) -- (113,815) -- (114,810) -- (116,808) -- (116,806) -- (118,802) -- (117,798) -- (121,793) -- (121,790) -- (122,789) -- (122,786) -- (125,780) -- (125,777) -- (128,772) -- (128,770) -- (129,769) -- (129,767) -- (132,760) -- (133,754) -- (137,746) -- (137,743) -- (138,742) -- (139,736) -- (140,735) -- (139,733) -- (142,730) -- (142,726) -- (145,723) -- (145,718) -- (147,714) -- (148,707) -- (150,704) -- (150,701) -- (153,696) -- (153,691) -- (154,690) -- (154,688) -- (156,686) -- (156,682) -- (158,680) -- (158,677) -- (160,675) -- (161,669) -- (163,666) -- (163,662) -- (166,658) -- (166,655) -- (170,647) -- (170,644) -- (172,640) -- (172,634) -- (174,632) -- (174,628) -- (176,626) -- (176,623) -- (178,619) -- (175,613) -- (176,612) -- (178,614) -- (180,613) -- (180,611) -- (181,610) -- (181,608) -- (184,601) -- (184,598) -- (185,597) -- (185,592) -- (188,588) -- (188,586) -- (191,579) -- (192,572) -- (194,569) -- (195,563) -- (198,557) -- (198,554) -- (200,551) -- (200,544) -- (204,537) -- (203,533) -- (205,531) -- (206,525) -- (208,523) -- (208,519) -- (211,512) -- (211,508) -- (208,506) -- (206,506) -- (205,510) -- (203,513) -- (203,518) -- (201,520) -- (201,523) -- (198,526) -- (199,529) -- (197,531) -- (197,535) -- (196,536) -- (197,540) -- (195,542) -- (192,548) -- (192,552) -- (193,553) -- (189,559) -- (189,563) -- (187,566) -- (185,575) -- (183,577) -- (183,581) -- (180,585) -- (180,592) -- (178,595) -- (178,597) -- (176,601) -- (176,604) -- (174,606) -- (174,611) -- (169,617) -- (170,619) -- (168,623) -- (169,627) -- (167,631) -- (167,634) -- (165,637) -- (164,644) -- (162,646) -- (160,650) -- (160,655) -- (158,657) -- (157,662) -- (155,664) -- (155,666) -- (154,667) -- (154,670) -- (155,672) -- (152,675) -- (152,679) -- (150,681) -- (149,686) -- (146,692) -- (146,696) -- (142,704) -- (141,715) -- (138,719) -- (137,727) -- (135,729) -- (135,731) -- (133,734) -- (133,737) -- (131,739) -- (129,747) -- (127,749) -- (127,751) -- (126,752) -- (126,756) -- (124,758) -- (124,763) -- (122,766) -- (122,768) -- (120,770) -- (119,776) -- (116,779) -- (115,789) -- (113,791) -- (113,794) -- (111,798) -- (111,802) -- (108,806) -- (108,810) -- (106,814) -- (107,817) -- (105,820) -- (105,822) -- (103,824) -- (103,828) -- (96,834) -- (98,841) -- (95,845) -- (95,848) -- (94,849) -- (93,854) -- (91,856) -- (89,867) -- (87,869) -- (86,874) -- (83,878) -- (83,883) -- (81,886) -- (80,897) -- (77,900) -- (77,905) -- (84,910) -- (88,910) -- (92,913) -- (97,914) -- (100,917) -- (108,919) -- (112,921) -- (114,923) -- (118,923) -- (122,926) -- (124,926) -- (129,929) -- (132,929) -- (135,933) -- (137,933) -- (140,936) -- (144,937) -- (152,942) -- (158,943) -- (161,946) -- (163,946) -- (171,950) -- (173,950) -- (178,954) -- (184,957) -- (186,957) -- (190,959) -- (192,961) -- (198,961) -- (203,965) -- (210,967) -- (212,969) -- (216,970) -- (218,972) -- (220,972) -- (222,974) -- (224,974) -- (230,977) -- (232,977) -- (237,981) -- (240,981) -- (241,982) -- (245,983) -- (250,987) -- (254,988) -- (257,991) -- (264,993) -- (271,997) -- (276,997) -- (281,1002) -- (287,1002) -- (295,1008) -- (298,1008) -- (302,1011) -- (306,1012) -- (309,1015) -- (314,1016) -- (318,1019) -- (321,1019) -- (324,1021) -- (331,1023) -- (339,1029) -- (341,1029) -- (344,1031) -- (350,1032) -- (354,1036) -- (356,1036) -- (360,1038) -- (363,1038) -- (367,1041) -- (369,1041) -- (374,1044) -- (379,1045) -- (381,1047) -- (384,1047) -- (388,1050) -- (394,1053) -- (396,1053) -- (404,1057) -- (406,1059) -- (408,1059) -- (410,1061) -- (412,1061) -- (414,1063) -- (417,1063) -- (422,1066) -- (426,1067) -- (428,1069) -- (435,1073) -- (440,1074) -- (448,1079) -- (451,1079) -- (459,1084) -- (464,1084) -- (467,1087) -- (470,1087) -- (474,1090) -- (477,1090) -- (481,1093) -- (483,1093) -- (487,1096) -- (492,1098) -- (494,1100) -- (496,1100) -- (498,1102) -- (501,1102) -- (505,1105) -- (509,1106) -- (515,1109) -- (518,1112) -- (520,1113) -- (523,1113) -- (526,1116) -- (529,1116) -- (533,1118) -- (535,1120) -- (539,1121) -- (544,1125) -- (550,1126) -- (551,1127) -- (554,1127) -- (559,1131) -- (564,1132) -- (569,1136) -- (571,1136) -- (573,1138) -- (575,1138) -- (577,1140) -- (581,1141) -- (583,1143) -- (585,1143) -- (588,1145) -- (595,1147) -- (597,1149) -- (599,1149) -- (602,1151) -- (607,1152) -- (612,1155) -- (618,1156) -- (619,1157) -- (619,1159) -- (618,1160) -- (608,1160) -- (607,1159) -- (604,1159) -- (603,1160) -- (601,1159) -- (576,1159) -- (575,1158) -- (549,1158) -- (548,1159) -- (546,1158) -- (544,1158) -- (543,1159) -- (527,1159) -- (526,1158) -- (524,1159) -- (517,1159) -- (516,1158) -- (514,1158) -- (513,1159) -- (510,1159) -- (509,1158) -- (507,1159) -- (499,1159) -- (498,1158) -- (461,1158) -- (460,1159) -- (451,1159) -- (450,1158) -- (444,1158) -- (443,1159) -- (436,1159) -- (435,1158) -- (421,1158) -- (420,1159) -- (418,1158) -- (402,1158) -- (401,1159) -- (399,1159) -- (398,1158) -- (382,1158) -- (381,1157) -- (379,1158) -- (373,1158) -- (372,1159) -- (349,1159) -- (348,1158) -- (343,1158) -- (342,1159) -- (339,1159) -- (338,1158) -- (326,1158) -- (325,1159) -- (323,1159) -- (322,1158) -- (317,1158) -- (316,1159) -- (310,1158) -- (309,1159) -- (307,1159) -- (306,1158) -- (290,1158) -- (289,1157) -- (287,1157) -- (284,1159) -- (279,1159) -- (278,1158) -- (276,1159) -- (263,1159) -- (262,1158) -- (260,1159) -- (256,1159) -- (255,1158) -- (245,1158) -- (244,1159) -- (239,1159) -- (238,1158) -- (226,1158) -- (225,1159) -- (220,1159) -- (219,1158) -- (197,1158) -- (196,1159) -- (186,1159) -- (185,1158) -- (177,1158) -- (176,1159) -- (174,1159) -- (173,1158) -- (170,1158) -- (169,1159) -- (152,1159) -- (150,1158) -- (149,1159) -- (134,1159) -- (132,1158) -- (131,1159) -- (129,1159) -- (128,1158) -- (120,1158) -- (119,1159) -- (107,1159) -- (106,1160) -- (97,1160) -- (96,1161) -- (94,1160) -- (86,1160) -- (85,1159) -- (68,1159) -- (67,1160) -- (65,1159) -- (62,1159) -- (61,1160) -- (52,1160) -- (51,1159) -- (49,1159) -- (45,1161) -- (43,1161) -- (42,1160) -- (40,1160) -- (39,1161) -- (35,1161) -- (34,1160) -- (32,1160) -- (31,1161) -- (24,1161) -- (21,1163) -- (21,1165) -- (22,1166) -- (28,1166) -- (29,1167) -- (31,1166) -- (43,1166) -- (44,1167) -- (68,1167) -- (69,1166) -- (71,1167) -- (73,1167) -- (74,1166) -- (95,1166) -- (96,1167) -- (99,1167) -- (100,1166) -- (105,1166) -- (106,1167) -- (109,1165) -- (112,1165) -- (113,1166) -- (119,1166) -- (120,1165) -- (126,1165) -- (127,1166) -- (133,1166) -- (134,1165) -- (143,1165) -- (144,1166) -- (149,1166) -- (150,1165) -- (167,1165) -- (168,1164) -- (170,1164) -- (171,1165) -- (174,1165) -- (175,1164) -- (181,1164) -- (182,1165) -- (191,1165) -- (192,1164) -- (196,1164) -- (197,1163) -- (201,1163) -- (202,1164) -- (214,1164) -- (215,1165) -- (222,1165) -- (223,1164) -- (229,1164) -- (230,1165) -- (232,1165) -- (236,1163) -- (249,1163) -- (250,1164) -- (268,1164) -- (269,1165) -- (274,1165) -- (275,1164) -- (277,1164) -- (278,1165) -- (284,1165) -- (285,1164) -- (293,1164) -- (294,1163) -- (296,1164) -- (306,1164) -- (307,1165) -- (320,1165) -- (321,1164) -- (323,1164) -- (324,1165) -- (328,1165) -- (329,1164) -- (331,1165) -- (368,1165) -- (369,1164) -- (372,1164) -- (373,1163) -- (375,1163) -- (376,1164) -- (389,1164) -- (390,1165) -- (394,1165) -- (395,1164) -- (403,1164) -- (404,1165) -- (406,1164) -- (414,1164) -- (415,1165) -- (419,1165) -- (420,1164) -- (425,1164) -- (426,1165) -- (428,1164) -- (430,1165) -- (435,1165) -- (436,1164) -- (438,1164) -- (439,1165) -- (450,1165) -- (451,1164) -- (460,1164) -- (461,1163) -- (463,1163) -- (464,1164) -- (467,1164) -- (468,1163) -- (470,1163) -- (471,1164) -- (484,1164) -- (485,1163) -- (487,1164) -- (519,1164) -- (520,1165) -- (529,1165) -- (530,1164) -- (540,1164) -- (541,1165) -- (548,1165) -- (549,1164) -- (552,1164) -- (553,1163) -- (556,1163) -- (557,1164) -- (559,1163) -- (561,1164) -- (563,1164) -- (564,1163) -- (579,1163) -- (580,1164) -- (591,1164) -- (592,1165) -- (597,1165) -- (598,1164) -- (600,1164) -- (601,1165) -- (612,1164) -- (613,1165) -- (614,1164) -- (618,1164) -- (619,1163) -- (621,1164) -- (623,1163) -- (626,1163) -- (627,1164) -- (637,1164) -- (638,1163) -- (642,1163) -- (643,1164) -- (645,1164) -- (646,1163) -- (650,1163) -- (651,1164) -- (681,1164) -- (682,1165) -- (684,1164) -- (699,1164) -- (700,1165) -- (704,1164) -- (705,1165) -- (717,1165) -- (718,1164) -- (730,1164) -- (731,1163) -- (733,1164) -- (761,1164) -- (762,1163) -- (764,1163) -- (765,1164) -- (781,1164) -- (782,1165) -- (813,1165) -- (814,1164) -- (824,1164) -- (825,1165) -- (828,1165) -- (829,1164) -- (830,1165) -- (831,1164) -- (834,1164) -- (835,1165) -- (839,1165) -- (840,1166) -- (842,1166) -- (843,1165) -- (848,1165) -- (849,1164) -- (853,1164) -- (854,1163) -- (856,1164) -- (860,1163) -- (861,1164) -- (886,1164) -- (887,1165) -- (889,1164) -- (897,1164) -- (898,1165) -- (901,1165) -- (902,1164) -- (910,1164) -- (912,1162) -- (916,1162) -- (917,1163) -- (926,1163) -- (927,1164) -- (938,1164) -- (939,1165) -- (941,1164) -- (955,1164) -- (957,1166) -- (964,1166) -- (965,1165) -- (972,1165) -- (973,1164) -- (983,1164) -- (984,1165) -- (986,1165) -- (990,1163) -- (995,1163) -- (996,1164) -- (1013,1164) -- (1014,1163) -- (1018,1163) -- (1021,1165) -- (1031,1165) -- (1032,1164) -- (1056,1164) -- (1057,1165) -- (1059,1165) -- (1060,1164) -- (1068,1164) -- (1069,1165) -- (1071,1165) -- (1073,1164) -- (1074,1165) -- (1076,1165) -- (1079,1163) -- (1081,1163) -- (1082,1164) -- (1084,1164) -- (1085,1163) -- (1092,1163) -- (1093,1164) -- (1098,1164) -- (1099,1165) -- (1111,1165) -- (1112,1164) -- (1123,1165) -- (1124,1164) -- (1126,1164) -- (1127,1165) -- (1138,1165) -- (1139,1166) -- (1166,1166) -- (1167,1165) -- (1186,1165) -- (1187,1164) -- (1189,1164) -- (1190,1165) -- (1211,1164) -- (1213,1166) -- (1215,1166) -- (1216,1165) -- (1221,1165) -- (1222,1166) -- (1229,1166) -- (1230,1165) -- (1232,1165) -- (1233,1166) -- (1240,1166) -- (1241,1165) -- (1245,1165) -- (1246,1166) -- (1249,1166) -- (1250,1165) -- (1254,1165) -- (1255,1164) -- (1269,1164) -- (1270,1165) -- (1273,1165) -- (1276,1163) -- (1278,1164) -- (1285,1164) -- (1286,1165) -- (1292,1165) -- (1293,1164) -- (1295,1164) -- (1296,1163) -- (1296,1161) -- (1292,1159) -- (1282,1159) -- (1281,1160) -- (1275,1160) -- (1274,1159) -- (1272,1160) -- (1270,1159) -- (1258,1159) -- (1257,1160) -- (1233,1159) -- (1232,1160) -- (1229,1160) -- (1228,1159) -- (1185,1159) -- (1184,1158) -- (1176,1158) -- (1175,1159) -- (1172,1159) -- (1171,1158) -- (1167,1158) -- (1166,1159) -- (1154,1159) -- (1153,1158) -- (1150,1158) -- (1149,1159) -- (1140,1159) -- (1139,1158) -- (1125,1158) -- (1124,1159) -- (1122,1158) -- (1077,1158) -- (1075,1160) -- (1072,1160) -- (1071,1159) -- (1057,1159) -- (1056,1158) -- (1029,1158) -- (1028,1157) -- (1025,1157) -- (1024,1158) -- (1004,1158) -- (1003,1159) -- (998,1159) -- (996,1157) -- (994,1157) -- (992,1159) -- (988,1159) -- (987,1158) -- (976,1158) -- (975,1157) -- (973,1158) -- (925,1158) -- (924,1157) -- (920,1157) -- (919,1158) -- (914,1158) -- (913,1157) -- (902,1157) -- (901,1158) -- (866,1158) -- (865,1157) -- (862,1157) -- (861,1158) -- (846,1158) -- (845,1157) -- (832,1157) -- (831,1158) -- (823,1158) -- (822,1157) -- (819,1157) -- (817,1158) -- (815,1157) -- (812,1159) -- (804,1159) -- (803,1158) -- (800,1158) -- (799,1159) -- (793,1159) -- (792,1158) -- (772,1158) -- (771,1159) -- (768,1159) -- (767,1158) -- (759,1158) -- (758,1157) -- (743,1157) -- (742,1158) -- (735,1157) -- (733,1159) -- (726,1159) -- (725,1160) -- (723,1160) -- (722,1159) -- (702,1159) -- (700,1160) -- (699,1159) -- (694,1159) -- (693,1158) -- (687,1158) -- (686,1159) -- (683,1159) -- (682,1158) -- (679,1158) -- (678,1159) -- (667,1159) -- (666,1158) -- (664,1158) -- (663,1159) -- (646,1159) -- (644,1157) -- (645,1155) -- (647,1155) -- (651,1152) -- (654,1151) -- (659,1146) -- (662,1145) -- (667,1141) -- (670,1140) -- (673,1137) -- (677,1136) -- (679,1134) -- (680,1134) -- (692,1124) -- (696,1123) -- (705,1115) -- (708,1115) -- (716,1108) -- (720,1107) -- (724,1103) -- (732,1098) -- (736,1094) -- (742,1091) -- (745,1088) -- (755,1082) -- (758,1079) -- (760,1079) -- (762,1077) -- (763,1077) -- (769,1071) -- (770,1071) -- (772,1069) -- (776,1068) -- (781,1063) -- (783,1063) -- (789,1057) -- (793,1055) -- (796,1052) -- (799,1051) -- (802,1048) -- (805,1047) -- (816,1037) -- (824,1033) -- (827,1030) -- (830,1029) -- (835,1024) -- (839,1023) -- (843,1019) -- (844,1019) -- (847,1016) -- (850,1015) -- (854,1011) -- (855,1011) -- (860,1007) -- (863,1006) -- (875,997) -- (877,997) -- (892,984) -- (899,981) -- (901,979) -- (900,977) -- (901,971) -- (905,966) -- (905,963) -- (907,961) -- (907,959) -- (910,954) -- (910,952) -- (915,946) -- (915,944) -- (914,943) -- (914,937) -- (916,934) -- (916,930) -- (918,928) -- (918,925) -- (921,919) -- (921,916) -- (923,914) -- (923,912) -- (924,911) -- (924,907) -- (926,905) -- (927,899) -- (929,896) -- (929,892) -- (931,890) -- (931,887) -- (933,885) -- (932,881) -- (935,876) -- (935,872) -- (937,869) -- (938,863) -- (942,858) -- (942,852) -- (944,850) -- (944,846) -- (945,845) -- (945,843) -- (947,841) -- (947,838) -- (949,836) -- (949,832) -- (950,831) -- (950,829) -- (954,824) -- (953,820) -- (955,817) -- (956,811) -- (960,805) -- (960,801) -- (962,797) -- (962,792) -- (965,786) -- (965,784) -- (968,779) -- (969,771) -- (972,767) -- (972,764) -- (975,758) -- (976,753) -- (979,749) -- (980,741) -- (983,737) -- (983,732) -- (986,728) -- (986,724) -- (988,721) -- (988,719) -- (992,711) -- (992,704) -- (994,702) -- (994,698) -- (995,697) -- (995,695) -- (997,692) -- (998,687) -- (1000,685) -- (1001,677) -- (1003,675) -- (1006,668) -- (1006,663) -- (1009,657) -- (1010,651) -- (1012,648) -- (1012,644) -- (1015,640) -- (1015,638) -- (1016,637) -- (1016,633) -- (1017,632) -- (1018,627) -- (1020,624) -- (1020,617) -- (1021,616) -- (1022,611) -- (1026,604) -- (1026,601) -- (1028,597) -- (1028,592) -- (1031,588) -- (1031,585) -- (1034,579) -- (1034,577) -- (1036,573) -- (1036,570) -- (1039,564) -- (1039,560) -- (1043,552) -- (1043,549) -- (1045,546) -- (1046,540) -- (1048,537) -- (1048,530) -- (1051,524) -- (1052,516) -- (1056,509) -- (1057,503) -- (1059,500) -- (1059,496) -- (1062,493) -- (1062,489) -- (1064,486) -- (1064,480) -- (1071,466) -- (1072,459) -- (1075,453) -- (1075,450) -- (1076,449) -- (1076,446) -- (1077,445) -- (1077,442) -- (1076,440) -- (1080,432) -- (1080,428) -- (1076,425) -- (1073,425) -- (1067,422) -- (1062,418) -- (1060,418) -- (1056,416) -- (1053,413) -- (1045,410) -- (1042,407) -- (1039,406) -- (1037,404) -- (1034,404) -- (1032,402) -- (1023,398) -- (1016,393) -- (1012,392) -- (1010,390) -- (1003,387) -- (1001,385) -- (997,383) -- (995,383) -- (987,379) -- (985,377) -- (982,377) -- (973,372) -- (971,372) -- (966,368) -- (962,367) -- (955,363) -- (953,363) -- (948,359) -- (944,359) -- (941,356) -- (935,353) -- (933,353) -- (931,351) -- (927,349) -- (925,349) -- (923,347) -- (921,347) -- (919,345) -- (915,343) -- (913,343) -- (909,340) -- (903,337) -- (901,337) -- (899,335) -- (897,335) -- (889,331) -- (885,328) -- (883,328) -- (879,326) -- (877,324) -- (873,323) -- (871,321) -- (868,320) -- (865,317) -- (863,317) -- (861,315) -- (855,312) -- (853,312) -- (851,310) -- (849,310) -- (844,306) -- (831,300) -- (829,298) -- (827,298) -- (825,296) -- (822,295) -- (815,290) -- (813,290) -- (810,288) -- (804,287) -- (802,285) -- (800,285) -- (793,280) -- (791,280) -- (783,276) -- (781,274) -- (777,273) -- (775,271) -- (773,271) -- (765,267) -- (763,265) -- (759,263) -- (757,263) -- (753,261) -- (751,259) -- (747,258) -- (744,255) -- (742,255) -- (739,252) -- (733,251) -- (728,246) -- (724,245) -- (721,242) -- (719,243) -- (715,243) -- (712,240) -- (707,239) -- (704,236) -- (702,236) -- (697,233) -- (695,233) -- (690,229) -- (686,228) -- (682,225) -- (678,224) -- (676,222) -- (671,220) -- (669,218) -- (665,217) -- (663,215) -- (661,215) -- (658,212) -- (656,211) -- (652,211) -- (647,207) -- (645,207) -- (640,204) -- (638,204) -- (632,201) -- (630,199) -- (630,196) -- (631,195) -- (631,192) -- (632,191) -- (632,185) -- (631,184) -- (631,181) -- (630,180) -- (630,165) -- (631,164) -- (631,157) -- (632,155) -- (630,153) -- (630,151) -- (631,150) -- (631,147) -- (632,146) -- (632,143) -- (633,142) -- (632,140) -- (632,134) -- (631,133) -- (632,132) -- (631,131) -- (631,122) -- (633,119) -- (632,117) -- (632,104) -- (631,102) -- (633,98) -- (632,97) -- (632,82) -- (631,81) -- (632,75) -- (631,74) -- (631,60) -- (630,59) -- (631,54) -- (630,53) -- (630,51) -- (632,48) -- (631,46) -- (631,38) -- (633,36) -- (633,30) -- (632,29) -- (632,26) -- (630,23) -- cycle
(633,537) -- (637,538) -- (639,540) -- (644,542) -- (647,545) -- (650,546) -- (652,548) -- (654,549) -- (657,549) -- (659,551) -- (667,555) -- (669,557) -- (673,558) -- (675,560) -- (680,561) -- (683,564) -- (690,566) -- (695,570) -- (699,572) -- (702,572) -- (706,575) -- (710,576) -- (720,582) -- (725,582) -- (728,585) -- (740,591) -- (742,591) -- (749,595) -- (751,595) -- (753,597) -- (755,597) -- (757,599) -- (768,604) -- (770,606) -- (773,606) -- (778,610) -- (782,611) -- (786,614) -- (788,614) -- (793,617) -- (798,618) -- (801,621) -- (800,625) -- (797,629) -- (797,631) -- (796,632) -- (796,636) -- (794,639) -- (794,642) -- (792,645) -- (792,649) -- (790,652) -- (790,654) -- (787,661) -- (787,664) -- (784,670) -- (784,673) -- (783,674) -- (782,679) -- (779,683) -- (779,688) -- (778,689) -- (778,692) -- (777,693) -- (777,695) -- (776,696) -- (776,698) -- (773,705) -- (773,711) -- (771,713) -- (771,716) -- (768,721) -- (768,724) -- (765,730) -- (765,734) -- (763,737) -- (763,741) -- (759,746) -- (759,752) -- (757,755) -- (756,762) -- (755,763) -- (754,768) -- (752,771) -- (751,777) -- (748,783) -- (748,787) -- (746,789) -- (746,791) -- (745,792) -- (745,798) -- (743,801) -- (743,805) -- (741,808) -- (741,812) -- (739,814) -- (739,818) -- (738,819) -- (738,821) -- (736,824) -- (736,826) -- (734,828) -- (734,832) -- (732,836) -- (732,839) -- (729,844) -- (729,847) -- (724,856) -- (724,863) -- (721,866) -- (721,873) -- (718,880) -- (718,883) -- (715,887) -- (704,882) -- (697,877) -- (692,877) -- (691,878) -- (691,881) -- (692,882) -- (696,882) -- (699,885) -- (703,887) -- (705,887) -- (709,890) -- (713,890) -- (715,892) -- (715,895) -- (712,899) -- (712,902) -- (710,906) -- (710,909) -- (707,913) -- (707,920) -- (704,925) -- (703,932) -- (701,934) -- (701,937) -- (699,940) -- (699,944) -- (696,949) -- (696,951) -- (694,955) -- (694,958) -- (689,970) -- (689,977) -- (687,981) -- (687,984) -- (685,987) -- (685,991) -- (682,996) -- (682,998) -- (681,999) -- (681,1003) -- (679,1005) -- (677,1013) -- (675,1016) -- (675,1020) -- (672,1027) -- (672,1030) -- (668,1036) -- (668,1041) -- (667,1042) -- (667,1045) -- (664,1052) -- (663,1060) -- (661,1063) -- (661,1068) -- (657,1076) -- (657,1081) -- (656,1082) -- (656,1084) -- (653,1087) -- (653,1091) -- (651,1094) -- (650,1100) -- (648,1104) -- (648,1109) -- (646,1111) -- (645,1116) -- (642,1122) -- (642,1125) -- (640,1128) -- (640,1130) -- (637,1137) -- (638,1139) -- (637,1141) -- (637,1146) -- (634,1149) -- (632,1147) -- (632,1134) -- (633,1133) -- (633,1129) -- (632,1128) -- (633,1126) -- (633,1097) -- (632,1095) -- (634,1093) -- (634,1088) -- (633,1087) -- (633,1085) -- (634,1084) -- (634,1082) -- (633,1081) -- (633,1074) -- (634,1072) -- (633,1071) -- (633,1065) -- (632,1064) -- (632,1056) -- (633,1055) -- (633,1053) -- (632,1052) -- (632,1048) -- (633,1047) -- (632,1045) -- (632,1043) -- (633,1042) -- (632,1040) -- (632,1037) -- (633,1036) -- (633,1033) -- (632,1031) -- (634,1029) -- (634,1027) -- (633,1026) -- (634,1024) -- (633,1022) -- (633,1014) -- (634,1013) -- (633,1012) -- (633,1007) -- (632,1006) -- (633,1004) -- (633,1002) -- (632,1001) -- (633,991) -- (631,989) -- (631,987) -- (633,985) -- (632,981) -- (634,979) -- (634,977) -- (632,973) -- (632,970) -- (633,969) -- (633,966) -- (632,965) -- (632,963) -- (633,962) -- (633,950) -- (631,948) -- (632,946) -- (631,940) -- (633,938) -- (633,933) -- (632,932) -- (632,928) -- (633,927) -- (633,925) -- (632,924) -- (632,916) -- (633,915) -- (633,907) -- (632,906) -- (633,904) -- (633,884) -- (632,883) -- (633,878) -- (632,877) -- (632,875) -- (633,874) -- (633,870) -- (632,869) -- (631,865) -- (632,864) -- (632,859) -- (633,858) -- (632,856) -- (632,853) -- (634,851) -- (641,854) -- (644,857) -- (648,858) -- (651,860) -- (654,860) -- (656,858) -- (656,857) -- (653,854) -- (651,854) -- (649,852) -- (647,852) -- (637,847) -- (633,848) -- (631,846) -- (632,845) -- (632,834) -- (633,833) -- (633,820) -- (632,818) -- (634,814) -- (634,811) -- (632,807) -- (632,804) -- (634,802) -- (634,799) -- (633,798) -- (633,795) -- (632,794) -- (633,784) -- (631,781) -- (631,779) -- (633,775) -- (633,766) -- (632,764) -- (633,763) -- (633,757) -- (634,756) -- (634,754) -- (632,750) -- (632,747) -- (634,745) -- (633,743) -- (634,739) -- (632,736) -- (632,732) -- (634,729) -- (634,725) -- (632,723) -- (632,722) -- (634,720) -- (634,718) -- (633,717) -- (633,710) -- (634,709) -- (633,708) -- (634,707) -- (633,706) -- (633,697) -- (632,696) -- (632,693) -- (633,692) -- (633,685) -- (634,684) -- (634,682) -- (633,681) -- (633,673) -- (634,672) -- (632,669) -- (633,667) -- (633,663) -- (634,662) -- (633,648) -- (632,647) -- (633,645) -- (633,641) -- (632,639) -- (634,635) -- (634,632) -- (632,629) -- (632,627) -- (633,626) -- (633,622) -- (632,621) -- (632,611) -- (633,610) -- (632,608) -- (632,590) -- (633,589) -- (632,587) -- (632,582) -- (633,581) -- (633,577) -- (632,576) -- (632,570) -- (631,569) -- (631,563) -- (633,561) -- (631,557) -- (631,548) -- (632,547) -- (631,546) -- (631,542) -- (632,541) -- (632,538) -- cycle
(662,863) -- (662,865) -- (665,868) -- (669,869) -- (675,874) -- (683,874) -- (684,873) -- (684,871) -- (683,870) -- (681,870) -- (676,866) -- (668,862) -- (663,862) -- cycle
(1071,436) -- (1073,438) -- (1073,439) -- (1072,440) -- (1072,444) -- (1071,445) -- (1070,450) -- (1066,458) -- (1066,464) -- (1061,473) -- (1061,476) -- (1059,478) -- (1058,485) -- (1056,487) -- (1056,490) -- (1055,491) -- (1055,493) -- (1054,494) -- (1054,496) -- (1053,497) -- (1053,499) -- (1050,506) -- (1050,509) -- (1048,512) -- (1044,526) -- (1042,529) -- (1041,538) -- (1039,540) -- (1039,542) -- (1038,543) -- (1038,547) -- (1036,549) -- (1036,554) -- (1033,558) -- (1033,561) -- (1031,565) -- (1028,568) -- (1029,569) -- (1029,573) -- (1028,574) -- (1028,577) -- (1025,582) -- (1025,585) -- (1022,592) -- (1021,599) -- (1020,600) -- (1019,605) -- (1017,608) -- (1017,610) -- (1015,614) -- (1014,624) -- (1010,631) -- (1010,634) -- (1007,640) -- (1007,643) -- (1005,645) -- (1004,653) -- (1002,656) -- (1002,658) -- (999,662) -- (999,664) -- (997,668) -- (997,671) -- (994,676) -- (994,679) -- (992,683) -- (992,688) -- (987,696) -- (986,704) -- (985,705) -- (985,709) -- (984,710) -- (984,713) -- (982,717) -- (981,723) -- (978,727) -- (977,733) -- (976,734) -- (974,742) -- (972,745) -- (972,747) -- (970,749) -- (970,751) -- (969,752) -- (969,757) -- (967,760) -- (967,763) -- (964,768) -- (964,771) -- (962,773) -- (962,776) -- (960,779) -- (960,782) -- (957,787) -- (957,789) -- (955,793) -- (955,799) -- (952,803) -- (952,807) -- (949,812) -- (949,815) -- (947,818) -- (947,821) -- (945,824) -- (944,830) -- (940,837) -- (940,839) -- (935,843) -- (933,847) -- (933,849) -- (930,854) -- (930,859) -- (932,861) -- (931,867) -- (928,872) -- (928,875) -- (925,881) -- (925,886) -- (924,887) -- (924,890) -- (920,896) -- (919,904) -- (917,907) -- (917,910) -- (915,913) -- (915,916) -- (912,922) -- (912,925) -- (909,930) -- (909,933) -- (905,942) -- (905,946) -- (903,948) -- (902,953) -- (899,959) -- (899,964) -- (897,967) -- (896,974) -- (895,975) -- (893,975) -- (891,973) -- (887,971) -- (883,971) -- (883,974) -- (884,975) -- (890,976) -- (891,977) -- (891,979) -- (885,984) -- (881,985) -- (878,988) -- (872,991) -- (866,996) -- (863,997) -- (860,1000) -- (844,1010) -- (837,1017) -- (834,1017) -- (825,1024) -- (824,1024) -- (819,1028) -- (816,1029) -- (807,1038) -- (805,1038) -- (803,1040) -- (800,1041) -- (796,1045) -- (795,1045) -- (787,1051) -- (784,1052) -- (779,1057) -- (777,1057) -- (775,1059) -- (772,1060) -- (766,1066) -- (762,1068) -- (759,1071) -- (756,1072) -- (752,1076) -- (747,1078) -- (744,1081) -- (737,1085) -- (733,1089) -- (731,1089) -- (729,1091) -- (728,1091) -- (721,1098) -- (715,1101) -- (711,1105) -- (710,1105) -- (702,1111) -- (700,1111) -- (687,1122) -- (685,1122) -- (682,1125) -- (677,1127) -- (673,1131) -- (670,1132) -- (666,1136) -- (664,1136) -- (655,1144) -- (651,1146) -- (648,1149) -- (646,1149) -- (643,1153) -- (642,1153) -- (640,1155) -- (638,1155) -- (637,1154) -- (638,1152) -- (638,1149) -- (642,1141) -- (642,1138) -- (644,1134) -- (644,1131) -- (646,1129) -- (646,1125) -- (650,1119) -- (650,1113) -- (652,1111) -- (652,1108) -- (655,1103) -- (655,1099) -- (658,1093) -- (659,1087) -- (663,1080) -- (664,1073) -- (666,1069) -- (667,1060) -- (669,1057) -- (669,1053) -- (673,1045) -- (674,1037) -- (677,1034) -- (677,1030) -- (679,1028) -- (679,1026) -- (682,1019) -- (682,1015) -- (683,1014) -- (685,1003) -- (687,1001) -- (688,994) -- (691,988) -- (691,983) -- (693,981) -- (693,979) -- (695,976) -- (694,969) -- (698,965) -- (698,963) -- (699,962) -- (699,960) -- (702,953) -- (702,950) -- (703,949) -- (703,947) -- (706,940) -- (706,936) -- (708,933) -- (708,929) -- (711,924) -- (711,921) -- (714,916) -- (714,913) -- (716,910) -- (717,902) -- (719,899) -- (719,896) -- (722,889) -- (722,883) -- (724,880) -- (724,875) -- (727,872) -- (727,870) -- (729,867) -- (729,865) -- (730,864) -- (730,862) -- (731,861) -- (731,859) -- (732,858) -- (732,856) -- (735,849) -- (735,845) -- (739,837) -- (739,834) -- (741,831) -- (741,828) -- (744,822) -- (744,818) -- (747,813) -- (747,810) -- (750,804) -- (750,799) -- (752,796) -- (752,791) -- (755,787) -- (756,779) -- (758,776) -- (758,773) -- (760,770) -- (761,764) -- (764,761) -- (766,748) -- (769,743) -- (770,737) -- (772,735) -- (772,731) -- (773,730) -- (774,725) -- (777,722) -- (777,717) -- (780,711) -- (780,704) -- (781,703) -- (781,700) -- (784,695) -- (784,690) -- (785,689) -- (785,687) -- (788,683) -- (788,678) -- (790,676) -- (790,673) -- (793,669) -- (794,660) -- (796,657) -- (796,654) -- (798,651) -- (798,648) -- (801,645) -- (801,639) -- (805,635) -- (804,634) -- (804,632) -- (805,631) -- (805,626) -- (806,625) -- (805,622) -- (807,618) -- (812,613) -- (813,613) -- (822,606) -- (828,604) -- (836,598) -- (841,596) -- (844,593) -- (847,592) -- (853,587) -- (855,587) -- (859,583) -- (862,582) -- (865,579) -- (869,577) -- (873,573) -- (875,573) -- (881,568) -- (883,568) -- (890,562) -- (892,562) -- (900,554) -- (903,553) -- (906,550) -- (908,550) -- (918,542) -- (920,542) -- (928,535) -- (931,534) -- (937,529) -- (938,529) -- (943,524) -- (947,522) -- (951,518) -- (961,512) -- (964,509) -- (969,507) -- (972,504) -- (980,500) -- (992,489) -- (998,486) -- (1000,486) -- (1009,478) -- (1016,476) -- (1017,475) -- (1017,473) -- (1019,471) -- (1022,470) -- (1030,464) -- (1033,463) -- (1041,455) -- (1045,455) -- (1050,451) -- (1052,452) -- (1053,451) -- (1053,449) -- (1056,446) -- (1060,445) -- (1065,439) -- cycle
(852,957) -- (852,959) -- (854,961) -- (856,961) -- (862,966) -- (866,966) -- (870,969) -- (872,969) -- (873,970) -- (876,970) -- (877,969) -- (877,966) -- (876,965) -- (874,965) -- (870,961) -- (867,960) -- (863,957) -- (861,957) -- (857,955) -- (855,955) -- cycle
(821,941) -- (821,942) -- (826,947) -- (830,949) -- (832,949) -- (838,953) -- (841,953) -- (843,951) -- (843,948) -- (839,945) -- (832,943) -- (828,940) -- (822,940) -- cycle
(788,925) -- (788,927) -- (791,930) -- (797,932) -- (801,936) -- (810,938) -- (812,935) -- (812,934) -- (809,931) -- (802,929) -- (795,924) -- (789,924) -- cycle
(757,910) -- (758,912) -- (759,912) -- (764,916) -- (768,917) -- (774,922) -- (777,922) -- (780,920) -- (780,918) -- (778,916) -- (762,908) -- (760,908) -- cycle
(727,895) -- (727,896) -- (731,900) -- (734,900) -- (736,902) -- (740,903) -- (743,906) -- (748,906) -- (750,904) -- (749,902) -- (742,899) -- (740,897) -- (734,894) -- (728,894) -- cycle
(630,209) -- (632,207) -- (636,207) -- (647,215) -- (649,215) -- (653,218) -- (655,218) -- (658,221) -- (661,221) -- (663,223) -- (666,223) -- (671,228) -- (677,229) -- (679,231) -- (683,233) -- (690,235) -- (693,238) -- (695,238) -- (699,241) -- (701,241) -- (704,244) -- (712,247) -- (718,252) -- (725,253) -- (729,256) -- (736,259) -- (738,261) -- (740,261) -- (745,264) -- (747,264) -- (751,267) -- (756,268) -- (758,270) -- (767,274) -- (769,276) -- (774,277) -- (779,281) -- (781,281) -- (783,283) -- (787,285) -- (789,285) -- (792,288) -- (796,289) -- (804,293) -- (806,295) -- (808,295) -- (811,297) -- (814,297) -- (827,305) -- (830,305) -- (836,310) -- (838,310) -- (841,313) -- (844,314) -- (846,316) -- (849,316) -- (854,319) -- (857,319) -- (861,323) -- (866,325) -- (868,327) -- (870,327) -- (876,330) -- (878,332) -- (880,332) -- (882,334) -- (885,334) -- (892,340) -- (896,340) -- (897,341) -- (899,341) -- (901,343) -- (904,344) -- (907,347) -- (914,349) -- (916,351) -- (918,351) -- (920,353) -- (924,354) -- (926,356) -- (928,356) -- (930,358) -- (936,360) -- (939,363) -- (942,363) -- (944,365) -- (955,370) -- (957,372) -- (959,372) -- (961,374) -- (963,374) -- (967,377) -- (971,378) -- (973,380) -- (977,381) -- (982,385) -- (987,386) -- (998,392) -- (1000,392) -- (1002,394) -- (1012,399) -- (1015,402) -- (1019,402) -- (1022,405) -- (1025,405) -- (1032,410) -- (1038,412) -- (1041,415) -- (1043,416) -- (1046,416) -- (1050,419) -- (1053,419) -- (1056,422) -- (1058,422) -- (1068,427) -- (1070,427) -- (1072,429) -- (1069,432) -- (1064,434) -- (1051,444) -- (1049,444) -- (1046,447) -- (1043,448) -- (1034,455) -- (1032,455) -- (1030,457) -- (1029,457) -- (1024,462) -- (1022,462) -- (1018,466) -- (1013,468) -- (1010,471) -- (1005,473) -- (1002,476) -- (999,477) -- (992,483) -- (988,485) -- (984,489) -- (981,490) -- (977,494) -- (972,496) -- (969,499) -- (966,500) -- (962,504) -- (954,508) -- (951,511) -- (944,515) -- (940,519) -- (938,519) -- (936,521) -- (932,523) -- (927,528) -- (925,528) -- (920,534) -- (917,535) -- (911,540) -- (905,543) -- (889,556) -- (881,560) -- (877,564) -- (864,572) -- (858,577) -- (854,578) -- (841,588) -- (836,590) -- (834,592) -- (832,592) -- (828,596) -- (825,597) -- (821,601) -- (817,603) -- (809,610) -- (808,610) -- (803,614) -- (799,614) -- (798,613) -- (796,613) -- (792,611) -- (790,609) -- (784,606) -- (782,606) -- (775,601) -- (773,601) -- (769,598) -- (767,598) -- (765,596) -- (761,595) -- (758,592) -- (755,592) -- (735,582) -- (731,579) -- (729,579) -- (720,574) -- (718,574) -- (714,572) -- (712,570) -- (710,570) -- (686,558) -- (684,556) -- (677,553) -- (675,551) -- (671,550) -- (664,546) -- (662,546) -- (654,542) -- (652,540) -- (650,540) -- (648,538) -- (645,538) -- (641,536) -- (639,534) -- (632,533) -- (630,531) -- (630,529) -- (631,528) -- (630,525) -- (632,523) -- (631,521) -- (631,519) -- (632,518) -- (632,516) -- (631,515) -- (632,509) -- (631,508) -- (631,500) -- (630,498) -- (631,497) -- (631,482) -- (632,481) -- (632,472) -- (631,471) -- (631,468) -- (632,467) -- (632,464) -- (633,463) -- (633,457) -- (631,454) -- (631,448) -- (632,446) -- (630,444) -- (630,433) -- (631,432) -- (631,430) -- (630,429) -- (630,427) -- (631,426) -- (631,421) -- (630,420) -- (630,418) -- (631,417) -- (631,415) -- (629,413) -- (630,409) -- (631,408) -- (631,405) -- (630,404) -- (630,402) -- (631,401) -- (631,395) -- (630,394) -- (630,392) -- (631,391) -- (631,372) -- (630,371) -- (630,369) -- (631,368) -- (631,356) -- (630,355) -- (631,353) -- (631,345) -- (630,344) -- (630,331) -- (629,330) -- (630,328) -- (630,321) -- (629,319) -- (631,316) -- (631,312) -- (630,311) -- (630,300) -- (631,299) -- (631,292) -- (630,291) -- (631,289) -- (630,287) -- (630,285) -- (631,284) -- (631,269) -- (630,268) -- (630,264) -- (629,263) -- (629,253) -- (630,252) -- (630,242) -- (631,241) -- (631,239) -- (630,238) -- (630,236) -- (631,235) -- (631,232) -- (630,231) -- (631,226) -- (629,223) -- (629,221) -- (630,220) -- cycle
(264,352) -- (266,350) -- (268,350) -- (272,347) -- (275,347) -- (277,346) -- (279,342) -- (284,337) -- (289,335) -- (299,326) -- (302,325) -- (304,323) -- (308,321) -- (313,316) -- (317,314) -- (328,305) -- (334,302) -- (343,294) -- (349,291) -- (354,285) -- (356,285) -- (362,280) -- (363,280) -- (369,274) -- (370,274) -- (372,272) -- (374,272) -- (385,261) -- (388,260) -- (390,258) -- (394,256) -- (398,252) -- (404,249) -- (409,244) -- (413,242) -- (417,238) -- (422,236) -- (433,226) -- (434,226) -- (439,222) -- (444,220) -- (447,217) -- (451,215) -- (458,208) -- (461,207) -- (464,204) -- (474,198) -- (477,195) -- (480,194) -- (490,185) -- (492,185) -- (495,182) -- (496,182) -- (509,170) -- (512,169) -- (525,158) -- (527,159) -- (527,162) -- (526,163) -- (525,169) -- (523,173) -- (524,174) -- (526,174) -- (528,172) -- (528,169) -- (531,163) -- (531,159) -- (533,157) -- (534,157) -- (537,160) -- (541,162) -- (544,162) -- (550,165) -- (552,167) -- (551,171) -- (552,172) -- (558,173) -- (562,178) -- (563,178) -- (567,174) -- (568,174) -- (575,179) -- (580,180) -- (582,182) -- (591,186) -- (593,188) -- (597,189) -- (603,192) -- (605,194) -- (609,196) -- (611,196) -- (613,198) -- (615,198) -- (616,199) -- (618,199) -- (620,201) -- (622,201) -- (626,205) -- (626,207) -- (625,208) -- (625,213) -- (624,214) -- (625,221) -- (624,222) -- (624,227) -- (623,228) -- (623,231) -- (624,232) -- (623,234) -- (623,237) -- (624,238) -- (624,243) -- (622,245) -- (623,251) -- (624,252) -- (624,255) -- (623,256) -- (623,260) -- (624,261) -- (624,263) -- (623,264) -- (623,268) -- (624,269) -- (624,287) -- (623,288) -- (624,291) -- (623,292) -- (623,296) -- (624,297) -- (624,300) -- (622,302) -- (623,304) -- (623,307) -- (622,308) -- (622,312) -- (623,313) -- (623,315) -- (622,316) -- (623,318) -- (623,325) -- (622,326) -- (622,333) -- (623,335) -- (622,336) -- (622,338) -- (623,339) -- (623,341) -- (621,343) -- (621,344) -- (622,345) -- (623,371) -- (624,372) -- (624,374) -- (623,375) -- (623,405) -- (624,407) -- (623,408) -- (623,411) -- (622,412) -- (623,413) -- (623,438) -- (622,439) -- (622,442) -- (624,445) -- (623,447) -- (623,450) -- (624,451) -- (624,456) -- (625,457) -- (625,459) -- (624,460) -- (624,494) -- (625,495) -- (625,498) -- (624,499) -- (625,510) -- (624,511) -- (624,513) -- (625,514) -- (625,524) -- (627,526) -- (624,529) -- (623,529) -- (615,525) -- (613,523) -- (608,521) -- (606,519) -- (602,517) -- (600,517) -- (598,515) -- (587,510) -- (585,508) -- (583,508) -- (571,501) -- (564,499) -- (559,495) -- (557,495) -- (553,492) -- (550,492) -- (548,490) -- (544,488) -- (541,488) -- (539,487) -- (535,483) -- (529,482) -- (524,478) -- (521,478) -- (519,476) -- (509,471) -- (507,471) -- (503,468) -- (499,466) -- (497,466) -- (495,464) -- (493,464) -- (491,462) -- (483,458) -- (481,458) -- (473,454) -- (471,452) -- (469,452) -- (458,446) -- (453,446) -- (448,442) -- (441,440) -- (439,438) -- (439,437) -- (442,433) -- (442,431) -- (444,429) -- (444,427) -- (446,424) -- (446,422) -- (447,421) -- (447,417) -- (446,416) -- (443,416) -- (441,418) -- (441,421) -- (439,423) -- (439,425) -- (438,426) -- (439,428) -- (438,430) -- (438,434) -- (435,437) -- (434,437) -- (432,435) -- (426,432) -- (424,432) -- (419,428) -- (415,427) -- (413,425) -- (407,422) -- (405,422) -- (403,420) -- (396,417) -- (394,415) -- (391,415) -- (386,411) -- (384,411) -- (380,408) -- (376,407) -- (371,403) -- (366,402) -- (363,400) -- (360,400) -- (358,398) -- (345,392) -- (343,390) -- (339,389) -- (335,386) -- (333,386) -- (325,382) -- (323,380) -- (319,378) -- (317,378) -- (315,376) -- (311,375) -- (309,373) -- (307,373) -- (305,371) -- (301,370) -- (299,368) -- (293,365) -- (289,364) -- (284,360) -- (277,359) -- (272,356) -- (266,355) -- (264,353) -- cycle
(457,380) -- (456,380) -- (452,384) -- (452,387) -- (450,390) -- (449,396) -- (447,399) -- (447,405) -- (448,406) -- (450,406) -- (452,404) -- (452,401) -- (455,395) -- (456,388) -- (458,386) -- (458,381) -- cycle
(466,347) -- (465,347) -- (462,350) -- (462,352) -- (460,356) -- (461,359) -- (459,362) -- (459,365) -- (457,367) -- (457,369) -- (458,370) -- (461,370) -- (463,369) -- (465,366) -- (465,363) -- (467,360) -- (467,354) -- (468,353) -- (468,349) -- cycle
(478,317) -- (475,317) -- (471,323) -- (469,330) -- (467,332) -- (467,338) -- (468,339) -- (471,339) -- (475,334) -- (476,326) -- (480,321) -- (480,319) -- cycle
(488,283) -- (487,283) -- (485,285) -- (483,292) -- (480,296) -- (479,304) -- (481,306) -- (483,306) -- (486,303) -- (486,301) -- (487,300) -- (488,293) -- (491,285) -- cycle
(500,251) -- (498,251) -- (495,254) -- (495,256) -- (493,259) -- (493,261) -- (490,268) -- (490,271) -- (489,273) -- (492,276) -- (493,276) -- (495,274) -- (496,270) -- (497,269) -- (497,264) -- (499,261) -- (499,259) -- (501,255) -- (501,252) -- cycle
(511,216) -- (510,216) -- (506,221) -- (506,226) -- (504,228) -- (503,234) -- (500,238) -- (500,240) -- (501,241) -- (504,241) -- (507,238) -- (507,235) -- (508,234) -- (508,232) -- (510,229) -- (510,227) -- (513,220) -- (513,218) -- cycle
(522,183) -- (520,184) -- (516,192) -- (516,195) -- (514,199) -- (512,201) -- (514,205) -- (517,205) -- (519,203) -- (519,201) -- (521,198) -- (521,196) -- (522,195) -- (522,191) -- (524,189) -- (524,185) -- cycle;
\node[inner sep=0pt,font=\fontsize{10}{10}\selectfont] at (528,124) {\scalebox{4.6}{$5$}};
\node[inner sep=0pt,font=\fontsize{10}{10}\selectfont] at (226,347) {\scalebox{4.6}{$4$}};
\node[inner sep=0pt,font=\fontsize{10}{10}\selectfont] at (1105,429) {\scalebox{4.6}{$6$}};
\node[inner sep=0pt,font=\fontsize{10}{10}\selectfont] at (826,634) {\scalebox{4.7}{$2$}};
\node[inner sep=0pt,font=\fontsize{10}{10}\selectfont] at (389,692) {\scalebox{4.6}{$7$}};
\node[inner sep=0pt,font=\fontsize{10}{10}\selectfont] at (48,920) {\scalebox{4.7}{$3$}};
\node[inner sep=0pt,font=\fontsize{10}{10}\selectfont] at (918,992) {\scalebox{4.4}{$1$}};
\node[inner sep=0pt,font=\fontsize{10}{10}\selectfont] at (630,1190) {\scalebox{4.4}{$0$}};
\end{tikzpicture}%
}
\end{minipage}\end{center}

[Bei dem Beispiel haben diese Zahlen beziehungsweise die Werthe:
\[
0,\quad 3+2i,\quad 2+6i,\quad -6+3i,\quad 5+8i,\quad -3+5i,\quad -1+11i,\quad -4+9i.]
\]
%
\clearpage
\SourceLeaf{355}

\begin{center}
\large\bfseries GEOMETRISCHE SEITE DER TERNÄREN FORMEN.
\end{center}



Ein Punkt im Raume $(0)$ sei als Anfangspunkt angenommen. Der Übergang von da zu drei
andern Punkten $\boldsymbol{P}$, $\boldsymbol{P}'$, $\boldsymbol{P}''$, die mit jenem nicht
in einer Ebene liegen, sei resp. $t$, $t'$, $t''$; wo, so oft keine Verwechslung möglich
ist, die Punkte $\boldsymbol{P}$, $\boldsymbol{P}'$, $\boldsymbol{P}''$ selbst durch $(t)$,
$(t')$, $(t'')$ bezeichnet werden mögen.

Es sei ferner allgemein $(t,t')$ das Product der Länge (kürzesten Entfernung) der beiden
Linien $t$, $t'$ in den Cosinus ihrer Neigung etc.

Man hat allgemein $(\alpha t + \alpha' t' + \alpha'' t'' + .\,.\,.,\ \beta u + \beta' u' +
\beta'' u'' + .\,.\,.)$, wenn man die Multiplication
\[
(\alpha t + \alpha' t' + \alpha'' t'' + .\,.\,.) \times (\beta u + \beta' u' + \beta'' u'' + .\,.\,.)
\]
ausführt und statt $tu$, $tu'$, $tu''$, $t'u'$, $t'u''$ u.\,s.\,w. $(t,u)$, $(t,u')$,
$(t,u'')$, $(t',u')$, $(t',u'')$ u.\,s.\,w. schreibt.

Jeder Punkt im Raume wird durch ein Trinomium
\[
(xt + x't' + x''t'')
\]
dargestellt werden können.

Für alle Punkte, die in einer bestimmten Ebene liegen, wird dann eine Gleichung
\[
\lambda x + \lambda' x' + \lambda'' x'' = L
\]
stattfinden, wo $\lambda$, $\lambda'$, $\lambda''$, $L$ bestimmte Zahlen bedeuten. Für eine
Ebene durch die drei Punkte $\mu t$, $\mu' t'$, $\mu'' t''$ ist
\[
\lambda\mu = \lambda'\mu' = \lambda''\mu'' = L.
\]
%
\SourceLeaf{356}

Schreibt man
{\small\[
(t,t) = a,\quad (t',t') = a',\quad (t'',t'') = a'',\quad (t',t'') = b,\quad (t,t'') = b',\quad (t,t') = b''
\]}
und
\[\begin{aligned}
a'a'' - bb &= A, & aa'' - b'b' &= A', & aa' - b''b'' &= A'', \\
b'b'' - ab &= B, & bb'' - a'b' &= B', & bb' - a''b'' &= B'',
\end{aligned}\]
\[
D = aa'a'' + 2bb'b'' - abb - a'b'b' - a''b''b'',
\]
so ist
\[
\begin{array}{ll}
 & \text{senkrecht gegen}\\
 T=At+B''t'+B't'' & t'\text{ und }t''\\
 T'=B''t+A't'+Bt'' & t\text{ und }t''\\
 T''=B't+Bt'+A''t'' & t\text{ und }t',
\end{array}
\]
und allgemein, wenn
\[
\lambda x + \lambda' x' + \lambda'' x'' = L
\]
die Gleichung einer Ebene ist, so wird die Linie
\[
\lambda T + \lambda' T' + \lambda'' T''
\]
gegen dieselbe senkrecht sein. $L =$ Werth der Form $\begin{pmatrix} A, A', A'' \\ B, B', B'' \end{pmatrix}$,
wenn die Unbestimmten $= \lambda$, $\lambda'$, $\lambda''$ gesetzt werden.

Es ist dann ferner
\[\begin{aligned}
aT + b''T' + b'T'' &= Dt \\
b''T + a'T' + bT'' &= Dt' \\
b'T + bT' + a''T'' &= Dt'',
\end{aligned}\]
und die Linien $t$, $t'$, $t''$ sind senkrecht gegen die Ebenen, deren Gleichungen
\[\begin{aligned}
ax + b''x' + b'x'' &= \text{Const.} \\
b''x + a'x' + bx'' &= \text{Const.} \\
b'x + bx' + a''x'' &= \text{Const.}
\end{aligned}\]

Der doppelte Flächeninhalt des Dreiecks durch die Punkte $mt$, $m't'$, $m''t''$ ist aequal
der Quadratwurzel aus dem Werthe der Form
\[
\begin{pmatrix} A, A', A'' \\ B, B', B'' \end{pmatrix} = F(m'm'',\, mm'',\, mm'),
\]
%
\SourceLeaf{357}

\noindent wenn substituirt wird $X = m'm''$, $X' = mm''$, $X'' = mm'$, während der sechsfache
Cubikinhalt der Pyramide, die sich dadurch mit dem Nullpunkte bildet, $= mm'm''\BareRoot{D}$
wird; folglich [ist] das Perpendikel
\[
= \BareRoot{\dfrac{D}{F\left(\dfrac{1}{m}, \dfrac{1}{m'}, \dfrac{1}{m''}\right)}};
\]

$T$, $T'$, $T''$ beziehen sich eben so auf die Form $\begin{pmatrix} AD, A'D, A''D \\ BD, B'D, B''D \end{pmatrix}$
wie $t$, $t'$, $t''$ auf $\begin{pmatrix} a, a', a'' \\ b, b', b'' \end{pmatrix}$.



\begin{center}
\bfseries BEMERKUNG.
\end{center}

Die vorstehende Notiz befindet sich in einem Handbuche und stammt aus dem Juli 1831. Da sie
zeigt, dass \textsc{Gauss} sich auch mit complexen Grössen von mehr als zwei Einheiten beschäftigt
hat, die er bereits in der Schlussbemerkung seiner Selbstanzeige der Theoria residuorum
biquadraticorum, Commentatio secunda vom 23. April 1831 (Bd. II. S. 169 bis 198) erwähnt
hatte, so hat sie hier noch einmal Platz gefunden, obwohl sie schon in Bd. II. S. 305 bis
307 wegen ihrer Beziehung zur Theorie der ternären Formen abgedruckt ist.
\begin{flushright}
\textsc{Stäckel.}
\end{flushright}


%
\clearpage\SourceLeaf{358}

\begin{center}
\large\bfseries [DIE KUGEL.]
\end{center}

\begin{center}
[1.]
\end{center}

\begin{center}
\bfseries ORTHOGRAPHISCHE PROJECTION DER KUGEL.
\end{center}

\begin{center}
Veränderte Erscheinung eines unendlich kleinen Theils der Kugelfläche,\\
wenn die Projectionsebene verändert wird.
\end{center}

Man wählt zur Axe der $x$ die Durchschnittslinie der beiden Projectionsebenen oder vielmehr als Axe der $y$ in der bleibenden Projectionsebene eine Parallele mit der geraden Linie, welche die beiden Projectionsplätze verbindet.

Man fragt, wie dasjenige Stück der Kugelfläche, welches in der ersten Projection wie ein Kreis erschien, in der zweiten erscheinen wird.

Coordinaten des Centrums dieser Fläche (Radius $= 1$):

\begin{center}
\begin{tabular}{ccccc}
in & der & ersten & Projection & $x,y,z$\\
,, & ,, & zweiten & ,, & $x',y',z'$.
\end{tabular}
\end{center}

Der Drehungswinkel $= \lambda$.

Es ist
\[
x' = x
\]
\[
\BareRoot{(1-xx)} = \rho
\]
\[
\BareRoot{(1-x'x')} = \rho
\]
\[
\frac{y}{\rho} = \sin u \qquad \frac{y'}{\rho} = \sin(u+\lambda)
\]
\[
dy' = \frac{z'}{z}\cdot dy - \frac{x\sin\lambda}{z}\cdot dx.
\]
%
\SourceLeaf{359}

Es erscheint also ein $\alpha+\beta.i$ wie $\alpha+\dfrac{\beta z' - \alpha\cdot x\sin\lambda}{z}\cdot i$ oder die beiden Radien des Kreises $1$ und $i$ in der ersten Projection erscheinen als conjugirte Halbmesser der Ellipse in der zweiten Projection $1-\dfrac{x\sin\lambda}{z}\cdot i$ und $\dfrac{z'}{z}\cdot i$, woraus sich nach p.\,164 [des Handbuches\HistoricalNote{*)}{[*)}{Dieser Band, S.\,339.]} die Dimensionen der Projectionsellipse leicht ableiten lassen.



\Needspace{20\baselineskip}
\begin{center}
[II.]
\end{center}

\begin{center}
\bfseries [STEREOGRAPHISCHE UND CENTRALE PROJECTION.]
\end{center}

\begin{center}
[1.]
\end{center}

\begin{center}
Zusammenhang der stereographischen und der Centralprojection, so dass\\
Massstab für jene im Pol $C$ halb so gross als für diese.
\end{center}

Der Kreis durch $AB$ stellt für erstere den Äquator, für letztere den Parallel von $45^0$ vor.

\begin{center}\resizebox{4.0cm}{!}{%
\begin{tikzpicture}[x=1pt,y=-1pt]
\path[use as bounding box] (0,0) rectangle (578,1044);
\path[fill=black,draw=none,even odd rule]
(299,78) -- (296,78) -- (295,79) -- (295,83) -- (294,84) -- (294,87) -- (291,94) -- (291,97) -- (289,99) -- (289,103) -- (287,105) -- (287,107) -- (284,112) -- (284,115) -- (280,123) -- (281,129) -- (278,133) -- (279,137) -- (275,142) -- (275,146) -- (273,148) -- (270,154) -- (270,156) -- (266,163) -- (266,167) -- (264,169) -- (265,171) -- (265,176) -- (263,178) -- (263,181) -- (261,184) -- (261,188) -- (259,190) -- (259,194) -- (257,198) -- (257,202) -- (254,206) -- (254,208) -- (252,211) -- (252,217) -- (251,218) -- (251,222) -- (250,223) -- (250,225) -- (246,230) -- (247,234) -- (245,237) -- (245,240) -- (242,246) -- (242,252) -- (241,253) -- (240,260) -- (238,263) -- (235,274) -- (233,276) -- (233,278) -- (230,283) -- (229,288) -- (225,295) -- (226,297) -- (226,301) -- (221,311) -- (222,317) -- (221,318) -- (221,320) -- (219,322) -- (219,325) -- (216,328) -- (216,331) -- (214,334) -- (214,337) -- (212,339) -- (211,343) -- (209,346) -- (210,348) -- (210,352) -- (209,353) -- (208,358) -- (206,361) -- (205,367) -- (203,369) -- (203,371) -- (200,378) -- (200,385) -- (198,387) -- (198,391) -- (197,392) -- (196,397) -- (193,402) -- (192,407) -- (190,409) -- (190,412) -- (188,415) -- (188,418) -- (189,420) -- (186,423) -- (185,431) -- (180,441) -- (180,447) -- (178,450) -- (178,453) -- (176,455) -- (176,457) -- (173,462) -- (173,471) -- (170,475) -- (170,477) -- (169,478) -- (169,482) -- (166,488) -- (166,490) -- (164,494) -- (164,497) -- (162,499) -- (162,502) -- (161,503) -- (161,508) -- (160,509) -- (160,514) -- (158,517) -- (157,523) -- (155,525) -- (155,531) -- (152,535) -- (152,537) -- (148,545) -- (148,549) -- (146,551) -- (146,556) -- (143,560) -- (143,566) -- (142,567) -- (142,570) -- (139,573) -- (138,582) -- (135,586) -- (134,589) -- (132,591) -- (132,598) -- (131,599) -- (131,601) -- (128,608) -- (128,611) -- (126,615) -- (125,623) -- (123,625) -- (123,628) -- (120,633) -- (120,639) -- (119,640) -- (119,643) -- (117,645) -- (118,648) -- (116,651) -- (115,657) -- (112,661) -- (112,664) -- (106,669) -- (105,672) -- (99,679) -- (98,683) -- (94,687) -- (94,692) -- (91,695) -- (89,702) -- (87,704) -- (85,711) -- (83,713) -- (82,719) -- (80,722) -- (80,727) -- (79,728) -- (79,730) -- (76,733) -- (76,737) -- (75,738) -- (74,745) -- (72,748) -- (72,756) -- (70,759) -- (70,765) -- (69,766) -- (69,768) -- (70,769) -- (70,773) -- (69,774) -- (69,783) -- (68,784) -- (68,786) -- (70,790) -- (70,794) -- (68,797) -- (69,799) -- (69,802) -- (68,803) -- (68,809) -- (69,810) -- (69,822) -- (71,824) -- (71,836) -- (72,837) -- (72,841) -- (73,842) -- (73,846) -- (74,847) -- (74,852) -- (76,854) -- (76,856) -- (77,857) -- (78,864) -- (80,867) -- (80,870) -- (82,873) -- (82,877) -- (83,878) -- (84,883) -- (87,888) -- (88,892) -- (92,896) -- (92,898) -- (95,901) -- (95,903) -- (97,905) -- (100,912) -- (103,915) -- (104,918) -- (106,920) -- (107,922) -- (107,925) -- (119,938) -- (119,939) -- (123,943) -- (124,946) -- (148,969) -- (150,969) -- (160,978) -- (162,978) -- (170,986) -- (174,987) -- (178,991) -- (180,991) -- (182,993) -- (184,993) -- (187,996) -- (189,996) -- (191,998) -- (193,998) -- (196,1001) -- (200,1003) -- (202,1003) -- (205,1005) -- (207,1005) -- (214,1009) -- (218,1009) -- (228,1014) -- (231,1013) -- (236,1016) -- (239,1016) -- (240,1017) -- (242,1017) -- (243,1018) -- (247,1018) -- (248,1019) -- (253,1019) -- (255,1021) -- (259,1021) -- (260,1020) -- (262,1021) -- (269,1021) -- (275,1024) -- (285,1024) -- (286,1023) -- (292,1023) -- (293,1024) -- (296,1023) -- (297,1024) -- (301,1024) -- (302,1023) -- (309,1023) -- (310,1022) -- (316,1022) -- (317,1021) -- (330,1021) -- (331,1020) -- (334,1020) -- (335,1019) -- (337,1020) -- (340,1020) -- (345,1017) -- (348,1017) -- (351,1015) -- (356,1015) -- (360,1013) -- (371,1012) -- (373,1010) -- (377,1010) -- (381,1007) -- (384,1007) -- (389,1003) -- (397,1001) -- (407,996) -- (409,994) -- (412,993) -- (414,991) -- (421,988) -- (424,985) -- (427,984) -- (430,981) -- (438,976) -- (443,971) -- (447,969) -- (457,959) -- (458,959) -- (468,949) -- (470,948) -- (471,945) -- (474,942) -- (474,941) -- (481,934) -- (481,932) -- (483,930) -- (483,927) -- (489,922) -- (490,918) -- (494,914) -- (494,912) -- (496,910) -- (497,906) -- (500,903) -- (504,895) -- (507,892) -- (508,890) -- (508,886) -- (511,881) -- (511,877) -- (514,874) -- (514,871) -- (515,870) -- (515,865) -- (517,863) -- (519,859) -- (519,852) -- (521,849) -- (521,840) -- (523,838) -- (522,835) -- (524,833) -- (524,825) -- (525,824) -- (525,796) -- (524,794) -- (525,793) -- (525,774) -- (526,773) -- (526,771) -- (525,770) -- (525,764) -- (523,761) -- (523,757) -- (522,756) -- (522,753) -- (521,752) -- (521,746) -- (520,745) -- (519,734) -- (516,727) -- (515,721) -- (512,718) -- (512,714) -- (509,711) -- (508,709) -- (508,706) -- (505,703) -- (505,698) -- (503,696) -- (501,692) -- (501,690) -- (499,688) -- (496,681) -- (491,675) -- (486,666) -- (483,663) -- (482,660) -- (477,655) -- (476,652) -- (461,636) -- (460,636) -- (455,631) -- (454,631) -- (446,622) -- (445,622) -- (430,609) -- (429,609) -- (424,605) -- (422,605) -- (418,602) -- (413,600) -- (408,596) -- (406,596) -- (400,593) -- (398,591) -- (390,588) -- (386,585) -- (384,585) -- (381,583) -- (377,583) -- (375,581) -- (372,581) -- (366,578) -- (363,578) -- (362,577) -- (356,576) -- (355,575) -- (346,574) -- (345,573) -- (342,573) -- (341,572) -- (324,571) -- (323,570) -- (319,570) -- (318,569) -- (301,569) -- (299,567) -- (299,564) -- (300,563) -- (300,560) -- (301,559) -- (301,553) -- (300,552) -- (300,549) -- (301,548) -- (301,540) -- (300,539) -- (301,537) -- (300,535) -- (300,533) -- (301,532) -- (301,524) -- (300,523) -- (300,521) -- (301,520) -- (301,513) -- (300,512) -- (300,500) -- (301,499) -- (300,497) -- (300,492) -- (301,491) -- (300,489) -- (300,479) -- (301,478) -- (301,471) -- (300,470) -- (300,464) -- (301,462) -- (299,460) -- (299,457) -- (300,456) -- (300,453) -- (299,451) -- (301,449) -- (301,441) -- (300,440) -- (300,438) -- (301,437) -- (301,424) -- (300,422) -- (301,421) -- (301,416) -- (302,415) -- (302,412) -- (301,411) -- (301,408) -- (302,407) -- (302,401) -- (301,400) -- (301,389) -- (302,388) -- (302,384) -- (300,382) -- (300,380) -- (301,379) -- (301,372) -- (302,371) -- (302,358) -- (301,357) -- (301,352) -- (302,351) -- (302,342) -- (300,340) -- (302,335) -- (301,333) -- (302,331) -- (301,323) -- (302,322) -- (302,316) -- (303,315) -- (302,314) -- (303,313) -- (302,312) -- (303,311) -- (303,305) -- (301,303) -- (303,299) -- (301,296) -- (302,295) -- (301,294) -- (301,283) -- (300,282) -- (302,278) -- (302,275) -- (301,274) -- (301,269) -- (302,268) -- (302,261) -- (301,260) -- (302,259) -- (301,258) -- (302,257) -- (301,256) -- (301,248) -- (302,247) -- (301,245) -- (301,235) -- (302,234) -- (302,232) -- (301,231) -- (301,226) -- (302,225) -- (302,217) -- (303,216) -- (303,214) -- (301,212) -- (301,209) -- (302,208) -- (301,206) -- (301,193) -- (302,191) -- (300,187) -- (302,184) -- (302,181) -- (303,180) -- (303,178) -- (302,177) -- (302,171) -- (301,170) -- (301,154) -- (302,153) -- (302,147) -- (301,146) -- (301,143) -- (302,142) -- (301,141) -- (301,133) -- (303,129) -- (300,126) -- (300,123) -- (301,122) -- (301,103) -- (302,102) -- (302,100) -- (300,96) -- (302,93) -- (302,86) -- (301,85) -- (301,81) -- cycle
(73,804) -- (75,802) -- (76,802) -- (86,809) -- (90,810) -- (100,819) -- (105,822) -- (107,822) -- (112,827) -- (121,832) -- (125,836) -- (129,838) -- (133,842) -- (134,842) -- (139,846) -- (141,846) -- (146,851) -- (148,851) -- (154,856) -- (156,856) -- (159,859) -- (163,861) -- (172,868) -- (174,868) -- (176,870) -- (179,871) -- (184,876) -- (191,880) -- (194,883) -- (197,884) -- (200,887) -- (203,888) -- (209,893) -- (211,893) -- (217,896) -- (223,902) -- (225,902) -- (228,905) -- (231,906) -- (234,909) -- (237,910) -- (241,914) -- (243,914) -- (247,918) -- (251,919) -- (257,925) -- (263,928) -- (268,932) -- (275,935) -- (279,940) -- (281,940) -- (288,945) -- (292,946) -- (295,949) -- (295,951) -- (294,952) -- (294,968) -- (295,969) -- (295,971) -- (294,972) -- (294,976) -- (296,978) -- (296,980) -- (295,981) -- (295,991) -- (296,993) -- (293,998) -- (293,1000) -- (295,1003) -- (295,1005) -- (294,1006) -- (294,1013) -- (295,1014) -- (295,1016) -- (293,1018) -- (291,1018) -- (289,1016) -- (280,1016) -- (279,1017) -- (277,1017) -- (276,1016) -- (269,1016) -- (268,1015) -- (262,1015) -- (261,1014) -- (256,1015) -- (252,1013) -- (247,1013) -- (246,1012) -- (241,1011) -- (238,1009) -- (227,1006) -- (224,1004) -- (221,1004) -- (219,1002) -- (215,1002) -- (213,1000) -- (208,1000) -- (205,997) -- (203,997) -- (200,994) -- (198,994) -- (194,991) -- (192,991) -- (189,988) -- (186,988) -- (184,986) -- (182,986) -- (180,984) -- (179,984) -- (176,981) -- (176,980) -- (172,979) -- (170,977) -- (168,977) -- (164,973) -- (162,973) -- (155,965) -- (154,965) -- (143,954) -- (142,954) -- (132,943) -- (132,941) -- (127,937) -- (126,937) -- (122,933) -- (121,929) -- (113,920) -- (110,913) -- (104,906) -- (103,902) -- (101,900) -- (101,898) -- (99,894) -- (97,892) -- (96,888) -- (94,886) -- (94,884) -- (90,876) -- (90,873) -- (87,869) -- (88,867) -- (84,861) -- (84,858) -- (82,855) -- (82,852) -- (80,849) -- (80,846) -- (79,845) -- (79,842) -- (78,841) -- (78,836) -- (77,835) -- (77,832) -- (76,831) -- (77,829) -- (77,826) -- (76,825) -- (76,820) -- (75,819) -- (76,817) -- (76,814) -- (73,809) -- cycle
(86,803) -- (88,801) -- (91,803) -- (93,803) -- (96,805) -- (99,805) -- (105,808) -- (108,808) -- (112,810) -- (119,811) -- (125,814) -- (129,814) -- (133,816) -- (137,816) -- (144,820) -- (150,821) -- (151,822) -- (156,823) -- (159,825) -- (167,826) -- (170,828) -- (179,829) -- (182,831) -- (188,832) -- (191,834) -- (193,834) -- (195,836) -- (200,836) -- (202,838) -- (205,838) -- (207,840) -- (223,842) -- (226,844) -- (229,844) -- (235,847) -- (239,847) -- (241,849) -- (245,849) -- (249,851) -- (254,851) -- (258,853) -- (261,853) -- (265,855) -- (268,855) -- (271,857) -- (276,857) -- (281,860) -- (284,860) -- (288,862) -- (293,862) -- (296,865) -- (296,867) -- (295,868) -- (295,874) -- (294,875) -- (294,882) -- (295,883) -- (295,893) -- (296,894) -- (296,899) -- (294,901) -- (295,903) -- (295,923) -- (296,924) -- (296,926) -- (295,927) -- (295,933) -- (297,935) -- (296,940) -- (295,941) -- (292,941) -- (290,940) -- (287,937) -- (281,934) -- (277,930) -- (275,930) -- (271,926) -- (268,925) -- (264,921) -- (246,910) -- (242,906) -- (228,897) -- (224,893) -- (221,893) -- (217,889) -- (211,887) -- (208,884) -- (205,883) -- (197,876) -- (194,875) -- (191,872) -- (183,868) -- (178,863) -- (173,860) -- (171,860) -- (162,853) -- (158,852) -- (150,845) -- (148,845) -- (146,843) -- (145,843) -- (134,833) -- (132,833) -- (130,831) -- (127,831) -- (108,817) -- (100,813) -- (97,810) -- (90,807) -- cycle
(96,800) -- (97,799) -- (106,799) -- (107,800) -- (121,800) -- (122,801) -- (124,801) -- (125,800) -- (132,800) -- (133,801) -- (135,800) -- (139,800) -- (140,801) -- (142,800) -- (144,800) -- (145,801) -- (171,801) -- (172,800) -- (174,801) -- (191,801) -- (192,800) -- (196,800) -- (197,801) -- (211,801) -- (212,802) -- (214,801) -- (225,801) -- (226,802) -- (228,801) -- (230,802) -- (232,802) -- (233,801) -- (241,801) -- (242,802) -- (256,802) -- (257,801) -- (263,801) -- (264,800) -- (276,800) -- (277,801) -- (280,801) -- (281,800) -- (283,800) -- (284,801) -- (286,801) -- (287,800) -- (292,800) -- (295,803) -- (295,807) -- (296,808) -- (296,823) -- (295,824) -- (295,827) -- (296,828) -- (295,830) -- (295,849) -- (297,852) -- (297,856) -- (295,858) -- (291,858) -- (288,856) -- (285,856) -- (284,855) -- (282,855) -- (281,854) -- (279,854) -- (278,853) -- (276,853) -- (269,850) -- (264,850) -- (261,848) -- (259,848) -- (258,847) -- (254,847) -- (253,846) -- (248,845) -- (245,843) -- (241,843) -- (238,841) -- (234,841) -- (231,839) -- (228,839) -- (226,837) -- (223,837) -- (220,835) -- (212,835) -- (209,833) -- (206,833) -- (204,831) -- (201,831) -- (200,830) -- (195,830) -- (189,827) -- (181,826) -- (173,822) -- (170,823) -- (167,821) -- (165,821) -- (161,819) -- (157,819) -- (156,818) -- (151,817) -- (148,815) -- (143,815) -- (141,813) -- (135,812) -- (131,810) -- (125,810) -- (124,809) -- (113,807) -- (110,805) -- (104,804) -- (102,802) -- (97,802) -- cycle
(519,799) -- (521,801) -- (522,808) -- (521,809) -- (521,815) -- (518,822) -- (518,824) -- (519,825) -- (519,829) -- (517,831) -- (517,833) -- (516,834) -- (517,838) -- (516,839) -- (515,848) -- (513,851) -- (512,860) -- (507,868) -- (507,872) -- (503,877) -- (503,880) -- (500,886) -- (499,892) -- (495,897) -- (494,901) -- (491,905) -- (490,908) -- (481,920) -- (481,922) -- (472,932) -- (471,936) -- (455,951) -- (454,954) -- (450,958) -- (448,958) -- (444,961) -- (442,961) -- (438,965) -- (434,967) -- (429,972) -- (429,973) -- (425,978) -- (423,978) -- (416,983) -- (410,985) -- (405,989) -- (401,991) -- (398,991) -- (396,993) -- (394,993) -- (388,996) -- (385,999) -- (380,1000) -- (378,1002) -- (372,1003) -- (367,1006) -- (364,1006) -- (360,1008) -- (352,1009) -- (351,1010) -- (347,1010) -- (343,1013) -- (336,1013) -- (335,1014) -- (332,1014) -- (331,1015) -- (328,1015) -- (327,1016) -- (311,1016) -- (310,1017) -- (307,1017) -- (306,1016) -- (301,1016) -- (300,1015) -- (301,1012) -- (299,1010) -- (299,1001) -- (301,998) -- (300,996) -- (300,984) -- (301,982) -- (300,981) -- (300,977) -- (299,976) -- (299,974) -- (300,973) -- (300,967) -- (299,966) -- (299,961) -- (300,960) -- (299,959) -- (299,953) -- (300,952) -- (300,946) -- (299,945) -- (300,943) -- (300,941) -- (299,940) -- (299,919) -- (300,918) -- (300,916) -- (298,914) -- (298,912) -- (300,910) -- (300,898) -- (301,897) -- (301,894) -- (300,893) -- (300,885) -- (299,883) -- (301,881) -- (301,876) -- (299,874) -- (299,868) -- (300,867) -- (300,865) -- (298,863) -- (300,858) -- (300,855) -- (299,854) -- (299,836) -- (300,835) -- (300,829) -- (299,828) -- (299,825) -- (300,824) -- (300,816) -- (299,815) -- (300,808) -- (299,807) -- (299,801) -- (301,799) -- (307,799) -- (308,798) -- (317,798) -- (318,799) -- (321,799) -- (323,798) -- (324,799) -- (330,799) -- (331,800) -- (336,800) -- (337,801) -- (341,801) -- (342,800) -- (356,800) -- (357,801) -- (383,801) -- (384,800) -- (395,800) -- (396,799) -- (400,799) -- (404,801) -- (412,801) -- (413,800) -- (416,800) -- (417,801) -- (426,801) -- (427,800) -- (440,801) -- (441,800) -- (447,800) -- (448,801) -- (450,801) -- (452,800) -- (456,802) -- (458,802) -- (459,801) -- (466,801) -- (467,802) -- (469,801) -- (485,801) -- (486,800) -- (490,800) -- (491,801) -- (510,801) -- (511,800) -- (517,800) -- cycle
(108,678) -- (109,679) -- (108,680) -- (107,688) -- (105,690) -- (105,692) -- (102,699) -- (102,702) -- (98,708) -- (98,712) -- (96,714) -- (94,726) -- (93,727) -- (93,729) -- (91,731) -- (91,733) -- (90,734) -- (90,738) -- (88,741) -- (88,744) -- (86,746) -- (86,750) -- (84,753) -- (84,756) -- (81,762) -- (81,765) -- (80,766) -- (81,769) -- (79,773) -- (77,775) -- (76,775) -- (75,774) -- (75,770) -- (76,769) -- (75,767) -- (75,764) -- (76,763) -- (76,757) -- (77,756) -- (78,747) -- (80,743) -- (80,739) -- (82,737) -- (83,731) -- (85,728) -- (85,722) -- (88,717) -- (88,714) -- (90,711) -- (90,709) -- (93,706) -- (94,702) -- (97,698) -- (97,695) -- (100,692) -- (101,689) -- (104,685) -- (104,681) -- (107,678) -- cycle
(294,574) -- (296,576) -- (296,584) -- (295,585) -- (295,594) -- (296,595) -- (296,598) -- (295,599) -- (295,604) -- (294,605) -- (295,607) -- (295,610) -- (294,611) -- (294,624) -- (293,625) -- (293,627) -- (294,628) -- (293,637) -- (295,639) -- (295,643) -- (294,644) -- (294,657) -- (293,658) -- (293,663) -- (294,664) -- (294,669) -- (293,670) -- (293,673) -- (294,674) -- (294,676) -- (293,677) -- (293,680) -- (294,681) -- (294,683) -- (293,684) -- (293,689) -- (294,690) -- (294,704) -- (293,705) -- (293,712) -- (294,713) -- (294,716) -- (293,717) -- (293,722) -- (295,724) -- (295,727) -- (294,728) -- (294,735) -- (295,736) -- (295,738) -- (294,739) -- (294,742) -- (293,744) -- (294,745) -- (294,752) -- (295,753) -- (294,755) -- (294,764) -- (295,765) -- (295,769) -- (294,770) -- (294,773) -- (295,774) -- (294,776) -- (294,791) -- (295,792) -- (295,794) -- (293,796) -- (289,796) -- (288,795) -- (275,795) -- (274,794) -- (257,794) -- (256,795) -- (243,795) -- (242,794) -- (214,794) -- (213,795) -- (192,795) -- (191,794) -- (189,795) -- (187,794) -- (183,794) -- (182,795) -- (179,795) -- (178,794) -- (172,794) -- (171,795) -- (148,795) -- (147,794) -- (129,794) -- (128,795) -- (101,795) -- (100,796) -- (93,794) -- (92,795) -- (87,795) -- (86,796) -- (77,796) -- (75,794) -- (75,792) -- (78,789) -- (78,787) -- (81,782) -- (81,780) -- (83,776) -- (83,773) -- (85,769) -- (86,762) -- (87,761) -- (88,756) -- (90,753) -- (91,748) -- (93,745) -- (93,743) -- (95,739) -- (95,734) -- (97,731) -- (97,728) -- (99,724) -- (99,721) -- (101,718) -- (101,715) -- (105,707) -- (106,700) -- (108,698) -- (108,694) -- (110,692) -- (110,686) -- (112,684) -- (112,681) -- (113,680) -- (113,677) -- (115,674) -- (115,670) -- (116,669) -- (116,667) -- (117,665) -- (126,656) -- (127,653) -- (131,649) -- (131,648) -- (134,645) -- (135,645) -- (140,640) -- (140,639) -- (153,627) -- (154,627) -- (161,621) -- (167,618) -- (172,613) -- (178,610) -- (181,607) -- (186,606) -- (188,604) -- (195,601) -- (198,598) -- (200,598) -- (201,597) -- (205,597) -- (208,594) -- (211,593) -- (214,590) -- (218,589) -- (220,587) -- (226,587) -- (227,586) -- (230,586) -- (237,583) -- (242,583) -- (247,580) -- (252,580) -- (254,578) -- (258,576) -- (272,576) -- (273,575) -- (280,575) -- (281,574) -- cycle
(301,573) -- (307,573) -- (308,574) -- (311,574) -- (312,575) -- (321,575) -- (322,576) -- (323,575) -- (324,576) -- (330,576) -- (331,577) -- (336,577) -- (337,578) -- (344,577) -- (347,580) -- (355,580) -- (357,582) -- (359,582) -- (362,584) -- (366,584) -- (368,586) -- (370,586) -- (372,588) -- (375,588) -- (377,590) -- (383,590) -- (387,593) -- (392,594) -- (394,596) -- (396,596) -- (400,599) -- (402,599) -- (411,606) -- (416,608) -- (421,613) -- (426,614) -- (437,624) -- (438,624) -- (446,632) -- (449,633) -- (452,636) -- (452,637) -- (461,646) -- (461,647) -- (471,657) -- (472,660) -- (483,674) -- (483,676) -- (488,681) -- (490,685) -- (490,687) -- (492,689) -- (492,691) -- (495,694) -- (497,698) -- (497,700) -- (499,702) -- (499,704) -- (501,707) -- (503,715) -- (508,721) -- (508,725) -- (509,726) -- (509,728) -- (512,732) -- (513,743) -- (514,744) -- (515,752) -- (517,756) -- (517,761) -- (519,763) -- (520,777) -- (521,778) -- (521,789) -- (522,790) -- (522,794) -- (520,796) -- (518,795) -- (513,795) -- (512,794) -- (507,797) -- (503,795) -- (498,795) -- (496,796) -- (495,795) -- (471,795) -- (470,794) -- (468,795) -- (464,795) -- (463,794) -- (433,794) -- (432,793) -- (427,793) -- (423,795) -- (421,794) -- (419,794) -- (418,795) -- (413,795) -- (412,794) -- (402,794) -- (401,793) -- (397,794) -- (396,795) -- (392,795) -- (391,794) -- (382,794) -- (381,795) -- (369,795) -- (368,794) -- (325,794) -- (324,795) -- (322,794) -- (307,794) -- (306,795) -- (301,795) -- (299,793) -- (300,787) -- (301,786) -- (301,783) -- (300,782) -- (300,776) -- (301,775) -- (301,765) -- (300,764) -- (300,758) -- (301,756) -- (300,755) -- (300,753) -- (301,752) -- (301,749) -- (302,748) -- (302,747) -- (300,745) -- (301,743) -- (301,737) -- (302,736) -- (302,734) -- (301,733) -- (301,728) -- (300,727) -- (301,724) -- (299,721) -- (301,718) -- (301,716) -- (300,715) -- (301,713) -- (300,711) -- (300,705) -- (302,702) -- (302,698) -- (301,697) -- (301,689) -- (302,688) -- (302,683) -- (300,681) -- (300,672) -- (301,671) -- (301,665) -- (300,664) -- (301,660) -- (300,659) -- (300,657) -- (301,656) -- (301,652) -- (302,651) -- (302,649) -- (300,645) -- (300,643) -- (301,642) -- (301,638) -- (300,637) -- (300,634) -- (299,633) -- (301,629) -- (301,627) -- (300,626) -- (300,620) -- (301,619) -- (300,610) -- (302,606) -- (302,604) -- (301,603) -- (301,597) -- (300,596) -- (300,587) -- (299,586) -- (299,582) -- (300,581) -- (299,575) -- cycle
(295,99) -- (297,101) -- (297,103) -- (296,104) -- (296,122) -- (295,123) -- (295,127) -- (297,131) -- (296,132) -- (296,144) -- (295,145) -- (295,185) -- (294,186) -- (295,188) -- (295,193) -- (296,194) -- (296,196) -- (295,197) -- (295,203) -- (294,204) -- (295,215) -- (296,216) -- (296,219) -- (295,220) -- (295,228) -- (294,229) -- (294,233) -- (296,236) -- (296,239) -- (294,243) -- (294,246) -- (296,248) -- (296,249) -- (294,251) -- (294,255) -- (295,256) -- (295,260) -- (296,261) -- (295,263) -- (296,267) -- (295,268) -- (295,276) -- (294,277) -- (294,283) -- (295,284) -- (295,286) -- (294,287) -- (294,301) -- (295,302) -- (295,305) -- (294,306) -- (295,308) -- (295,310) -- (294,311) -- (294,313) -- (295,314) -- (294,322) -- (295,323) -- (295,325) -- (293,328) -- (293,334) -- (294,335) -- (293,339) -- (295,342) -- (294,343) -- (294,354) -- (295,356) -- (294,357) -- (294,362) -- (293,363) -- (294,365) -- (293,367) -- (293,370) -- (294,371) -- (293,381) -- (295,383) -- (295,385) -- (294,386) -- (294,395) -- (295,396) -- (295,399) -- (294,400) -- (294,420) -- (293,421) -- (293,425) -- (294,426) -- (294,431) -- (295,432) -- (295,435) -- (294,436) -- (293,441) -- (294,442) -- (294,445) -- (295,446) -- (294,448) -- (294,456) -- (295,457) -- (294,466) -- (296,468) -- (296,471) -- (295,472) -- (295,478) -- (294,479) -- (294,482) -- (295,483) -- (295,485) -- (294,486) -- (294,510) -- (295,511) -- (295,517) -- (296,518) -- (295,520) -- (295,524) -- (294,525) -- (294,529) -- (295,530) -- (295,533) -- (294,534) -- (295,536) -- (294,538) -- (294,547) -- (295,548) -- (295,552) -- (296,553) -- (295,561) -- (298,566) -- (294,570) -- (272,570) -- (271,571) -- (265,570) -- (264,571) -- (259,571) -- (258,572) -- (255,572) -- (252,574) -- (249,574) -- (248,575) -- (240,575) -- (237,577) -- (235,577) -- (228,580) -- (225,580) -- (224,581) -- (213,584) -- (209,586) -- (207,588) -- (204,588) -- (199,592) -- (192,594) -- (190,596) -- (186,597) -- (180,602) -- (178,602) -- (176,604) -- (171,606) -- (163,613) -- (156,617) -- (152,621) -- (149,622) -- (139,632) -- (139,633) -- (133,639) -- (133,640) -- (128,645) -- (128,646) -- (126,648) -- (124,648) -- (122,646) -- (124,642) -- (124,638) -- (126,635) -- (126,632) -- (130,624) -- (130,621) -- (133,615) -- (133,611) -- (135,609) -- (136,604) -- (138,601) -- (138,595) -- (139,594) -- (140,589) -- (142,586) -- (142,582) -- (144,580) -- (145,575) -- (147,572) -- (147,568) -- (149,566) -- (150,564) -- (150,561) -- (152,558) -- (152,551) -- (153,550) -- (153,548) -- (155,546) -- (155,542) -- (157,540) -- (158,534) -- (161,528) -- (161,525) -- (162,524) -- (163,519) -- (166,515) -- (166,508) -- (167,507) -- (167,504) -- (169,501) -- (169,499) -- (172,493) -- (172,488) -- (176,483) -- (176,478) -- (179,473) -- (179,467) -- (181,463) -- (181,459) -- (183,457) -- (183,454) -- (185,451) -- (185,448) -- (188,444) -- (188,441) -- (190,439) -- (190,437) -- (192,433) -- (192,429) -- (194,425) -- (195,416) -- (198,412) -- (198,409) -- (200,407) -- (200,401) -- (204,393) -- (204,389) -- (205,388) -- (205,386) -- (207,384) -- (208,376) -- (211,371) -- (211,365) -- (213,363) -- (213,360) -- (216,355) -- (216,351) -- (219,347) -- (219,345) -- (220,344) -- (219,340) -- (222,333) -- (222,329) -- (224,327) -- (225,321) -- (227,318) -- (227,316) -- (230,313) -- (229,311) -- (229,308) -- (231,306) -- (233,302) -- (235,291) -- (238,285) -- (238,282) -- (240,278) -- (241,272) -- (243,269) -- (243,266) -- (246,263) -- (246,259) -- (247,258) -- (246,252) -- (249,247) -- (249,242) -- (252,239) -- (252,234) -- (254,232) -- (257,225) -- (257,221) -- (259,218) -- (259,212) -- (261,209) -- (261,205) -- (264,200) -- (264,196) -- (265,195) -- (266,189) -- (271,180) -- (271,177) -- (273,173) -- (273,168) -- (275,166) -- (276,156) -- (278,153) -- (278,151) -- (280,149) -- (280,146) -- (281,145) -- (281,143) -- (284,136) -- (284,132) -- (285,131) -- (285,124) -- (288,118) -- (288,116) -- (290,113) -- (290,107) -- (292,104) -- (292,102) -- cycle;
\node[inner sep=0pt,font=\fontsize{10}{10}\selectfont] at (299,51) {\scalebox{4.9}{$D$}};
\node[inner sep=0pt,font=\fontsize{10}{10}\selectfont] at (36,795) {\scalebox{4.6}{$A$}};
\node[inner sep=0pt,font=\fontsize{10}{10}\selectfont] at (546,798) {\scalebox{4.5}{$B$}};
\node[inner sep=0pt,font=\fontsize{10}{10}\selectfont] at (319,821) {\scalebox{4.3}{$C$}};
\node[inner sep=0pt,font=\fontsize{10}{10}\selectfont] at (318,865) {\scalebox{4.3}{$P$}};
\node[inner sep=0pt,font=\fontsize{10}{10}\selectfont] at (321,938) {\scalebox{4.5}{$Q$}};
\end{tikzpicture}%
}
\end{center}
\indent $P$ und $Q$ Darstellung eines und desselben Punkts der Kugel in den beiden Projectionen; $ACB$ gegen $CPQ$ normal, Winkel $CAP = PAQ$.

$AD$ normal gegen $PA$ und $D$ auf der verlängerten $QPC$, wird $D$ die stereographische Darstellung des jenem Kugelpunkte gegenüberliegenden.

Ein Kreis um den Mittelpunkt $Q$ mit dem Radius $QA$ ist die stereographische Darstellung desjenigen grössten Kreises, der jenen Kugelpunkt zum Pole hat.

Ist $a+bi$ die complexe Zahl für centrale Projection, $r$ Halbmesser, so ist die complexe Zahl für die stereographische Projection

\[
\frac{(a+ib)r}{aa+bb}\left\{\BareRoot{(aa+bb+rr)}-r\right\}.
\]


%
\SourceLeaf{360}

\Needspace{17\baselineskip}
\begin{center}
[2.]
\end{center}

\begin{center}
\bfseries Centrale Projection.
\end{center}

Sind $a+bi$, $a'+b'i$ die complexen Zahlen für die Projectionen zweier Punkte, so ist von ihrer Bogendistanz

\[
\text{cosinus} = \frac{rr+aa'+bb'}{\Delta}
\]
\[
\text{sinus} = \BareRoot{\left\{[(a'-a)^2+(b'-b)^2]rr+(ab'-ba')^2\right\}}:\Delta,
\]
\[
\Delta = \BareRoot{\left\{(rr+aa+bb)(rr+a'a'+b'b')\right\}}.
\]

\begin{center}
[III.]
\end{center}

\begin{center}
\bfseries LEHRSATZ.
\end{center}

\begin{center}
\bfseries [ÜBER DIE STEREOGRAPHISCHE PROJECTION.]
\end{center}

Die wirkliche (geradlinige) Distanz zweier Punkte auf der Kugelfläche findet sich aus der Distanz ihrer stereographischen Projectionen, wenn man letztere mit dem Product der beiden Cosinus derjenigen Neigungen multiplicirt, welche die von jenen Projectionen nach dem Endpunkt einer im Centrum errichteten Normale aequal dem Durchmesser der Kugel (Halbmesser des Äquators in der Projection) gezogenen Geraden gegen die Normale haben.

\[
\text{Wirkliche Distanz} = PP'\cdot\frac{CN}{PN}\cdot\frac{CN}{P'N}.
\]

\begin{center}\resizebox{3.7cm}{!}{%
\begin{tikzpicture}[x=1pt,y=-1pt]
\path[use as bounding box] (0,0) rectangle (478,755);
\path[fill=black,draw=none,even odd rule]
(146,693) -- (149,696) -- (160,696) -- (162,695) -- (165,692) -- (165,691) -- (164,690) -- (151,690) -- (150,691) -- (147,691) -- cycle
(116,692) -- (116,693) -- (118,695) -- (123,695) -- (127,697) -- (129,696) -- (134,696) -- (137,694) -- (137,692) -- (134,690) -- (126,690) -- (125,691) -- (123,690) -- (120,690) -- cycle
(199,692) -- (203,695) -- (207,695) -- (208,694) -- (210,695) -- (213,695) -- (214,694) -- (219,694) -- (222,691) -- (222,690) -- (221,689) -- (201,689) -- (199,691) -- cycle
(347,690) -- (347,691) -- (349,693) -- (367,693) -- (368,692) -- (368,689) -- (367,688) -- (349,688) -- cycle
(321,692) -- (323,694) -- (334,694) -- (335,693) -- (337,693) -- (340,690) -- (340,689) -- (339,688) -- (327,688) -- (326,689) -- (322,689) -- (321,690) -- cycle
(294,689) -- (294,691) -- (296,693) -- (304,693) -- (305,694) -- (307,694) -- (308,693) -- (310,693) -- (312,691) -- (312,690) -- (309,688) -- (295,688) -- cycle
(231,690) -- (231,691) -- (234,693) -- (241,693) -- (242,694) -- (244,694) -- (245,693) -- (249,693) -- (250,692) -- (250,690) -- (248,688) -- (235,688) -- cycle
(175,691) -- (175,693) -- (177,695) -- (191,695) -- (192,694) -- (192,691) -- (189,689) -- (186,689) -- (185,688) -- (178,689) -- cycle
(284,689) -- (281,687) -- (279,687) -- (277,689) -- (272,689) -- (271,688) -- (268,688) -- (267,689) -- (264,689) -- (262,691) -- (262,692) -- (263,693) -- (268,693) -- (269,694) -- (277,694) -- (278,693) -- (282,693) -- (284,691) -- cycle
(369,645) -- (368,646) -- (368,649) -- (373,654) -- (373,656) -- (378,661) -- (381,661) -- (382,660) -- (382,658) -- (378,654) -- (375,649) -- (371,645) -- cycle
(117,645) -- (116,645) -- (112,650) -- (111,654) -- (107,659) -- (107,663) -- (109,663) -- (113,659) -- (115,654) -- (118,651) -- (118,646) -- cycle
(350,621) -- (349,622) -- (349,624) -- (352,628) -- (353,631) -- (360,637) -- (362,637) -- (363,636) -- (363,634) -- (360,631) -- (358,627) -- (354,623) -- cycle
(133,613) -- (131,613) -- (128,616) -- (127,623) -- (123,628) -- (123,630) -- (124,631) -- (126,631) -- (130,627) -- (131,623) -- (134,618) -- (134,614) -- cycle
(330,598) -- (329,600) -- (336,608) -- (336,609) -- (340,613) -- (343,613) -- (344,612) -- (344,610) -- (340,606) -- (337,601) -- (334,598) -- cycle
(149,582) -- (146,583) -- (144,585) -- (144,588) -- (143,590) -- (140,593) -- (139,597) -- (135,601) -- (136,603) -- (139,604) -- (143,600) -- (147,592) -- (147,589) -- (150,585) -- (150,583) -- cycle
(314,562) -- (312,561) -- (310,563) -- (310,565) -- (309,566) -- (309,574) -- (301,578) -- (299,582) -- (292,582) -- (290,584) -- (288,584) -- (286,586) -- (282,588) -- (280,588) -- (270,595) -- (266,597) -- (264,597) -- (261,600) -- (258,600) -- (255,603) -- (251,603) -- (249,604) -- (246,607) -- (235,612) -- (233,614) -- (229,616) -- (227,616) -- (224,619) -- (222,619) -- (215,625) -- (212,625) -- (210,627) -- (204,629) -- (201,632) -- (198,632) -- (193,636) -- (184,640) -- (181,643) -- (179,643) -- (176,645) -- (174,645) -- (168,648) -- (157,656) -- (155,656) -- (151,659) -- (148,660) -- (146,662) -- (127,671) -- (125,673) -- (122,674) -- (115,679) -- (113,679) -- (111,681) -- (108,682) -- (104,686) -- (101,686) -- (100,685) -- (104,677) -- (103,675) -- (101,675) -- (100,679) -- (96,683) -- (96,685) -- (94,687) -- (93,690) -- (90,693) -- (90,695) -- (93,698) -- (95,697) -- (104,697) -- (108,694) -- (108,692) -- (105,690) -- (105,689) -- (107,687) -- (109,687) -- (114,683) -- (118,683) -- (121,680) -- (130,676) -- (133,673) -- (135,673) -- (140,670) -- (142,670) -- (148,666) -- (151,666) -- (154,663) -- (156,663) -- (160,661) -- (165,656) -- (167,656) -- (171,653) -- (177,651) -- (179,649) -- (182,648) -- (184,646) -- (186,646) -- (188,644) -- (191,644) -- (193,643) -- (196,640) -- (201,638) -- (203,636) -- (209,633) -- (211,633) -- (215,631) -- (220,627) -- (224,626) -- (227,623) -- (231,622) -- (238,616) -- (241,616) -- (247,611) -- (249,611) -- (251,609) -- (256,608) -- (261,604) -- (263,604) -- (267,601) -- (271,601) -- (281,593) -- (288,591) -- (294,586) -- (297,586) -- (307,580) -- (312,580) -- (321,590) -- (324,590) -- (325,589) -- (325,587) -- (319,581) -- (318,578) -- (313,573) -- (313,564) -- cycle
(163,554) -- (160,554) -- (158,556) -- (156,564) -- (152,568) -- (152,571) -- (154,573) -- (155,573) -- (158,570) -- (159,567) -- (161,565) -- (162,559) -- (163,558) -- cycle
(318,531) -- (315,531) -- (314,532) -- (314,537) -- (313,538) -- (314,542) -- (311,548) -- (311,550) -- (312,551) -- (314,551) -- (317,547) -- (317,540) -- (319,536) -- (319,532) -- cycle
(180,519) -- (178,519) -- (176,521) -- (171,531) -- (166,537) -- (166,539) -- (170,540) -- (174,537) -- (174,535) -- (177,531) -- (177,528) -- (181,523) -- (181,520) -- cycle
(322,500) -- (320,501) -- (320,506) -- (319,507) -- (318,511) -- (316,513) -- (316,522) -- (317,523) -- (320,523) -- (322,520) -- (323,512) -- (324,511) -- (325,503) -- cycle
(192,493) -- (190,493) -- (190,494) -- (186,499) -- (183,506) -- (181,508) -- (182,510) -- (186,510) -- (189,507) -- (190,505) -- (190,501) -- (192,499) -- (192,497) -- (193,496) -- (193,494) -- cycle
(327,472) -- (324,472) -- (323,473) -- (323,476) -- (324,478) -- (321,483) -- (321,488) -- (323,490) -- (325,490) -- (327,488) -- (327,484) -- (329,480) -- (329,475) -- cycle
(209,460) -- (207,460) -- (204,463) -- (201,471) -- (197,476) -- (197,478) -- (199,480) -- (200,480) -- (203,477) -- (203,476) -- (207,471) -- (207,469) -- (210,465) -- (210,461) -- cycle
(334,446) -- (331,446) -- (329,448) -- (329,453) -- (326,458) -- (326,460) -- (328,462) -- (331,462) -- (332,461) -- (332,459) -- (333,458) -- (333,452) -- (334,451) -- cycle
(223,429) -- (222,429) -- (218,433) -- (218,435) -- (216,437) -- (214,441) -- (213,445) -- (215,449) -- (216,449) -- (218,447) -- (219,443) -- (222,438) -- (222,436) -- (224,434) -- (224,430) -- cycle
(337,414) -- (335,414) -- (334,415) -- (334,418) -- (333,419) -- (333,421) -- (332,422) -- (332,426) -- (331,427) -- (331,430) -- (332,431) -- (335,431) -- (337,429) -- (337,424) -- (339,420) -- (339,416) -- cycle
(240,397) -- (238,397) -- (236,399) -- (236,400) -- (232,405) -- (232,408) -- (227,414) -- (227,416) -- (229,418) -- (231,418) -- (235,414) -- (235,412) -- (238,409) -- (238,405) -- (240,403) -- (240,401) -- (241,400) -- (241,398) -- cycle
(342,384) -- (340,384) -- (339,385) -- (340,388) -- (337,395) -- (337,399) -- (339,401) -- (341,401) -- (342,400) -- (342,397) -- (343,396) -- (343,393) -- (344,392) -- (344,387) -- cycle
(257,361) -- (255,361) -- (253,363) -- (253,364) -- (251,366) -- (251,368) -- (244,378) -- (244,381) -- (246,383) -- (248,383) -- (250,379) -- (253,376) -- (254,372) -- (258,366) -- (258,362) -- cycle
(348,351) -- (347,351) -- (344,355) -- (344,361) -- (343,362) -- (343,364) -- (342,365) -- (342,369) -- (341,370) -- (341,371) -- (343,373) -- (345,373) -- (349,366) -- (348,364) -- (348,361) -- (349,360) -- (349,352) -- cycle
(273,331) -- (271,331) -- (268,334) -- (265,341) -- (263,343) -- (261,347) -- (261,350) -- (262,351) -- (264,351) -- (266,349) -- (266,347) -- (269,344) -- (272,337) -- (274,335) -- (274,332) -- cycle
(353,319) -- (352,319) -- (350,321) -- (350,327) -- (349,328) -- (349,334) -- (348,335) -- (348,338) -- (349,339) -- (352,339) -- (354,337) -- (354,327) -- (355,326) -- (355,321) -- cycle
(287,301) -- (285,301) -- (282,304) -- (280,308) -- (280,310) -- (275,316) -- (275,317) -- (277,319) -- (280,319) -- (282,317) -- (283,313) -- (288,305) -- (288,302) -- cycle
(359,287) -- (358,287) -- (356,289) -- (356,294) -- (355,295) -- (355,299) -- (354,300) -- (354,303) -- (353,304) -- (353,309) -- (354,310) -- (357,310) -- (359,306) -- (359,302) -- (360,301) -- (360,297) -- (361,296) -- (361,289) -- cycle
(302,269) -- (301,269) -- (299,271) -- (299,275) -- (296,278) -- (290,288) -- (290,289) -- (292,291) -- (296,289) -- (296,286) -- (301,281) -- (301,276) -- (303,274) -- (303,270) -- cycle
(365,259) -- (363,259) -- (361,262) -- (361,266) -- (358,271) -- (359,275) -- (360,276) -- (364,276) -- (365,275) -- (365,267) -- (366,266) -- (366,260) -- cycle
(320,234) -- (318,234) -- (314,240) -- (314,243) -- (310,247) -- (310,250) -- (307,253) -- (307,255) -- (309,257) -- (310,257) -- (315,251) -- (316,248) -- (318,246) -- (318,244) -- (320,242) -- (320,240) -- (322,237) -- (322,236) -- cycle
(370,228) -- (368,228) -- (367,229) -- (366,238) -- (364,241) -- (364,244) -- (366,246) -- (367,246) -- (369,244) -- (369,242) -- (371,238) -- (371,229) -- cycle
(336,202) -- (334,202) -- (332,204) -- (330,209) -- (328,211) -- (327,213) -- (327,216) -- (324,220) -- (324,223) -- (325,224) -- (328,224) -- (330,222) -- (333,216) -- (333,212) -- (336,208) -- (336,206) -- (337,205) -- (337,203) -- cycle
(376,196) -- (374,197) -- (372,201) -- (372,205) -- (371,206) -- (371,209) -- (369,213) -- (371,215) -- (373,215) -- (374,214) -- (374,212) -- (376,209) -- (376,202) -- (377,201) -- (377,197) -- cycle
(352,171) -- (351,171) -- (347,175) -- (347,177) -- (345,179) -- (343,184) -- (340,187) -- (340,189) -- (341,190) -- (344,190) -- (346,186) -- (350,181) -- (353,174) -- (353,172) -- cycle
(382,163) -- (380,164) -- (380,167) -- (379,168) -- (379,171) -- (378,172) -- (378,176) -- (376,178) -- (376,180) -- (378,182) -- (379,182) -- (382,179) -- (382,172) -- (383,171) -- (383,167) -- (384,166) -- (384,165) -- cycle
(367,142) -- (365,142) -- (363,144) -- (360,151) -- (358,153) -- (356,157) -- (356,160) -- (357,161) -- (359,161) -- (361,159) -- (363,152) -- (367,147) -- (367,145) -- (368,144) -- cycle
(389,134) -- (387,134) -- (385,136) -- (385,139) -- (382,146) -- (383,148) -- (385,148) -- (387,145) -- (387,138) -- cycle
(382,113) -- (380,113) -- (377,116) -- (377,118) -- (374,124) -- (371,127) -- (371,130) -- (372,131) -- (374,131) -- (377,128) -- (377,125) -- (383,115) -- cycle
(399,79) -- (398,79) -- (394,83) -- (391,91) -- (388,94) -- (386,98) -- (386,100) -- (388,101) -- (389,100) -- (389,98) -- (391,96) -- (392,96) -- (398,101) -- (396,103) -- (396,106) -- (394,108) -- (393,107) -- (391,108) -- (390,110) -- (390,115) -- (389,116) -- (389,119) -- (391,120) -- (393,118) -- (395,120) -- (395,137) -- (396,138) -- (396,147) -- (395,148) -- (395,150) -- (396,151) -- (396,157) -- (397,158) -- (397,162) -- (396,163) -- (396,170) -- (397,171) -- (397,173) -- (396,174) -- (396,176) -- (394,178) -- (395,184) -- (396,185) -- (396,187) -- (395,188) -- (395,198) -- (396,199) -- (396,202) -- (395,203) -- (395,208) -- (396,209) -- (396,211) -- (395,213) -- (397,216) -- (395,218) -- (395,221) -- (397,223) -- (397,225) -- (396,226) -- (396,235) -- (395,236) -- (395,240) -- (396,241) -- (396,252) -- (395,253) -- (395,258) -- (397,260) -- (397,261) -- (396,262) -- (395,271) -- (397,274) -- (397,277) -- (398,278) -- (398,280) -- (397,281) -- (397,291) -- (396,292) -- (397,300) -- (395,302) -- (395,304) -- (396,305) -- (396,311) -- (397,313) -- (396,314) -- (395,321) -- (396,322) -- (396,326) -- (398,330) -- (398,335) -- (397,336) -- (397,343) -- (398,344) -- (397,350) -- (398,351) -- (398,353) -- (397,354) -- (397,362) -- (396,363) -- (396,365) -- (397,366) -- (397,373) -- (398,374) -- (398,378) -- (397,379) -- (397,384) -- (398,385) -- (398,388) -- (397,389) -- (397,391) -- (399,394) -- (398,400) -- (397,401) -- (398,402) -- (398,408) -- (399,409) -- (398,411) -- (398,429) -- (397,430) -- (397,432) -- (398,433) -- (398,438) -- (399,439) -- (399,479) -- (398,480) -- (398,483) -- (399,484) -- (399,488) -- (398,489) -- (398,491) -- (399,492) -- (399,498) -- (398,499) -- (401,502) -- (399,506) -- (400,508) -- (400,521) -- (399,522) -- (399,530) -- (400,531) -- (400,539) -- (399,540) -- (399,543) -- (400,544) -- (400,547) -- (399,548) -- (399,550) -- (400,551) -- (400,559) -- (399,561) -- (401,563) -- (401,567) -- (400,568) -- (401,570) -- (401,574) -- (400,575) -- (400,578) -- (401,579) -- (401,594) -- (402,595) -- (401,597) -- (401,605) -- (402,606) -- (402,608) -- (401,609) -- (401,614) -- (403,616) -- (403,619) -- (402,620) -- (402,634) -- (403,635) -- (403,646) -- (402,647) -- (403,650) -- (402,651) -- (402,653) -- (403,654) -- (403,658) -- (402,659) -- (402,661) -- (403,662) -- (403,666) -- (404,667) -- (403,669) -- (403,672) -- (404,673) -- (404,680) -- (403,681) -- (401,681) -- (389,668) -- (386,668) -- (385,669) -- (386,672) -- (395,679) -- (395,681) -- (400,685) -- (400,687) -- (398,689) -- (397,689) -- (395,687) -- (392,687) -- (391,688) -- (386,688) -- (385,689) -- (383,688) -- (375,688) -- (373,690) -- (373,691) -- (375,693) -- (390,693) -- (395,690) -- (397,690) -- (404,693) -- (407,693) -- (410,690) -- (410,688) -- (409,687) -- (409,681) -- (408,680) -- (409,676) -- (410,675) -- (410,673) -- (408,671) -- (408,669) -- (409,668) -- (409,663) -- (408,662) -- (408,659) -- (410,656) -- (409,654) -- (409,647) -- (408,646) -- (408,640) -- (409,639) -- (409,628) -- (408,627) -- (408,624) -- (410,622) -- (410,618) -- (407,615) -- (409,611) -- (407,608) -- (407,606) -- (408,605) -- (408,602) -- (407,601) -- (407,596) -- (408,595) -- (408,591) -- (407,590) -- (407,588) -- (408,587) -- (408,581) -- (406,579) -- (406,578) -- (408,576) -- (408,572) -- (407,571) -- (407,569) -- (408,568) -- (409,564) -- (408,563) -- (408,559) -- (406,557) -- (406,555) -- (407,554) -- (407,548) -- (406,547) -- (406,543) -- (407,542) -- (407,534) -- (406,533) -- (406,529) -- (407,528) -- (407,526) -- (405,523) -- (405,521) -- (407,517) -- (407,511) -- (406,510) -- (407,505) -- (406,504) -- (406,501) -- (405,500) -- (406,498) -- (406,489) -- (405,488) -- (405,484) -- (406,483) -- (406,475) -- (405,474) -- (405,468) -- (406,467) -- (406,454) -- (407,453) -- (407,451) -- (406,450) -- (406,444) -- (405,443) -- (406,437) -- (404,433) -- (404,430) -- (405,429) -- (404,424) -- (406,421) -- (405,419) -- (405,402) -- (404,401) -- (404,393) -- (405,392) -- (405,390) -- (404,389) -- (404,387) -- (405,386) -- (405,384) -- (404,383) -- (404,376) -- (405,375) -- (404,370) -- (405,369) -- (405,365) -- (406,364) -- (406,362) -- (404,359) -- (404,348) -- (403,347) -- (403,341) -- (404,340) -- (404,330) -- (402,327) -- (402,324) -- (403,323) -- (403,315) -- (402,314) -- (402,312) -- (403,311) -- (403,306) -- (402,305) -- (402,300) -- (404,297) -- (404,294) -- (402,292) -- (403,273) -- (401,271) -- (401,269) -- (402,268) -- (402,263) -- (400,261) -- (401,259) -- (401,256) -- (402,255) -- (402,252) -- (401,251) -- (402,249) -- (402,244) -- (401,243) -- (402,241) -- (402,224) -- (401,223) -- (402,221) -- (401,219) -- (402,217) -- (402,215) -- (401,214) -- (401,204) -- (402,203) -- (401,201) -- (402,199) -- (401,197) -- (401,191) -- (402,190) -- (401,188) -- (401,183) -- (402,182) -- (402,179) -- (400,175) -- (400,170) -- (403,167) -- (401,163) -- (402,157) -- (400,153) -- (401,151) -- (400,149) -- (400,139) -- (401,138) -- (400,127) -- (401,126) -- (401,124) -- (404,121) -- (405,117) -- (401,114) -- (403,112) -- (403,110) -- (400,107) -- (400,100) -- (399,99) -- (399,97) -- (401,95) -- (404,98) -- (405,98) -- (406,97) -- (406,92) -- (400,86) -- (400,83) -- (401,81) -- cycle;
\node[inner sep=0pt,font=\fontsize{10}{10}\selectfont] at (397,50) {\scalebox{5.8}{$N$}};
\node[inner sep=0pt,font=\fontsize{10}{10}\selectfont] at (277,555) {\scalebox{5.8}{$P'$}};
\node[inner sep=0pt,font=\fontsize{10}{10}\selectfont] at (45,702) {\scalebox{6.0}{$P$}};
\node[inner sep=0pt,font=\fontsize{10}{10}\selectfont] at (436,706) {\scalebox{5.8}{$C$}};
\end{tikzpicture}%
}
\end{center}
\indent Man erhält so die Chorde oder eine dem Sinus des halben Bogens proportionale Grösse. Besser ist's, entweder die Projectionen durch bicomplexe Grössen $t$, $t'$ auszudrücken, wo die Distanz durch [die] Tangente des halben Bogens, gemessen durch

\[
\frac{t'-t}{1+tt'},
\]

\noindent dargestellt wird; oder sich auch der tricomplexen Grössen dabei zu bedienen.


%
\SourceLeaf{361}

\Needspace{13\baselineskip}
\begin{center}
[IV.]
\end{center}

\begin{center}
\bfseries STEREOGRAPHISCHE PROJECTION DER KUGELFLÄCHE.
\end{center}

Conjugirte complexe Zahlen sollen hier immer durch kleine und grosse gleichnamige Buchstaben bezeichnet werden.

Halbmesser der Kugel wird $=1$ gesetzt.

$n$: lineares Vergrösserungsverhältniss.

1. In dem durch die complexe Zahl $a$ bezeichneten Punkte ist
\[
n = 1+aA.
\]

2. Der gegenüberliegende Punkt wird durch die complexe Zahl $-\dfrac{1}{A}$ bezeichnet; in demselben ist also
\[
n = 1+\frac{1}{aA} = \frac{1+aA}{aA}.
\]

3. Der bicomplexen Zahl $a$ entspricht die tricomplexe
\[
\frac{2a+(1-aA)\eta}{1+aA}.
\]

4. Allgemeine Verwandlungsformel:
\[
t' = \frac{at-b}{ct+d},
\]

wo $t$, $t'$ sich auf Einerlei Punkt in der ersten oder zweiten Darstellungsform beziehen.

Hier wird
\[
n' = \text{Mod.}\ \frac{ad+bc}{(ct+d)^2}(1+tT).
\]

\Needspace{11\baselineskip}Der Minimumwerth findet statt für
\[
t = \frac{C}{D}
\]

\noindent und ist
\[
\text{Mod.}\ \frac{ad+bc}{cC+dD}\cdot\frac{D}{d} = \text{Mod.}\ \frac{ad+bc}{cC+dD}.
\]
%
\SourceLeaf{362}
\noindent Soll dieser für $t'=0$ gelten, so muss
\[
aC=bD
\]
sein; wodurch
\[
n'=\text{Mod.}\,\frac{a}{D}=\text{Mod.}\,\frac{b}{C}
\]
wird. Soll dieser zugleich $=1$ werden, so kann man, wenn $\varepsilon$ eine beliebige
Grösse, deren Modulus $=1$, ist, setzen
\[
\varepsilon{}C=b,\qquad \varepsilon{}D=a
\]
und
\[
c=\varepsilon{}B,\qquad d=\varepsilon{}A.
\]
Schreibt man anstatt $a,b,c,d$: $\sqrt{\varepsilon}.a$,
$\sqrt{\varepsilon}.b$, $\sqrt{\varepsilon}.c$, $\sqrt{\varepsilon}.d$,
wodurch die Allgemeinheit nicht vermindert wird, so wird schlechthin
\[
c=B,\qquad d=A
\]
oder die Transformationsformel
\[
t'=\frac{at-b}{Bt+A},\qquad
t=\frac{At'+b}{-Bt'+a}.
\]

5. Die einfachste Transformationsformel ist
\[
t'=\frac{t-a}{At+1}.
\]

6. Zwischen zweien Punkten $a,b$ liegt in der Mitte der Punkt
\[
\frac{\BareRoot{(1+aB)(1+bB)}-\BareRoot{(1+aA)(1+bA)}}
{-A\BareRoot{(1+aB)(1+bB)}+B\BareRoot{(1+aA)(1+bA)}}
\]
oder, wenn die den $a,b$ entgegengesetzten Punkte $-\dfrac{1}{A},-\dfrac{1}{B}$
mit $a',b'$ bezeichnet werden,
\[
\frac{a'\BareRoot{(a-b')(b-b')}-b'\BareRoot{(a-a')(b-a')}}
{\BareRoot{(a-b')(b-b')}-\BareRoot{(a-a')(b-a')}}.
\]
%
\SourceLeaf{363}
\clearpage
\begin{center}[V.]\end{center}
\begin{center}\bfseries ZUR THEORIE DER DREHUNG DER KUGELFLÄCHE IN SICH SELBST.\end{center}

Die Aufgabe, zwei Drehungen in Eine zu vereinigen, kann auf doppelte Weise aufgelöset werden.

I. Man soll die beiden Drehungen, die erste $a$ um den Punkt $A$ und die zweite $\beta$ um den Punkt $B$, in Eine vereinigen.

\begin{center}
\begin{minipage}{4.4cm}\resizebox{4.4cm}{!}{%
\begin{tikzpicture}[x=1pt,y=-1pt]
\path[use as bounding box] (0,0) rectangle (734,539);
\path[fill=black,draw=none,even odd rule]
(51,106) -- (53,108) -- (71,108) -- (74,105) -- (74,103) -- (56,103) -- (55,104) -- (52,104) -- cycle
(83,102) -- (83,104) -- (87,106) -- (89,106) -- (90,105) -- (100,105) -- (103,103) -- (103,101) -- (101,99) -- (88,99) -- (87,100) -- (85,100) -- cycle
(134,99) -- (132,97) -- (127,96) -- (124,98) -- (116,98) -- (115,99) -- (113,99) -- (112,100) -- (112,102) -- (117,105) -- (126,105) -- (127,104) -- (130,104) -- (132,103) -- (134,101) -- cycle
(149,98) -- (149,99) -- (152,102) -- (155,102) -- (156,103) -- (159,103) -- (160,102) -- (166,102) -- (170,99) -- (169,97) -- (167,97) -- (164,95) -- (153,95) -- cycle
(187,95) -- (187,97) -- (188,98) -- (191,98) -- (192,99) -- (204,99) -- (207,97) -- (207,95) -- (205,93) -- (191,93) -- cycle
(224,93) -- (225,95) -- (227,95) -- (231,97) -- (233,96) -- (243,96) -- (245,94) -- (245,92) -- (242,90) -- (232,90) -- (231,91) -- (227,91) -- cycle
(260,93) -- (261,94) -- (266,94) -- (267,95) -- (269,95) -- (270,94) -- (273,94) -- (274,93) -- (281,93) -- (284,91) -- (284,89) -- (281,87) -- (271,87) -- (268,89) -- (264,89) -- (262,90) -- (260,92) -- cycle
(298,88) -- (298,90) -- (300,92) -- (312,92) -- (317,89) -- (317,88) -- (314,86) -- (301,86) -- cycle
(666,64) -- (663,62) -- (652,62) -- (649,64) -- (640,64) -- (639,65) -- (615,65) -- (614,66) -- (610,66) -- (609,65) -- (607,65) -- (603,67) -- (565,68) -- (564,69) -- (547,70) -- (546,71) -- (535,71) -- (534,72) -- (532,71) -- (529,71) -- (528,72) -- (522,71) -- (521,72) -- (513,72) -- (512,73) -- (509,72) -- (505,74) -- (485,74) -- (484,75) -- (472,74) -- (469,76) -- (460,76) -- (459,77) -- (450,77) -- (449,78) -- (442,77) -- (441,78) -- (416,78) -- (413,80) -- (394,80) -- (393,81) -- (374,81) -- (373,82) -- (366,81) -- (365,82) -- (362,82) -- (361,83) -- (349,83) -- (348,84) -- (337,84) -- (334,86) -- (334,89) -- (333,90) -- (333,98) -- (331,100) -- (331,105) -- (328,110) -- (328,114) -- (327,115) -- (326,122) -- (325,123) -- (324,128) -- (322,131) -- (321,137) -- (319,139) -- (319,141) -- (320,142) -- (319,152) -- (318,153) -- (318,156) -- (316,160) -- (316,164) -- (314,168) -- (314,171) -- (312,173) -- (311,181) -- (309,184) -- (309,193) -- (307,196) -- (307,200) -- (305,202) -- (305,207) -- (303,211) -- (302,218) -- (300,221) -- (300,226) -- (298,228) -- (298,231) -- (297,232) -- (297,241) -- (295,243) -- (295,248) -- (293,251) -- (293,256) -- (291,259) -- (291,262) -- (290,263) -- (289,269) -- (287,273) -- (287,278) -- (286,279) -- (285,288) -- (283,291) -- (283,296) -- (282,297) -- (282,300) -- (279,306) -- (279,309) -- (277,313) -- (276,323) -- (274,325) -- (274,327) -- (273,328) -- (273,332) -- (272,333) -- (272,340) -- (271,341) -- (270,346) -- (267,350) -- (267,354) -- (264,358) -- (263,373) -- (262,374) -- (262,377) -- (261,378) -- (262,384) -- (260,386) -- (260,388) -- (258,391) -- (258,395) -- (256,398) -- (256,402) -- (255,403) -- (255,411) -- (253,413) -- (252,415) -- (252,418) -- (250,421) -- (251,424) -- (248,427) -- (248,429) -- (247,430) -- (247,432) -- (249,434) -- (249,437) -- (245,441) -- (246,443) -- (245,445) -- (245,452) -- (246,453) -- (250,453) -- (252,452) -- (259,445) -- (260,445) -- (267,438) -- (268,438) -- (285,420) -- (286,420) -- (299,408) -- (300,408) -- (313,395) -- (315,395) -- (320,390) -- (320,388) -- (324,384) -- (327,383) -- (327,381) -- (330,378) -- (331,378) -- (343,367) -- (344,367) -- (349,361) -- (350,361) -- (360,352) -- (361,352) -- (382,329) -- (387,326) -- (397,316) -- (398,316) -- (406,308) -- (407,308) -- (415,300) -- (415,299) -- (420,294) -- (423,293) -- (434,281) -- (437,280) -- (447,270) -- (450,269) -- (471,248) -- (472,248) -- (479,241) -- (480,241) -- (481,239) -- (489,231) -- (494,228) -- (505,217) -- (506,217) -- (510,213) -- (510,212) -- (514,207) -- (515,207) -- (525,198) -- (530,195) -- (532,193) -- (532,192) -- (535,189) -- (536,189) -- (543,181) -- (547,179) -- (552,173) -- (555,172) -- (557,168) -- (566,159) -- (567,159) -- (580,147) -- (582,147) -- (584,145) -- (585,145) -- (587,143) -- (589,139) -- (597,131) -- (602,128) -- (610,119) -- (611,119) -- (618,112) -- (619,112) -- (623,108) -- (624,108) -- (630,101) -- (633,100) -- (637,96) -- (638,96) -- (639,94) -- (661,72) -- (662,72) -- (666,68) -- cycle
(656,70) -- (647,79) -- (646,79) -- (638,87) -- (637,87) -- (634,90) -- (634,92) -- (632,94) -- (630,94) -- (625,99) -- (624,99) -- (619,104) -- (618,104) -- (618,105) -- (613,110) -- (608,113) -- (602,120) -- (601,120) -- (597,124) -- (596,124) -- (588,132) -- (587,132) -- (582,137) -- (581,137) -- (573,146) -- (572,146) -- (566,152) -- (565,152) -- (547,170) -- (546,170) -- (540,176) -- (539,176) -- (503,210) -- (502,210) -- (496,217) -- (495,217) -- (489,223) -- (484,226) -- (479,231) -- (479,232) -- (476,235) -- (475,235) -- (461,249) -- (461,250) -- (458,253) -- (457,253) -- (451,259) -- (447,261) -- (443,265) -- (438,268) -- (438,269) -- (433,274) -- (428,277) -- (427,279) -- (424,282) -- (423,282) -- (416,289) -- (416,290) -- (402,304) -- (400,305) -- (400,306) -- (398,308) -- (396,308) -- (384,320) -- (383,320) -- (371,332) -- (370,332) -- (369,334) -- (356,347) -- (355,347) -- (338,363) -- (337,363) -- (336,365) -- (333,368) -- (332,368) -- (328,373) -- (324,375) -- (321,378) -- (319,383) -- (317,385) -- (315,385) -- (311,390) -- (310,390) -- (290,409) -- (288,409) -- (281,417) -- (278,418) -- (254,443) -- (252,441) -- (255,436) -- (255,432) -- (256,431) -- (256,426) -- (257,425) -- (257,422) -- (260,418) -- (260,415) -- (261,414) -- (262,408) -- (263,407) -- (264,398) -- (265,397) -- (265,395) -- (267,392) -- (267,389) -- (268,388) -- (268,384) -- (269,383) -- (268,380) -- (270,377) -- (271,369) -- (272,368) -- (272,366) -- (275,359) -- (275,353) -- (276,352) -- (276,350) -- (279,343) -- (280,336) -- (281,335) -- (280,331) -- (282,328) -- (282,323) -- (283,322) -- (283,319) -- (285,316) -- (286,306) -- (288,303) -- (289,297) -- (290,296) -- (291,285) -- (293,282) -- (293,278) -- (295,276) -- (295,273) -- (296,272) -- (296,268) -- (297,267) -- (298,258) -- (301,251) -- (301,245) -- (303,243) -- (303,241) -- (304,240) -- (303,236) -- (305,233) -- (305,228) -- (308,223) -- (308,217) -- (310,214) -- (310,210) -- (312,206) -- (312,203) -- (313,202) -- (313,200) -- (315,198) -- (316,192) -- (317,191) -- (317,186) -- (316,184) -- (318,181) -- (319,173) -- (320,172) -- (321,166) -- (322,165) -- (321,162) -- (323,160) -- (322,157) -- (325,151) -- (325,146) -- (326,145) -- (327,140) -- (329,137) -- (329,130) -- (330,129) -- (331,122) -- (333,119) -- (334,111) -- (335,110) -- (335,106) -- (336,104) -- (339,101) -- (340,101) -- (342,98) -- (338,94) -- (338,92) -- (341,89) -- (343,88) -- (345,89) -- (350,89) -- (351,88) -- (365,88) -- (368,86) -- (376,86) -- (378,87) -- (379,86) -- (391,86) -- (392,85) -- (413,85) -- (415,83) -- (432,83) -- (433,84) -- (435,84) -- (436,83) -- (439,83) -- (440,82) -- (451,82) -- (452,81) -- (460,82) -- (463,80) -- (481,80) -- (482,79) -- (484,80) -- (495,80) -- (496,79) -- (505,79) -- (506,78) -- (517,78) -- (518,77) -- (523,77) -- (524,76) -- (545,76) -- (546,75) -- (556,75) -- (557,74) -- (568,74) -- (569,73) -- (572,73) -- (574,74) -- (575,73) -- (586,73) -- (587,72) -- (604,72) -- (605,71) -- (616,71) -- (617,70) -- (634,70) -- (635,69) -- (651,69) -- (652,68) -- (654,68) -- cycle;
\node[inner sep=0pt,font=\fontsize{10}{10}\selectfont] at (339,47) {\scalebox{5.8}{$C$}};
\node[inner sep=0pt,font=\fontsize{10}{10}\selectfont] at (700,66) {\scalebox{5.9}{$B$}};
\node[inner sep=0pt,font=\fontsize{10}{10}\selectfont] at (47,77) {\scalebox{6.1}{$D$}};
\node[inner sep=0pt,font=\fontsize{10}{10}\selectfont] at (228,492) {\scalebox{6.0}{$A$}};
\end{tikzpicture}%
}
\end{minipage}\hfill
\begin{minipage}{8.1cm}\raggedright
Man verbinde $AB$; setze
\[
\begin{aligned}
BAC&=-\tfrac12a\quad\text{oder}\quad CAB=+\tfrac12a,\\
ABC&=+\tfrac12\beta,
\end{aligned}
\]
und construire das Dreieck $ABC$. Dann ist $2ACD$ die aus der Verbindung der beiden Drehungen zusammengesetzte.
\end{minipage}
\end{center}

\begin{center}\resizebox{3.2cm}{!}{%
\begin{tikzpicture}[x=1pt,y=-1pt]
\path[use as bounding box] (0,0) rectangle (539,539);
\path[fill=black,draw=none,even odd rule]
(446,54) -- (434,54) -- (433,55) -- (430,55) -- (429,56) -- (424,56) -- (423,57) -- (420,57) -- (419,58) -- (411,58) -- (408,60) -- (406,59) -- (403,59) -- (402,60) -- (395,61) -- (392,63) -- (390,62) -- (386,62) -- (380,65) -- (377,65) -- (376,66) -- (369,66) -- (368,67) -- (363,67) -- (360,69) -- (359,69) -- (357,67) -- (357,65) -- (356,64) -- (351,64) -- (350,65) -- (352,68) -- (350,70) -- (337,71) -- (334,73) -- (331,73) -- (330,74) -- (325,74) -- (324,75) -- (318,76) -- (317,77) -- (309,77) -- (308,78) -- (303,78) -- (300,80) -- (289,81) -- (286,83) -- (280,83) -- (279,84) -- (276,84) -- (273,86) -- (268,86) -- (267,87) -- (262,87) -- (261,88) -- (258,88) -- (256,87) -- (255,88) -- (251,88) -- (250,89) -- (240,90) -- (237,92) -- (233,92) -- (228,96) -- (228,97) -- (226,99) -- (226,105) -- (225,106) -- (225,109) -- (221,114) -- (221,118) -- (219,121) -- (219,123) -- (215,128) -- (215,132) -- (211,137) -- (211,139) -- (209,142) -- (209,147) -- (208,148) -- (208,151) -- (204,154) -- (204,159) -- (200,162) -- (201,164) -- (201,168) -- (197,176) -- (197,179) -- (195,181) -- (195,183) -- (193,187) -- (193,190) -- (190,196) -- (190,200) -- (187,205) -- (187,208) -- (186,209) -- (185,214) -- (183,216) -- (181,223) -- (179,225) -- (179,228) -- (176,232) -- (176,237) -- (173,242) -- (173,248) -- (169,255) -- (169,258) -- (167,260) -- (167,262) -- (165,265) -- (164,270) -- (162,272) -- (162,276) -- (160,280) -- (158,282) -- (158,286) -- (156,289) -- (156,293) -- (154,297) -- (154,301) -- (152,303) -- (151,308) -- (148,312) -- (148,314) -- (144,322) -- (144,325) -- (142,328) -- (142,330) -- (139,334) -- (139,336) -- (138,337) -- (138,342) -- (136,344) -- (136,346) -- (135,347) -- (135,351) -- (132,355) -- (132,357) -- (130,360) -- (130,363) -- (128,366) -- (127,371) -- (122,380) -- (122,384) -- (119,388) -- (119,395) -- (117,397) -- (117,401) -- (115,403) -- (115,406) -- (110,416) -- (110,419) -- (106,426) -- (106,428) -- (104,430) -- (104,433) -- (102,436) -- (102,439) -- (101,440) -- (101,444) -- (99,447) -- (99,450) -- (97,453) -- (97,457) -- (96,458) -- (96,460) -- (94,463) -- (94,465) -- (92,467) -- (90,471) -- (90,475) -- (91,476) -- (95,476) -- (99,473) -- (100,470) -- (104,467) -- (105,464) -- (124,444) -- (124,443) -- (129,438) -- (129,437) -- (133,433) -- (133,432) -- (146,419) -- (147,415) -- (160,401) -- (160,400) -- (166,395) -- (169,389) -- (172,386) -- (173,386) -- (177,381) -- (178,381) -- (180,379) -- (184,378) -- (186,376) -- (185,372) -- (188,369) -- (188,368) -- (194,363) -- (196,358) -- (209,344) -- (212,338) -- (221,329) -- (221,328) -- (226,323) -- (226,322) -- (231,318) -- (231,317) -- (239,308) -- (240,305) -- (244,301) -- (245,298) -- (250,293) -- (250,292) -- (261,281) -- (263,277) -- (268,272) -- (268,271) -- (273,266) -- (273,265) -- (278,260) -- (278,259) -- (287,251) -- (288,248) -- (296,239) -- (298,235) -- (309,224) -- (309,223) -- (312,220) -- (313,217) -- (317,214) -- (317,213) -- (321,209) -- (321,208) -- (337,190) -- (338,187) -- (349,176) -- (353,176) -- (354,175) -- (353,169) -- (366,158) -- (366,157) -- (370,153) -- (372,147) -- (380,140) -- (383,135) -- (383,133) -- (409,104) -- (410,100) -- (430,80) -- (431,76) -- (438,69) -- (439,66) -- (446,60) -- (447,58) -- (447,55) -- cycle
(436,59) -- (437,60) -- (437,62) -- (427,73) -- (427,74) -- (418,84) -- (418,85) -- (406,97) -- (401,105) -- (397,109) -- (396,109) -- (394,111) -- (393,114) -- (383,123) -- (383,124) -- (376,131) -- (375,134) -- (362,147) -- (362,149) -- (360,151) -- (357,151) -- (357,154) -- (355,156) -- (355,157) -- (352,160) -- (352,161) -- (348,164) -- (346,169) -- (340,174) -- (339,177) -- (331,186) -- (331,188) -- (327,191) -- (327,192) -- (324,195) -- (324,198) -- (322,200) -- (320,200) -- (317,203) -- (313,210) -- (310,213) -- (309,213) -- (308,215) -- (298,225) -- (298,226) -- (292,231) -- (292,232) -- (289,235) -- (289,237) -- (282,244) -- (281,247) -- (266,263) -- (266,265) -- (262,269) -- (261,269) -- (251,280) -- (248,285) -- (242,291) -- (241,295) -- (226,311) -- (226,312) -- (221,317) -- (221,318) -- (216,322) -- (215,325) -- (198,344) -- (193,352) -- (187,358) -- (185,358) -- (184,362) -- (178,368) -- (176,368) -- (174,370) -- (174,371) -- (171,374) -- (170,378) -- (168,380) -- (167,380) -- (161,387) -- (160,390) -- (155,395) -- (154,398) -- (148,403) -- (147,407) -- (144,410) -- (144,411) -- (139,416) -- (138,416) -- (135,419) -- (134,422) -- (122,433) -- (122,435) -- (118,440) -- (118,441) -- (107,453) -- (107,454) -- (103,458) -- (102,458) -- (101,457) -- (101,455) -- (109,439) -- (109,435) -- (110,434) -- (110,432) -- (113,428) -- (114,423) -- (119,413) -- (120,407) -- (122,405) -- (122,402) -- (124,398) -- (126,396) -- (126,393) -- (129,390) -- (129,385) -- (130,384) -- (130,380) -- (132,377) -- (132,375) -- (134,373) -- (137,373) -- (138,372) -- (138,370) -- (137,370) -- (135,368) -- (135,367) -- (138,362) -- (138,360) -- (140,358) -- (141,353) -- (145,345) -- (145,340) -- (146,339) -- (146,336) -- (149,330) -- (149,328) -- (151,325) -- (151,322) -- (156,313) -- (156,311) -- (157,310) -- (157,306) -- (160,302) -- (160,299) -- (162,296) -- (162,292) -- (166,284) -- (166,281) -- (173,267) -- (173,264) -- (176,260) -- (177,253) -- (179,251) -- (179,248) -- (180,247) -- (180,243) -- (183,238) -- (183,234) -- (185,232) -- (186,227) -- (188,225) -- (188,221) -- (191,218) -- (192,211) -- (195,206) -- (195,203) -- (197,200) -- (197,195) -- (199,191) -- (201,189) -- (201,187) -- (203,184) -- (203,181) -- (205,177) -- (205,174) -- (207,171) -- (207,168) -- (211,163) -- (211,159) -- (212,158) -- (213,153) -- (215,151) -- (215,146) -- (218,141) -- (219,136) -- (222,131) -- (223,124) -- (226,120) -- (227,114) -- (229,112) -- (229,107) -- (231,105) -- (231,101) -- (233,99) -- (237,99) -- (240,97) -- (244,97) -- (245,96) -- (248,96) -- (249,95) -- (258,94) -- (259,93) -- (262,93) -- (263,92) -- (273,91) -- (274,90) -- (278,90) -- (279,89) -- (286,89) -- (288,87) -- (300,86) -- (303,84) -- (307,84) -- (308,83) -- (316,83) -- (317,82) -- (322,82) -- (323,81) -- (332,81) -- (335,79) -- (338,79) -- (339,78) -- (349,77) -- (350,76) -- (355,76) -- (356,75) -- (363,75) -- (366,73) -- (380,72) -- (384,69) -- (390,69) -- (391,68) -- (394,68) -- (395,67) -- (401,67) -- (402,66) -- (411,65) -- (412,64) -- (415,64) -- (416,63) -- (420,64) -- (423,62) -- (429,61) -- (431,59) -- cycle;
\node[inner sep=0pt,font=\fontsize{10}{10}\selectfont] at (480,48) {\scalebox{6.1}{$C$}};
\node[inner sep=0pt,font=\fontsize{10}{10}\selectfont] at (189,75) {\scalebox{6.1}{$B$}};
\node[inner sep=0pt,font=\fontsize{10}{10}\selectfont] at (48,489) {\scalebox{6.0}{$A$}};
\end{tikzpicture}%
}
\end{center}
II. Anstatt Drehungen um Punkte betrachte man Schiebungen längs grössten Kreisen.

Dann ist die Schiebung $2AC$ aus $2AB$ und $2BC$ zusammengesetzt.

\Needspace{10\baselineskip}
\begin{center}\bfseries BEMERKUNGEN.\end{center}

Die vorstehenden Notizen [I], [II] und [III] stehen in einem Handbuche. Die folgende Notiz [IV] befindet sich auf einem einzelnen Zettel ohne Datirung. Sie zeigt, dass \textsc{Gauss} den Zusammenhang der linearen Substitutionen einer complexen Veränderlichen mit den Drehungen der Kugelfläche in sich selbst erkannt hatte. Aus diesem Grunde ist die Notiz [V], die sich auf denselben Gegenstand bezieht, jedoch in rein geometrischer Behandlungsweise, hinzugefügt worden; sie findet sich in einem andern Handbuche und ist vermuthlich schon vor 1819 niedergeschrieben worden.
\begin{flushright}\textsc{Stäckel.}\end{flushright}
%
\clearpage
\SourceLeaf{364}
\begin{center}\bfseries [MUTATIONEN DES RAUMES.]\end{center}
\begin{center}[I.]\end{center}
\begin{center}\bfseries [TRANSFORMATIONEN DES RAUMES.]\end{center}
\begin{center}[1.]\end{center}

Man setze
\[
\begin{aligned}
\BareRoot{\frac{1-p-q-r}{2}}&=a\\
\BareRoot{\frac{1-p+q+r}{2}}&=b\\
\BareRoot{\frac{1+p-q+r}{2}}&=c\\
\BareRoot{\frac{1+p+q-r}{2}}&=d,
\end{aligned}
\]
so ist
\[
\begin{array}{ccc}
p&ad+bc&ac-bd\\
bc-ad&q&ab+cd\\
-ac-bd&cd-ab&r
\end{array}
\]
das allgemeine Transformationsschema der Raumcoordinaten, oder, bedeuten $a,b,c,d$ was sie wollen[, so sind die 9 Grössen des Schemas] proportional folgenden:
\[
\begin{array}{ccc}
\tfrac12(cc+dd-aa-bb)&ad+bc&ac-bd\\
bc-ad&\tfrac12(bb+dd-aa-cc)&ab+cd\\
-ac-bd&cd-ab&\tfrac12(bb+cc-aa-dd).
\end{array}
\]
%
\SourceLeaf{365}
\Needspace{12\baselineskip}
\begin{center}[2.]\end{center}

Besser auf folgende Weise:
\[
\begin{aligned}
\BareRoot{\frac{1+p+q+r}{2}}&=a\\
\BareRoot{\frac{1+p-q-r}{2}}&=b\\
\BareRoot{\frac{1-p+q-r}{2}}&=c\\
\BareRoot{\frac{1-p-q+r}{2}}&=d,
\end{aligned}
\]
\[
\begin{array}{ccc}
\tfrac12(aa+bb-cc-dd)&-ad-bc&-ac+bd\\
+ad-bc&\tfrac12(aa-bb+cc-dd)&-ab-cd\\
+ac+bd&+ab-cd&\tfrac12(aa-bb-cc+dd).
\end{array}
\]

\Needspace{13\baselineskip}
\begin{center}[3.]\end{center}

Aus der Verbindung zweier Transformationen,
\[
\left.\begin{array}{llll}
\text{deren}&\text{erster}&\text{die Scale}&a,\ b,\ c,\ d\\
\text{,,}&\text{zweiter}&\text{,,\quad,,}&\alpha,\ \beta,\ \gamma,\ \delta
\end{array}\right\}\quad\text{entspricht,}
\]
entsteht eine neue, deren Scale:
\[
\begin{gathered}
a\alpha-b\beta-c\gamma-d\delta\\
a\beta+b\alpha-c\delta+d\gamma\\
a\gamma+b\delta+c\alpha-d\beta\\
a\delta-b\gamma+c\beta+d\alpha.
\end{gathered}
\]

\Needspace{10\baselineskip}
\begin{center}[4.]\end{center}

Sind die Coordinaten des ruhenden Punkts $\xi,\eta,\zeta$, Vergrösserung $=nn$, Drehung $=\lambda$, so kann man setzen
%
\SourceLeaf{366}
\[
\begin{aligned}
a&=n\cos\tfrac12\lambda\\
b&=n\xi\sin\tfrac12\lambda\\
c&=n\eta\sin\tfrac12\lambda\\
d&=n\zeta\sin\tfrac12\lambda.
\end{aligned}
\]

\Needspace{14\baselineskip}
\begin{center}[5.]\end{center}

Schreiben wir
\[
\begin{aligned}
a\alpha-b\beta-c\gamma-d\delta&=A\\
a\beta+b\alpha-c\delta+d\gamma&=B\\
a\gamma+b\delta+c\alpha-d\beta&=C\\
a\delta-b\gamma+c\beta+d\alpha&=D;\\
aa+bb+cc+dd&=m\\
\alpha\alpha+\beta\beta+\gamma\gamma+\delta\delta&=\mu,
\end{aligned}
\]
so ist:
\[
\begin{aligned}
\frac{C+Di}{A+Bi}&\equiv\frac{c+di}{a+bi}&&\text{(mod. }m\text{)},\\[3pt]
\frac{C-Di}{A+Bi}&\equiv\frac{\gamma-\delta i}{\alpha+\beta i}&&\text{(mod. }\mu\text{)};\\[12pt]
\frac{B+Di}{A+Ci}&\equiv\frac{\beta+\delta i}{\alpha+\gamma i}&&\text{(mod. }m\text{)},\\[3pt]
\frac{B-Di}{A+Ci}&\equiv\frac{b-di}{a+ci}&&\text{(mod. }\mu\text{)};\\[12pt]
\frac{B+Ci}{A+Di}&\equiv\frac{b+ci}{a+di}&&\text{(mod. }m\text{)},\\[3pt]
\frac{B-Ci}{A+Di}&\equiv\frac{\beta-\gamma i}{\alpha+\delta i}&&\text{(mod. }\mu\text{)}.
\end{aligned}
\]

\Needspace{6\baselineskip}
\begin{center}[6.]\end{center}

Wir bezeichnen allgemein die Combination $a,b,c,d$ durch $(a,b,c,d)$ und schreiben
\[
(a,b,c,d)(\alpha,\beta,\gamma,\delta)=(A,B,C,D).
\]
%
\SourceLeaf{367}%
\noindent Es ist also $(a,b,c,d)(\alpha,\beta,\gamma,\delta)$ nicht mit $(\alpha,\beta,\gamma,\delta)(a,b,c,d)$ zu verwechseln.

Man hat dann
\[
(a,b,c,d)(a,-b,-c,-d)=(aa+bb+cc+dd,0,0,0).
\]
Ferner bezeichne man die Combination $(a,b,c,d)$ durch einen Buchstaben, z.~B. $g$, und dann die Combination $(a,-b,-c,-d)$ durch $g'$.

Es ist also
\[
gg'=g'g=(aa+bb+cc+dd,0,0,0)=(m,0,0,0).
\]
Ferner ist $(\alpha,\beta,\gamma,\delta)=h$, $\alpha\alpha+\beta\beta+\gamma\gamma+\delta\delta=\mu$ gesetzt:
\[
ghg'=(\alpha m,-\beta m,-\gamma m,-\delta m)=h'gg'=gg'h'.
\]

\Needspace{15\baselineskip}
\begin{center}[7.]\end{center}

Man kann auch die Scale $a,b,c,d$ als auf diejenige Versetzung der Coordinatenaxen sich beziehend betrachten, die der Substitution
\[
\begin{array}{ccc}
+a+\dfrac{bb}{n+a},&+d+\dfrac{bc}{n+a},&-c+\dfrac{bd}{n+a},\\[6pt]
-d+\dfrac{bc}{n+a},&+a+\dfrac{cc}{n+a},&+b+\dfrac{cd}{n+a},\\[6pt]
+c+\dfrac{bd}{n+a},&-b+\dfrac{cd}{n+a},&+a+\dfrac{dd}{n+a}
\end{array}
\]
entspricht, wo
\[
\begin{aligned}
n&=\BareRoot{(aa+bb+cc+dd)}&&\text{Vergrösserung,}\\
a&=n\cos\lambda,&\BareRoot{(bb+cc+dd)}&=n\sin\lambda,\\
\lambda&\text{ ganze Drehung.}
\end{aligned}
\]

\Needspace{9\baselineskip}
\begin{center}[8.]\end{center}

Die Transformationsscale $a,b,c,d$ bedeutet\\
eine Vergrösserung $=(aa+bb+cc+dd)$\\
nebst Drehung $=\lambda$ um den Punkt, dessen Coordinaten
\[
\frac ba\operatorname{cotg}\tfrac12\lambda,\qquad
\frac ca\cot\tfrac12\lambda,\qquad
\frac da\operatorname{cotg}\tfrac12\lambda.
\]
%
\par
\SourceLeaf{368}%
Geben drei auf einander folgende Scalen
\[
M,0,0,0
\]
und sind die betreffenden Drehungspunkte $P,P',P''$,\\
\hspace*{5em}die Winkel $\lambda,\lambda',\lambda''$,\\
so sind $\tfrac12\lambda,\tfrac12\lambda',\tfrac12\lambda''$ die Winkel des sphärischen Dreiecks zwischen $P,P',P''$.

\Needspace{14\baselineskip}
\begin{center}[II.]\end{center}
\begin{center}\bfseries MUTATIONEN DES RAUMES VON DREI DIMENSIONEN.\end{center}
\begin{center}1.\end{center}

Seit langer Zeit habe ich das Verfahren gebraucht, jede Mutation eines Raumes (worunter ich verstehe seine Verrückung in einem andern Raume verbunden mit einer allgemeinen Vergrösserung oder Verkleinerung des erstern, unbeschadet der Conservation der Ähnlichkeit) durch vier Grössen
\[
a,b,c,d
\]
auszudrücken, deren Complex ich die Mutationsscale nenne. Es wird zugleich angenommen, dass Ein Punkt des beweglichen Raumes in dem andern als fest oder absolut betrachteten immer unbewegt bleibe, welcher Punkt mit Nullpunkt bezeichnet werden mag.

\Needspace{9\baselineskip}
\begin{center}2.\end{center}

Die Bedeutung dieser vier Grössen ist hierbei folgende. In dem festen Raume werden drei von $0$ ausgehende, rechte Winkel mit einander machende gerade Linien angenommen $01,02,03$. Man setze
\[
\BareRoot{(bb+cc+dd)}=\rho,
\]
und lasse
\[
\frac b\rho,\ \frac c\rho,\ \frac d\rho
\]
%
\SourceLeaf{369}%
\noindent die Cosinus der drei Winkel sein, welche eine Gerade $0P$ resp. mit $01,02,03$ macht. Diese gerade Linie hat in dem bewegten Raume nach der Mutation dieselbe Lage wie vor derselben; [wird ferner]
\[
\BareRoot{(aa+\rho\rho)}=k
\]
gesetzt, oder
\[
k=\BareRoot{(aa+bb+cc+dd)},
\]
ist $k$ die linearische Vergrösserung des mutirten Raumes.

Endlich
\[
a=k\cos\theta,\qquad \rho=k\sin\theta
\]
gesetzt, ist $2\theta$ diejenige Drehung um $0P$, durch welche, verbunden mit der Vergrösserung $k$, die Mutation erzeugt werden kann.

\Needspace{12\baselineskip}
\begin{center}\bfseries BEMERKUNG.\end{center}

Die unter [I] zusammengefassten Notizen stehen in einem Handbuche. Die Nummern 1 bis 7 stimmen in Schrift und Tinte genau mit astronomischen Aufzeichnungen aus dem Jahre 1819 überein, die sich auf den vorhergehenden Seiten befinden. Verschieden von ihnen sieht die als Nummer 8 bezeichnete Bemerkung aus. Sie trägt sachlich den Character eines spätern Zusatzes, und diese Annahme wird auch dadurch bestätigt, dass augenscheinlich gleichzeitig mit ihr niedergeschriebene astronomische Notizen aus der Zeit von 1822 bis 1823 stammen. Die Notiz [II] füllt einen einzelnen Zettel ohne Datirung. Nach der Handschrift zu urtheilen, stammt sie aus noch späterer Zeit.
\begin{flushright}\textsc{Stäckel.}\end{flushright}
%
\end{document}
